% Options for packages loaded elsewhere
\PassOptionsToPackage{unicode}{hyperref}
\PassOptionsToPackage{hyphens}{url}
\documentclass[
  ukrainian,
]{article}
\usepackage{xcolor}
\usepackage[margin=2.5cm]{geometry}
\usepackage{amsmath,amssymb}
\setcounter{secnumdepth}{-\maxdimen} % remove section numbering
\usepackage{iftex}
\ifPDFTeX
  \usepackage[T1]{fontenc}
  \usepackage[utf8]{inputenc}
  \usepackage{textcomp} % provide euro and other symbols
\else % if luatex or xetex
  \usepackage{unicode-math} % this also loads fontspec
  \defaultfontfeatures{Scale=MatchLowercase}
  \defaultfontfeatures[\rmfamily]{Ligatures=TeX,Scale=1}
\fi
\usepackage{lmodern}
\ifPDFTeX\else
\setmainfont{DejaVu Serif}
\setsansfont{DejaVu Sans}
\setmonofont{DejaVu Sans Mono}
\setmathfont{STIX Two Math}
\fi
\ifPDFTeX\else
  % xetex/luatex font selection
\fi
% Use upquote if available, for straight quotes in verbatim environments
\IfFileExists{upquote.sty}{\usepackage{upquote}}{}
\IfFileExists{microtype.sty}{% use microtype if available
  \usepackage[]{microtype}
  \UseMicrotypeSet[protrusion]{basicmath} % disable protrusion for tt fonts
}{}
\makeatletter
\@ifundefined{KOMAClassName}{% if non-KOMA class
  \IfFileExists{parskip.sty}{%
    \usepackage{parskip}
  }{% else
    \setlength{\parindent}{0pt}
    \setlength{\parskip}{6pt plus 2pt minus 1pt}}
}{% if KOMA class
  \KOMAoptions{parskip=half}}
\makeatother
\usepackage{longtable,booktabs,array}
\newcounter{none} % for unnumbered tables
\usepackage{calc} % for calculating minipage widths
% Correct order of tables after \paragraph or \subparagraph
\usepackage{etoolbox}
\makeatletter
\patchcmd\longtable{\par}{\if@noskipsec\mbox{}\fi\par}{}{}
\makeatother
% Allow footnotes in longtable head/foot
\IfFileExists{footnotehyper.sty}{\usepackage{footnotehyper}}{\usepackage{footnote}}
\makesavenoteenv{longtable}
\usepackage{graphicx}
\graphicspath{{/Users/kalinicenkovitalijmikolajovic/.codex/worktrees/5c15/Script_for_translate_book/Output_Final/}}
\setkeys{Gin}{width=0.66\linewidth,height=0.70\textheight,keepaspectratio}
\makeatletter
\newsavebox\pandoc@box
\newcommand*\pandocbounded[1]{% scales image to fit in text height/width
  \sbox\pandoc@box{#1}%
  \Gscale@div\@tempa{\textheight}{\dimexpr\ht\pandoc@box+\dp\pandoc@box\relax}%
  \Gscale@div\@tempb{\linewidth}{\wd\pandoc@box}%
  \ifdim\@tempb\p@<\@tempa\p@\let\@tempa\@tempb\fi% select the smaller of both
  \ifdim\@tempa\p@<\p@\scalebox{\@tempa}{\usebox\pandoc@box}%
  \else\usebox{\pandoc@box}%
  \fi%
}
% Set default figure placement to htbp
\def\fps@figure{htbp}
\makeatother
\ifLuaTeX
\usepackage[bidi=basic,shorthands=off,provide=*]{babel}
\else
\usepackage[bidi=default,shorthands=off,provide=*]{babel}
\fi
\ifLuaTeX
  \usepackage{selnolig} % disable illegal ligatures
\fi
\setlength{\emergencystretch}{3em} % prevent overfull lines
\providecommand{\tightlist}{%
  \setlength{\itemsep}{0pt}\setlength{\parskip}{0pt}}
\usepackage{bookmark}
\IfFileExists{xurl.sty}{\usepackage{xurl}}{} % add URL line breaks if available
\urlstyle{same}
\hypersetup{
  pdflang={uk},
  hidelinks,
  pdfcreator={LaTeX via pandoc}}

\author{}
\date{}

\usepackage{float}
\usepackage[section]{placeins}
\raggedbottom
\clubpenalty=10000
\widowpenalty=10000
\displaywidowpenalty=10000
\brokenpenalty=10000
\emergencystretch=2em
\allowdisplaybreaks[2]
\setlength{\textfloatsep}{12pt plus 2pt minus 4pt}
\setlength{\floatsep}{10pt plus 2pt minus 2pt}
\setlength{\intextsep}{10pt plus 2pt minus 2pt}
\makeatletter
\def\fps@figure{htbp}
\makeatother
\begin{document}

ЧАСТКОВО ДИФЕРЕНЦІАЛ РІВНЯННЯ для Вчені та Інженери

Стенлі Дж. Фарлоу, професор математики, Університет Мену

\section{С. Фарлоу}\label{ux441.-ux444ux430ux440ux43bux43eux443}

\section{РІВНЯННЯ З ПОХІДНИМИ ПОХІДНИМИ ДЛЯ НАУКОВЦІВ І
ІНЖЕНЕРІВ}\label{ux440ux456ux432ux43dux44fux43dux43dux44f-ux437-ux43fux43eux445ux456ux434ux43dux438ux43cux438-ux43fux43eux445ux456ux434ux43dux438ux43cux438-ux434ux43bux44f-ux43dux430ux443ux43aux43eux432ux446ux456ux432-ux456-ux456ux43dux436ux435ux43dux435ux440ux456ux432}

Перекладено з англійської \textbf{A.I.} Плісом

Редактор

С. І. Похожаєва

5\textbar{} 66K 22.161.6 φ24 υ ΥΠΚ 517.2

\section{Фарлоу С.}\label{ux444ux430ux440ux43bux43eux443-ux441.}

F 24 Рівняння з частими похідними для науковців та інженерів. Москва,
видавництво Мир, 1985, 384 с.

King's American Mathematician, який є підручником з теорії
диференціальних рівнянь у похідних похідних. Він характеризується
компактністю, ясністю та ясністю подачі, а також неформальним підходом
до подачі матеріалу. Вона містить багато ілюстрацій, графіків і діаграм,
але замість суворих доказів часто даються роздуми, засновані на інтуїції
або аналогіях.

Для інженерів і нематематиків --- біологів, диміків, а також студентів
університету.
\(\Phi = \frac{1702050000-380}{041(01)-85} 26-85, \text{ q. 1}\) BBK
22.161.6

Редакційна колегія журналу Literature on Mathematical Sciences

© переклад російською, Мир, 1985

© John Wiley \& Sons, Inc.~1982. Всі права захищені. Авторизований
переклад з англомовного видання, опублікований John Wiley \& Sons, Inc.

\section{З РЕДАКТОРА
ПЕРЕКЛАДУ}\label{ux437-ux440ux435ux434ux430ux43aux442ux43eux440ux430-ux43fux435ux440ux435ux43aux43bux430ux434ux443}

Книга, запропонована для читачів, слід розглядати як підручник для
курсів: «Диференціальні рівняння в похідних похідних» та «Рівняння
математичної фізики» для студентів технічних університетів. Вона також
буде корисною для інженерів різних спеціальностей. У книзі представлено
традиційні методи розв'язання диференціальних рівнянь у похідних
похідних, включаючи деякі чисельні методи. Водночас особлива увага
приділяється методу розділення змінних і методу інтегральних
перетворень, і цю частину книги можна використовувати на практичних
заняттях в університетах

Орієнтація книги на широку «нематематичну» аудиторію визначила стиль
викладання розглянутих питань: інтуїтивний підхід і велика увага до
фізичного значення не лише самих рівнянь, а й меж і початкових умов для
різних задач. Вважається, що читачі знайомі лише з основами
диференціального та інтегрального числення та елементами теорії
звичайних диференціальних рівнянь.

Фахівці можуть не задовольнитися інтуїтивним описом певних понять і
методів автором, але більш суворий виклад неминуче призведе до
ускладнення книги і звузить коло її читачів, що, очевидно, не було
наміром автора.

Більшим запереченням може бути відсутність загальних результатів щодо
задач для диференціальних рівнянь у похідних з змінними коефіцієнтами,
включно з рівняннями вищого порядку. Однак читач має чітко розуміти, що
ця книга є вступним курсом з теорії диференціальних рівнянь у часткових
похідних. Питання лінійної теорії диференціальних рівнянь, які автор не
включив, можна знайти в літературі, наведеній наприкінці книги.

Матеріал книги подається автором з великою методологічною майстерністю.
Однак не все було однаково успішним, і винний читач знайде це,
наприклад, у вивченні принципу Дюамеля, сферичних заповітів та деяких
інших уривків у книзі.

Без сумніву, запропонована книга С. Фарлоу знайде читача не лише серед
студентів і інженерів, а й серед викладачів, які можуть багато чого
навчитися з конкретного способу презентації та розробки спеціальних
математичних курсів.

\section{ПЕРЕДМОВА}\label{ux43fux435ux440ux435ux434ux43cux43eux432ux430}

За останні кілька років спостерігається значне зростання кількості
учнів, які починають вивчати диференціальні рівняння з похідними у своїх
одинадцятих класах. Крім того, багато з них спеціалізуються на сферах,
де інтуїція важливіша за математичну строгість. Тому, пишучи цю книгу, я
прагнув створити текст, який стимулює інтуїтивне мислення читачів, не
втрачаючи при цьому значної математичної строгості. Можливо, на одному
полюсі, подати предмет на високому рівні «епсилон-дельта», але тоді, як
це зазвичай буває, студенти не знатимуть, що робити з отриманою
інформацією. Інша крайність --- повне ігнорування математичних тонкощів,
але тоді ні учні, ні вчителі не знатимуть, як діяти. Я намагався знайти
розумний компроміс між цими крайнощами.

Ця книга виникла з курсу лекцій, які я читав протягом останніх п'яти
років. Розташування матеріалу не є загальноприйнятим --- вони майже
незалежні один від одного

Друг лекцій.

Основна увага приділяється двом найважливішим аналітичним методам:
розділенню змінних та інтегральним перетворенням. Також представлені
деякі нестандартні розділи, такі як метод Монте-Карло, варіаційне
числення, теорія керування, теорія потенціалів, інтегральні рівняння. Це
пов'язано з тим, що більшості студентів доведеться стикатися з цими
питаннями під час навчання, і якщо вони не ознайомляться з ними зараз,
навряд чи коли-небудь систематично їх вивчатимуть.

Цю книгу можна використовувати як навчальний посібник для одно- або
двосеместрового курсу як для початківців, так і для старших студентів.
Читачі повинні знати лише диференціальне та інтегральне числення та
звичайні диференціальні рівняння. Матеріал більшості лекцій подається в
одному-двох уроках. Типовий односеместровий курс може включати лекції
1-13, 15-17, 19-20, 22-23, 25-27, 30-32, 37-39. Усі 47 лекцій легко
читати протягом двох семестрів. Водночас часу буде достатньо для
вирішення проблем.

Автор вдячний Вайлі за пропозицію написати книгу, а також рецензентам
--- професорам Крісу Рорресу та М. Куршиду Алі, які зробили низку
корисних коментарів до тексту книги. Будь-які пропозиції щодо подальшого
покращення цього підручника будуть вдячні як від учнів, так і від
викладачів. Я також вдячний Дороті, Сюзан, Александер і Дезі Фарлоу.

Стенлі Дж. Фарлоу

\section{ВСТУП ДО ТЕОРІЇ ДИФЕРЕНЦІАЛЬНИХ РІВНЯНЬ У ПОХІДНИХ
ПОХІДНИХ}\label{ux432ux441ux442ux443ux43f-ux434ux43e-ux442ux435ux43eux440ux456ux457-ux434ux438ux444ux435ux440ux435ux43dux446ux456ux430ux43bux44cux43dux438ux445-ux440ux456ux432ux43dux44fux43dux44c-ux443-ux43fux43eux445ux456ux434ux43dux438ux445-ux43fux43eux445ux456ux434ux43dux438ux445}

МЕТА ЛЕКЦІЇ: Показати, що таке диференціальні рівняння в похідних
похідних, де і як вони виникають і як їх розв'язують. Коротко розгляньте
методи класифікації рівнянь і надайте список базових понять, які будуть
детально вивчені на наступних лекціях.

Більшість фізичних явищ у таких сферах, як гідродинаміка, електрика та
магнетизм, механіка, оптика та теплопередача, загалом можна описати за
допомогою рівнянь у похідних похідних (PNP). Більшість рівнянь у
математичній фізиці є диференціальними рівняннями з похідними похідними.
Правда, з деякими спрощеними припущеннями ці рівняння зводяться до
звичайних диференціальних рівнянь, але повний опис таких систем неминуче
призводить до використання диференціальних рівнянь у похідних похідних.

\section{Що таке часткові диференціальні
рівняння?}\label{ux449ux43e-ux442ux430ux43aux435-ux447ux430ux441ux442ux43aux43eux432ux456-ux434ux438ux444ux435ux440ux435ux43dux446ux456ux430ux43bux44cux43dux456-ux440ux456ux432ux43dux44fux43dux43dux44f}

Рівняння з частковими похідними --- це рівняння, яке містить
диференціали у похідних похідних. На відміну від звичайних
диференціальних рівнянь (TAC), де невідома функція залежить лише від
однієї змінної, у рівняннях часткових похідних вона залежить від кількох
змінних (наприклад, температура u(x, t) залежить від координати x і часу
t).

Тепер наведемо деякі з найважливіших диференціальних рівнянь у похідних
похідних. Для спрощення запису ми використаємо наступні нотації

\[u_t \equiv \frac{\partial u}{\partial t}\],
\(u_x \equiv \frac{\partial u}{\partial x}\) ,
\(u_{xx} \equiv \frac{\partial^2 u}{\partial x^2}\) , \ldots{}

\section{Деякі диференціальні рівняння в похідних
похідних}\label{ux434ux435ux44fux43aux456-ux434ux438ux444ux435ux440ux435ux43dux446ux456ux430ux43bux44cux43dux456-ux440ux456ux432ux43dux44fux43dux43dux44f-ux432-ux43fux43eux445ux456ux434ux43dux438ux445-ux43fux43eux445ux456ux434ux43dux438ux445}

\[u_t = u_{xx}\]

(одновимірне рівняння теплопровідності), \(u_t = u_{xx} + u_{yy}\)
(двовимірне рівняння теплопровідності),
\(u_{rr} + \frac{1}{r}u_r + \frac{1}{r^2}u_{00} = 0\) (рівняння Лапласа
за полярними координатами), \(u_{tt} = u_{xx} + u_{yy} + u_{zz}\)
(тривимірне хвильове рівняння), \(u_{tt} = u_{xx} + zu_t + \beta u\)
(телеграфне рівняння).

\section{Примітки щодо наведених
прикладів}\label{ux43fux440ux438ux43cux456ux442ux43aux438-ux449ux43eux434ux43e-ux43dux430ux432ux435ux434ux435ux43dux438ux445-ux43fux440ux438ux43aux43bux430ux434ux456ux432}

У всіх наведених вище прикладах функція невідома і залежить від більше
ніж однієї змінної. Змінна і (яку ми диференціюємо) називається залежною
змінною. Змінні, за якими відбувається диференціювання, називаються
незалежними змінними. Наприклад, у рівнянні.

\[u_t = u_{xx}\]

Залежна змінна u(x, t) є функцією двох незалежних змінних x і t; та в
рівнянні

\[u_t = u_{rr} + \frac{1}{r} u_r + \frac{1}{r^2} u_{00}\]

змінна \(u(r, \theta, t)\) залежить від \(r, \theta, t\) .

\section{Чому необхідно вивчати диференціальні рівняння у похідних
похідних?}\label{ux447ux43eux43cux443-ux43dux435ux43eux431ux445ux456ux434ux43dux43e-ux432ux438ux432ux447ux430ux442ux438-ux434ux438ux444ux435ux440ux435ux43dux446ux456ux430ux43bux44cux43dux456-ux440ux456ux432ux43dux44fux43dux43dux44f-ux443-ux43fux43eux445ux456ux434ux43dux438ux445-ux43fux43eux445ux456ux434ux43dux438ux445}

Більшість фізичних законів природи можна сформулювати через рівняння з
похідними похідними. Прикладами є рівняння Максвелла, закон
теплопередачі Ньютона, рівняння Нав'є-Стокса, рівняння руху Ньютона та
рівняння Шредінгера в кватичній механіці. У всіх цих рівняннях фізичні
явища описуються мовою просторових і часових похідних. Похідні
з'являються в рівняннях, оскільки вони описують найважливіші фізичні
величини (такі як швидкість, прискорення, сила, тертя, потік, струм
тощо). Таким чином, виникають диференціальні рівняння в частинних
похідних, які містять невідому функцію, яку потрібно визначити.

Мета цієї книги --- показати читачеві:

\begin{enumerate}
\def\labelenumi{\arabic{enumi}.}
\item
  Як сформулювати фізичну задачу у формі диференціального рівняння у
  похідних похідних (побудова математичної моделі).
\item
  Як розв'язати диференціальне рівняння в похідних (з урахуванням
  початкових і граничних умов).
\end{enumerate}

Перш ніж перейти до побудови математичних моделей, давайте коротко
розглянемо методи розв'язання рівнянь у похідних похідних.

\section{Як розв'язати диференціальні рівняння в похідних
похідних?}\label{ux44fux43a-ux440ux43eux437ux432ux44fux437ux430ux442ux438-ux434ux438ux444ux435ux440ux435ux43dux446ux456ux430ux43bux44cux43dux456-ux440ux456ux432ux43dux44fux43dux43dux44f-ux432-ux43fux43eux445ux456ux434ux43dux438ux445-ux43fux43eux445ux456ux434ux43dux438ux445}

Гарне питання. Виявляється, існує цілий арсенал методів, придатних для
практичного використання. Найважливішими є ті, у яких диференціальні
рівняння з похідними похідними зводяться до звичайних диференціальних
рівнянь. Перелічимо десять методів розв'язання диференціальних рівнянь у
похідних похідних:

\begin{enumerate}
\def\labelenumi{\arabic{enumi}.}
\tightlist
\item
  Метод розділення змінних. Рівняння з факторами, отриманими з n
  незалежними змінними, зводиться до n звичайних-
\end{enumerate}

диференціальнірівняння 1).

\begin{enumerate}
\def\labelenumi{\arabic{enumi}.}
\setcounter{enumi}{1}
\item
  Метод інтегральних перетворень. Рівняння в похідних з n незалежними
  змінними зводиться до рівняння з похідними у (n-1) незалежних змінних;
  Відповідно, рівняння з частковими похідними з двома незалежними
  змінними можна звести до звичайного диференціального рівняння \(^2\)
  ).
\item
  Метод перетворення координат. Початкове диференціальне рівняння в
  похідних похідних зводиться до звичайного рівняння з частковими
  похідними або до іншого, простішого рівняння у похідних похідних
  шляхом відповідного перетворення координат (наприклад, обертання осей
  координат тощо).
\item
  Перетворення залежної змінної. Початкове диференціальне рівняння в
  похідних похідних перетворюється у таке рівняння з похідними похідними
  для іншої невідомої функції, яка
\end{enumerate}

легше розв'язати, ніж початкові 3).

\begin{enumerate}
\def\labelenumi{\arabic{enumi}.}
\setcounter{enumi}{4}
\tightlist
\item
  Чисельні методи. Початкове диференціальне рівняння в похідних похідних
  зводиться до системи різницей, яку розв'язують методом ітерацій на
  комп'ютері. У багатьох випадках це єдиний спосіб розв'язати
  диференціальне рівняння з похідними похідними. Окрім різницьких
  методів для розв'язання рівнянь з частковими про-
\end{enumerate}

\begin{enumerate}
\def\labelenumi{\alph{enumi})}
\setcounter{enumi}{4}
\tightlist
\item
  Ефективніше використовувати загальні перетворення незалежних і
  залежних змінних, включаючи похідні. Ед.
\end{enumerate}

4) Це стосується повного розділення змінних. При частковому розділенні
змінних одне диференціальне рівняння в похідних похідних зводиться до
кількох рівнянь у часткових похідних з меншою кількістю незалежних
змінних. --- Ед. Ед.

2) Це відбувається у випадку одномірного інтегрального перетворення. У
випадку k-вимірного інтегрального перетворення LCP з n незалежними
змінними зводиться до рівняння часткових з n---k незалежними змінними.
--- Ед. Ед.

Існують також інші чисельні методи, зокрема ті, що базуються на
наближенні розв'язків за поліноміальними поверхнями (наближення за
сплайнами).

\begin{enumerate}
\def\labelenumi{\arabic{enumi}.}
\setcounter{enumi}{5}
\tightlist
\item
  Метод теорії збурень. Початкова нелінійна задача зводиться до
  послідовності лінійних задач, апрокенн-
\end{enumerate}

нелінійна задача 1).

\begin{itemize}
\tightlist
\item
  \begin{enumerate}
  \def\labelenumi{\arabic{enumi}.}
  \setcounter{enumi}{6}
  \tightlist
  \item
    Метод функцій Гріна. Початкові та граничні умови замінюються
    системою найпростіших джерел, і задача розв'язується для кожного
    найпростішого джерела. Повне розв'язання початкової задачі
    отримується шляхом підсумовування розв'язків елементарних джерел.
  \end{enumerate}
\item
  \begin{enumerate}
  \def\labelenumi{\arabic{enumi}.}
  \setcounter{enumi}{7}
  \tightlist
  \item
    Метод інтегральних рівнянь. Рівняння з частковими похідними
    зводяться до інтегрального рівняння (рівняння, в якому невідома
    функція знаходиться під знаком інтеграла). Існує багато різних
    методів розв'язання інтегральних рівнянь.
  \end{enumerate}
\item
  \begin{enumerate}
  \def\labelenumi{\arabic{enumi}.}
  \setcounter{enumi}{8}
  \tightlist
  \item
    Варіаційні методи. Замість рівняння з частковими похідними
    диференціальними розв'язують задачу мінімізації. Виявляється, що
    функція, яка дає мінімум певному виразу (наприклад, загальній
    енергії системи), одночасно є розв'язком початкового рівняння
    похідниху похідних 2.
  \end{enumerate}
\item
  \begin{enumerate}
  \def\labelenumi{\arabic{enumi}.}
  \setcounter{enumi}{9}
  \tightlist
  \item
    Метод розкладу за власними функціями. Розв'язок рівняння часткових
    похідних шукається як послідовність власних функцій. Ці власні
    функції знаходяться як розв'язки так званої задачі власних значень,
    що відповідає початковій задачі для рівняння часткових похідних.
  \end{enumerate}
\end{itemize}

\section{Типи рівнянь у похідних
похідних}\label{ux442ux438ux43fux438-ux440ux456ux432ux43dux44fux43dux44c-ux443-ux43fux43eux445ux456ux434ux43dux438ux445-ux43fux43eux445ux456ux434ux43dux438ux445}

Рівняння в похідних похідних можна класифікувати за багатьма
характеристиками. Класифікація рівнянь важлива, оскільки, як виявилося,
кожен клас має свою загальну теорію та методи розв'язання рівнянь.

Тут ми наведемо шість основних методів класифікації

рівняння.

\begin{enumerate}
\def\labelenumi{\arabic{enumi}.}
\tightlist
\item
  Порядок рівняння. Порядок рівняння --- це найвищий порядок часткових
  похідних, включених до рівняння. Наприклад,
\end{enumerate}

Примітка. Ед.

1) Точніше, це метод лінеаризації нелінійних задач. Загальна теорія
збурень містить лінійні та нелінійні теорії збурень, кожна з яких, у
свою чергу, є регулярною та особливою. --- Ед. Ед.

Ці методи застосовні до спеціальних задач, що зумовлено існуванням
відповідних виразів --- функціоналів для цих задач.

Застосовувався раніше

\[u_t = u_{xx}\]

(рівняння другого порядку), \(u_t = u_x\) (рівняння першого порядку),
\(u_1 = uu_{xxx} + \sin x\) (рівняння трьох порядків).

\begin{enumerate}
\def\labelenumi{\arabic{enumi}.}
\setcounter{enumi}{1}
\tightlist
\item
  Кількість змінних. Кількість змінних --- це кількість незалежних
  змінних. Наприклад,
\end{enumerate}

\[u_t = u_{xx}\]

(рівняння з двома змінними \(x\) і \(t\) ),
\(u_t = u_{rr} + \frac{1}{r} u_r + \frac{1}{r^2} u_{\theta\theta}\)
(рівняння з трьома змінними \(r\) ,

\begin{enumerate}
\def\labelenumi{\arabic{enumi}.}
\setcounter{enumi}{2}
\tightlist
\item
  Лінійність. Рівняння з частковими похідними можуть бути лінійними та
  нелінійними. У лінійних рівняннях залежна змінна та всі її часткові
  похідні включаються лінійно, зокрема, вони не множаться одна на одну,
  не множаться в квадраті тощо). Точніше, лінійне рівняння другого
  порядку з двома незалежними змінними є рівнянням вигляду
\end{enumerate}

\[(1.1) Au_{xx} + Bu_{xy} + Cu_{yy} + Du_x + Eu_y + Fu = G\]

де A, B, C, D, E, F і G --- це константи або задані функції
пояснювальних змінних x і y.

\[u_{tt} = e^{-t}u_{xx} + \sin t\]

(лінійне рівняння), \(uu_{xx} + u_t = 0\) (нелінійне рівняння),
\(u_{xx} + yu_{yy} = 0\) (лінійне рівняння), \(xu_x + yu_y + u^2 = 0\)
(нелінійне рівняння).

\begin{itemize}
\tightlist
\item
  \begin{enumerate}
  \def\labelenumi{\arabic{enumi}.}
  \setcounter{enumi}{3}
  \tightlist
  \item
    Однорідність. Рівняння (1.1) називають однорідним, якщо права
    частина G(x, y) ідентична нулю для всіх x і y. Якщо G(x, y) не
    ідентичне нулю, то рівняння називають неоднорідним.
  \end{enumerate}
\item
  \begin{enumerate}
  \def\labelenumi{\arabic{enumi}.}
  \setcounter{enumi}{4}
  \tightlist
  \item
    Типи шансів. Якщо коефіцієнти A, B, C, D, E і F рівняння (1.1) є
    сталими, то це рівняння називається рівнянням зі сталими
    коефіцієнтами (інакше рівнянням зі змінними коефіцієнтами).
  \end{enumerate}
\item
  \begin{enumerate}
  \def\labelenumi{\arabic{enumi}.}
  \setcounter{enumi}{5}
  \tightlist
  \item
    Три основні типи лінійних рівнянь. Усі лінійні рівняння з частковими
    похідними другого порядку виду (1.1) належать до одного з трьох
    типів: а) параболічних, б) гіперболічних, в) еліптичних.
  \end{enumerate}
\end{itemize}

Параболічний тип. Рівняння параболічного типу описують процеси
теплопровідності та дифузії і визначаються умовою \(B^2 - 4AC = 0\) .

Гіперболічний тип. Рівняння гіперболічного типу описів

осцилюючі системи та хвильові рухи визначаються умовою \(B^2-4AC>0\) .

Еліпічний тип. Рівняння еліптичного типу описують усталені процеси і
визначаються умовним \(B^2-4AC<0\) . Приклади

\begin{enumerate}
\def\labelenumi{\alph{enumi})}
\tightlist
\item
\end{enumerate}

\[u_i = u_{xx}\], \(B^2 - 4AC = 0\) (параболічна), (гіперболічна),
(гіперболічна), (гіперболічна \(y > 0\) ), (гіперболічна), (еліптична),
(еліптична), параболічна на \(y = 0\) , гіперболічна на \(y < 0\) .

(У випадку змінних коефіцієнтів тип рівняння може змінюватися від точки
до точки.)

\section{ЗАУВАЖЕННЯ}\label{ux437ux430ux443ux432ux430ux436ux435ux43dux43dux44f}

\begin{itemize}
\tightlist
\item
  \begin{enumerate}
  \def\labelenumi{\arabic{enumi}.}
  \tightlist
  \item
    Загалом, значення \(B^2-4AC\) є функцією незалежних змінних.
    Відповідно, базовий тип рівняння може змінюватися в межах визначення
    рівняння, хоча це не є обов'язковим.
  \end{enumerate}
\item
  \begin{enumerate}
  \def\labelenumi{\arabic{enumi}.}
  \setcounter{enumi}{1}
  \tightlist
  \item
    У рівнянні (1.1) незалежними змінними є x і y. У багатьох задачах
    однією з двох незалежних змінних є час, і рівняння (1.1) можна
    записати через x і t.
  \end{enumerate}
\end{itemize}

\pandocbounded{\includegraphics[keepaspectratio]{images/_page_10_Figure_9.jpeg}}

PHC 1.1. Діаграма класифікації диференціальних рівнянь у похідних
похідних.

\begin{enumerate}
\def\labelenumi{\arabic{enumi}.}
\setcounter{enumi}{2}
\tightlist
\item
  Рис. 1.1. Представлено діаграму класифікації для диференціальних
  рівнянь у похідних похідних.
\end{enumerate}

\section{ЗАВДАННЯ}\label{ux437ux430ux432ux434ux430ux43dux43dux44f}

\begin{itemize}
\tightlist
\item
  \begin{enumerate}
  \def\labelenumi{\arabic{enumi}.}
  \tightlist
  \item
    Класифікуємо наступні рівняння відповідно до всіх характеристик,
    зазначених у днаграмі на рис. 1.1:
  \end{enumerate}
\item
  \begin{enumerate}
  \def\labelenumi{\alph{enumi})}
  \tightlist
  \item
    \(u_t = u_{xx} + 2u_x + u\) ,
  \end{enumerate}
\end{itemize}

\begin{enumerate}
\def\labelenumi{\arabic{enumi})}
\setcounter{enumi}{5}
\tightlist
\item
  \(u_t = u_{xx} + e^{-t}\) , B) \(u_{xx} + 3u_{xy} + u_{yy} = \sin x\)
  , \(u_{tt} = uu_{xxx} + e^{-t}\) .
\end{enumerate}

\begin{itemize}
\tightlist
\item
  \begin{enumerate}
  \def\labelenumi{\arabic{enumi}.}
  \setcounter{enumi}{1}
  \tightlist
  \item
    Скільки розв'язків \(u_1 = u_{xx}\) рівняння існує? Спробуйте знайти
    рішення, як \(u(x, t) = e^{a\lambda + bt}\) .
  \end{enumerate}
\item
  \begin{enumerate}
  \def\labelenumi{\arabic{enumi}.}
  \setcounter{enumi}{2}
  \tightlist
  \item
    Якщо функції \(u_1(x, t)\) і \(u_2(x, t)\) задовольняють рівняння
    (1.1), чи задовольняє сума цих функцій його? Доведи це.
  \end{enumerate}
\item
  \begin{enumerate}
  \def\labelenumi{\arabic{enumi}.}
  \setcounter{enumi}{3}
  \tightlist
  \item
    Ймовірно, найпростіше диференціальне рівняння в похідних ---
  \end{enumerate}
\end{itemize}

\[\frac{\partial u(x, y)}{\partial x} = 0.\]

Чи можете ви розв'язати це рівняння? (Знайдіть усі функції, що
задовольняють це рівняння.)

\begin{enumerate}
\def\labelenumi{\arabic{enumi}.}
\setcounter{enumi}{4}
\tightlist
\item
  Що можна сказати про рівняння
\end{enumerate}

\[\frac{\partial^2 u(x, y)}{\partial x \partial y} = 0\]

?

Чи можете ви знайти всі розв'язки цього рівняння? Скільки їх? Порівняйте
з кількістю розв'язків звичайного диференціального рівняння

\[\frac{d^2y}{dx^2} = 0.\]

\section{ЗАДАЧІ ТИПУ ДИФУЗІЇ (ПАРАБОЛІЧНІ
РІВНЯННЯ)}\label{ux437ux430ux434ux430ux447ux456-ux442ux438ux43fux443-ux434ux438ux444ux443ux437ux456ux457-ux43fux430ux440ux430ux431ux43eux43bux456ux447ux43dux456-ux440ux456ux432ux43dux44fux43dux43dux44f}

МЕТА ЛЕКЦІЇ: Показати, як параболічні рівняння використовуються для
розв'язання задач теплопровідності та дифузії. Наведено кілька прикладів
параболічних рівнянь, а також пояснюється фізичне значення різних членів
\((u_1, u_2, u_3, u_4, u_6)\) , включених у ці рівняння.

На конкретному прикладі розглядається основна ідея зміщеної задачі для
параболічного рівняння. Одна з головних цілей цієї лекції --- надати
читачеві інтуїтивне уявлення

Розуміння параболічних проблем.

Почнемо з розгляньмо просту фізичну задачу і те, як її можна описати
математичною моделлю з рівнянням у похідних похідних.

Потім ми ускладнимо задачу і покажемо, як нові фізичні умови призводять
до появи нових рівнянь у похідних похідних. У цій лекції диференціальні
рівняння не виводяться і не розв'язуються. Про це ми поговоримо пізніше.

\section{Простий експеримент з
теплопровідністю}\label{ux43fux440ux43eux441ux442ux438ux439-ux435ux43aux441ux43fux435ux440ux438ux43cux435ux43dux442-ux437-ux442ux435ux43fux43bux43eux43fux440ux43eux432ux456ux434ux43dux456ux441ux442ux44e}

Припустимо, що ми провели простий експеримент, який розбили на такі
кроки:

КРОК 1. Візьміть достатньо довгий (скажімо, \(L=2\,\mathrm{m}\) )
стрижень (наприклад, мідний) діаметром \(2\,\mathrm{cm}\) , бічна
поверхня (але не кінці) якої термічно ізольована. Замість стрижня можна
використовувати мідну трубку, покриту теплоізоляцією зовні та зсередини.

Іншими словами, тепло може проходити лише через кінці стрижня і не може
проходити через його бічну поверхню.

КРОК 2. Тепер помістимо цей стрижень у пристрій із фіксованою
температурою \(T_{\mathfrak{o}}\) (градуси Цельсія) на достатньо довгий
час, щоб температура всередині стрижня стала

Так само, як і температура пристрою. Для простоти припустимо, що
температура пристрою \(T_0 = 10^{\circ}\) C.

КРОК 3. Давайте витягнемо стрижень з нашого пристрою у певний момент
часу, що зручно взяти як нуль t = 0, і приєднаємо до нього дві термопари
з кінців. Завдання цих елементів --- підтримувати фіксовані температури
\(T_1\) і \(T_2\) на кінцях (наприклад, \(T_1=0^{\circ}\) C і
\(T_2=50^{\circ}\) C). Іншими словами, два термостати постійно
контролюють температуру на кінцях стрижня, і якщо температура
відхиляється від встановлених значень \(T_1\) і \(T_2\) , активуються
потужні нагрівальні або охолоджувальні елементи (пристрої) для
відповідного регулювання температури. Наш експеримент ілюструє RPS, 2 1,

\pandocbounded{\includegraphics[keepaspectratio]{images/_page_13_Figure_4.jpeg}}

Рис. 2.1. Схема експерименту: a --- це термопара, яка підтримує
встановлену температуру \(T_1\) на лівому кінці стрижня; \(\delta\) ---
це термопара, яка підтримує встановлену температуру \(T_2\) на правому
кінці стрижня.

КРОК 4. Тепер ми будемо слідкувати за температурним профілем у стрижні
за допомогою пристрою візуалізації. Чому ми хочемо провести такий
експеримент? Це інше питання, про яке ми поговоримо пізніше.

Це кінець опису нашого експерименту.

Головна мета цієї лекції --- показати, як цю фізичну проблему та її
модифікації можна змоделювати мовою параболічних рівнянь.

\section{Математична модель
теплопровідності}\label{ux43cux430ux442ux435ux43cux430ux442ux438ux447ux43dux430-ux43cux43eux434ux435ux43bux44c-ux442ux435ux43fux43bux43eux43fux440ux43eux432ux456ux434ux43dux43eux441ux442ux456}

Опис нашої фізичної проблеми вимагає трьох типів відносин.

\begin{itemize}
\tightlist
\item
  \begin{enumerate}
  \def\labelenumi{\arabic{enumi}.}
  \tightlist
  \item
    Рівняння з частковими похідними, які описують фізичне явище
    теплопровідності.
  \end{enumerate}
\item
  \begin{enumerate}
  \def\labelenumi{\arabic{enumi}.}
  \setcounter{enumi}{1}
  \tightlist
  \item
    Граничні умови, що описують процеси теплопередачі на межах.
  \end{enumerate}
\item
  \begin{enumerate}
  \def\labelenumi{\arabic{enumi}.}
  \setcounter{enumi}{2}
  \tightlist
  \item
    Початкові твердження, які описують стан системи на початку процесу.
  \end{enumerate}
\end{itemize}

\section{Рівняння
теплопровідності}\label{ux440ux456ux432ux43dux44fux43dux43dux44f-ux442ux435ux43fux43bux43eux43fux440ux43eux432ux456ux434ux43dux43eux441ux442ux456}

Базове одновимірне рівняння теплопровідності записується у вилці

(2.1)

\[u_t = \alpha^2 u_{xx}, \quad 0 < x < L, \quad 0 < t < \infty.\]

Це рівняння пов'язує значення \(u_t\) --- швидкість зміни температури з
часом (виміряна в градусах/с)

та величину \(u_{xx}\) --- це увігнутий профілю температури u(x, t)
(який є мірою різниці між температурою в заданій точці та температурою в
сусідніх точках).

\pandocbounded{\includegraphics[keepaspectratio]{images/_page_14_Figure_14.jpeg}}

Рис. 2.2. Стрілки показують зміну температури відповідно до рівняння
\(u_t = \alpha^2 u_{xx}\) .

Це рівняння буде виведено з базового рівняння збереження тепла в
наступних лекціях, але наразі ми розглянемо його заданим. Це рівняння
каже, що температура \(u\left(x,\ t\right)\) (у певній точці t у певній
точці x стрижня) зростає \((u_t>0)\) або зменшується \((u_t<0)\) залежно
від того, чи є друга похідна \(u_{xx}\) позитивною чи негативною. Фіг.
Рисунок 2.2 ілюструє зміну температури в різних точках на елементі.

Давайте подивимось, як значення \(u_{xx}\) можна інтерпретувати мовою
теплопровідності. Припустимо, ми наблизимо значення \(u_{xx}\) фінальних
постів

\[u_{xx}(x, t) \cong \frac{1}{\Delta x^2} [u(x + \Delta x, t) - 2u(x, t) + u(x - \Delta x, t)].\]

Це співвідношення можна переписати у вигляді

\[u_{\lambda x}(x, t) \simeq \frac{2}{\Delta x^2} \left[ \frac{u(x + \Delta x, t) + u(x - \Delta x, t)}{2} - u(x, t) \right].\]

Тепер ми можемо дати наступну інтерпретацію значення \(u_{xx}\) .

\begin{enumerate}
\def\labelenumi{\arabic{enumi}.}
\item
  Якщо температура u(x, t) менша за середнє значення температури у двох
  сусідніх точках, то \(u_{xx} > 0\) (тут загальний тепловий потік по
  осі x додатний).
\item
  Якщо температура u(x, t) дорівнює середньому значенню температур у
  двох сусідніх точках, то \(u_{xx} = 0\) (тут загальний тепловий потік
  по осі x дорівнює нулю).
\item
  Якщо температура \(u\left(x,\ t\right)\) вища за середню температуру у
  двох сусідніх точках, то \(u_{xx}<0\) (тут --- повний потік
\end{enumerate}

тепла вздовж осі x є від'ємною).

Усе це ілюструється на рис. 2.2. Іншими словами, якщо температура в
точці x більша за середню температуру у двох сусідніх точках
\(x - \Delta x\) і \(x + \Delta x\) , то температура в точці x
зменшиться. Відповідно, точна швидкість зниження температури, а саме
величина \(u_t\) , пропорційна цій різниці. Коефіцієнт пропорційності
\(\alpha^2\) визначається властивостями матеріалу. Ми розглянемо цю
константу детальніше у майбутніх лекціях.

\section{Граничні умови
(ГІ)}\label{ux433ux440ux430ux43dux438ux447ux43dux456-ux443ux43cux43eux432ux438-ux433ux456}

Усі фізичні проблеми характеризуються наявністю певних меж. Отже, щоб
рухатися вперед, ми повинні включити межі в математичну модель, щоб
адекватно описати фізичну проблему. У нашому експерименті граничні умови
(ГІ) дуже прості. Оскільки температура на кінцях x=0 і x=L завжди
залишається сталою і дорівнює \(T_1\) та \(T_2\) , можливо записати

(2.2)

\[\left\{ \begin{array}{l} u\left(0,\ t\right) = T_{t}, \\ u\left(L,\ t\right) = T_{2}, \end{array} \right. \quad 0 < t < \infty,\]

\section{Початкові умови
(NU)}\label{ux43fux43eux447ux430ux442ux43aux43eux432ux456-ux443ux43cux43eux432ux438-nu}

Усі фізичні процеси мають починатися у певний момент часу (зазвичай
приймається як нуль t=0), тому ми повинні встановити фізичні умови саме
в цю точку. Відтоді, як ми почали контролювати температуру з моменту,
коли паличка мала постійну температуру \(T_{\mathrm{ }o}\) , ми

(2.3)

\[u(x, 0) = T_0, \quad 0 \le x \le L\]

Тепер ми створили математичну модель експерименту. Система співвідношень
(2.1), (2.2) та (2.3) називається задачею зміщення для рівняння
теплопровідності і зазвичай записується як

(2.4)

\[\begin{aligned} & (\text{YUII}) \quad u_t = \alpha^2 u_{xx}, & 0 < x < L, \ 0 < t < \infty, \\ & \left( \Gamma \text{Y} \right) \quad \begin{cases} u \ (0, \ t) = T_1, \\ u \ (L, \ t) = T_2, \end{cases} & 0 < t < \infty, \\ & (\text{HY}) \quad u \ (x, \ 0) = T_0, \quad 0 \leqslant x \leqslant L. \end{aligned}\]

Цікаво, що це зовсім не очевидно: чи існує одна функція u(x, t), яка
задовольняє умови (2.4), і чи описує ця функція температуру в стрижні?
Отже, наше безпосереднє завдання --- знайти цю функцію u(x, t) з (2.4).

Перш ніж завершити цю лекцію, давайте розглянемо деякі варіації основної
проблеми. Почнемо з невеликих змін рівняння теплопровідності.

\section{Деякі рівняння дифузійного
типу}\label{ux434ux435ux44fux43aux456-ux440ux456ux432ux43dux44fux43dux43dux44f-ux434ux438ux444ux443ux437ux456ux439ux43dux43eux433ux43e-ux442ux438ux43fux443}

Теплопередача через бічну поверхню пропорційна різниці температур

Рівняння

\[u_t = \alpha^2 u_{xx} - \beta (u - u_0)\]

Описує теплопровідність елемента, враховуючи не лише дифузію теплових
\(\alpha^2 u_{xx}\) вздовж елемента, а й теплопередачу через бічну
поверхню елемента. Відтік тепла ( \(\beta > 0\) ) або його приплив (
\(\beta < 0\) ) пропорційний різниці між температурою u стрижня та
температурою навколишнього середовища \(u_0\) ( \(\beta\) --- це сталий
коефіцієнт пропорційності). Якщо коефіцієнт \(\beta\) великий порівняно
з \(\alpha^2\) , то теплопотік уздовж стрижня буде малим порівняно з

Протікання через бічну поверхню, і тому тепло буде проходити через неї
відповідно до закону \(u_t = -\beta (u - u_a)\) .

У хімії, де величина відіграє роль концентрації, рівняння

\[u_1 = \alpha^2 u_{xx} - \beta (u - u_0)\]

означає, що швидкість зміни \(u_i\) кількості речовини залежить як від
дифузії \(\alpha^2 u_{xx}\) (у напрямку осі x), так і від появи (
\(\beta < 0\) ) або розпаду ( \(\beta > 0\) ) речовини в хімічній
реакції, і пропорційна різниці між двома концентраціями u і
\(u_{\theta}\) .

\section{Внутрішнє джерело
тепла}\label{ux432ux43dux443ux442ux440ux456ux448ux43dux454-ux434ux436ux435ux440ux435ux43bux43e-ux442ux435ux43fux43bux430}

Неоднородне рівняння

\[u_t = \alpha^2 u_{xx} + f(x, t)\]

відповідає випадку, коли всередині елемента є джерело тепла (розташоване
вздовж усього елемента і діє в будь-який момент t). Наприклад, може бути
так, що дріт з електричним струмом проходить через стрижень і опір є
постійним джерелом тепла f(x,t) = K.

\section{Рівняння конвективної
дифузії}\label{ux440ux456ux432ux43dux44fux43dux43dux44f-ux43aux43eux43dux432ux435ux43aux442ux438ux432ux43dux43eux457-ux434ux438ux444ux443ux437ux456ux457}

Припустимо, що домішка поширюється вздовж потоку, що рухається зі
швидкістю v. Очевидно, що концентрація \(u\left(x,\,t\right)\) домішок
змінюється залежно від x (позитивний напрямок осі x обирається вздовж
потоку) і часу t. Швидкість зміни концентрації \(u_t\) визначається
рівнянням конвективної дифузії

\[u_t = \alpha^2 u_{xx} - v u_x.\]

Термін \(\alpha^2 u_{xx}\) описує внесок дифузії, а \(vu_x\) ---
конвективний компонент. Що важливіше, дифузія або конвекція залежить від
співвідношення \(\alpha^2\) до v. Димові частинки піднімаються
конвективно через рух гарячого повітря і одночасно розсіюють через
вихорові рухи повітря. Окрім цих модифікацій рівняння теплопровідності,
можливо змінювати граничні умови на кінцях стрижня відповідно до
фізичної ситуації. Деякі з цих змін ми обговоримо у третій лекції.

\section{ЗАУВАЖЕННЯ}\label{ux437ux430ux443ux432ux430ux436ux435ux43dux43dux44f-1}

Рівняння теплопровідності \(u_t = \alpha^2(x) u_{xx}\) з змінним
коефіцієнтом \(\alpha(x)\) відповідає задачі, в якій дифузія залежить
від x (тобто матеріал є гетерогенним). Наприклад, мідь і сталь---

\pandocbounded{\includegraphics[keepaspectratio]{images/_page_18_Picture_4.jpeg}}

PHC. 2.3.

Пластини контактують, як показано на рис. 2.3. Припустимо, що
температура \(u(0,t)=0^{\circ}\) C підтримується на лівому боці міді, а
температура \(u(L,t)=20^{\circ}\) C --- з правого боку сталі, тоді
рівняння, що описує теплопровідність у цій системі, виглядає так:

\[u_t = \alpha^2(x) \, u_{_{\! A \, X}}, \quad 0 < x < L,\]

де

\[
\alpha(x) = \begin{cases}
\alpha_1, & \text{коефіцієнт дифузії в міді}, \quad 0 < x < L/2, \\
\alpha_2, & \text{коефіцієнт дифузії в сталі}, \quad L/2 < x < L.
\end{cases}
\]

\section{ЗАВДАННЯ}\label{ux437ux430ux432ux434ux430ux43dux43dux44f-1}

\begin{enumerate}
\def\labelenumi{\arabic{enumi}.}
\tightlist
\item
  Якщо початкова температура в стрижні
\end{enumerate}

\[u(x, 0) = \sin \pi x, \quad 0 \leqslant x \leqslant 1,\]

а граничні умови такі:

\[
u(0, t) = 0, \qquad u(1, t) = 0.
\]

Тоді як поводитиметься температура в стрижні \(U(x, t)\) при \(t > 0\) ?
2. Припустимо, що в стрижні є постійне внутрішнє джерело тепла, так що
рівняння теплопровідності в стрижні записується як

\[u_t = \alpha^2 u_{xx} + 1, \quad 0 < x < 1.\]

Припустимо, що ми зафіксували температуру на межі з значеннями u(0, t) =
0 і u(1, t) = 1. Якою буде стаціонарна температура у стрижні? Іншими
словами, який c-незалежний розподіл U(x, t) перешкоджує u(x, t)?

\begin{enumerate}
\def\labelenumi{\arabic{enumi}.}
\setcounter{enumi}{2}
\tightlist
\item
  Припустимо, що металевий стрижень втрачає тепло через бічну поверхню
  згідно з рівнянням
\end{enumerate}

\[u_1 = \alpha^2 u_{xx} - \beta u\], \(0 < x < 1\) .

І припустимо, що ми підтримуємо кінці стрижня при температурах U(0, T) =
1, U(1, T) = 1. Знайдіть стаціонарний розподіл температури по стрижню.
Побудуйте графік цього розподілу. У яких моментах відбувається тепловий
потік?

\begin{enumerate}
\def\labelenumi{\arabic{enumi}.}
\setcounter{enumi}{3}
\tightlist
\item
  Припустимо, що термоізольований металевий стрижень довжини L = 1 має
  початкову температуру \(\sin(3\pi x)\) , а фіксовану температуру --- 0
  і \(10\,^{\circ}\) C на лівому та правому кінцях.
\end{enumerate}

\section{Лекція 3}\label{ux43bux435ux43aux446ux456ux44f-3}

\section{ГРАНИЧНІ УМОВИ В ЗАДАЧАХ ДИФУЗІЙНОГО
ТИПУ}\label{ux433ux440ux430ux43dux438ux447ux43dux456-ux443ux43cux43eux432ux438-ux432-ux437ux430ux434ux430ux447ux430ux445-ux434ux438ux444ux443ux437ux456ux439ux43dux43eux433ux43e-ux442ux438ux43fux443}

МЕТА ЛЕКЦІЇ: Показати, як виникають різні граничні умови у задачах
теплопровідності та дифузії, а також ввести важливе поняття теплового
потоку.

Розглядаються три найважливіші типи граничних умов:

\begin{itemize}
\tightlist
\item
  \begin{enumerate}
  \def\labelenumi{\arabic{enumi}.}
  \tightlist
  \item
    u = g(t) (температура задається на межі),
  \end{enumerate}
\item
  \(2. \ \frac{\partial u}{\partial n} + \lambda u = g\left(t\right)\)
  (встановлення температури навколишнього середовища);
\end{itemize}

n --- вектор зовнішньої нормалі до межі),
\(3. \frac{\partial u}{\partial n} = g(t)\) (визначений тепловий потік
через межу).

У задачах теплопровідності зазвичай використовують три основні типи
граничних умов. Ця лекція наводить приклади того, як різні фізичні умови
експерименту призводять до різних типів граничних умов.

\section{Граничні умови першого виду (температура задається на
межі)}\label{ux433ux440ux430ux43dux438ux447ux43dux456-ux443ux43cux43eux432ux438-ux43fux435ux440ux448ux43eux433ux43e-ux432ux438ux434ux443-ux442ux435ux43cux43fux435ux440ux430ux442ux443ux440ux430-ux437ux430ux434ux430ux454ux442ux44cux441ux44f-ux43dux430-ux43cux435ux436ux456}

Розглянемо тепловий потік в одномірному стрижні, показаний на рис. 3 1,
і припустимо, що температура кінців змінюється відповідно до законів
\(g_1(t)\) і \(g_2(t)\) .

\pandocbounded{\includegraphics[keepaspectratio]{images/_page_20_Picture_4.jpeg}}

Фіг. 31. Температура буде відновлена на межі.

Як згадувалося в попередній лекції, кожен пристрій, що підтримує
встановлену температуру на кінцях, складається з термостата та
нагрівального елемента для підтримки відповідних теплових умов. Проблеми
з граничними умовами першого виду дуже поширені. У деяких випадках
складність полягає у пошуку контролерів граничних умов, тобто температур
\(g_1(t)\) та \(g_2(t)\) меж, які призводять до певної зміни температури
всередині стрижня. Наприклад, у металургії часто необхідно обирати
граничні умови контролю, щоб температура металу всередині печі
змінювалася з часом, а температурний градієнт був невеликим.

\pandocbounded{\includegraphics[keepaspectratio]{images/_page_20_Picture_7.jpeg}}

Фіг. 3.2. Зниження температури на межі.

Граничні умови першого виду також виникають у багатовимірних задачах
теплопровідності. Як приклад можна навести цікаву задачу визначення
температурного поля всередині кругового диска з радіусом \emph{R} (див.
рис. 3.2), коли гранична температура задається у полярних координатах
формулою \(u(R, \theta, t) = \cos t \sin \theta\) .

Звісно, щоб розв'язати задачу, потрібно знати початкову температуру, але
в цьому випадку початкові умови перестають впливати на розв'язок через
дуже короткий час, і отримана температура всередині диска визначається
лише значеннями температури на межі.

\section{Граничні умови другого виду (задана температура навколишнього
середовища)}\label{ux433ux440ux430ux43dux438ux447ux43dux456-ux443ux43cux43eux432ux438-ux434ux440ux443ux433ux43eux433ux43e-ux432ux438ux434ux443-ux437ux430ux434ux430ux43dux430-ux442ux435ux43cux43fux435ux440ux430ux442ux443ux440ux430-ux43dux430ux432ux43aux43eux43bux438ux448ux43dux44cux43eux433ux43e-ux441ux435ux440ux435ux434ux43eux432ux438ux449ux430}

Припустимо, що ми знову розглядаємо термічно ізольований мідний
стержень, але тепер відмовляємося від заданого температурного режиму на
кінцях стрижня. Натомість давайте наблизимо кінці стрижня до контакту з
двома середовищами. Нехай температура одного з них \(g_1(t)\) , а іншого
\(g_2(t)\) . Іншими словами, припустимо, що лівий кінець стрижня
поміщений у контейнер із рідиною, температура якої змінюється відповідно
до закону \(g_1(t)\) . А правий кінець поміщається в іншу рідину з
температурою \(g_2(t)\) (рис. 3, 3)

\pandocbounded{\includegraphics[keepaspectratio]{images/_page_21_Picture_6.jpeg}}

Рис. 3.3. Конвективна теплопередача через межі, a --- це рідина при
температурі \(g_1(t)\) , \(\delta\) --- рідина при температурі
\(g_2(t)\) .

Задаючи граничні умови такого типу, ми не можемо вважати, що граничні
температури стрижня збігаються з температурами рідини. Але ми знаємо
закон Ньютона: якщо температура одного з кінців стрижня менша за
температуру відповідної рідини, тепло потече в стрижень зі швидкістю,
пропорційною різницям температур. Іншими словами, для одномірного
стрижня з межами x=0 і x=1 закон теплопередачі Ньютона формулюється так:

(3.1)

\[\{\]

вихідний тепловий потік на \(x=0\) дорівнює
\(h\left[u\left(0,\,t\right)-g_{1}\left(t\right)\right]\) , вихідний
тепловий потік при \(x=L\) дорівнює
\(h\left[u\left(L,\,t\right)-g_{2}\left(t\right)\right]\) ,

де h --- коефіцієнт теплопередачі, який показує, скільки калорій
проходить через межу за одну секунду при різниці температур у один
градус. Вихідний потік дорівнює кількості кало-

Пройшов крізь палочний спис за одну секунду. Зверніть увагу, що вихідний
тепловий потік буде позитивним на кінці стрижня, де температура стрижня
вища за температуру навколишнього середовища. Рівняння (3.1) разом із
відомим законом теплопровідності Фур'є тепер можна використовувати для
отримання Р и с

\pandocbounded{\includegraphics[keepaspectratio]{images/_page_22_Figure_3.jpeg}}

. 3.4. Ілюстрація закону Фур'є: \(\frac{\partial u}{\partial n} > 0\)
означає, що тепло надходить у цю область;
\(\frac{\partial u}{\partial n} < 0\) означає, що тепло виходить із цієї
зони; n --- напрямок зовнішньої нормалі;
\(\frac{\partial u}{\partial n} -\) Зміна температури в напрямку N;
\(\frac{\partial u}{\partial n} = -\frac{\partial u}{\partial (-n)}\) є
похідною у напрямку внутрішньої нормальної --- це похідна, взята
\(\mathbf{e}\) протилежному знаку протилежного знака пормалі.

граничні умови. Закон Фур'є дає друге представлення (перше міститься у
(3.1)) для вихідного теплового потоку. Прирівнюючи ці два вирази до
потоку, ми отримуємо граничні умови, які шукаємо. Сформулюємо закон
Фур'є (встановлений експериментально):

(3.2) Тепловий потік, що проходить через межу площі, пропорційний
нормальній похідній температури у напрямку внутрішньої норми.

Цей закон стверджує, що якщо температура швидко підвищується у напрямку
зовнішньої температури до межі D (рисунок 3.4), то потік буде текти з
навколишнього середовища до області D.

У нашій одновимірній задачі закон Фур'є має вигляд:

(3.3)

\[\begin{cases} \text{вытекающий поток тепла при } x=0 \text{ равен } k \frac{\partial u}{\partial x}, \\ \text{вытекающий поток тепла при } x=' \text{ равен } -k \frac{\partial u}{\partial x}, \end{cases}\]

де k --- теплопровідність матеріалу, яка є мірою того, наскільки добре
матеріал проводить тепло. Погано провідні матеріали мають коефіцієнт
близький до нуля в одиницях GHS, тоді як мідь та алюміній мають значення
значно вище нуля.

\pandocbounded{\includegraphics[keepaspectratio]{images/_page_23_Figure_5.jpeg}}

Фіг. 3.5. Ще одна ілюстрація закону Чотирьох.

Закон Фур'є (3.3) фактично описує теплопровідність всюди всередині
стрижня, а не лише на його межі; Наприклад, (див. рис. 3.5).

(3.4) Потік, що проходить через точку

\[x\]

зліва направо \(=-k\frac{\partial n}{\partial x}\) Закон Фур'є (3.4)
стверджує: якщо \(u_x < 0\) , то в точці \(x_0\) тепло тече зліва
направо, якщо \(u_x > 0\) , то в точці \(x_0\) тепло тече справа наліво.
(Тепло завжди переходить від вищих температур до нижчих.)

Нарешті, якщо використовувати вирази (3.1) і (3.3) для теплового потоку,
то отримуємо бажані граничні умови для задачі, описаної на рис. 3.3, у
чисто математичній формі

\[(\Gamma Y) \quad \begin{cases} \frac{\partial u\left(0,\,t\right)}{\partial x} = \frac{h}{k}\left[u\left(0,\,t\right) - g_{1}\left(t\right)\right], \\ \frac{\partial u\left(L,\,t\right)}{\partial x} = -\frac{h}{k}\left[u\left(L,\,t\right) - g_{2}\left(t\right)\right]. \end{cases} \quad 0 < t < \infty,\]

Дуже часто константа h/k позначається як \(\lambda\) , а граничні умови
теплового потоку на межі записуються у вигляді

(3.5)

\[u_{x}(0, t) = \lambda [u(0, t) - g_{1}(t)], u_{x}(L, t) = -\lambda [u(L, t) - g_{2}(t)].\]

Для випадку вищих розмірностей отримують подібні граничні умови.
Наприклад, якщо межа кругового диска осяяна рухомою рідиною з
температурою \(g(\theta,t)\) , то граничні умови записуються як

\[\frac{\partial u}{\partial r}(R, \theta, t) = -\frac{h}{k} [u(R, \theta, t) - g(\theta, t)].\]

Тут \(\frac{\partial u}{\partial r}(R, \theta, t)\) позначає зовнішню
нормальну похідну (у позитивному напрямку r-осі) значення u, обчислену в
точці \((R, \theta)\) межі. Ми назвемо такі граничні умови лінійними
(оскільки вони лінійні в U і \(u_r\) ), але неоднорідними, оскільки
функція \(g(\theta, t)\) знаходиться праворуч.

\section{Граничні умови третього виду (враховуючи потік; зокрема,
враховується випадок термічно ізольованих
меж)}\label{ux433ux440ux430ux43dux438ux447ux43dux456-ux443ux43cux43eux432ux438-ux442ux440ux435ux442ux44cux43eux433ux43e-ux432ux438ux434ux443-ux432ux440ux430ux445ux43eux432ux443ux44eux447ux438-ux43fux43eux442ux456ux43a-ux437ux43eux43aux440ux435ux43cux430-ux432ux440ux430ux445ux43eux432ux443ux454ux442ux44cux441ux44f-ux432ux438ux43fux430ux434ux43eux43a-ux442ux435ux440ux43cux456ux447ux43dux43e-ux456ux437ux43eux43bux44cux43eux432ux430ux43dux438ux445-ux43cux435ux436}

Жоден потік не проходить через термоізольовані межі, тому нормальна
похідна (зовнішня чи внутрішня) має зрівнятися з нулем на межі, оскільки
потік пропорційний нормальній похідній. У випадку одномірного елемента з
ізольованими кінцями x=0 і x=L, граничні умови мають вигляд

\[u_{\lambda}(0, t) = 0,\] \(u_{\lambda}(L, t) = 0.\)
\(0 < t < \infty,\) У двовимірній області ізоляція межі означає, що
нормальна похідна температури на межі перетворюється на нуль.
Припустимо, наприклад, що круглий диск термічно ізольований уздовж межі.
Тоді граничну умову слід записати так: \(u_r(R, \theta, t) = 0\) для
всіх \(0 \le \theta < 2\pi\) і всіх \(0 < t < \infty\) .

З іншого боку, якщо вказати кількість тепла, що проходить через межу
нашого диска, то граничні умови набувають вигляду

\[u_r(R, \theta, t) = f(\theta, t),\]

де f(0, t) --- це кількість тепла, що переходить у круглий диск із
зовнішнього джерела.

Тепер проілюструємо різні типи граничних умов.

\section{Типові граничні умови для задачі одномірної
теплопровідності}\label{ux442ux438ux43fux43eux432ux456-ux433ux440ux430ux43dux438ux447ux43dux456-ux443ux43cux43eux432ux438-ux434ux43bux44f-ux437ux430ux434ux430ux447ux456-ux43eux434ux43dux43eux43cux456ux440ux43dux43eux457-ux442ux435ux43fux43bux43eux43fux440ux43eux432ux456ux434ux43dux43eux441ux442ux456}

Припустимо, що у нас є мідний стрижень довжиною 200 см, бікова поверхня
якого термічно ізольована, а початкова температура становить 0 °C.
Припустимо також, що верхній кінець стрижня x=0 термічно ізольований, а
нижній кінець x=200 змивається рухомою водою, яка має постійну
температуру \(g_2(t)=20\) °C (див. рис. 3.6).

\pandocbounded{\includegraphics[keepaspectratio]{images/_page_25_Picture_6.jpeg}}

Фіг. 3.6. Змішане завдання.

Математична модель цієї задачі описується наступними чотирма
відношеннями:

\[(3.6)
\begin{aligned}
(\text{УЧП})\quad & u_t = \alpha^2 u_{xx}, && 0 < x < 200,\ 0 < t < \infty, \\
(\text{ГУ})\quad & u_x(0,t) = 0, && 0 < t < \infty, \\
& u_x(200,t) = -\frac{h}{k}\left[u(200,t) - 20\right], && 0 < t < \infty, \\
(\text{НУ})\quad & u(x,0) = 0^\circ\mathrm{C}, && 0 \leqslant x \leqslant 200.
\end{aligned}
\]

де \(\alpha^2 = 1,16\) см 2/с --- теплова дифузія міді;

k = 0,93 кал/см·с·°C --- теплопровідність міді;

h --- коефіцієнт теплопередачі. Визначення цього співвідношення ---
велика проблема. Він вимірюється швидкістю виходу тепла від другого
кінця стрижня до води. Це залежить від швидкості потоку води, стану
поверхні, яку слід змивати тощо.

Читач може спробувати провести експеримент, щоб визначити

ділення значення цього коефіцієнта.

\section{ЗАУВАЖЕННЯ}\label{ux437ux430ux443ux432ux430ux436ux435ux43dux43dux44f-2}

\begin{enumerate}
\def\labelenumi{\arabic{enumi}.}
\tightlist
\item
  Типова задача внутрішньої теплопровідності для квадрата показана на
  рис. 3.7.
\end{enumerate}

У цій задачі початкова температура \(u\left(x,\,y,\,t\right)\) при t = 0
задається всередині квадрата, і рівняння з похідними з частковими має
вигляд Р и с

\pandocbounded{\includegraphics[keepaspectratio]{images/_page_26_Figure_9.jpeg}}

. 3.7. Типові граничні умови для задачі дифузії в квадраті: a ---
кількість тепла, \(f_1(t)\) проходить через межу; b --- термоізольована
межа; a --- це навколишня температура Равії \(10^{\circ}\) C; z ---
температура підтримується рівною \(10^{\circ}\) C.

а граничні умови будуть виконані на \(0 < t < \infty\) і визначатимуть
температуру в будь-який момент. Якою б не була ця температура, вона
повинна відповідати граничним умовам, зазначеним на рис. 3.7.

\begin{enumerate}
\def\labelenumi{\arabic{enumi}.}
\setcounter{enumi}{1}
\tightlist
\item
  Зверніть увагу, що гранична умова
\end{enumerate}

\[u_r(R, \theta, t) = -\frac{h}{k} [u(R, \theta, t) - g(\theta, t)]\]

На колі не означає, що гранична температура дорівнює \(g(\theta,t)\) ,
але якщо коефіцієнт теплопередачі h великий, то ці умови фактично
еквівалентні вимогі \(g(\theta,t)\) температури на межі.

\section{ЗАВДАННЯ}\label{ux437ux430ux432ux434ux430ux43dux43dux44f-2}

\begin{enumerate}
\def\labelenumi{\arabic{enumi}.}
\item
  Намалюйте контури графіків для розв'язання змішаної задачі (3.6) для
  різних моментів часу. Чи відповідають ваші графи визначеним граничним
  умовам? Яка стаціонарна температура стрижня? Очевидно, ваша відповідь
  базується на інтуїції?
\item
  Як би ви інтерпретували наступне змішане завдання?
\end{enumerate}

\[
\begin{aligned}
(\text{УЧП})\quad & u_t = \alpha^2 u_{xx}, && 0 < x < 1,\ 0 < t < \infty, \\
(\text{ГУ})\quad & u(0,t) = 0,\quad u_x(1,t) = 1, && 0 < t < \infty, \\
(\text{НУ})\quad & u(x,0) = \sin(\pi x), && 0 \leqslant x \leqslant 1.
\end{aligned}
\]

Чи можете ви графічно показати рішення цієї проблеми в різні моменти
часу? Чи буде це рішення, як правило, стаціонарним? Це очевидно?

\begin{enumerate}
\def\labelenumi{\arabic{enumi}.}
\setcounter{enumi}{2}
\tightlist
\item
  Яку фізичну інтерпретацію ви можете дати проблемі
\end{enumerate}

\[
\begin{aligned}
(\text{УЧП})\quad & u_t = \alpha^2 u_{xx}, && 0 < x < 1,\ 0 < t < \infty, \\
(\text{ГУ})\quad & u_x(0,t) = 0,\quad u_x(1,t) = 0, && 0 < t < \infty, \\
(\text{НУ})\quad & u(x,0) = \sin(\pi x), && 0 \leqslant x \leqslant 1.
\end{aligned}
\]

Чи можете ви графічно показати рішення цієї проблеми в різні моменти
часу?

Що можна сказати про стабільну температуру?

\begin{enumerate}
\def\labelenumi{\arabic{enumi}.}
\setcounter{enumi}{3}
\tightlist
\item
  Припустимо, що бічна поверхня металевого стрижня не ізольована і має
  початкову температуру 20°C, але температура одного кінця стає 50°C.
  Потім стрижень промивають рідиною з температурою 30°C. Як виглядає
  змішана задача, що відповідає цьому випадку?
\end{enumerate}

\section{ВИВЕДЕННЯ РІВНЯННЯ
ТЕПЛОПРОВІДНОСТІ}\label{ux432ux438ux432ux435ux434ux435ux43dux43dux44f-ux440ux456ux432ux43dux44fux43dux43dux44f-ux442ux435ux43fux43bux43eux43fux440ux43eux432ux456ux434ux43dux43eux441ux442ux456}

МЕТА ЛЕКЦІЇ: Показати, як можна отримати одномірне рівняння
теплопровідності

\[u_t = \alpha^2 u_{xx} + f(x, t),\]

спираючись лише на фундаментальний закон збереження кількості тепла.
Розглядається залежність швидкості теплопередачі від основних
термофізичних параметрів: теплопровідності, теплоємності та густини
матеріалу. Розглядаються різні модифікації базового рівняння
теплопровідності.

У кожній галузі знань існує набір базових принципів, правильність яких
очевидна. Усі інші твердження мають бути виведені з цих базових
принципів. Звісно, те, що очевидно одній людині, може здатися сумнівним
для іншої. Історія будь-якої науки полягає у постійному пошуку більш
фундаментальних принципів, що лежать в основі цієї науки. Основні
принципи мають бути настільки універсальними, щоб ніхто не сумнівався.

? --- Припущення A --- Припущення B --- Припущення C --- Припущення D
--- ? Рис. 4.1. Аксіоматичний метод.

Наприклад, одна людина може вважати, що всі факти певної науки можна
вивести з одного базового припущення, яке ми позначимо як B. На основі
припущення B вона або вона може довести теорему C, що призведе до
теореми D, а яка, у свою чергу, призведе до доведення багатьох інших
теорем (див. рисунок 4.1). Ось як виглядає прогрес у науці --- дослідник
може отримувати все більше нових результатів. Найвидатніші фізики,
хіміки та біологи або саме на цьому шляху.

З іншого боку, замість доведення нових теорем можна спробувати, якщо
можливо, знайти новий принцип (назвемо його A), який є більш
фундаментальним за B. Припущення A природно вважається більш
фундаментальним, ніж B.

Якщо B можна вивести з A. Отже, якщо можна знайти новий принцип A, тоді
межа наших знань зміщується до сфери більш фундаментальних ідей про
певну науку. У теорії теплопровідності основним принципом є закон
збереження енергії (теплова енергія). Усі інші твердження випливають із
цього базового принципу (рис. 4.2).

Закон збереження енергії

\[u_t = \alpha^2 u_{xx} + f(x,t)\]

Інші властивості теплопровідності Рис. 4.2. Енергозбереження: наріжний
камінь теплопровідності.

Звісно, ми могли б забути цю лекцію і використати рівняння
теплопровідності як фундаментальне поняття (дехто може так подумати,
якщо це очевидно), але тоді ми обдуримо більш серйозних студентів, адже
закон збереження вважається найфундаментальнішим принципом у науці.
Науковці, моделюючи свої задачі, часто записують відношення, що
виражають закони збереження, а потім переписують їх у вигляді рівнянь у
похідних похідних.

Повернемося до головної мети нашої лекції --- вивести рівняння
теплопровідності з рівняння збереження енергії.

\section{Виведення рівняння
теплопровідності}\label{ux432ux438ux432ux435ux434ux435ux43dux43dux44f-ux440ux456ux432ux43dux44fux43dux43dux44f-ux442ux435ux43fux43bux43eux43fux440ux43eux432ux456ux434ux43dux43eux441ux442ux456-1}

Розглянемо однорідний стрижень довжини L, щодо якого зробимо такі
припущення:

\begin{enumerate}
\def\labelenumi{\arabic{enumi}.}
\tightlist
\item
  Стрижень виготовлений з одного однорідного провідного матеріалу.
\end{enumerate}

\pandocbounded{\includegraphics[keepaspectratio]{images/_page_29_Figure_10.jpeg}}

Фіг. 4.3. Тонкий тепловий стержень.

\begin{enumerate}
\def\labelenumi{\arabic{enumi}.}
\setcounter{enumi}{1}
\item
  Бікова поверхня стрижня термічно ізольована (тепло може поширюватися
  лише вздовж осі x).
\item
  Стрижень тонка, тобто температура всіх точок у кожному поперечному
  перерізі стрижня є сталою. Якщо розглядати частину вудки на сегменті
  \([x, x+\Delta x]\) і використовувати
\end{enumerate}

закон збереження кількості тепла, тоді можна написати:

(4.1) Загальна зміна кількості тепла на сегменті \([x, x + \Delta x] =\)
дорівнює загальній кількості тепла, що пройшло через межі + Нульова
кількість тепла, утвореного всередині сегмента \([x, x + \Delta x]\) .

Оскільки загальна кількість тепла (у калоріях) всередині сегмента
\(\{x, x + \Delta x\}\) в будь-який момент часу t обчислюється формулою:

Загальна кількість тепла всередині сегмента
\([x, x+\Delta x] = \int\limits_{-\infty}^{x+\Delta x} coAu\) (s, t) ds,

де c --- питома теплоємність матеріалу (показує здатність матеріалу
накопичувати тепло),

Щільність матеріалу,

A --- це площа поперечного перерізу стрижня.

\textbf{va}con збереження енергії (4.1) можна надати у такій
математичній формі:

(4.2)

\[\frac{d}{dt} \int_{x}^{x+\Delta x} c\rho Au(s, t) ds = c\rho A \int_{x}^{x+\Delta x} u_{t}(s, t) ds = eA \left[ u_{x}(x+\Delta x, t) - u_{x}(x, t) \right] + A \int_{x+\Delta x}^{x+\Delta x} f(s, t) ds.\]

\textbf{де} k --- теплопровідність матеріалу (вказує на здатність
матеріалу проводити тепло), f(x, t) --- об'ємна потужність зовнішнього
джерела тепла (калорії/см·с).

Задача полягає в тому, щоб записати рівняння (4.2) у формі, яка не
містить інтегралів. Щоб розв'язати цю задачу, нагадаємо читачеві теорему
про середнє значення з ходу інтегральної нечисельності.

\section{Теорема про
значення}\label{ux442ux435ux43eux440ux435ux43cux430-ux43fux440ux43e-ux437ux43dux430ux447ux435ux43dux43dux44f}

Якщо функція f(x) є неперервною на сегменті {[}a, b{]}, то існує
принаймні одна точка \(\xi \in [a, b]\) така, що

\[\int_{a}^{b} f(x) dx = f(\xi) (b-a).\]

Застосовуючи цей результат до рівняння (4.2), ми отримуємо наступний
результат

\[
coA\,u_t(\xi, t)\,\Delta x = kA\left[u_x(x + \Delta x, t) - u_x(x, t)\right] + Af(\xi, t)\,\Delta x,
\]

де \(x < \xi < x + \Delta x\), тобто

\[
u_t(\xi, t) = \frac{k}{c\rho}\left[\frac{u_x(x + \Delta x, t) - u_x(x, t)}{\Delta x}\right] + \frac{1}{c\rho} f(\xi, t).
\]

Спрямовуючи \(\Delta x\) до 0, отримуємо потрібне рівняння

(4.3)

\[u_t(x, t) = \alpha^2 u_{xx}(x, t) + F(x, t).\]

де

\[\alpha^{2}=\frac{k}{c\rho}\]

--- коефіцієнт теплової дифузії,
\(F\left(x,\ t\right)=\frac{1}{c\rho}f\left(x,\ t\right)\) --- густина
джерел тепла.

Ми могли б зупинитися на цьому. Однак ми також розглянемо випадок, коли
бічна поверхня стрижня не термічно ізольована. Припустимо, що кількість
теплового потоку через бічну поверхню стрижня в цьому випадку
пропорційна різниці між температурою стрижня u(x,t) і температурою
навколишнього середовища, яка залишається сталою і дорівнює нулю. У
цьому випадку закон збереження кількості тепла веде до рівняння

\$\$

(4.4) u\_t = \alpha\^{}2 u\_\{xx\} - \beta u + F(x, t),

\$\$

де \(\beta\) --- коефіцієнт пропорційності потоку через бічну поверхню.

\section{ЗАУВАЖЕННЯ}\label{ux437ux430ux443ux432ux430ux436ux435ux43dux43dux44f-3}

\begin{enumerate}
\def\labelenumi{\arabic{enumi}.}
\tightlist
\item
  Як зазначалося вище, константа k називається теплопровідністю речовини
  і чисельно дорівнює кількості тепла (у калоріях), що проходить за
  секунду через пластину товщиною 1 см і площею поперечного перерізу 1
  см³ при різниці температур на протилежних гранях пластини рівною 1°C.
  Числові значення теплопровідності для різних матеріалів можна знайти в
  довідниках з фізики або хімії. Типові значення теплопровідності
  варіюються від 1 (для міді) до значень, близьких до нуля для хороших
  теплоізоляторів. Якщо стрижень зроблений із однорідного матеріалу, то
  значення k не залежить від x. Для деяких матеріалів значення k
  залежить від температури, і тому рівняння теплопровідності в цьому
  випадку стає нелінійним:
\end{enumerate}

\[u_{t} = \frac{1}{c\rho} \frac{\partial}{\partial x}\left(k(u)\,u_{x}\right).\]

Однак найчастіше теплопровідність k дуже слабо впливає на температуру u,
і цю нелінійність можна ігнорувати.

\begin{enumerate}
\def\labelenumi{\arabic{enumi}.}
\setcounter{enumi}{1}
\item
  Значення c зазвичай називають питомою теплоємністю речовини. Він
  характеризує здатність речовини зберігати теплову енергію. Наприклад,
  запечена картопля має високу теплоємність, оскільки може зберігати
  велику кількість тепла на одиницю маси (саме тому вона зберігає високу
  температуру так довго). Чисельна теплоємність дорівнює кількості
  калорій тепла, які потрібно передати 1 граму речовини, щоб підвищити
  її температуру на 1°C. Для більшості розглянутих задач значення c
  можна вважати сталим, не записуючи ні x, ні i. Значення питомої
  теплоємності для конкретних матеріалів можна знайти у фізичних
  довідниках.
\item
  Одиниці вимірювання основних величин, включених у рівняння
\end{enumerate}

теплопровідність (у системі GHS) наступна:

i-температура (градуси Цельсія), \(u_t\) --- швидкість зміни температури
(°C/s), \(u_x\) --- нахил температурної кривої (°C/см), \(u_{xx}\) ---
опуклість температурної кривої (°C/см²),

c = щільність тепла (калор/г \(\cdot\) °C)

k --- теплопровідність (кал/см \(\cdot\) при \(\cdot\) °C), \(\rho\) ---
щільність (г/см3), \(\alpha^2\) --- це теплова дифузія (см²/с).

\begin{enumerate}
\def\labelenumi{\arabic{enumi}.}
\setcounter{enumi}{3}
\tightlist
\item
  Слід зазначити, що теплова дифузія матеріалу
  \(\alpha^2 = \frac{k}{\epsilon \rho}\) прямо пропорційна
  теплопровідності матеріалу і обернено пропорційна густині \(\rho\) та
  питому теплоємності \(\epsilon\) ; Ми сподіваємося, що цей зв'язок
  відповідає інтуїтивним ідеям читача.
\end{enumerate}

\section{ЗАВДАННЯ}\label{ux437ux430ux432ux434ux430ux43dux43dux44f-3}

\begin{enumerate}
\def\labelenumi{\arabic{enumi}.}
\tightlist
\item
  Підставити розмірність кожної з \(u,\ u_t,\ \dots\) величин у рівняння
\end{enumerate}

\$\$

u\_t = \alpha\^{}2 u\_\{xx\} - \beta u

\$\$

і переконайтеся, що всі терміни мають однаковий розмір
\({}^{\circ}\mathrm{C}/c\) .

\begin{enumerate}
\def\labelenumi{\arabic{enumi}.}
\setcounter{enumi}{1}
\tightlist
\item
  Підставити розмірність кожної з величин у рівняння
\end{enumerate}

\$\$

u\_t = \alpha\^{}2 u\_\{xx\} - v u\_x,

\$\$

де \(\upsilon\) вимірюється в одиницях швидкості, і переконайтеся, що
всі терми мають однакову розмірність.

\begin{enumerate}
\def\labelenumi{\arabic{enumi}.}
\setcounter{enumi}{2}
\tightlist
\item
  Отримаємо рівняння теплопровідності
\end{enumerate}

\[
u_t = \frac{1}{c\rho}\frac{\partial}{\partial x}\left[h(x)u_x\right] + f(x,t).
\]

для випадку, коли коефіцієнт теплопровідності залежить від координати
\(x\).

\begin{enumerate}
\def\labelenumi{\arabic{enumi}.}
\setcounter{enumi}{3}
\tightlist
\item
  Припустимо, що \(u(x,t)\) --- концентрація речовини в потоці, який
  рухається зі швидкістю \(v\). Якщо концентрація змінюється внаслідок
  дифузії та конвекції, то відповідне рівняння має вигляд
\end{enumerate}

\[
u_t = \alpha^2 u_{xx} - v u_x.
\]

Виходячи з того, що в будь-який момент часу в певному
\([x, x + \Delta x]\) сегменті матерії не з'являється і не зникає.

ПРИМІТКА. Закон збереження матерії виглядає так:

Зміни маси речовини на сегменті \([x, x + \Delta x] =\) = Зміна
внаслідок дифузії через межі + Зміна внаслідок переносу через межі.

Лекція 5

\section{ЗМІННЕ
РОЗДІЛЕННЯ}\label{ux437ux43cux456ux43dux43dux435-ux440ux43eux437ux434ux456ux43bux435ux43dux43dux44f}

МЕТА ЛЕКЦІЇ: Познайомити читача з потужним методом розділення змінних і
показати, як цей метод можна використати для розв'язання добре відомої
проблеми дифузії. Оскільки цей метод створює труднощі для деяких учнів
(головним чином через громіздкі алгебранічні перетворення), ми
супроводжуємо презентацію інтуїтивними міркуваннями.

Основна ідея методу полягає в тому, щоб розкласти початкову умову на
найпростіші компоненти, знайти відгук системи на кожен найпростіший
компонент і підсумувати всі відповіді. Таким чином можна знайти
відповідь на довільну початкову умову.

Існуючий метод поділу змінних певною мірою зменшує можливість такої
інтерпретації методу, але все ж не виключає його.

Метод розділення змінних є одним із найвідоміших методів розв'язання
зсувних задач і використовується, якщо:

\begin{itemize}
\tightlist
\item
  \begin{enumerate}
  \def\labelenumi{\arabic{enumi}.}
  \tightlist
  \item
    Рівняння є лінійним і однорідним (не обов'язково з постійними
    коефіцієнтами).
  \end{enumerate}
\item
  \begin{enumerate}
  \def\labelenumi{\arabic{enumi}.}
  \setcounter{enumi}{1}
  \tightlist
  \item
    Граничні умови визначаються як
  \end{enumerate}
\end{itemize}

\$\$

\alpha u\_x(0, t) + \beta u(0, t) = 0, \qquad \gamma u\_x(1, t) +
\delta u(1, t) = 0.

\$\$

де \(\alpha\), \(\beta\), \(\gamma\) і \(\delta\) --- константи
(граничні умови, задані в такій формі, називаються лінійними однорідними
граничними умовами).

Метод був створений за часів Фур'є (зазвичай відомий як метод Фур'є) і,
здається, є найпопулярнішим (де це застосовно).

\pandocbounded{\includegraphics[keepaspectratio]{images/_page_34_Figure_8.jpeg}}

Рис. 5.1. Схематичне представлення задачі дифузії.

Замість того, щоб вивчати метод у загальному випадку, давайте спочатку
розглянемо конкретну задачу (загальний випадок обговоримо пізніше).
Розглянемо змішану задачу типу дифузії: щоб знайти розв'язок

\[
\begin{aligned}
(\text{УЧП})\quad & u_t = \alpha^2 u_{xx}, && 0 < x < 1,\quad 0 < t < \infty, \\
(\text{ГУ})\quad & u(0,t)=0,\quad u(1,t)=0, && 0 < t < \infty, \\
(\text{НУ})\quad & u(x,0)=\varphi(x), && 0 \le x \le 1.
\end{aligned}
\]

Перш ніж почати розділяти змінні, давайте дамо фізичну інтерпретацію
нашої проблеми. Отже, існує стрижень скінченної довжини, кінці якого
підтримуються при сталій температурі, рівній нулю (насправді кінці можна
підтримувати при значно вищій температурі, значення якої вважається
початковою точкою). Додаткові дані про завдання подаються як початкова
умова. Наша мета --- знайти розподіл температури \(u\left(x,t\right)\) в
наступний час.

\section{Загальні принципи методу розділення
змінних}\label{ux437ux430ux433ux430ux43bux44cux43dux456-ux43fux440ux438ux43dux446ux438ux43fux438-ux43cux435ux442ux43eux434ux443-ux440ux43eux437ux434ux456ux43bux435ux43dux43dux44f-ux437ux43cux456ux43dux43dux438ux445}

Для найпростішого диференціального рівняння в похідних похідних
розділення змінних є пошуком розв'язків у вигляді

\$\$

u(x, t) = X(x)T(t),

\$\$

де X(x) --- це функція, що залежить лише від змінної x, а T(t) ---
функція, що залежить лише від t. Це розв'язок у певному сенсі є
найпростішим, оскільки температура u(x, t), представлена у такій формі,
зберігатиме «форму» профілю у різний час часу t (див. рисунок 5.2).

\pandocbounded{\includegraphics[keepaspectratio]{images/_page_35_Figure_6.jpeg}}

РИСУНОК 5.2 Графік функції X(x)T(t) у різних точках часу t.

Загальна ідея полягає в тому, щоб знайти нескінченну кількість таких
часткових розв'язків, які задовольняють граничні умови

\[
u(0,t)=0, \qquad u(1,t)=0.
\]

Для цього підставляємо розв'язки (5.1) у ці граничні умови. У результаті
отримуємо

\[
\begin{aligned}
u(0,t) &= B e^{-\lambda^2 \alpha^2 t} = 0 \Rightarrow B = 0, \\
u(1,t) &= A e^{-\lambda^2 \alpha^2 t}\sin \lambda = 0 \Rightarrow \sin \lambda = 0.
\end{aligned}
\]

Отже,

\[
\lambda_n = n\pi, \qquad n = 1, 2, \ldots
\]

Зверніть увагу, що друга гранична умова могла б виконуватися і при
\(A=0\) , але тоді розв'язок (5.1) був би тривіально нульовим.

\pandocbounded{\includegraphics[keepaspectratio]{images/_page_38_Figure_3.jpeg}}

Фіг. 5.3. Фундаментальне рішення
\(n_n(x, t) = A_n e^{-(\pi n \alpha)^{n/t}} \sin{(\pi n x)}\) .

Тепер, коли ми зробили другий крок, маємо нескінченний набір ознак

(5.2)

\[u_n(x, t) = A_n e^{-(n\pi\alpha)^2 t} \sin(n\pi x), \quad n = 1, 2 \dots\]

кожна з яких задовольняє рівняння з частковими похідними та граничними
умовами). Отримані функції --- це будівельні блоки, з яких ми будемо
створювати рішення нашої проблеми. Це рішення буде певною сумою

2) Зверніть увагу, що функції \(u_n\) і \(u_{-n}\) відрізняються лише
знаком.

цих найпростіших функцій. Конкретний тип суми залежить від початкового
стану. Фіг. Рисунок 5.3 показує графіки кількох таких фундаментальних
функцій

КРОК 3. (Пошук розв'язку, який задовольняє рівняння, межу та початкові
умови.)

Останній (і, ймовірно, найцікавіший з математичної точки зору) крок ---
знайти таку суму фундаментальних розв'язків

\[u(x, t) = \sum_{n=1}^{\infty} A_n e^{-(n\pi\alpha)^n t} \sin(n\pi x),\]

тобто при виборі таких коефіцієнтів \(A_n\) , що функція задовольнятиме
початкову умову

\[u\left( x,\ 0\right) =\varphi\left( x\right) .\]

Підставляючи суму в початковому стані, отримуємо

(5.3)

\[q(x) = \sum_{n=1}^{\infty} A_n \sin(n\pi x).\]

Це рівняння підводить нас до цікавого питання: чи можна розкласти
початкову температуру \(\phi(x)\) на ряд елементарних функцій вигляду

\[A_1 \sin(\pi x) + A_2 \sin(2\pi x) + A_3 \sin(3\pi x) + \dots\]

?

Позитивну відповідь на це питання дав французький математик Жозеф Фур'є.
Виявилося, що для достатньо хороших функцій така декомпозиція можлива
\(^1\) ). Тоді виникає нове питання: знайти коефіцієнти розкладу \(A_n\)
?

Насправді це дуже просто. Для цього використаємо властивість системи
функцій

\[\{\sin(n\pi x)\}_{n=1}^{\infty}.\]

відомий як ортогональність.

У нашому випадку (див. завдання 2) ці функції задовольняють умови

\[\int_{0}^{1} \sin(m\pi x) \sin(n\pi x) dx = \begin{cases} 0, & m \neq n, \\ 1/2, & m = n. \end{cases}\]

Умови розкладу функцій у рядах Фур'є наведені в лекції 11. Для глибшого
ознайомлення з цим питанням можна звернутися до численних освітніх
літератур (див., наприклад: Г. П. Толстов, Серія Фур'є. --- Москва:
Наука, 1980). --- Ед. Ед.

Ця властивість чітко показана на рис. 5.4, яка показує графіки деяких
функцій із цієї системи.

Отже, ми хочемо знайти коефіцієнти у розкладі

\[\varphi(x) = A_1 \sin(\pi x) + A_2 \sin(2\pi x) + A_3 \sin(3\pi x) + \dots\]

\pandocbounded{\includegraphics[keepaspectratio]{images/_page_40_Figure_5.jpeg}}

РИСУНОК 5.4. Послідовність ортогональних функцій

Помножте обидві сторони цього співвідношення на \(\sin{(m\pi x)}\) (m
--- довільне ціле число) і інтегруйте від нуля до одиниці. У результаті
ми отримуємо

\[\int_{0}^{1} \varphi(x) \sin(m\pi x) \, dx = A_{m} \int_{0}^{1} \sin^{2}(m\pi x) \, dx = \frac{A_{m}}{2}\]

(усі інші члени зрівнялися з нулем через ортогональність). Розв'язуючи
рівняння відносно \(A_m\) , отримуємо

\[A_m = 2 \int_0^1 \varphi(x) \sin(m\pi x) dx.\]

Отже, отримано, що розв'язок записується у вигляді

(5.4)

\[u(x, t) = \sum_{n=1}^{\infty} A_n e^{-(n\pi\alpha)^2 t} \sin(\pi n x),\]

де коефіцієнти \(A_n\) визначаються формулами

\[A_n = 2\int_0^1 \varphi(x) \sin(\pi nx) dx\]

Ви можете переконатися, що отримане рішення задовольняє всі умови
початкової задачі. На цьому завершується третій крок.

Багато студентів розчаровані, коли бачать, наскільки складним виглядає
рішення, і майже не намагаються дивитися на нього знову (виглядає дуже
погано). Однак це рішення не здаватиметься таким складним, якщо ви
витратите час на його аналіз. Насправді, складніша форма розв'язку
робить його більш інформативним. Нижче наведено кілька нотаток, які
допоможуть вам інтерпретувати це рішення.

\section{ЗАУВАЖЕННЯ}\label{ux437ux430ux443ux432ux430ux436ux435ux43dux43dux44f-4}

Зверніть увагу, що єдина різниця між синусоїдальним розкладом функції
φ(x) і розв'язком (5.4) --- це наявність часового фактора e-(nπα)31 у
кожному члені ряду. Отже, якщо початкова умова має просту декомпозицію,
наприклад, вигляду

\[\varphi(x) = \sin(\pi x) + \frac{1}{2}\sin(3\pi x),\]

тоді рішення можна негайно записати

\[u(x, t) = e^{-(\pi\alpha)^2 t} \sin(\pi x) + \frac{1}{2} e^{-(3\pi\alpha)^2 t} \sin(3\pi x).\]

У цьому випадку очевидно, що якщо ми розкладемо \(\phi(x)\) на ряд Фур'є
вздовж синусів, то отримуємо

\[A_1 = 1,\] \(A_2 = 0,\) \(A_3 = 1/2,\) \(A_4 = A_5 = \dots = 0.\) - 2.
Розв'язок (5.4) можна інтерпретувати так: ми представляємо початкову
температуру \(\varphi(x)\) як суму найпростіших функцій виду
\(A_n \sin(\pi nx)\) ; кожна така функція породжує відповідь
\(A_n e^{-(\pi n\alpha)^2 t} \sin(\pi nx)\) . Сумуючи всі такі
відповіді, отримуємо розв'язок, який відповідає початковій умові
\(u(x, 0) = \varphi(x)\) . - 3. Кожен член у декомпозиції

\[u(x, t) = A_1 e^{-(\pi\alpha)^2 t} \sin(\pi x) + A_2 e^{-(2\pi\alpha)^2 t} \sin(2\pi x) + \cdots\]

є функцією від x і t. Зверніть увагу, що внесок членів з великими
числами в t\textgreater0 дуже малий через множник
\(e^{-(n\pi\alpha)^2 t}\) . Відповідно, після закінчення досить великого
Р и с

\pandocbounded{\includegraphics[keepaspectratio]{images/_page_42_Figure_2.jpeg}}

. 5.5. У задачах дифузії вищі порядки терміни швидко затухають;
\(\phi(x)\) --- початкова температура. Слід зазначити, що по-перше, у
розчині зникають високі частоти --- по-перше, зникли найбільші коливання
міді.

Повний розв'язок приблизно збігається з першим членом, який є півхвилі
синусоїди, що з часом згасає (рис. 1). 5.5.)

\section{ЗАВДАННЯ}\label{ux437ux430ux432ux434ux430ux43dux43dux44f-4}

\begin{enumerate}
\def\labelenumi{\arabic{enumi}.}
\tightlist
\item
  Показати, що функції
\end{enumerate}

\[u(x, t) = e^{-\lambda^2 \alpha' t} \left[ A \sin(\lambda x) + B \cos(\lambda x) \right]\]

задовольняють рівняння \(u_l = \alpha^2 u_{xx}\) для довільних значень
A, B і \(\lambda\) .

\begin{enumerate}
\def\labelenumi{\arabic{enumi}.}
\setcounter{enumi}{1}
\tightlist
\item
  Покажи це
\end{enumerate}

\[\int_0^1 \sin(m\pi x) \sin(n\pi x) dx = \begin{cases} 0, & m \neq n, \\ 1/2, & m = n. \end{cases}\]

ПРИМІТКА. Використання ідентичності

\[\sin(mx)\sin(nx) = \frac{1}{2}[\cos(m-n)x - \cos(m+n)x].\]

\begin{enumerate}
\def\labelenumi{\arabic{enumi}.}
\setcounter{enumi}{2}
\tightlist
\item
  Знайти розклад у рядах Фур'є синусами функції \(\varphi(x) = 1\) на
  сегменті {[}0, 1{]}. Побудуйте перші три терміни розкладу.
\end{enumerate}

Використовуючи результати завдання 3, знайдіть розв'язок наступної
змішаної задачі:

(PPP)

\[u_t = u_{xx}\] \(0 < x < 1\) , (GI)

\[\begin{cases} u(0, t) = 0, \\ u(1, t) = 0, \end{cases}\]
\(0 < t < \infty\) , (НУ) \(u(x, 0) = 1, \quad 0 \le x \le 1\) .

(Зверніть увагу, що ця проблема фізично не має значення, оскільки
припускається, що температура на кінцях стрижня миттєво знижується з
одного до нуля \(\phi(x)\) .

\begin{enumerate}
\def\labelenumi{\arabic{enumi}.}
\setcounter{enumi}{4}
\tightlist
\item
  Знайти розв'язок проблеми 4, якщо початкова умова встановлена у
  вигляді
\end{enumerate}

\[u(x, 0) = \sin(2\pi x) + \frac{1}{3}\sin(4\pi x) + \frac{1}{5}\sin(6\pi x).\]

\begin{enumerate}
\def\labelenumi{\arabic{enumi}.}
\setcounter{enumi}{5}
\tightlist
\item
  Розв'яжіть задачу 4, якщо початкова умова
  \(u\left( x,\ 0\right) =x-x^{2},\quad 0< x<1.\) Лекція 6
\end{enumerate}

\section{ПЕРЕТВОРЕННЯ НЕОДНОРІДНИХ ГРАНИЧНИХ УМОВ У ОДНОРІДНІ ГРАНИЧНІ
УМОВИ}\label{ux43fux435ux440ux435ux442ux432ux43eux440ux435ux43dux43dux44f-ux43dux435ux43eux434ux43dux43eux440ux456ux434ux43dux438ux445-ux433ux440ux430ux43dux438ux447ux43dux438ux445-ux443ux43cux43eux432-ux443-ux43eux434ux43dux43eux440ux456ux434ux43dux456-ux433ux440ux430ux43dux438ux447ux43dux456-ux443ux43cux43eux432ux438}

МЕТА ЛЕКЦІЇ: Показати, як змішану задачу з неоднорідними граничними
умовами можна звести до задачі з однорідними граничними умовами.

\[
\begin{aligned}
(\text{УЧП})\quad & u_t - \alpha^2 u_{xx} = f(x,t), \\
(\text{ГУ})\quad & \alpha_1 u_x(0,t) + \beta_1 u(0,t) = g_1(t), \\
& \alpha_2 u_x(L,t) + \beta_2 u(L,t) = g_2(t), \\
(\text{НУ})\quad & u(x,0) = \varphi(x).
\end{aligned}
\]

Після відповідної заміни невідомої функції ця задача часто зводиться до
задачі вигляду

\[
\begin{aligned}
U_t - \alpha^2 U_{xx} &= F(x,t), \\
\alpha_1 U_x(0,t) + \beta_1 U(0,t) &= 0, \\
\alpha_2 U_x(L,t) + \beta_2 U(L,t) &= 0, \\
U(x,0) &= \overline{\varphi}(x),
\end{aligned}
\]

тобто до задачі з нульовими граничними умовами. Якщо \(F(x,t)=0\), нову
задачу можна розв'язувати методом розділення змінних; якщо ж
\(F(x,t)\neq 0\), доводиться використовувати інші методи, зокрема
інтегральні перетворення або розклади за власними функціями.

Щоб метод розділення змінних був застосовним, граничні умови мають бути
лінійними та однорідними, тобто мати вигляд

(6.1)

\[
\begin{aligned}
\alpha_1 u_x(0,t) + \beta_1 u(0,t) &= 0, \\
\alpha_2 u_x(L,t) + \beta_2 u(L,t) &= 0.
\end{aligned}
\]

Мета цієї лекції --- пояснити, як задачі з неоднорідними граничними
умовами форми

(6.2)

\[
\begin{aligned}
(\text{УЧП})\quad & u_t = \alpha^2 u_{xx}, \\
(\text{ГУ})\quad & \alpha_1 u_x(0,t) + \beta_1 u(0,t) = g_1(t), \\
& \alpha_2 u_x(L,t) + \beta_2 u(L,t) = g_2(t), \\
(\text{НУ})\quad & u(x,0) = \varphi(x),
\end{aligned}
\]

можна звести до задач з однорідними граничними умовами.

Почнемо з найпростішого прикладу, коли кінці стрижня підтримуються при
сталих температурах.

\section{Перетворення неоднорідних граничних умов у однорідні граничні
умови}\label{ux43fux435ux440ux435ux442ux432ux43eux440ux435ux43dux43dux44f-ux43dux435ux43eux434ux43dux43eux440ux456ux434ux43dux438ux445-ux433ux440ux430ux43dux438ux447ux43dux438ux445-ux443ux43cux43eux432-ux443-ux43eux434ux43dux43eux440ux456ux434ux43dux456-ux433ux440ux430ux43dux438ux447ux43dux456-ux443ux43cux43eux432ux438-1}

Розглянемо проблему поширення тепла в термоізольованому стрижні, кінці
якого підтримуються при сталих температурах \(k_1\) і \(k_2\), тобто

\[
\begin{aligned}
(\text{УЧП})\quad & u_t = \alpha^2 u_{xx}, && 0 < x < L,\ 0 < t < \infty, \\
(6.3)\ (\text{ГУ})\quad & u(0,t) = k_1,\quad u(L,t) = k_2, && 0 < t < \infty, \\
(\text{НУ})\quad & u(x,0) = \varphi(x), && 0 \leqslant x \leqslant L.
\end{aligned}
\]

Складність цієї задачі полягає в тому, що неоднорідні граничні умови не
дозволяють безпосередньо застосувати метод розділення змінних. Однак
очевидно, що при \(t \to \infty\) розв'язок має наближатися до
стаціонарного лінійного профілю температури. Тому природно шукати
розв'язок у вигляді суми стаціонарної та перехідної частин:

\[
u(x,t)=\left[k_1+\frac{x}{L}(k_2-k_1)\right]+U(x,t).
\]

У цьому випадку достатньо знайти перехідну температуру \(U(x,t)\).

\pandocbounded{\includegraphics[keepaspectratio]{images/_page_45_Figure_7.jpeg}}

Рис. 6.1. Розв'язання задачі (6.3) у різні моменти часу: а) початкова
температура \(u(x,0)=\varphi(x)\); б) стаціонарна температура
\(u(x,\infty)=k_1+\frac{x}{L}(k_2-k_1)\); в) перехідна температура.

Після цієї заміни задача (6.3) зводиться до нової задачі для функції
\(U(x,t)\):

(6.4)

\[
\begin{aligned}
U_t &= \alpha^2 U_{xx}, && 0 < x < L, \\
U(0,t) &= 0,\quad U(L,t) = 0, && 0 < t < \infty, \\
U(x,0) &= \varphi(x) - \left[k_1 + \frac{x}{L}(k_2 - k_1)\right] = \overline{\varphi}(x).
\end{aligned}
\]

де \(\overline{\varphi}(x)\) --- нова, але відома початкова умова.

Це вже задача з однорідними граничними умовами, отже її можна
розв'язувати методом розділення змінних. Розв'язок має вигляд

(6.5)

\[
U(x,t)=\sum_{n=1}^{\infty} a_n e^{-(n\pi\alpha/L)^2 t}\sin\!\left(\frac{n\pi x}{L}\right),
\]

де

\[
a_n=\frac{2}{L}\int_0^L \overline{\varphi}(\xi)\sin\!\left(\frac{n\pi \xi}{L}\right)\,d\xi.
\]

Це все для сердечників із фіксованою температурою на кінцях. А як щодо
більш реалістичних GU з правою стороною, залежною від часу? Основні ідеї
тут ті ж, що й у попередній задачі, але вони набагато складніші.

\section{Перетворення граничних умов, залежних від часу, у
нульові}\label{ux43fux435ux440ux435ux442ux432ux43eux440ux435ux43dux43dux44f-ux433ux440ux430ux43dux438ux447ux43dux438ux445-ux443ux43cux43eux432-ux437ux430ux43bux435ux436ux43dux438ux445-ux432ux456ux434-ux447ux430ux441ux443-ux443-ux43dux443ux43bux44cux43eux432ux456}

Розглянемо типову задачу

\[
\begin{aligned}
(\text{PPP})\quad & u_t = \alpha^2 u_{xx}, && 0 < x < L,\quad 0 < t < \infty, \\
(\text{ГУ})\quad & u(0,t)=g_1(t),\quad u_x(L,t)+h\,u(L,t)=g_2(t), && 0 < t < \infty, \\
(\text{НУ})\quad & u(x,0)=\varphi(x), && 0 \le x \le L.
\end{aligned}
\] Щоб перетворити ці граничні умови на нульові умови, ми, після деяких
спроб і помилок, зупинилися на такій формі розв'язку:

(6.7)

\[
u(x, t)=A(t)\left(1-\frac{x}{L}\right)+B(t)\frac{x}{L}+U(x, t).
\]

де функції \(A\left(t\right)\) і \(B\left(t\right)\) обираються так, що
«квазі-стаціонарна» частина розв'язку (6.7)

(6.8)

\[
S(x, t)=A(t)\left(1-\frac{x}{L}\right)+B(t)\frac{x}{L}.
\]

задовольняє граничні умови початкової задачі. У цьому випадку функція
\(U\left(x,\ t\right)\) задовольнятиме однорідні граничні умови.

Підставляючи функцію S(x, t) у граничні умови

\[
S(0, t)=g_1(t), \qquad S_x(L, t)+hS(L, t)=g_2(t).
\]

одержуємо два рівняння для \(A(t)\) і \(B(t)\). У результаті маємо

(6.9)

\[
A(t)=g_1(t), \qquad B(t)=\frac{g_1(t)+Lg_2(t)}{1+Lh}.
\]

Отже,

\[
u(x, t)=g_1(t)\left(1-\frac{x}{L}\right)+\frac{g_1(t)+Lg_2(t)}{1+Lh}\frac{x}{L}+U(x, t).
\]

Якщо підставити ці вирази для \(u(x,t)\) у початкову задачу, то
отримаємо нову задачу для невідомої функції \(U(x,t)\):

\[
\begin{aligned}
(\text{УЧП})\quad & U_t = \alpha^2 U_{xx} - S_t(x,t), && 0 < x < L,\quad 0 < t < \infty, \\
(\text{ГУ})\quad & U(0,t) = 0,\quad U_x(L,t) + hU(L,t) = 0, && 0 < t < \infty, \\
(\text{НУ})\quad & U(x,0) = \varphi(x) - S(x,0), && 0 \leqslant x \leqslant L.
\end{aligned}
\]

Тепер ми стикаємося з новою проблемою однорідних граничних умов (на
жаль, рівняння стало гетерогенним). Ми не можемо розв'язати цю задачу
методом розділення змінних, але якщо читач погоджується чекати, то за
кілька лекцій він ознайомиться з розв'язком цієї задачі методом цілісних
перетворень і розкладання за власними функціями.

\section{ЗАУВАЖЕННЯ
.}\label{ux437ux430ux443ux432ux430ux436ux435ux43dux43dux44f-.}

\begin{enumerate}
\def\labelenumi{\arabic{enumi}.}
\item
  Як зазначалося раніше, мета цієї лекції --- навчити читача, як
  перетворювати задачі з неунірольними граничними умовами на задачі з
  нульовими граничними умовами. Після цього може статися, що нове
  диференціальне рівняння в похідних похідних стане однорідним. Якщо це
  станеться, ми можемо вважати себе щасливцями, оскільки в цьому випадку
  проблему можна розв'язати методом розділення змінних (див. перший
  приклад з цієї лекції). Якщо перетворене рівняння виявиться
  гетерогенним, тоді нову задачу потрібно розв'язати іншим методом.
\item
  Лінійні неоднорідні граничні умови найзагальної форми
\end{enumerate}

\[
\alpha_1 u_x(0,t) + \beta_1 u(0,t) = g_1(t), \qquad \alpha_2 u_x(L,t) + \beta_2 u(L,t) = g_2(t).
\]

Ви також можете перетворити ці умови на нульові граничні умови,
використовуючи метод, описаний у другому прикладі. Звісно, нове рівняння
при цьому, як правило, буде гетерогенним.

\begin{itemize}
\tightlist
\item
  \begin{enumerate}
  \def\labelenumi{\arabic{enumi}.}
  \setcounter{enumi}{2}
  \tightlist
  \item
    Деякі методи розв'язку не містять вимог щодо унікальності граничних
    умов і, отже, не потребують попередніх перетворень. Пізніше, коли ми
    ознайомимося з перетворенням Лапласа, побачимо, що нульові граничні
    умови є необов'язковими (у випадку нульових граничних умов розв'язок
    дещо простіший, ніж в інших випадках).
  \end{enumerate}
\item
  \begin{enumerate}
  \def\labelenumi{\arabic{enumi}.}
  \setcounter{enumi}{3}
  \tightlist
  \item
    Для граничних умов типу
  \end{enumerate}
\end{itemize}

\[
u(0, t)=g_1(t), \qquad u(L, t)=g_2(t).
\]

Метод, описаний у другому прикладі, призводить до наступної форми
розв'язання задачі:

\[
u(x, t)=g_1(t)+\frac{x}{L}\left[g_2(t)-g_1(t)\right]+U(x, t).
\]

\section{ЗАВДАННЯ}\label{ux437ux430ux432ux434ux430ux43dux43dux44f-5}

\begin{enumerate}
\def\labelenumi{\arabic{enumi}.}
\tightlist
\item
  Знайти розв'язок наступної змішаної задачі:
\end{enumerate}

\[
\begin{aligned}
(\text{УЧП})\quad & u_t = \alpha^2 u_{xx} + x\cos t, && 0 < x < 1,\quad 0 < t < \infty, \\
(\text{ГУ})\quad & u(0,t) = 1,\quad u_x(1,t) + hu(1,t) = 1, && 0 < t < \infty, \\
(\text{НУ})\quad & u(x,0) = \sin(\pi x), && 0 \leqslant x \leqslant 1.
\end{aligned}
\]

попередньо перетворюючи його на задачу з нульовими граничними умовами.
Чи відповідає рішення, яке ви отримуєте, вашим задумам щодо цієї
проблеми?

\begin{enumerate}
\def\labelenumi{\arabic{enumi}.}
\setcounter{enumi}{1}
\tightlist
\item
  Перетворіть у задачу з нульовими граничними умовами і розв'яжіть
  задачу
\end{enumerate}

\[
\begin{aligned}
(\text{УЧП})\quad & u_t = u_{xx}, && 0 < x < 1,\quad 0 < t < \infty, \\
(\text{ГУ})\quad & u(0,t)=0,\quad u(1,t)=1, && 0 < t < \infty, \\
(\text{НУ})\quad & u(x,0)=x^2, && 0 \le x \le 1.
\end{aligned}
\]

Як виглядатиме розв'язок у різні моменти часу? Чи відповідає він вашим
інтуїтивним уявленням? Що таке стаціонарний розв'язок і як виглядає
перехідна частина?

\begin{enumerate}
\def\labelenumi{\arabic{enumi}.}
\setcounter{enumi}{2}
\tightlist
\item
  Завдання
\end{enumerate}

\[
\begin{aligned}
(\text{УЧП})\quad & u_t = u_{xx}, && 0 < x < 1,\ 0 < t < \infty, \\
(\text{ГУ})\quad & u_x(0,t) = 0,\quad u_x(1,t) + h u(1,t) = 1, && 0 < t < \infty, \\
(\text{НУ})\quad & u(x,0) = \sin(\pi x), && 0 \leqslant x \leqslant 1.
\end{aligned}
\]

Перетворимо на задачу з нульовими граничними умовами. Чи буде нове
рівняння однорідним?

Лекція 7

\section{РОЗВ'ЯЗАННЯ СКЛАДНІШИХ ЗАДАЧ ШЛЯХОМ РОЗДІЛЕННЯ
ЗМІННИХ}\label{ux440ux43eux437ux432ux44fux437ux430ux43dux43dux44f-ux441ux43aux43bux430ux434ux43dux456ux448ux438ux445-ux437ux430ux434ux430ux447-ux448ux43bux44fux445ux43eux43c-ux440ux43eux437ux434ux456ux43bux435ux43dux43dux44f-ux437ux43cux456ux43dux43dux438ux445}

МЕТА ЛЕКЦІЇ: Показати, як більш складні проблеми теплопровідності можна
розв'язати методом розділення змінних. У цій лекції ми головним чином
розглядаємо такі задачі, розв'язання яких дозволить читачеві краще
ознайомитися з методом ділення змінних. Сподіваємося, що читач зможе
екстраполювати викладені тут ідеї для власних завдань.

У тій же лекції читач ознайомиться з задачею за її власними функціями,
або, як її ще називають, проблемою Штурма-Ліувіля. Обговорюються деякі
властивості проблеми Штурма-Ліувіля.

Як згадувалося вище, мета цієї лекції --- розв'язати змішану задачу
шляхом розділення змінних. Можливо, саме таким методом читач вирішить
свою проблему. Водночас ми сподіваємося, що він зможе перенести метод,
розглянутий у цій лекції, для власних завдань.

Почнемо з одномірної задачі теплопровідності, в якій одна з граничних
умов містить похідну.

\section{Задача теплопровідності з похідною в граничній
умові}\label{ux437ux430ux434ux430ux447ux430-ux442ux435ux43fux43bux43eux43fux440ux43eux432ux456ux434ux43dux43eux441ux442ux456-ux437-ux43fux43eux445ux456ux434ux43dux43eux44e-ux432-ux433ux440ux430ux43dux438ux447ux43dux456ux439-ux443ux43cux43eux432ux456}

Розглянемо пристрій, показаний на рис. 7.1. Припустимо, що верхній
кінець стрижня підтримується при нульовій температурі, тобто u(0, t) =
0, а нижній кінець занурений у воду, температура якої також стала і
дорівнює нулю (у цьому випадку --- нуль Р и с

\pandocbounded{\includegraphics[keepaspectratio]{images/_page_50_Picture_4.jpeg}}

. 7.1. Схематичне представлення зсувної задачі.

\pandocbounded{\includegraphics[keepaspectratio]{images/_page_50_Figure_6.jpeg}}

Фіг. 7.2. Природа шахраїв, як бути задоволеними

\[
(\text{ГУ})\quad
\begin{cases}
u(0,t)=0,\\
u_x(1,t)=-h\,u(1,t).
\end{cases}
\]

відповідає певній еталонній температурі). Теплообмін між стрижнем і
водою відбувається згідно із законом Ньютона, який записується у x=1 як
наступна гранична умова:

\[
u_x(1, t) = -hu(1, t).
\]

Тепер припустимо, що початкова температура стрижня дорівнює u(x, 0) = x,
але одразу при t \textgreater{} 0 включаємо граничні умови. Щоб
визначити температуру в стрижні, необхідно розв'язати проблему зміщення

\[
\begin{aligned}
(\text{УЧП})\quad & u_t = \alpha^2 u_{xx}, && 0 < x < 1,\ 0 < t < \infty, \\
(7.1)\ (\text{ГУ})\quad & u(0,t) = 0,\quad u_x(1,t) + h u(1,t) = 0, && 0 < t < \infty, \\
(\text{НУ})\quad & u(x,0) = x, && 0 \leq x \leq 1.
\end{aligned}
\]

Щоб використати метод розділення змінних, розбиємо процес розв'язання на
кілька етапів.

\textbf{КРОК} 1. (Перетворення диференціального рівняння в похідних у
два звичайних диференціальних рівняння)

Підставляючи \(u(x,t)=X(x)T(t)\) у рівняння (7.1), отримуємо

\[
XT'=\alpha^2X''T.
\]

Поділивши обидві частини на \(\alpha^2XT\), маємо

\[
\frac{T'}{\alpha^2 T} = \frac{X''}{X} .
\]

Оскільки ліва частина останнього рівняння залежить лише від часу, а
права --- лише від x (і оскільки x і t незалежні), обидві сторони
рівняння мають бути константами. Прирівнюючи обидві частини до однієї
сталої \(\mu\) , отримуємо

(7.2)

\[
T' - \mu \alpha^{2} T = 0, X'' - \mu X = 0.
\]

Процес розщеплення змінних завершено.

КРОК 2. (Визначення константи відокремлення)

По-перше, константа і не повинна бути додатною, інакше \(^{1}\) )
функція \(T\left( t\right)\) буде зростати експоненційно, що призведе до
нескінченного підвищення температури u=XT. Очевидно, що таке рішення не
має фізичного значення.

По-друге, припустимо, що \(\mu = 0\) . Тоді

\[
X'' = 0
\]

і, отже.

\[
X(x) = A + \hat{B}x.
\]

Однак, оскільки граничні умови вимагають, що

\[
u(0, t) = X(0) T(t) = 0, \qquad u(1,t)=\bigl(X'(1)+hX(1)\bigr)T(t)=0.
\]

Ми одразу дійшли висновку, що

\[
X(0)=0\Rightarrow A=0,\qquad X'(1)+hX(1)=0\Rightarrow B=0.
\]

Це означає, що \(u(x, t) \equiv 0\) . Іншими словами, лише нульове
розв'язання відповідає нульовій константі розділення. Ми відкинемо його,
оскільки нас цікавлять ненульові розв'язки.

Якщо, нарешті, \(\mu < 0\) , можна поставити \(\mu = -\lambda^2\) і
переписати рівняння (7.2) у вигляді

\[
T' + \lambda^2 \alpha^2 T = 0, \qquad X'' + \lambda^2 X = 0.
\]

1) Розв'язок X(x) відповідної задачі з граничним значенням буде лише
тривіальним, X(x) = 0. --- Ед. Ед.

Загальні розв'язки цих рівнянь легко знайти:

\[
T(t) = Ae^{-(\lambda \alpha)^2 t},
\qquad
X(x) = B\sin(\lambda x) + C\cos(\lambda x).
\]

Отже, будь-яка функція провідника

(7.3)

\[
u(x, t) = e^{-(\lambda \alpha)^2 t} [A \sin(\lambda x) + B \cos(\lambda x)]
\]

при будь-яких значеннях \(\lambda\) A і B задовольняє рівняння (читач
може перевірити це самостійно простою заміною).

Тепер нам потрібно з'ясувати, які з цих функцій задовольняють граничним
умовам

(7.4)

\[
u(0, t) = 0, u_x(1, t) + hu(1, t) = 0.
\]

Підстановка розв'язків виду (7.3) у граничні умови (7.4) дозволяє знайти
відношення, яким мають задовольняти значення \(\lambda\) , \(\Lambda\)
та B:

\[
Be^{-(\lambda\alpha)^2 t} = 0 \Rightarrow B = 0,
\]

\(A\lambda e^{-(\lambda\alpha)^2 t} \cos \lambda + hAe^{-(\lambda\alpha)^2 t} \sin \lambda = 0.\)
Після простих алгебранхічних перетворень останнє рівняння зводиться до
вигляду

\[
\tan \lambda = -\lambda/h
\]

.

Іншими словами, щоб визначити \(\lambda\) , потрібно знайти точки
перетину графів функцій tg \(\lambda\) і --- \(\lambda/h\) (рисунок
7.3).

Числові значення значень \(\lambda_1\) , \(\lambda_2\) , \ldots{} можна
знайти за допомогою комп'ютера, якщо задано значення h Р и с

\pandocbounded{\includegraphics[keepaspectratio]{images/_page_52_Figure_15.jpeg}}

. 7.3. Перетин графів функцій \(tg(\lambda)\) і \(-\lambda/h\) .

викликати власні значення задачі граничного значення

(7.5)

\[
X''+\lambda^{2}X=0,\qquad
\begin{cases}
X(0)=0,\\
X'(1)+hX(1)=0.
\end{cases}
\]

Іншими словами, з цими значеннями параметра \(\lambda\) існують
нетривіальні розв'язки задачі (7.5). Власні значення задачі (7.5)
визначаються як корені рівняння

\[
\tan\lambda = -\frac{\lambda}{h}.
\]

Їхні числові значення знаходилися на комп'ютері для випадку \(h=1\).
Таблиця 7.1 показує результати розрахунків перших кількох власних
значень.

{\def\LTcaptype{none} % do not increment counter
\begin{longtable}[]{@{}llll@{}}
\toprule\noalign{}
n & \(\lambda_B\) & rt . & \(\lambda_n\) \\
\midrule\noalign{}
\endhead
\bottomrule\noalign{}
\endlastfoot
123 & 2,024,917,98 & 45 & 11,0814,20 \\
\end{longtable}
}

Розв'язки задачі (7.5), що відповідають власним значенням \(\lambda_n\)
, називаються власними функціями \(X_n(x)\) . Для розглянутого завдання
власні функції виглядають так:

\[
X_n(x) = \sin(\lambda_n x)
\]

і зображені на рис. 7.4.

\pandocbounded{\includegraphics[keepaspectratio]{images/_page_53_Figure_10.jpeg}}

Рис. 7.4. Власні функції \(X_n(x)\) завдання (7.5) при h=1; Усі власні
функції задовольняють граничні умови.

КРОК 3, (Пошук фундаментальних рішень)

Тепер у нас є нескінченний набір функцій (фундаментальних розв'язків) у
нашому розпорядженні

\[
u_n(x, t) = X_n(x) T_n(t) = e^{-(\lambda_n\alpha)^2 t} \sin(\lambda_n x),
\]

кожна з яких задовольняє рівняння та граничні умови. Залишився ще один
крок --- загалом

\[
u(x, t) = \sum_{n=1}^{\infty} a_n X_n(x) T_n(t)
= \sum_{n=1}^{\infty} a_n e^{-(\lambda_n \alpha)^2 t}\sin(\lambda_n x)
\]

Вибрати коефіцієнти \(a_n\) так, щоб функція u(x, t) задовольняла
початкову умову u(x, 0) = x (у цьому випадку функція u(x, t)
задовольнятиме рівняння та граничні умови, оскільки, по-перше, кожен
член суми їх задовольняє, а по-друге, і рівняння, і граничні умови є
лінійними та однорідними).

Отже, ймовірність \(a_n\) слід визначати з цієї умови

(7 6)

\[
u(x, 0) = x = \sum_{n=1}^{\infty} a_n \sin(\lambda_n x).
\]

Ми маємо зробити останній крок.

КРОК 4. (Розклад початкової умови на ряд за власними функціями)

Щоб знайти коефіцієнти \(a_n\) у розкладі (7,6) на власні функції,
помножте обидві частини цього співвідношення на \(\sin(\lambda_m x)\) і
інтегруйте від 0 до 1. У результаті ми отримуємо

\[
\int_{0}^{1} x \sin(\lambda_m x)\,dx = \sum_{n=1}^{\infty} a_{n} \int_{0}^{1} \sin(\lambda_{n} \xi) \sin(\lambda_{m} \xi) d\xi
= a_{m} \int_{0}^{1} \sin^{2}(\lambda_{m} \xi) d\xi
= a_{m} \left( \frac{\lambda_{m} - \sin \lambda_{m} \cos \lambda_{m}}{2\lambda_{m}} \right).
\]

Розв'язавши останнє відношення відносно \(a_m\) і замінивши m на n,
отримуємо бажаний результат:

(7.7)

\[
a_n = \frac{2\lambda_n}{(\lambda_n - \sin \lambda_n \cos \lambda_n)} \int_0^1 x \sin(\lambda_n x)\,dx.
\]

Таблиця 7.2.

\section{\texorpdfstring{Шанси \(a_n\) у формулі
(7.8)}{Шанси a\_n у формулі (7.8)}}\label{ux448ux430ux43dux441ux438-a_n-ux443-ux444ux43eux440ux43cux443ux43bux456-7.8}

{\def\LTcaptype{none} % do not increment counter
\begin{longtable}[]{@{}llll@{}}
\toprule\noalign{}
n & an & n & \(a_n\) \\
\midrule\noalign{}
\endhead
\bottomrule\noalign{}
\endlastfoot
123 & 0,240,22-0,03 & 45 & -0.11-0.09  \\
\end{longtable}
}

Отже, розв'язок задачі (7.1) записується у вигляді

(7.8)

\[
u(x, t) = \sum_{n=1}^{\infty} a_n e^{-(\lambda_n \alpha)^2 t} \sin(\lambda_n x),
\]

де коефіцієнти \(a_n\) задаються формулами (7.7). Перші п'ять
коефіцієнтів \(a_n\) нашої задачі були знайдені чисельно. Їхні значення
наведені в Таблиці 1. 7.2.

\pandocbounded{\includegraphics[keepaspectratio]{images/_page_55_Figure_8.jpeg}}

Рис. 7.5. Графічне представлення функції (7.8); межа умовно
задовольняється у всі моменти часу.

Відповідно, перші три члени розв'язання змішаної задачі

(7.9)

\[
\begin{aligned}
(\text{УЧП})\quad & u_t=u_{xx}, && 0<x<1,\ 0<t<\infty,\\
(\text{ГУ})\quad & u(0,t)=0,\quad u_x(1,t)+h\,u(1,t)=0, && 0<t<\infty,\\
(\text{НУ})\quad & u(x,0)=x, && 0\le x\le1.
\end{aligned}
\]

мають вигляд

\[
u(x, t) = = 0.24e^{-4t} \sin(2.02x) + 0.22e^{-24t} \sin(4.9x) - 0.03e^{-48.3t} \sin(7.98x) + \dots
\]

Фіг. На рисунку 7-5 показані графіки цього розв'язку для кількох точок
часу. Рекомендуємо читачеві відповісти на такі питання: Чи відповідає
тип рішення, яке ви отримуєте, з вашими інтуїтивними ідеями? Чи
виконуються граничні умови?

\section{ЗАУВАЖЕННЯ}\label{ux437ux430ux443ux432ux430ux436ux435ux43dux43dux44f-5}

Задача власних значень (7.5) є особливим випадком більш загальної
задачі, яку зазвичай називають проблемою Штурма-Ліувіля:

(7.10)

\[
\begin{aligned}
(\text{ОДУ})\quad & [p(x)y']' - q(x)y + \lambda r(x)y = 0,\qquad 0<x<1,\\
(\text{ГУ})\quad & \alpha_1 y(0) + \beta_1 y'(0)=0,\\
& \alpha_2 y(1) + \beta_2 y'(1)=0.
\end{aligned}
\]

Якщо розв'язати задачу з однорідними граничними умовами методом
розділення змінних, то функції X(x) завжди знаходяться як розв'язки
деякої задачі Штурма--Ліувіля. Очевидно, що задача власних значень (7.5)
є особливим випадком задачі (7.10).

Перелічимо основні властивості задачі (7.10), які були встановлені
Штурмом і Ліувілем за певних природних обмежень.

функції p(x), q(x) та r(x).

\begin{enumerate}
\def\labelenumi{\arabic{enumi}.}
\tightlist
\item
  Існує нескінченна послідовність власних значень, що задовольняють
  нерівності
\end{enumerate}

\[
\lambda_1 < \lambda_2 < \lambda_3 < \ldots < \lambda_n < \ldots \to \infty
\]

. 2. Кожне власне значення \(\lambda_n\) відповідає одній (точній до
константного множника) ненульовій власній функції \(y_n(x)\) . 3. Якщо
\(y_n(x)\) і \(y_m(x)\) дві різні власні функції, що відповідають різним
власним значенням \(\lambda_n\) і \(\lambda_m\) , то вони є
ортогональними з вагою r(x) на сегменті {[}0, 1{]}, тобто вони
задовольняють це відношення

\[
\int_{0}^{1} r(x) y_{n}(x) y_{m}(x) dx = 0.
\]

Більше деталей про проблему Штурма-Ліувіля можна знайти в літературі,
наведеній наприкінці книги. 1) Див., наприклад, Тіхоног А. Н.,
Самарський А. А. Рівняння математичної фізики. --- Москва: Наука, 1972.
--- ред. Ед.

\section{ЗАВДАННЯ}\label{ux437ux430ux432ux434ux430ux43dux43dux44f-6}

\begin{enumerate}
\def\labelenumi{\arabic{enumi}.}
\tightlist
\item
  Розв'язати наступну задачу для рівняння теплопровідності методом
  розкладу (PPP)
\end{enumerate}

\[
u_1 = u_{xx}, \quad 0 < x < 1, \quad 0 < t < \infty,
\]

(GI)

\[
\begin{cases} u(0, t) = 0, \\ u_x(1, t) = 0, \\ u(x, 0) = x, \quad 0 \le x \le 1. \end{cases}
\]

Чи відповідає рішення, яке ви знаходите, з інтонативними
представленнями? Що таке стаціонарне рішення? 2. Знайдіть власні
цінності та функції проблеми Ассолта-Ліувіля (ODY)

\[
X' + \lambda X = 0, \quad 0 < x < 1,
\]

(FY) \(\begin{cases} X(0) = 0, \\ X'(1) = 0. \end{cases}\) Задайте
специфічні вирази для функцій p(x), q(x) та r(x) із загальної задачі
Штурма--Ліувіля для цього випадку. 3. Розв'язати задачу теплопровідності
для планки з термоізольованими межами (PPP)

\[
u_t = u_{xx}, \quad 0 < x < 1, \quad 0 < t < \infty,
\]

(GI)

\[
\begin{cases} u_x (0, t) = 0, \\ u_x (1, t) = 0, \end{cases} \quad 0 < t < \infty,
\]

(ГУ) \(u(x, 0) = x, \quad 0 \le x \le 1.\) Чи відповідає рішення, яке ви
отримуєте, з вашою інтерпретацією цієї проблеми? Що таке стаціонарне
рішення? Яке його фізичне значення? 4. Знайдіть власні цінності та
функції завдань (ОДА)

\[
X'' + \lambda X = 0, \quad 0 < x < 1,
\]

(GI) \(\begin{cases} X'(0) = 0, \\ X'(1) = 0. \end{cases}\)

\section{ПЕРЕТВОРЕННЯ СКЛАДНИХ РІВНЯНЬ У ПРОСТІШУ
ФОРМУ}\label{ux43fux435ux440ux435ux442ux432ux43eux440ux435ux43dux43dux44f-ux441ux43aux43bux430ux434ux43dux438ux445-ux440ux456ux432ux43dux44fux43dux44c-ux443-ux43fux440ux43eux441ux442ux456ux448ux443-ux444ux43eux440ux43cux443}

МЕТА ЛЕКЦІЇ: Показати, як рівняння теплопровідності з додатковими
членами можна звести до стандартного рівняння для нової функції
\(w(x,t)\).

У цій лекції ми розглядаємо два типові приклади:

\[
u_t = \alpha^2 u_{xx} - \beta u,
\]

та

\[
u_t = \alpha^2 u_{xx} - v u_x, \qquad v=\text{const}.
\]

Обидва рівняння можна перетворити до найпростішого рівняння
теплопровідності

\[
w_t = \alpha^2 w_{xx}
\]

за допомогою відповідних замін:

\[
\begin{aligned}
u(x,t)&=e^{-\beta t}w(x,t),\\
u(x,t)&=\exp\!\left[\frac{v}{2\alpha^2}\!\left(x-\frac{vt}{2}\right)\right]w(x,t).
\end{aligned}
\]

Розв'язавши просте рівняння для \(w(x,t)\), ми повертаємося до
початкової функції \(u(x,t)\) за цими формулами. Те саме перетворення,
звісно, треба послідовно застосовувати і до початкових, і до граничних
умов.

Після прочитання двох останніх лекцій читач може створити враження, що
метод розділення змінних застосовується \emph{лише} до рівняння

\[
u_t = \alpha^2 u_{xx}.
\]

Насправді рівняння теплопотоку є найпростішим, але не єдиним
параболічним рівнянням, яке можна розв'язати методом розділення змінних.
Як уже згадувалося, якщо рівняння є лінійним і однорідним, його можна
розв'язати, поділивши змінні. Наприклад, двовимірна задача
теплопровідності в колі описується рівнянням

\[
u_t = \alpha^2 \left[ u_{rr} + \frac{1}{r} u_r + \frac{1}{r^2} u_{\theta\theta} \right],
\]

І хоча це рівняння з \emph{змінними коефіцієнтами}, його можна звести до
трьох звичайних диференціальних рівнянь. У цій лекції ми покажемо
читачеві, що іноді, замість того, щоб безпосередньо розв'язувати
диференціальне рівняння в похідних, краще перетворити його на простіше
рівняння, яке можна розв'язати методом розділення змінних або іншим
методом. Наступний приклад ілюструє цей підхід.

\section{Перетворення проблеми з перенесенням бічної поверхні у проблему
з термічно ізольованою бічною
поверхнею}\label{ux43fux435ux440ux435ux442ux432ux43eux440ux435ux43dux43dux44f-ux43fux440ux43eux431ux43bux435ux43cux438-ux437-ux43fux435ux440ux435ux43dux435ux441ux435ux43dux43dux44fux43c-ux431ux456ux447ux43dux43eux457-ux43fux43eux432ux435ux440ux445ux43dux456-ux443-ux43fux440ux43eux431ux43bux435ux43cux443-ux437-ux442ux435ux440ux43cux456ux447ux43dux43e-ux456ux437ux43eux43bux44cux43eux432ux430ux43dux43eux44e-ux431ux456ux447ux43dux43eux44e-ux43fux43eux432ux435ux440ux445ux43dux435ux44e}

Розглянемо наступне завдання: (8.1)

\[
\begin{aligned}
u_{t} &= \alpha^{2}u_{xx} - \beta u, \quad 0 < x < 1, \quad 0 < t < \infty,\\
u(0, t) &= 0, \quad u(1, t) = 0, \quad 0 < t < \infty,\\
u(x, 0) &= \varphi(x), \quad 0 \le x \le 1.
\end{aligned}
\]

де термін \(-\beta u\) описує потік тепла через бічну поверхню стрижня
(рис. 8.1).
\pandocbounded{\includegraphics[keepaspectratio]{images/_page_59_Figure_7.jpeg}}
Рис. 8.1. Задача теплопровідності, описана рівнянням
\(u_t=\alpha^2u_{xx}-\beta u\); член \(-\beta u\) відповідає втратам
тепла через бічну поверхню стрижня. Наша мета --- ввести нову
температурну \(w\left(x,\ t\right)\) замість \(u\left(x,\ t\right)\) ,
щоб рівняння для \(w\left(x,\ t\right)\) стало простішим за початкове
рівняння

\[
u_t = \alpha^2 u_{xx} - \beta u
\]

.

Це досить загальний метод розв'язання диференціальних рівнянь у
частинних похідних, але вибір відповідного перетворення базується на
інтуїтивних уявленнях про поведінку розв'язків початкового рівняння.
Наприклад, у задачі (8.1) температура u(x, t) у кожній точці \(x_0\)
змінюється внаслідок двох факторів:

\begin{itemize}
\tightlist
\item
  \begin{enumerate}
  \def\labelenumi{\arabic{enumi})}
  \tightlist
  \item
    дифузія тепла вздовж стрижня (через компонентний \(\alpha^2 u_{xx}\)
    ),
  \end{enumerate}
\item
  \begin{enumerate}
  \def\labelenumi{\arabic{enumi})}
  \setcounter{enumi}{1}
  \tightlist
  \item
    теплопередача через бічну поверхню стрижня (через член
    \(-\beta u\)).
  \end{enumerate}
\end{itemize}

Найважливішим елементом нашого міркування є наступне: якщо не було
дифузії вздовж стрижня ( \(\alpha = 0\) ),

тоді температура в кожній точці стрижня експоненційно знижується до нуля
згідно із законом

\[
u(x_0,t)=u(x_0,0)e^{-\beta t}.
\]

. Виходячи з цього спостереження, запитаємо, чи можна розв'язок проблеми
(8.1) представити як утворення двох факторів

\[
(8.2)\quad u(x,t)=e^{-\beta t}w(x,t).
\]

де \(w(x,t)\) описує розподіл температури, що виникає лише через
дифузію. Після підстановки цього виразу у задачу (8.1) отримуємо
перетворену задачу без члена \(-\beta u\). (8.4)

\[
\begin{aligned}
w(x,t)&=\sum_{n=1}^{\infty} a_n e^{-(n\pi\alpha)^2 t}\sin(n\pi x),\\
a_n&=2\int_0^1 \varphi(\xi)\sin(n\pi \xi)\,d\xi.
\end{aligned}
\]

і, отже, розв'язок початкової задачі (8.1) має вигляд

\[
u(x, t) = e^{-\beta t}w(x, t).
\]

Нижче, у розділі «Нотатки», наведено ще один приклад проблеми,
розв'язаної цим методом.

\section{ПРИМІТКИ}\label{ux43fux440ux438ux43cux456ux442ux43aux438}

\begin{enumerate}
\def\labelenumi{\arabic{enumi}.}
\tightlist
\item
  Вирішити проблему
\end{enumerate}

\[
\begin{aligned}
u_t &= u_{xx} - u, && 0 < x < 1,\ 0 < t < \infty, \\
u(0,t) &= 0,\quad u(1,t) = 0, && 0 < t < \infty, \\
u(x,0) &= \sin(\pi x) + 0.5\sin(3\pi x), && 0 \leqslant x \leqslant 1.
\end{aligned}
\]

За методом, описаним вище, ми - a) тимчасово відкинути конвективний член
u у початковому рівнянні. - b) записати розв'язок змішаної задачі без
конвективного члена

\[
u(x, t) = e^{-\pi^2 t} \sin(\pi x) + 0.5e^{-(3\pi)^2 t} \sin(3\pi x),
\]

\begin{enumerate}
\def\labelenumi{\alph{enumi})}
\setcounter{enumi}{2}
\tightlist
\item
  помножимо це розв'язання на конвективний фактор
  \(e^{-\beta t}=e^{-t}\) отримаємо остаточне розв'язання початкової
  задачі
\end{enumerate}

\[
u(x, t) = e^{-t} \left[ e^{-\pi^2 t} \sin(\pi x) + 0.5e^{-(3\pi)^2 t} \sin(3\pi x) \right].
\]

\begin{enumerate}
\def\labelenumi{\arabic{enumi}.}
\setcounter{enumi}{1}
\tightlist
\item
  Рівняння конвективної дифузії
\end{enumerate}

\[
u_t = \alpha^2 u_{xx} - v u_x
\]

(коефіцієнт і є константою) також можна перетворити у вигляд

\[
w_t = \alpha^2 w_{xx}.
\]

Відповідне перетворення має вигляд

\[
u(x, t) = e^{v(x-vt/2)/2\alpha^2}w(x, t).
\]

Це перетворення підкреслює експоненціальний фактор, який повністю
зумовлений рухом середовища. Зверніть увагу, що цей множник описує
хвилю, що поширюється зліва направо зі швидкістю \(\upsilon/2\) . Читач
матиме можливість скористатися цією трансформацією.

\section{ЗАВДАННЯ}\label{ux437ux430ux432ux434ux430ux43dux43dux44f-7}

\begin{enumerate}
\def\labelenumi{\arabic{enumi}.}
\tightlist
\item
  Розв'язати задачу конвективної дифузії
\end{enumerate}

\[
\begin{aligned}
(\text{УЧП})\quad & u_t = u_{xx} - u_x, && 0 < x < 1,\ 0 < t < \infty, \\
(\text{ГУ})\quad & u(0,t) = 0,\quad u(1,t) = 0, && 0 < t < \infty, \\
(\text{НУ})\quad & u(x,0) = e^{x/2}, && 0 \leqslant x \leqslant 1.
\end{aligned}
\]

методом перетворення до простішого рівняння. На що схожий розв'язок?
Можна дати таку інтерпретацію цієї задачі: \(u(x,t)\) --- концентрація
деякої речовини в рухомому середовищі, яке рухається зліва направо зі
швидкістю \(v = 1\); початковий розподіл концентрації задається функцією
\(e^{x/2}\), а на межах концентрація підтримується нульовою. Чи
відповідає ваш розв'язок такій інтерпретації? 2. Розв'язати задачу

\[
\begin{aligned}
(\text{УЧП})\quad & u_t = u_{xx} - u + x, && 0 < x < 1,\ 0 < t < \infty, \\
(\text{ГУ})\quad & u(0,t) = 0,\quad u(1,t) = 1, && 0 < t < \infty, \\
(\text{НУ})\quad & u(x,0) = 0, && 0 \leqslant x \leqslant 1.
\end{aligned}
\]

\begin{itemize}
\tightlist
\item
  а) перетворивши неоднорідні граничні умови на однорідні;
\item
  б) перетворивши вихідне рівняння на нове рівняння без члена \(-u\);
\item
  в) розв'язавши задачу, що виникає після цих перетворень.
\end{itemize}

\begin{enumerate}
\def\labelenumi{\arabic{enumi}.}
\setcounter{enumi}{2}
\tightlist
\item
  Розв'язати задачу
\end{enumerate}

\[
\begin{aligned}
(\text{УЧП})\quad & u_t = u_{xx} - u, && 0 < x < 1,\ 0 < t < \infty, \\
(\text{ГУ})\quad & u(0,t) = 0,\quad u(1,t) = 0, && 0 < t < \infty, \\
(\text{НУ})\quad & u(x,0) = \sin(\pi x), && 0 \leqslant x \leqslant 1.
\end{aligned}
\]

застосовуючи метод розділення змінних безпосередньо до вихідного
рівняння, не виконуючи жодних попередніх перетворень. Чи збігається
отриманий розв'язок з розв'язком, знайденим раніше після заміни

\[
u(x,t) = e^{-t}w(x,t)
\]

?

\section{РОЗВ'ЯЗАННЯ НЕОДНОРІДНИХ УЧП МЕТОДОМ РОЗКЛАДУ ЗА ВЛАСНИМИ
ФУНКЦІЯМИ}\label{ux440ux43eux437ux432ux44fux437ux430ux43dux43dux44f-ux43dux435ux43eux434ux43dux43eux440ux456ux434ux43dux438ux445-ux443ux447ux43f-ux43cux435ux442ux43eux434ux43eux43c-ux440ux43eux437ux43aux43bux430ux434ux443-ux437ux430-ux432ux43bux430ux441ux43dux438ux43cux438-ux444ux443ux43dux43aux446ux456ux44fux43cux438}

ЦІЛЬ ЛЕКЦІЇ: Навчитися розв'язувати змішані задачі вигляду

\[
\begin{aligned}
(\text{УЧП})\quad & u_t = \alpha^2 u_{xx} + f(x,t), && 0 < x < 1,\ 0 < t < \infty, \\
(\text{ГУ})\quad & \alpha_1 u_x(0,t) + \beta_1 u(0,t) = 0, && 0 < t < \infty, \\
& \alpha_2 u_x(1,t) + \beta_2 u(1,t) = 0, && 0 < t < \infty, \\
(\text{НУ})\quad & u(x,0) = \varphi(x), && 0 \leqslant x \leqslant 1.
\end{aligned}
\]

Розв'язок неоднорідного УЧП можна шукати у вигляді ряду

\[
u(x,t) = \sum_{n=1}^{\infty} T_n(t) X_n(x).
\]

де \(X_n(x)\) --- власні функції відповідної однорідної задачі

\[
\begin{aligned}
(\text{УЧП})\quad & u_t = \alpha^2 u_{xx}, && 0 < x < 1,\ 0 < t < \infty, \\
(\text{ГУ})\quad & \alpha_1 u_x(0,t) + \beta_1 u(0,t) = 0, && 0 < t < \infty, \\
& \alpha_2 u_x(1,t) + \beta_2 u(1,t) = 0, && 0 < t < \infty, \\
(\text{НУ})\quad & u(x,0) = \varphi(x), && 0 \leqslant x \leqslant 1.
\end{aligned}
\]

А функції \(T_n(t)\) визначаються як розв'язки деяких звичайних
диференціальних рівнянь. На лекції 6 ми ознайомилися з тим, як
перетворювати неоднорідні граничні умови на однорідні. На жаль, рівняння
залишається гетерогенним, і ми стоїмо перед завданням

\[
\begin{aligned}
u_{t} &= \alpha^{2} u_{xx} + f(x, t), \quad 0 < x < 1, \quad 0 < t < \infty,\\
\alpha_{1} u_{x}(0, t) + \beta_{1} u(0, t) &= 0, \\
\alpha_{2} u_{x}(1, t) + \beta_{2} u(1, t) &= 0, \\
u(x, 0) &= \varphi(x), \quad 0 \le x \le 1.
\end{aligned}
\]

Мета нашої лекції --- опанувати метод розв'язання таких задач. Цей метод
аналогічний методу варіації довільних констант. теорія звичайних
диференціальних рівнянь називається методом розкладу за власними
функціями. Ідея методу надзвичайно проста. Відомо, що розв'язок задачі
(9.1) при f(x, t) = 0 (так звана відповідна однорідна задача)
записується як ряд

\[
u\left(x,\ t\right)=\sum_{n=1}^{\infty}a_{n}e^{-\left(\lambda_{n}\alpha\right)^{2}}/X_{n}\left(x\right),
\]

де \(X_n(x)\) --- власні функції задачі Штурма-Ліувіля, (9.2)

\[
X'' + \lambda^2 X = 0, \\ \alpha_1 X'(0) + \beta_1 X(0) = 0, \\ \alpha_2 X'(1) + \beta_2 X(1) = 0.
\]

Запитаймо себе, чи можна також шукати розв'язок гетерогенної проблеми
(9.1) у вигляді серії власних функцій, але більш загальної форми

\[
u(x, t) = \sum_{n=1}^{\infty} T_n(t) X_n(x).
\]

Існують певні фізичні підстави для такого припущення. Насправді терміни
f(x,t) описують джерело тепла, розташоване всередині стрижня. Це
означає, що залежність рішення від часу не буде виражена показниками

\[
e^{-(\lambda_n\alpha)^2 t}
\]

,

Це стосувалося однорідного завдання, але з деякими іншими функціями.

Щоб показати цей метод у дії і не ускладнювати подання деталями, ми
застосуємо його до розв'язання однієї дуже простої задачі.

\section{Построзв'язання методом демаркації за власними
функціями,}\label{ux43fux43eux441ux442ux440ux43eux437ux432ux44fux437ux430ux43dux43dux44f-ux43cux435ux442ux43eux434ux43eux43c-ux434ux435ux43cux430ux440ux43aux430ux446ux456ux457-ux437ux430-ux432ux43bux430ux441ux43dux438ux43cux438-ux444ux443ux43dux43aux446ux456ux44fux43cux438}

Розглянемо гетерогенну проблему

(9.3) (PY)

\[
\begin{aligned} u_t &= \alpha^2 u_{xx} + f(x, t), & 0 < x < 1, & 0 < t < \infty, \\ u(0, t) &= 0, & 0 < t < \infty, \\ u(1, t) &= 0, & 0 < t < \infty, \\ u(x, 0) &= \varphi(x), & 0 \leqslant x \leqslant 1. \end{aligned}
\]

Процес розв'язання задачі подається у вигляді послідовності наступних
кроків. КРОК 1. Основна ідея методу полягає в розкладанні густини
джерела f(x, t) на ряд власними функціями

\[
f(x, t) = f_1(t) X_1(x) + f_2(t) X_2(x) + \dots + f_n(t) X_n(x) + \dots
\]

а також визначити реакції системи \(u_n\left(x,\,t\right)\) на вплив
кожного компонента \(f_n\left(t\right)X_n\left(x\right)\) . Якщо знайти
всі відповіді, рішення проблеми виглядатиме так:

\[
u(x, t) = \sum_{n=1}^{\infty} u_n(x, t).
\]

Головна складність цього методу --- розклад густини джерела на
компоненти \(f_n(t)\,X_n(x)\) . Виявляється, що \(X_n(x)\) факторами в
цій задачі є власними векторами системи Штурма--Ліувіля, які виникають
при розв'язанні відповідної однорідної задачі методом розділення
змінних, а саме

\[
\begin{array}{c} u_t = \alpha^2 u_{xx} \quad \text{(отметим, что } f(x, \ t) = 0), \\ u(0, \ t) = 0, \\ u(1, \ t) = 0, \\ u(x, \ 0) = \varphi(x). \end{array}
\]

Завдання Штурма-Ліувіля в цьому випадку виглядає так:

\[
X'' + \lambda^2 X = 0,
\]

\(X(0) = 0,\) \(X(1) = 0,\) а його розв'язки є функціями

\[
X_n(x) = \sin(n\pi x), \quad n = 1, 2, 3 \dots
\]

Отже, розклад на густину джерела представлений як ряд

\[
(9.5) \ f(x,t) = f_1(t) \sin(\pi x) + f_2(t) \sin(2\pi x) + \dots + f_n(t) \sin(n\pi x) + \dots
\]

\(\mathcal{U}\) нарешті, щоб знайти функції \(f_n(t)\) , помножте обидві
сторони останнього відношення на \(\sin(m\pi x)\) і інтегруйте від 0 до
1 за координатою x.

У результаті ми отримуємо

\[
\int_{0}^{1} f(x, t) \sin(m\pi x)\,dx = \sum_{n=1}^{\infty} f_{n}(t) \int_{0}^{1} \sin(m\pi x) \sin(n\pi x)\,dx = \frac{1}{2} f_{m}(t)
\]

або (після заміни t на p)

(9.6)

\[
f_{n}(t) = 2 \int_{0}^{1} f(x, t) \sin(n\pi x)\,dx.
\]

Останнє співвідношення встановлює співвідношення між коефіцієнтами
\(f_n(t)\) і часткою джерела тепла f(x, t).

КРОК 2. (Пошук \(u_n(x, t) = X_n(x) T_n(t)\) реакції на вхідний вплив
\(f_n(t) X_n(x)\) .)

Замінимо густину джерела на його розклад у рядок

\[f(x, t) = \sum_{n=1}^{\infty} f_n(t) \sin(n\pi x)\]

і намагатися знайти індивідуальні відповіді, тобто функції,
\(T_{\mathbf{n}}(t)\) у

\[u(x, t) = \sum_{n=1}^{\infty} T_n(t) \sin(n\pi x).\]

Оскільки розв'язок задачі має вигляд

\[u(x, t) = \sum_{n=1}^{\infty} T_n(t) \sin(n\pi x),\]

потім підставимо цей вираз у початкову задачу

\[
u_{t} = \alpha^{2} u_{xx} + \sum_{n=1}^{\infty} f_{n}(t) \sin(n\pi x),
\quad u(0, t) = 0, \quad u(1, t) = 0, \quad u(x, 0) = \varphi(x),
\]

У результаті ми отримаємо

\[
\sum_{n=1}^{\infty} T'_n(t) \sin(n\pi x) = -\alpha^2 \sum_{n=1}^{\infty} (n\pi)^2 T_n(t) \sin(n\pi x) + + \sum_{n=1}^{\infty} f_n(t) \sin(\pi n\lambda),
\]

(9.7)

\[
\sum_{n=1}^{\infty} T_n(t) \sin 0 = 0
\]

(задовольняється ідентично),
\(\sum_{n=1}^{\infty} T_n(t) \sin (n\pi) = 0\) (задовольняється
ідентично), \(\sum_{n=1}^{\infty} T_n(0) \sin (n\pi x) = \phi(x)\) .

Переписування рівняння та початкової умови

(PPP)

\[
\sum_{n=1}^{\infty} [T'_n + (n\pi\alpha)^2 T_n - f_n(t)] \sin(n\pi x) = 0,
\]

(НУ)

\[
\sum_{n=1}^{\infty} T_n(0) \sin(n\pi x) = \varphi(x),
\]

можна побачити, що функції \(T_n(\ell)\) є розв'язками задачі Коші

(9.8)

\[
\begin{aligned}
T_n' + (n\pi\alpha)^2 T_n &= f_n(t), \\
T_n(0) &= 2 \int_0^1 \varphi(\xi)\sin(n\pi \xi)\,d\xi = a_n.
\end{aligned}
\]

Усі ці TAC легко розв'язати, а їхні рішення записуються у вигляді

(9.9)

\[
T_n(t) = a_n e^{-(n\pi\alpha)^2 t} + \int_0^t e^{-(n\pi\alpha)^2 (t-\tau)} f_n(\tau)\,d\tau.
\]

З того місця, де легко отримати повне розв'язання проблеми (9.3):

(9.10)

\[
\begin{aligned}
u(x, t) &= \sum_{n=1}^{\infty} T_n(t)\sin(n\pi x) \\
&= \sum_{n=1}^{\infty} a_n e^{-(n\pi\alpha)^2 t}\sin(n\pi x) \\
&\quad + \sum_{n=1}^{\infty} \sin(n\pi x)\int_0^t e^{-(n\pi\alpha)^2 (t-\tau)} f_n(\tau)\,d\tau.
\end{aligned}
\]

Перехідна частина (обумовлена стаціонарною частиною (умовлена початковою
умовою) правою частиною f(x, t)).

Розв'язок (9.10) показує, що температурна реакція складається з двох
частин: перша частина зумовлена початковими умовами, а друга ---
густиною джерела f(x,t). Зверніть увагу, що термін «стаціонарна частина»
не є найкращим описом другої частини температурної реакції, оскільки він
не означає, що ця частина не залежить від часу (стаціонарний стан може
бути періодичним, якщо f(x, t) є періодичною функцією часу t).

Проблема повністю вирішена. Перш ніж перейти до наступної теми, давайте
покажемо, як використовувати описаний вище метод для розв'язання

Конкретне завдання.

\section{Розв'язання конкретної задачі методом декомпозиції за власними
функціями}\label{ux440ux43eux437ux432ux44fux437ux430ux43dux43dux44f-ux43aux43eux43dux43aux440ux435ux442ux43dux43eux457-ux437ux430ux434ux430ux447ux456-ux43cux435ux442ux43eux434ux43eux43c-ux434ux435ux43aux43eux43cux43fux43eux437ux438ux446ux456ux457-ux437ux430-ux432ux43bux430ux441ux43dux438ux43cux438-ux444ux443ux43dux43aux446ux456ux44fux43cux438}

Розглянемо просту задачу

(9.11)

\[
\begin{aligned}
(\text{УЧП})\quad & u_t = \alpha^2 u_{xx} + \sin(3\pi x), && 0 < x < 1,\quad 0 < t < \infty, \\
(\text{ГУ})\quad & u(0,t)=0,\quad u(1,t)=0, && 0 < t < \infty, \\
(\text{НУ})\quad & u(x,0)=\sin(\pi x), && 0 \le x \le 1.
\end{aligned}
\]

Наша мета --- обчислити коефіцієнти \(T_n(t)\) при розкладі

\[
u\left(x,\,t\right)=\sum_{n=1}^{\infty}T_n(t)\sin(n\pi x)
\]

(Власні функції цієї задачі \(X_n(x)\) збігатися з власними функціями
попереднього). Якщо підставити цю декомпозицію на початкову задачу,
отримаємо наступну задачу Коші для визначення функцій \(T_n(t)\) :

\[
\begin{aligned}
T_n' + (n\pi\alpha)^2 T_n &= f_n(t), \\
f_n(t) &= 2 \int_0^1 \sin(3\pi x)\sin(n\pi x)\,dx
= \begin{cases}
1, & n = 3, \\
0, & n \ne 3,
\end{cases} \\
T_n(0) &= 2 \int_0^1 \sin(\pi \xi)\sin(n\pi \xi)\,d\xi
= \begin{cases}
1, & n = 1, \\
0, & n \ne 1.
\end{cases}
\end{aligned}
\]

Перепишем эти уравнения последовательно для n=1, 2\ldots{}

\[
\begin{aligned}
(n=1)\quad & T_1' + (\pi\alpha)^2 T_1 = 0, && T_1(0) = 1, && T_1(t) = e^{-(\pi\alpha)^2 t}, \\
(n=2)\quad & T_2' + (2\pi\alpha)^2 T_2 = 0, && T_2(0) = 0, && T_2(t) = 0, \\
(n=3)\quad & T_3' + (3\pi\alpha)^2 T_3 = 4, && T_3(0) = 0, && T_3(t) = \frac{1}{(3\pi\alpha)^2}\left[1 - e^{-(3\pi\alpha)^2 t}\right], \\
(n\geqslant 4)\quad & T_n' + (n\pi\alpha)^2 T_n = 0, && T_n(0) = 0, && T_n(t) = 0.
\end{aligned}
\]

Отже, розв'язок задачі можна записати у вигляді

\[u(x, t) = e^{-(\pi\alpha)^2 t} \sin(\pi x) + \frac{1}{(3\pi\alpha)^2} [1 - e^{-(3\pi\alpha)^2 t}] \sin(3\pi x).\]

(9.12)

\section{ЗАУВАЖЕННЯ}\label{ux437ux430ux443ux432ux430ux436ux435ux43dux43dux44f-6}

\begin{enumerate}
\def\labelenumi{\arabic{enumi}.}
\item
  Метод розкладу власних функцій є одним із найпотужніших методів
  розв'язання неоднорідних диференціальних рівнянь у похідних похідних.
  Пізніше, коли ми ознайомимося з інтегральними перетвореннями,
  побачимо, що існують й інші методи розв'язання задач такого типу.
\item
  Власні функції \(X_n(x)\) , за допомогою якої виконується розклад,
  залежать від задачі, яку потрібно розв'язати, точніше від рівнянь і
  граничних умов. Читача просять розв'язати задачу 4 наприкінці лекції,
  щоб перевірити, чи дійсно він навчився знаходити власні функції
  \(X_n(x)\) .
\item
  Якщо читач пам'ятає основні факти з теорії звичайних диференціальних
  рівнянь, то він знає, що розв'язки, що відповідають неоднорідним
  членам вигляду
\end{enumerate}

\[
P_n(x) e^{\alpha x} \cdot
\begin{cases} \sin{(\beta x)}, \\ \cos{(\beta x)}, \end{cases}
\]

можна знайти методом невизначених коефіцієнтів. Те саме і тут. Задача
(9.11) може бути розв'язана таким самим способом.

\section{ЗАДАЧІ}\label{ux437ux430ux434ux430ux447ux456}

\section{1. Решение
задачи}\label{ux440ux435ux448ux435ux43dux438ux435-ux437ux430ux434ux430ux447ux438}

\[
\begin{aligned}
(\text{УЧП}) \quad & u_t = u_{xx} + \sin(3\pi x), && 0 < x < 1, \\
(\Gamma \text{У}) \quad & u(0,t) = 0,\quad u(1,t) = 0, && 0 < t < \infty, \\
(\text{НУ}) \quad & u(x,0) = \sin(\pi x), && 0 \leqslant x \leqslant 1.
\end{aligned}
\]

Отже, розв'язок задається формулою (9.12). Чи відповідає це рішення
вашим інтуїтивним уявленням? Що саме означає отриманий розв'язок? 2.
Решите задачу

\[
\begin{aligned}
(\text{УЧП})\quad & u_t = u_{xx} + \sin(\pi x) + \sin(2\pi x), && 0 < x < 1,\ 0 < t < \infty,\\
(\text{ГУ})\quad & u(0,t)=0,\quad u(1,t)=0, && 0 < t < \infty,\\
(\text{НУ})\quad & u(x,0)=0, && 0 \le x \le 1.
\end{aligned}
\]

\section{3. Розв'яжіть
задачу}\label{ux440ux43eux437ux432ux44fux436ux456ux442ux44c-ux437ux430ux434ux430ux447ux443}

(PPP)

\[
u_t = u_{xx} + \sin(\pi x), \quad 0 < x < 1, \quad 0 < t < \infty,
\]

(GI)

\[
\begin{cases} u(0, t) = 0, \\ u(1, t) = 0, \\ u(x, 0) = 1, \quad 0 < x < 1, \end{cases}
\]

методом розкладу за власними функціями. 4. Найдите решение задачи

\[
\begin{aligned}
(\text{УЧП})\quad & u_t = u_{xx} + \sin(\lambda_1 x), && 0 < x < 1,\ 0 < t < \infty, \\
(\text{ГУ})\quad & u(0,t) = 0,\quad u_x(1,t) + u(1,t) = 0, && 0 < t < \infty, \\
(\text{НУ})\quad & u(x,0) = 0, && 0 \leqslant x \leqslant 1.
\end{aligned}
\]

методом разложения по собственным функциям, если \(\lambda_1\) ---
первый корень уравнения \(\operatorname{tg}\lambda = -\lambda\). Каковы
собственные функции этой задачи?

\section{5. Решите
задачу}\label{ux440ux435ux448ux438ux442ux435-ux437ux430ux434ux430ux447ux443}

\[
\begin{aligned}
(\text{УЧП})\quad & u_t = u_{xx}, && 0 < x < 1,\ 0 < t < \infty, \\
(\text{ГУ})\quad & u(0,t) = 0,\quad u(1,t) = \cos t, && 0 < t < \infty, \\
(\text{НУ})\quad & u(x,0) = \varphi(x), && 0 \leqslant x \leqslant 1.
\end{aligned}
\]

Спочатку перетворивши граничні умови в нуль. Отримана задача
розв'язується методом розкладу за власними функціями.

\section{ІНТЕГРАЛЬНІ ПЕРЕТВОРЕННЯ (СИНУСОЇДАЛЬНІ ТА КОСИНУСОВІ
ПЕРЕТВОРЕННЯ)}\label{ux456ux43dux442ux435ux433ux440ux430ux43bux44cux43dux456-ux43fux435ux440ux435ux442ux432ux43eux440ux435ux43dux43dux44f-ux441ux438ux43dux443ux441ux43eux457ux434ux430ux43bux44cux43dux456-ux442ux430-ux43aux43eux441ux438ux43dux443ux441ux43eux432ux456-ux43fux435ux440ux435ux442ux432ux43eux440ux435ux43dux43dux44f}

МЕТА ЛЕКЦІЇ: Ввести концепцію інтегрального перетворення та показати, як
використання інтегральних перетворень дозволяє перейти від CPP з n
змінними до нового рівняння з n---1 змінними. Лекція показує, що
інтегральні перетворення можна інтерпретувати як розклад вхідної дії на
елементарні (частотний розклад), визначення реакції системи на
елементарні впливи та подальше підсумовування всіх отриманих
результатів. У інтегральних перетвореннях операція диференціювання
переходить у операцію множення, і, відповідно, деякі часткові похідні
переходять у алгебраїчні вирази. Для розв'язання задач дифузії на
необмежених інтервалах вводять синусоїдальні та косинусні перетворення.
Розв'язок, отриманий у цьому випадку, цікавий тим, що містить функцію,
яка є додатковою до ймовірнісного інтеграла. Інтегральне перетворення
зазвичай називається перетворенням, яке кожній функції f(t) призначає
нову функцію F(s) відповідно до формули

\[
F(s) = \int_{A}^{B} K(s, t) f(t) dt.
\]

Зверніть увагу, що початкова функція залежить від t, а отримана функція
--- від s. Функція K(s,t) називається ядром перетворення. Це головний
елемент, який відрізняє одну трансформацію від іншої. Зазвичай ядро
обирається так, щоб перетворення мало певні задані властивості. Межі
інтегрування також залежать від типу трансформації. Голосною, що лежить
в основі інтегральних перетворень, є здатність позбутися часткових
похідних однієї з незалежних змінних. Це означає, що перетворене
рівняння матиме на одну змінну менше, ніж початкове. Наприклад, якщо
застосувати інтегральне перетворення до рівняння Щоб позбутися похідної
часу, отримуємо звичайне диференціальне рівняння з незалежною змінною X.

\[
u_{xx} + u_{cy} + u_{zz} = 0
\]

застосувати перетворення Фур'є до змінної x, тоді похідна \(u_{xx}\)
буде усунена, і в новому рівнянні залишаться лише дві незалежні змінні y
та z. Тоді отримуємо звичайне диференціальне рівняння з однією
незалежною змінною. Іншими словами, інтегральні перетворення спрощують
початкове завдання. Потім розв'язується перетворення задачі, і,
застосовуючи обернене перетворення до її розв'язку, отримується
розв'язок початкового перетворення. Схема описаного алгоритму показана
на рис. 10.1
\pandocbounded{\includegraphics[keepaspectratio]{images/_page_72_Figure_5.jpeg}}
Рис. 10.1, Загальна схема застосування інтегрованих перетворень. Фіг. На
рисунку 10.1 видно, що кожне перетворення прямого інтеграла пов'язане з
оберненим перетворенням, яке відновлює початкову функцію з
трансформованої. Прямі та обернені перетворення зазвичай називають парою
перетворень. Таблиця 1. Рисунок 10.1 показує кілька пар перетворень, які
ми будемо використовувати при розв'язуванні рівнянь із факторами. Слід
зазначити, що в цих перетвореннях використовуються два типи символів.
Наприклад, у випадку перетворення Лапласа позначення \(\mathcal{L}[f]\)
вказує, що функція f перетворена, тоді як позначення F(s) означає, що
перетворення породжує функцію, яка залежить від s. У цій лекції ми не
прагнемо вивчати всі пари перетворень. Ми зосередимося лише на
перетвореннях Фур'є синуса та косинуса (пари 1 і 2). Трохи пізніше ми
вивчаємо іншу з цих пар. Ознайомляючись із трансформаціями, ми
дізнаємося про їхні переваги та недоліки, дізнаємося, як Таблиця 10.1.
Деякі пари інтегральних перетворень 1. Синус-перетворення Фур'є:

\[
\mathcal{F}_s[f](\omega) = \frac{2}{\pi}\int_0^\infty f(t)\sin(\omega t)\,dt,
\]

\[
\mathcal{F}_s^{-1}[F](t) = \int_0^\infty F(\omega)\sin(\omega t)\,d\omega.
\]

\begin{enumerate}
\def\labelenumi{\arabic{enumi}.}
\setcounter{enumi}{1}
\tightlist
\item
  Косинус-перетворення Фур'є:
\end{enumerate}

\[
\mathcal{F}_c[f](\omega) = \frac{2}{\pi}\int_0^\infty f(t)\cos(\omega t)\,dt,
\]

\[
\mathcal{F}_c^{-1}[F](t) = \int_0^\infty F(\omega)\cos(\omega t)\,d\omega.
\]

\begin{enumerate}
\def\labelenumi{\arabic{enumi}.}
\setcounter{enumi}{2}
\tightlist
\item
  Перетворення Фур'є:
\end{enumerate}

\[
\mathcal{F}[f](\omega) = \frac{1}{\sqrt{2\pi}}\int_{-\infty}^{\infty} f(x)e^{-i\omega x}\,dx,
\]

\[
\mathcal{F}^{-1}[F](x) = \frac{1}{\sqrt{2\pi}}\int_{-\infty}^{\infty} F(\omega)e^{i\omega x}\,d\omega.
\]

\begin{enumerate}
\def\labelenumi{\arabic{enumi}.}
\setcounter{enumi}{3}
\tightlist
\item
  Скінченне синус-перетворення:
\end{enumerate}

\[
S_n = \frac{2}{L}\int_0^L f(x)\sin\!\left(\frac{n\pi x}{L}\right) dx,
\]

\[
f(x) = \sum_{n=1}^{\infty} S_n \sin\!\left(\frac{n\pi x}{L}\right).
\]

\begin{enumerate}
\def\labelenumi{\arabic{enumi}.}
\setcounter{enumi}{4}
\tightlist
\item
  Скінченне косинус-перетворення:
\end{enumerate}

\[
C_n = \frac{2}{L}\int_0^L f(x)\cos\!\left(\frac{n\pi x}{L}\right) dx,
\]

\[
f(x) = \frac{C_0}{2} + \sum_{n=1}^{\infty} C_n \cos\!\left(\frac{n\pi x}{L}\right).
\]

\begin{enumerate}
\def\labelenumi{\arabic{enumi}.}
\setcounter{enumi}{5}
\tightlist
\item
  Перетворення Лапласа:
\end{enumerate}

\[
\mathcal{L}[f](s) = \int_0^\infty e^{-st}f(t)\,dt,
\]

\[
\mathcal{L}^{-1}[F](t) = \frac{1}{2\pi i}\int_{\gamma-i\infty}^{\gamma+i\infty} e^{st}F(s)\,ds.
\]

\begin{enumerate}
\def\labelenumi{\arabic{enumi}.}
\setcounter{enumi}{6}
\tightlist
\item
  Перетворення Ханкеля:
\end{enumerate}

\[
H\{f\}(\xi) = F_n(\xi) = \int_0^\infty r J_n(\xi r) f(r)\,dr,
\]

\[
H^{-1}[F_n](r) = \int_0^\infty \xi J_n(\xi r) F_n(\xi)\,d\xi.
\]

Однак перед тим, як перейти до вивчення інтегральних перетворень, буде
корисно зрозуміти, що зазвичай називають спектром функції (або
спектральним розкладом функції).

\section{Спектральна
функція}\label{ux441ux43fux435ux43aux442ux440ux430ux43bux44cux43dux430-ux444ux443ux43dux43aux446ux456ux44f}

Інтегральні перетворення та спектр функції тісно пов'язані. Насправді
інтегральне перетворення можна інтерпретувати як розклад функції на
певний спектр компонент. Цей спектр залежить від обраного перетворення,
але кожна функція має певний спектр відносно цього перетворення Р и с
\pandocbounded{\includegraphics[keepaspectratio]{images/_page_74_Picture_7.jpeg}}
. 10 2 Розклад періодичної функції на ряд синусами та косинусами;
Послідовність квадратних імпульсів апроксимується синусами та косинусами
Наприклад, розглянемо розклад періодичної функції на ряд синусами та
косинусами (ряд Фур'є) 1:

\[
f(x) = \sum_{n=0}^{\infty} \left[ A_n \cos(nx) + B_n \sin(nx) \right].
\]

(див. рис. 10.2). Тут коефіцієнти \(A_n\) і \(B_n\) позначають внесок
\(\cos(nx)\) і \(\sin(nx)\) у функцію f(x), тоді як величина

\[
\sqrt{A_n^2 + B_n^2}
\]

1) Фур'є-ряди детальніше розглядаються у лекції 11, (називається
спектром функції) показує частку амплітуди функції f(x), що припадає на
частоту n. Наприклад, якщо функція f(x) --- це просто сума синуси та
косинуси

\[
f(x) = 1 + \sin x + \frac{1}{5}\sin(3x) + \cos x + \frac{1}{2}\cos(2x) + \frac{1}{4}\cos(4x).
\]

тоді його спектр (дискретний) виглядатиме так, як показано на рис. 10.3.
Зчитуючи значення значення \(\sqrt{A_n^2 + B_n^2}\) з графіка, можна
визначити амплітудне значення компоненти функції f(x), яка має частоту
n.
\pandocbounded{\includegraphics[keepaspectratio]{images/_page_75_Figure_8.jpeg}}
Рис. 10 3. Дискретний спектр функції f(r). Періодичні функції
розкладаються на нескінченні ряди (вони мають дискретний спектр), тоді
як неперіодичні функції повинні мати неперервний спектр. Якщо функція
визначена лише на скінченному інтервалі, її можна періодично розширювати
на всю вісь і розкладати на ряд Фур'є, який представлятиме нашу функцію
в межах визначення області. І хоча неперіодичну функцію не можна
розкласти на нескінченний ряд синусами та косинусами, можна спробувати
записати неперервний аналог ряду Фур'є:

\[
f(x) = \int_{-\infty}^{+\infty} \left[C(\omega)\cos(\omega x) + S(\omega)\sin(\omega x)\right]\,d\omega.
\]

де функції \(S(\omega)\) і \(C(\omega)\) слугують мірою спектральної
густини синуса і косинуса у функції \(f(x)\), а величина

\[
\sqrt{S^2(\omega) + C^2(\omega)}
\]

слугує мірою загального внеску частоти \(\omega\) у функцію \(f(x)\) і
називається неперервним спектром функції \(f(x)\). Виходячи з наведеного
вище інтуїтивного визначення спектра функції, давайте проаналізуємо
інтегральні перетворення «гвинт за гвинтом». Спершу запишемо кілька
властивостей цих перетворень, які ми будемо широко використовувати.

\section{Синусне та косинусне перетворення
похідних}\label{ux441ux438ux43dux443ux441ux43dux435-ux442ux430-ux43aux43eux441ux438ux43dux443ux441ux43dux435-ux43fux435ux440ux435ux442ux432ux43eux440ux435ux43dux43dux44f-ux43fux43eux445ux456ux434ux43dux438ux445}

Для подальших застосувань нам знадобляться такі властивості:

\[
\mathcal{F}_s[f'] = -\omega \mathcal{F}_c[f],
\qquad
\mathcal{F}_s[f''] = \frac{2}{\pi}\omega f(0) - \omega^2 \mathcal{F}_s[f],
\]

\[
\mathcal{F}_c[f'] = -\frac{2}{\pi}f(0) + \omega \mathcal{F}_s[f],
\qquad
\mathcal{F}_c[f''] = -\frac{2}{\pi}f_c(0) - \omega^2 \mathcal{F}_c[f].
\]

Деякі інші властивості прямих і обернених синусних та косинусних
перетворень наведено в таблиці наприкінці книги. Тепер покажемо, як
синусне перетворення можна використати для розв'язання важливої змішаної
задачі.

\section{Розв'язання задачі дифузії на
напівпрямій}\label{ux440ux43eux437ux432ux44fux437ux430ux43dux43dux44f-ux437ux430ux434ux430ux447ux456-ux434ux438ux444ux443ux437ux456ux457-ux43dux430-ux43dux430ux43fux456ux432ux43fux440ux44fux43cux456ux439}

Задача, яку ми хочемо розв'язати, формулюється так: (PPP)

\[
\begin{aligned}
(\text{УЧП})\quad & u_t = \alpha^2 u_{xx}, && 0 < x < \infty,\quad 0 < t < \infty, \\
(\text{ГУ})\quad & u(0,t) = A, && 0 < t < \infty, \\
(\text{НУ})\quad & u(x,0) = 0, && 0 \le x < \infty.
\end{aligned}
\]

\pandocbounded{\includegraphics[keepaspectratio]{images/_page_76_Figure_14.jpeg}}
Рис. 10.4. Задача дифузії в напівнескінченній області: \(u(0,t)=A\),
\(u(x,0)=0\). КРОК 1. Застосуємо синусне перетворення Фур'є за змінною
\(x\) до обох частин рівняння. Тоді дістанемо

\[
\mathcal{F}_s[u_t] = \alpha^2 \mathcal{F}_s[u_{xx}].
\]

Оскільки перетворення виконується за змінною \(x\), змінна \(t\)
поводиться як параметр. Тому

\[
\mathcal{F}_s[u_t]
= \frac{2}{\pi}\int_{0}^{\infty} u_t(x,t)\sin(\omega x)\,dx
= \frac{d}{dt}\mathcal{F}_s[u]
= \frac{dU}{dt},
\]

де \(U=U(\omega,t)=\mathcal{F}_s[u]\). Для другого члена, використовуючи
властивості синусного перетворення та граничну умову \(u(0,t)=A\), маємо

\[
\mathcal{F}_s[u_{xx}] = \frac{2}{\pi}\omega u(0,t) - \omega^2 \mathcal{F}_s[u]
= \frac{2A\omega}{\pi} - \omega^2 U(\omega,t).
\]

Підставляючи ці вирази у рівняння \(u_t=\alpha^2 u_{xx}\), отримуємо
звичайне диференціальне рівняння

\[
\frac{dU}{dt} = \alpha^2\left[-\omega^2 U + \frac{2A\omega}{\pi}\right].
\]

Початкова умова для \(U\) випливає з умови \(u(x,0)=0\):

\[
\mathcal{F}_s[u(x,0)] = U(\omega,0) = 0.
\]

Отже, замість початкової змішаної задачі ми отримали задачу Коші для
звичайного диференціального рівняння: (10.2)

\[
\begin{aligned}
\frac{dU}{dt} + \omega^2 \alpha^2 U &= \frac{2A\omega \alpha^2}{\pi}, \\
U(\omega,0) &= 0.
\end{aligned}
\]

КРОК 2. Щоб розв'язати отриману задачу Коші, можна використати будь-який
елементарний метод розв'язання звичайних диференціальних рівнянь. У
результаті маємо

\[
U(\omega,t)=\frac{2A}{\pi\omega}\left(1-e^{-\omega^2\alpha^2 t}\right).
\]

Отже, перед нами результат застосування синусного перетворення до
шуканої функції \(u(x,t)\). Останній крок полягає в застосуванні
оберненого перетворення:

\[
u(x,t)=\mathcal{F}_s^{-1}[U].
\]

КРОК 3. Скориставшись таблицями перетворень, отримуємо

\[
u(x,t)=A\operatorname{erfc}\!\left(\frac{x}{2\alpha\sqrt{t}}\right),
\]

де \(\operatorname{erfc}(x)\) --- доповнююча функція помилки:

\[
\operatorname{erfc}(x)=\frac{2}{\sqrt{\pi}}\int_x^\infty e^{-s^2}\,ds.
\]

Графік цієї функції показано на рис. 10.5. Вона не виражається через
елементарні функції, але широко табульована в довідниках.
\pandocbounded{\includegraphics[keepaspectratio]{images/_page_79_Figure_4.jpeg}}
Фіг. 10.5. Графіки функцій \(\operatorname{erf}(x)\) та
\(\operatorname{erfc}(x)\).

\section{Інтерпретація
рішення}\label{ux456ux43dux442ux435ux440ux43fux440ux435ux442ux430ux446ux456ux44f-ux440ux456ux448ux435ux43dux43dux44f}

Розв'язок

\[
u(x,t)=A\operatorname{erfc}\!\left(\frac{x}{2\alpha\sqrt{t}}\right)
\]

дає багато їжі для роздумів. На рисунку 10.6 показані графіки
ймовірнісної комплементарної функції на різних значеннях часу t.
\pandocbounded{\includegraphics[keepaspectratio]{images/_page_79_Figure_10.jpeg}}
Рис. 10.6. Розподіл температури в напівнескінченному стрижні, на кінці
якого підтримується постійна температура. Вони відрізняються один від
одного лише характеристичною шкалою вздовж осі x.

\section{ЗАВДАННЯ}\label{ux437ux430ux432ux434ux430ux43dux43dux44f-8}

\begin{itemize}
\tightlist
\item
  \begin{enumerate}
  \def\labelenumi{\arabic{enumi}.}
  \tightlist
  \item
    Довести ідентичності (10.1).
  \end{enumerate}
\item
  \begin{enumerate}
  \def\labelenumi{\arabic{enumi}.}
  \setcounter{enumi}{1}
  \tightlist
  \item
    Розв'яжіть задачу Коші (10.2).
  \end{enumerate}
\item
  \begin{enumerate}
  \def\labelenumi{\arabic{enumi}.}
  \setcounter{enumi}{2}
  \tightlist
  \item
    Використайте синусоїдальне або косинусне перетворення для
    розв'язання задачі
  \end{enumerate}
\end{itemize}

\[
\begin{aligned}
(\text{УЧП})\quad & u_t = \alpha^2 u_{xx}, && 0 < x < \infty,\ 0 < t < \infty, \\
(\text{ГУ})\quad & u_x(0,t) = 0, && 0 < t < \infty, \\
(\text{НУ})\quad & u(x,0) = H(1-x), && 0 \leqslant x < \infty.
\end{aligned}
\]

де \(H(x)\) --- функція Хевісайду:

\[
H(x)=
\begin{cases}
0, & x<0,\\
1, & x\ge 0.
\end{cases}
\]

Вхідна до початкової умови функція зображена на рисунку.
\pandocbounded{\includegraphics[keepaspectratio]{images/_page_80_Picture_7.jpeg}}
Як виглядатиме графік розв'язку в різні моменти часу? 5. Знайдіть
прообраз константи при синусному перетворенні, тобто знайдіть \(f(x)\),
якщо \(F(\omega)=c=\mathrm{const}\).

\section{РЯДИ І ПЕРЕТВОРЕННЯ
ФУР'Є}\label{ux440ux44fux434ux438-ux456-ux43fux435ux440ux435ux442ux432ux43eux440ux435ux43dux43dux44f-ux444ux443ux440ux454}

МЕТА ЛЕКЦІЇ: Ввести поняття ряду Фур'є і показати, як періодичну функцію
\(f(x)\) можна подати рядом за синусами і косинусами. Якщо ж функція
\(f(x)\) є неперіодичною і визначена на всій дійсній осі, то ряд Фур'є
природно переходить у перетворення Фур'є. Важлива роль рядів Фур'є в
теорії рівнянь з частинними похідними пов'язана з тим, що періодичні
функції, а також функції, визначені на скінченному відрізку, можна
розкладати в ряди за синусами і косинусами. Це часто дозволяє звести
початкову задачу до задач для окремих гармонік. Наприклад, пилкоподібну
функцію

\[
f(x)=x, \qquad -L < x < L, \qquad f(x+2L)=f(x),
\]

можна подати рядом Фур'є (11.1)

\[
f(x) = \frac{a_0}{2} + \sum_{n=1}^{\infty} \left[a_n \cos\!\left(\frac{n\pi x}{L}\right) + b_n \sin\!\left(\frac{n\pi x}{L}\right)\right].
\]

Коефіцієнти Фур'є \(a_n\) і \(b_n\) знаходяться за формулами Ейлера. Для
цієї функції маємо

\[
a_n = \frac{1}{L}\int_{-L}^{L} f(x)\cos\!\left(\frac{n\pi x}{L}\right)\,dx = 0, \qquad n=0,1,2,\dots,
\]

\[
b_n = \frac{1}{L}\int_{-L}^{L} f(x)\sin\!\left(\frac{n\pi x}{L}\right)\,dx
= -(-1)^n \frac{2L}{n\pi}, \qquad n=1,2,\dots.
\]

\pandocbounded{\includegraphics[keepaspectratio]{images/_page_82_Figure_7.jpeg}}
Рис. 11.1. Розклад пилкоподібної функції в ряд Фур'є. Тому розклад цієї
функції набирає вигляду (11.3)

\[
f(x) = \frac{2L}{\pi}\left[\sin\!\left(\frac{\pi x}{L}\right) - \frac{1}{2}\sin\!\left(\frac{2\pi x}{L}\right) + \frac{1}{3}\sin\!\left(\frac{3\pi x}{L}\right) - \dots \right].
\]

Кожен доданок цього ряду називається гармонікою, а всі гармоніки мають
частоти, кратні основній частоті. У теорії рядів Фур'є важливу роль
відіграє теорема Діріхле: якщо обмежена періодична функція на кожному
періоді має лише скінченне число максимумів, мінімумів і точок розриву,
то її ряд Фур'є збігається до самої функції в точках неперервності і до
півсуми односторонніх границь у точках розриву. Перед переходом до
перетворення Фур'є корисно ввести поняття дискретного частотного спектра
періодичної функції. Якщо

\[
c_n = \sqrt{a_n^2 + b_n^2}, \qquad n=0,1,2,\dots,
\]

то числа \(c_n\) показують внесок відповідних частотних компонент у
функцію \(f(x)\). Для пилкоподібної функції

\[
c_n = \frac{2L}{n\pi}, \qquad n=1,2,\dots, \qquad c_0 = 0.
\]

\pandocbounded{\includegraphics[keepaspectratio]{images/_page_84_Figure_8.jpeg}}
Рис. 11.2. Дискретний частотний спектр пилкоподібної функції: a ---
основна частота (період \(2L\)), b --- \(n\)-та гармоніка (період
\(2L/n\)). Послідовність \(\{c_n\}\) нагадує спектральний розклад білого
світла на кольори. Введемо тепер поняття перетворення Фур'є. Для
неперіодичних функцій, визначених на всій дійсній осі, роль ряду Фур'є
відіграє інтеграл Фур'є (11.4)

\[
f(x) = \int_{0}^{\infty} a(\xi)\cos(\xi x)\,d\xi + \int_{0}^{\infty} b(\xi)\sin(\xi x)\,d\xi,
\]

де (11.5)

\[
a(\xi) = \frac{1}{\pi}\int_{0}^{\infty} f(x)\cos(\xi x)\,dx, \qquad
b(\xi) = \frac{1}{\pi}\int_{0}^{\infty} f(x)\sin(\xi x)\,dx.
\]

За аналогією з рядом Фур'є визначимо неперервний частотний спектр

\[
C(\xi) = \sqrt{a^2(\xi) + b^2(\xi)}.
\]

Частотний спектр є мірою внеску частоти \(\xi\) у функцію \(f(x)\).
\pandocbounded{\includegraphics[keepaspectratio]{images/_page_85_Figure_7.jpeg}}
Рис. 11.3. Частотні спектри деяких функцій. Тепер перейдемо до
перетворення Фур'є.

\section{Преобразование
Фурье}\label{ux43fux440ux435ux43eux431ux440ux430ux437ux43eux432ux430ux43dux438ux435-ux444ux443ux440ux44cux435}

Як уже зазначалося раніше, неперіодичні функції, задані на
\((-\infty,+\infty)\), не розкладаються у звичайні ряди Фур'є. Проте для
деяких таких функцій можна побудувати аналог ряду Фур'є у вигляді
інтеграла Фур'є, тобто розкладу за неперервним спектром частот: (11.4)

\[
f(x) = \int_{0}^{\infty} a(\xi)\cos(\xi x)\,d\xi + \int_{0}^{\infty} b(\xi)\sin(\xi x)\,d\xi,
\]

де (11.5)

\[
a(\xi) = \frac{1}{\pi}\int_{0}^{\infty} f(x)\cos(\xi x)\,dx, \qquad
b(\xi) = \frac{1}{\pi}\int_{0}^{\infty} f(x)\sin(\xi x)\,dx.
\]

З формули (11.4) видно, що інтеграл Фур'є розкладає функцію \(f(x)\) у
спектр за всіма частотами \(0<\xi<\infty\), а не лише за частотами,
кратними основній, як це було для періодичних функцій. За аналогією з
рядами Фур'є визначимо частотний спектр формулою

\[
C(\xi) = \sqrt{a^2(\xi) + b^2(\xi)}.
\]

Частотний спектр є мірою внеску частоти \(\xi\) у функцію \(f(x)\). На
рис. 11.3 зображено кілька функцій та їхні спектри. Функції з різкими
зламами мають широкий частотний спектр, тоді як для простої гармоніки
\(f(x)=\sin(\xi_0 x)\) спектр зосереджується поблизу однієї частоти
\(\xi=\xi_0\). Тепер ми готові перейти до експоненціальної форми
перетворення Фур'є. Скориставшись формулою Ейлера

\[
e^{i\theta} = \cos\theta + i\sin\theta,
\]

після нескладних перетворень приходимо до інтегрального представлення
Фур'є (11.6)

\[
f(x) = \frac{1}{\sqrt{2\pi}} \int_{-\infty}^{+\infty}
\left[\frac{1}{\sqrt{2\pi}} \int_{-\infty}^{+\infty} f(s)e^{-i\xi s}\,ds \right]
e^{i\xi x}\,d\xi.
\]

Звідси отримуємо пару інтегральних перетворень (11.7)

\[
\mathcal{F}[f] = F(\xi) = \frac{1}{\sqrt{2\pi}} \int_{-\infty}^{+\infty} f(x)e^{-i\xi x}\,dx
\]

(пряме перетворення Фур'є) та

\[
\mathcal{F}^{-1}[F] = f(x) = \frac{1}{\sqrt{2\pi}} \int_{-\infty}^{+\infty} F(\xi)e^{i\xi x}\,d\xi
\]

(обернене перетворення Фур'є). У наступній лекції ми вивчатимемо
властивості цих перетворень.

\section{ЗАУВАЖЕННЯ}\label{ux437ux430ux443ux432ux430ux436ux435ux43dux43dux44f-7}

\begin{enumerate}
\def\labelenumi{\arabic{enumi}.}
\tightlist
\item
  Образ Фур'є \(F(\xi)\) функції \(f(x)\) може бути комплекснозначною
  функцією. Наприклад, для функції
\end{enumerate}

\[
f(x)=
\begin{cases}
0, & x \le 0,\\
e^{-x}, & x > 0,
\end{cases}
\]

образ Фур'є має вигляд

\[
F(\xi) = \frac{1}{\sqrt{2\pi}} \frac{1-i\xi}{1+\xi^2}.
\]

\begin{enumerate}
\def\labelenumi{\arabic{enumi}.}
\setcounter{enumi}{1}
\tightlist
\item
  Абсолютне значення образу Фур'є називається частотним спектром функції
  \(f(x)\). Наприклад, якщо позначити \(C(\xi)=|F(\xi)|\), то маємо
\end{enumerate}

\[
C(\xi) = \sqrt{\frac{1}{2\pi(1+\xi^2)}}.
\]

\begin{enumerate}
\def\labelenumi{\arabic{enumi}.}
\setcounter{enumi}{2}
\tightlist
\item
  Не всі функції допускають перетворення Фур'є, тобто інтеграл (11.7)
  може не існувати. Наприклад, функції \(f(x)=c\), \(\sin x\) і \(e^x\)
  не мають класичного образу Фур'є. У прикладних задачах перетворення
  Фур'є зазвичай застосовують до температурних, хвильових та інших
  фізичних полів, які прямують до нуля при \(|x|\to\infty\).
\end{enumerate}

\section{ЗАВДАННЯ}\label{ux437ux430ux432ux434ux430ux43dux43dux44f-9}

\begin{enumerate}
\def\labelenumi{\arabic{enumi}.}
\tightlist
\item
  Розкласти у ряд Фур'є за синусами функцію
\end{enumerate}

\[
f(x)=
\begin{cases}
-1, & -1 < x < 0,\\
1, & 0 \le x < 1,
\end{cases}
\qquad
f(x+2)=f(x).
\]

Побудуйте суми перших 2, 3 і 4 членів ряду, щоб простежити його
збіжність до \(f(x)\). Побудуйте також графіки частотного спектра цієї
функції. 2. Показати, що якщо почленно диференціювати ряд Фур'є (11.3)
для пилкоподібної функції, то отримаємо послідовність, яка не збігається
до похідної цієї функції. 3. Побудувати частотний спектр таких
періодичних функцій: - a) \(f(x)=\sin x\), - б) \(f(x)=\sin x+\cos 2x\),
- в) \(f(x)=\sin x+\cos x+0.5\sin 3x\). 4. Знайти образ Фур'є \(F(\xi)\)
та частотний спектр \(C(\xi)=|F(\xi)|\) для функції

\[
f(x)=
\begin{cases}
1, & -1 < x < 1,\\
0, & \text{в інших точках.}
\end{cases}
\]

\begin{enumerate}
\def\labelenumi{\arabic{enumi}.}
\setcounter{enumi}{4}
\tightlist
\item
  Показати, що абсолютне значення функції \(F(\xi)=\dfrac{1}{1+i\xi}\)
  дорівнює
\end{enumerate}

\[
\left|F(\xi)\right| = \sqrt{\frac{1}{1+\xi^2}}.
\]

ПРИМІТКА. Спочатку помножте чисельник і знаменник на \(1-i\xi\), а потім
позбудьтеся уявної одиниці \(i\) у знаменнику. 6. Перевірити
ортогональність функцій \(\{\sin(n\pi x/L)\}\) та \(\{\cos(n\pi x/L)\}\)
на відрізку \([-L,L]\):

\[
\int_{-L}^{L} \sin\!\left(\frac{m\pi x}{L}\right)\sin\!\left(\frac{n\pi x}{L}\right)\,dx =
\begin{cases}
0, & m \neq n,\\
L, & m = n,
\end{cases}
\]

\[
\int_{-L}^{L} \cos\!\left(\frac{m\pi x}{L}\right)\cos\!\left(\frac{n\pi x}{L}\right)\,dx =
\begin{cases}
0, & m \neq n,\\
L, & m = n,
\end{cases}
\]

\[
\int_{-L}^{L} \sin\!\left(\frac{m\pi x}{L}\right)\cos\!\left(\frac{n\pi x}{L}\right)\,dx = 0
\quad \text{для всіх } m,n=1,2,\dots
\]

\section{Лекція 12}\label{ux43bux435ux43aux446ux456ux44f-12}

\section{ПЕРЕТВОРЕННЯ ФУР'Є ТА ЙОГО ЗАСТОСУВАННЯ ДО РОЗВ'ЯЗАННЯ
ДИФЕРЕНЦІАЛЬНИХ РІВНЯНЬ У ПОХІДНИХ
ПОХІДНИХ}\label{ux43fux435ux440ux435ux442ux432ux43eux440ux435ux43dux43dux44f-ux444ux443ux440ux454-ux442ux430-ux439ux43eux433ux43e-ux437ux430ux441ux442ux43eux441ux443ux432ux430ux43dux43dux44f-ux434ux43e-ux440ux43eux437ux432ux44fux437ux430ux43dux43dux44f-ux434ux438ux444ux435ux440ux435ux43dux446ux456ux430ux43bux44cux43dux438ux445-ux440ux456ux432ux43dux44fux43dux44c-ux443-ux43fux43eux445ux456ux434ux43dux438ux445-ux43fux43eux445ux456ux434ux43dux438ux445}

МЕТА ЛЕКЦІЇ: Проілюструвати основні властивості перетворення Фур'є та
показати, як ці властивості застосовуються до розв'язання рівнянь з
частинними похідними. Зокрема, ми побачимо, як під дією перетворення
Фур'є операція диференціювання замінюється множенням, а диференціальні
рівняння переходять в алгебраїчні. (12.1)

\[
\mathcal{F}[f] = F(\xi) = \frac{1}{\sqrt{2\pi}} \int_{-\infty}^{+\infty} f(x)e^{-i\xi x}\,dx.
\]

Отже, якщо функція \(f(x)\) визначена на всій дійсній осі, то
підставляючи її у формулу (12.1), ми отримуємо нову функцію \(F(\xi)\),
визначену для \(-\infty < \xi < \infty\). Таблиця 1. 12.1 показує деякі
функції та їхні фур'є-образи. Більш розгорнута таблиця функцій та їхніх
фур'є-образів наведена в додатку. З наведеної вище таблиці видно, що
образ Фур'є функції може виявитися комплекснозначною функцією аргументу
\(\xi\) . У першому прикладі цієї таблиці функція \(F(\xi)\) містить як
множник уявну одиницю i, і таким чином є комплекснозначною функцією
екзистенційного аргументу ( \(\xi\) змінюється від \(-\infty\) до
\(+\infty\) ). Іншими словами, аргумент \(\xi\) є реальним, а значення
функції є складними. Перетворення Фур'є, як і більшість інших
інтегральних перетворень, дуже корисне, оскільки воно замінює операцію
диференціювання операцією множення і тим самим перетворює диференціальні
рівняння на алгебраїчні. Існує багато властивостей, які роблять його
зручним інструментом для розв'язання задач. Тут ми перелічимо лише
найважливіші з них. Таблиця 12.1. Приклади перетворень Фур'є:

\begin{enumerate}
\def\labelenumi{\arabic{enumi}.}
\tightlist
\item
  Для непарної функції
\end{enumerate}

\[
f(x)=
\begin{cases}
e^{-x}, & x\ge 0,\\
-e^{x}, & x<0,
\end{cases}
\qquad
F(\xi)=-i\sqrt{\frac{2}{\pi}}\frac{\xi}{1+\xi^2}.
\]

\begin{enumerate}
\def\labelenumi{\arabic{enumi}.}
\setcounter{enumi}{1}
\tightlist
\item
  Для прямокутного імпульсу
\end{enumerate}

\[
f(x)=
\begin{cases}
1, & |x|<1,\\
0, & |x|\ge 1,
\end{cases}
\qquad
F(\xi)=\sqrt{\frac{2}{\pi}}\frac{\sin\xi}{\xi}.
\]

\begin{enumerate}
\def\labelenumi{\arabic{enumi}.}
\setcounter{enumi}{2}
\tightlist
\item
  Для гауссової функції
\end{enumerate}

\[
f(x)=e^{-x^2},
\qquad
F(\xi)=\frac{1}{\sqrt{2}}e^{-\xi^2/4}.
\]

\section{Основні властивості перетворення
Фур'є}\label{ux43eux441ux43dux43eux432ux43dux456-ux432ux43bux430ux441ux442ux438ux432ux43eux441ux442ux456-ux43fux435ux440ux435ux442ux432ux43eux440ux435ux43dux43dux44f-ux444ux443ux440ux454}

\subsection{Властивість 1 (Пара перетворень
Фур'є)}\label{ux432ux43bux430ux441ux442ux438ux432ux456ux441ux442ux44c-1-ux43fux430ux440ux430-ux43fux435ux440ux435ux442ux432ux43eux440ux435ux43dux44c-ux444ux443ux440ux454}

Після застосування перетворення Фур'є до функції \(f(x)\), визначеної
для \(-\infty < x < \infty\), отримуємо нову функцію \(F(\xi)\) за
формулою

\[
\mathcal{F}[f] = F(\xi) = \frac{1}{\sqrt{2\pi}} \int_{-\infty}^{+\infty} f(x)e^{-i\xi x}\,dx.
\]

Якщо тепер застосувати обернене перетворення Фур'є до функції
\(F(\xi)\), визначеної для \(-\infty < \xi < \infty\), то

\[
\mathcal{F}^{-1}[F] = f(x) = \frac{1}{\sqrt{2\pi}} \int_{-\infty}^{+\infty} F(\xi)e^{i\xi x}\,d\xi.
\]

то отримаємо початкову функцію \(f(x)\).
\pandocbounded{\includegraphics[keepaspectratio]{images/_page_90_Figure_7.jpeg}}
Рис. 12.1. Графіки функції та її образу Фур'є. Наприклад, (див. рис.
12.1)

\[
e^{-|x|} \xrightarrow{\mathcal{F}} \sqrt{\frac{2}{\pi}}\frac{1}{1+\xi^2} \xrightarrow{\mathcal{F}^{-1}} e^{-|x|}.
\]

\subsection{Властивість 2
(Лінійність)}\label{ux432ux43bux430ux441ux442ux438ux432ux456ux441ux442ux44c-2-ux43bux456ux43dux456ux439ux43dux456ux441ux442ux44c}

Перетворення Фур'є --- це лінійне перетворення, тобто

\[
\mathcal{F}[af+bg] = a\mathcal{F}[f] + b\mathcal{F}[g].
\]

Це легко перевірити безпосередньо. Наприклад, для виразу

\[
e^{-|x|} + 3e^{-x^2}
\]

має вигляд

\[
\mathcal{F}\left[e^{-|x|} + 3e^{-x^2}\right]
= \mathcal{F}\left[e^{-|x|}\right] + 3\mathcal{F}\left[e^{-x^2}\right].
\]

\subsection{Властивість 3 (Перетворення
похідних)}\label{ux432ux43bux430ux441ux442ux438ux432ux456ux441ux442ux44c-3-ux43fux435ux440ux435ux442ux432ux43eux440ux435ux43dux43dux44f-ux43fux43eux445ux456ux434ux43dux438ux445}

Якщо нас цікавить, як перетворення Фур'є діє на похідні, то насамперед
треба звернути увагу на змінні, за якими диференціюється функція. Якщо
перетворення Фур'є виконується за змінною \(x\), то маємо

\[
\begin{aligned}
\mathcal{F}[u_x] &= \frac{1}{\sqrt{2\pi}} \int_{-\infty}^{+\infty} u_x(x,t)e^{-i\xi x}\,dx = i\xi \mathcal{F}[u], \\
\mathcal{F}[u_{xx}] &= \frac{1}{\sqrt{2\pi}} \int_{-\infty}^{+\infty} u_{xx}(x,t)e^{-i\xi x}\,dx = -\xi^2 \mathcal{F}[u].
\end{aligned}
\]

З іншого боку, якщо перетворювати часткову похідну \(u_t(x,t)\),
залишаючи змінною інтегрування ту саму змінну \(x\), то отримаємо

\[
\begin{aligned}
\mathcal{F}[u_t] &= \frac{1}{\sqrt{2\pi}} \int_{-\infty}^{+\infty} u_t(x,t)e^{-i\xi x}\,dx = \frac{\partial}{\partial t}\mathcal{F}[u], \\
\mathcal{F}[u_{tt}] &= \frac{1}{\sqrt{2\pi}} \int_{-\infty}^{+\infty} u_{tt}(x,t)e^{-i\xi x}\,dx = \frac{\partial^2}{\partial t^2}\mathcal{F}[u].
\end{aligned}
\]

\subsection{Властивість 4
(Згортка)}\label{ux432ux43bux430ux441ux442ux438ux432ux456ux441ux442ux44c-4-ux437ux433ux43eux440ux442ux43aux430}

Кожне інтегральне перетворення має так звану властивість згортки. Образ
Фур'є добутку двох функцій, взагалі кажучи, не дорівнює добутку їхніх
образів, тобто

\[
\mathcal{F}[f(x)g(x)] \neq \mathcal{F}[f]\mathcal{F}[g].
\]

Однак у теорії перетворень існує операція, яка називається згорткою
функцій \(f\) і \(g\) і позначається через \(f*g\). Для неї виконується
співвідношення (12.2)

\[
\mathcal{F}[f*g] = \mathcal{F}[f]\mathcal{F}[g].
\]

Що це за загадкова операція під назвою згортка? Вона визначається
формулою (12.3)

\[
(f*g)(x) = \frac{1}{\sqrt{2\pi}} \int_{-\infty}^{+\infty} f(x-\xi)g(\xi)\,d\xi.
\]

Неважко показати, що співвідношення (12.2) справді виконується. З
означення видно, що згортка двох функцій \(f(x)\) і \(g(x)\) є новою
функцією \((f*g)(x)\) аргументу \(x\).

\section{Приклад обчислення згортки двох
функцій}\label{ux43fux440ux438ux43aux43bux430ux434-ux43eux431ux447ux438ux441ux43bux435ux43dux43dux44f-ux437ux433ux43eux440ux442ux43aux438-ux434ux432ux43eux445-ux444ux443ux43dux43aux446ux456ux439}

Нехай задано дві функції

\[
f(x) = x, \qquad g(x) = e^{-x^2}.
\]

Тоді їхня згортка дорівнює

\[
(f*g)(x) = \frac{1}{\sqrt{2\pi}} \int_{-\infty}^{+\infty} (x-\xi)e^{-\xi^2}\,d\xi = \frac{x}{\sqrt{2}}.
\]

Тут ми використали відому формулу

\[
\int_{-\infty}^{+\infty} e^{-\xi^2}\,d\xi = \sqrt{\pi}.
\]

Важливість концепції згортки (12.3) полягає в тому, що дуже часто,
розв'язуючи диференціальні рівняння, ми стикаємося з необхідністю
застосовувати обернене перетворення Фур'є до виразу, який є добутком
двох образів Фур'є, тобто потрібно знайти (12.4)

\[
\mathcal{F}^{-1}\left\{\mathcal{F}[f]\mathcal{F}[g]\right\}.
\]

Застосовуючи обернене перетворення до обох частин (12.2), отримуємо

\[
f * g = \mathcal{F}^{-1}\left\{\mathcal{F}[f]\mathcal{F}[g]\right\}.
\]

Отже, щоб визначити (12.4), потрібно знайти прототипи кожного з
факторів, тобто функцій f і g, а потім обчислити їхню згортку. Тепер ми
готові розв'язати одну важливу проблему: теорію диференціальних рівнянь
у похідних похідних.

\section{Розв'язання проблеми
Коші}\label{ux440ux43eux437ux432ux44fux437ux430ux43dux43dux44f-ux43fux440ux43eux431ux43bux435ux43cux438-ux43aux43eux448ux456}

Розглянемо проблему поширення тепла в \emph{нескінченному} стрижні, якщо
задано початкову температуру \(u(x, 0) = \varphi(x)\) . Цю задачу
зазвичай називають задачою Коші для рівняння теплопровідності або
задачою початкових умов

\[
\begin{aligned}
(\text{УЧП})\quad & u_t=\alpha^2u_{xx}, && -\infty<x<\infty,\ 0<t<\infty,\\
(\text{НУ})\quad & u(x,0)=q(x), && -\infty<x<\infty.
\end{aligned}
\]

Розв'язання задачі можна поділити на три основні етапи. КРОК 1
(Конверсія завдань). Оскільки просторова змінна x змінюється між
\(-\infty\) і \(+\infty\) , ми піддамо рівняння та початкову умову
(12,6) перетворенню Фур'є на змінній \(x^{(1)}\) (x --- інтегрована
змінна у перетворенні). Після цього ми отримуємо

\[
\mathcal{F}[u_t] = \alpha^2 \mathcal{F}[u_{xx}],\qquad
\mathcal{F}[u(x,0)] = \mathcal{F}[\varphi(x)].
\]

Використовуючи властивості перетворення Фур'є, отримуємо задачу Коші

\[
\begin{cases}
\dfrac{dU}{dt}=-\alpha^2\xi^2U,\\
U(0)=\Phi(\xi).
\end{cases}
\]

де \(U=U(\xi,t)=\mathcal{F}[u(x,t)]\) та
\(\Phi(\xi)=\mathcal{F}[\varphi(x)]\). ШАГ 2 (Решение преобразованной
задачи). Поскольку новая переменная играет роль гонстанты в задаче
(12.7), ее решение легко находится: (12.8)

\[
U(\xi,t)=\Phi(\xi)e^{-\alpha^2\xi^2 t}.
\]

ШАГ 3 (Нахождение обратного преобразования). Искомое решение \(u(x,t)\)
находится по формуле

\[
u(x, t) = \mathcal{F}^{-1}\!\left[U(\xi, t)\right]
= \mathcal{F}^{-1}\!\left[\Phi(\xi)e^{-\alpha^2 \xi^2 t}\right].
\]

Припускається, що існує перетворення Фур'є цих функцій. Примітка. Ед.
Вот здесь нас выручит теорема о свертке (12.5). Используя эту теорему,
получаем (12.9)

\[
\begin{aligned}
u(x, t)
&= \mathcal{F}^{-1}\!\left[\Phi(\xi)e^{-\alpha^{2}\xi^{2} t}\right] \\
&= \mathcal{F}^{-1}\!\left[\Phi(\xi)\right] * \mathcal{F}^{-1}\!\left[e^{-\alpha^{2}\xi^{2} t}\right] \\
&= \Phi(x) * \left[\frac{1}{2\alpha \sqrt{\pi t}} e^{-x^{2}/4\alpha^{2} t}\right] \\
&= \frac{1}{2\alpha \sqrt{\pi t}} \int_{-\infty}^{+\infty} \Phi(\xi)e^{-(x-\xi)^{2}/4\alpha^{2} t}\,d\xi.
\end{aligned}
\]

Дело сделано! Формула (12.9) дает решение нашей задачи. Прежде чем
поставить точку, давайте проанализируем полученный результат. Отметим,
что подынтегральная функция состоит из двух сомножителей: 1. Начальной
температуры q(x), 2. Функции
\(G(x, t) = \frac{1}{2\alpha \sqrt{\pi t}} e^{-(x-\xi)^2/4c^2t}\) ,
которая называется функцией Грина или функцией источника. Можно
показать, что функция источника G(x, t) описывает отклик системы на
начальный температурный единичный импульс Р И С
\pandocbounded{\includegraphics[keepaspectratio]{images/_page_94_Figure_9.jpeg}}
. 12.2. Отклик G(x,t) на температурный импульс в точке \(x=\xi\) ,
графчк G(x,t) похож на график нормального распределения, при этом время
игриег роль среднеквадратического уклонения: a---функция G(x,t) при
пебольших значениях t, \(\delta\) ---функция G(x,t) сущестьенно
уменьши\ldots ась. в ючке \(x=\xi\) . Другими словами, функция
\(G(\kappa,t)\) описывает распределение температуры в стержие в мочент
времени t, если на точку \(x=\xi\) подействовал единичный тепловой
импульс (см. рис. 12.2). Теперь формуле (12.9) можно дать следующую
интерпретацию: начальную температуру
\(u\left(x,\,0\right)=\phi\left(x\right)\) можно представлять как
конининуальное множество точечных импульсов каждой величины \(q(\xi)\)
(в точке \(x = \xi\) ). Каждый точечный импульс дает распределение
температуры \(\phi(\xi)\) G(x, t). Результирующее распределение
находится суммированием (интегрированием) температур точечных источников
по формуле (12.9). Позже мы увидим, что это одно из проявлений общего
принципа суперпозиции.

\section{ЗАМЕЧАНИЕ}\label{ux437ux430ux43cux435ux447ux430ux43dux438ux435}

Главный недостаток преобразования Фурье состоит в том, что его нельзя
применять ко всем функциям. Например, функции

\[
f(x)=\text{const}, \qquad f(x)=e^x, \qquad f(x)=\sin x
\]

не имеют фурье-образа, так как интеграл

\[
\mathcal{F}[f] = \frac{1}{\sqrt{2\pi}} \int_{-\infty}^{+\infty} f(x)e^{-i\xi x}\,dx
\]

расходится. Фурье-образ имеют только те функции, которые достаточно
быстро стремятся к нулю при \(|x| \to \infty\). Кроме того, нельзя
применять преобразование Фурье по времени для решения задач с начальными
условиями, так как \(0 < t < \infty\).

\section{ЗАДАЧІ}\label{ux437ux430ux434ux430ux447ux456-1}

\begin{enumerate}
\def\labelenumi{\arabic{enumi}.}
\tightlist
\item
  Найдите фурье-образ функции
\end{enumerate}

\[
f(x)=
\begin{cases}
0, & x \le 0, \\
e^{-x}, & x > 0.
\end{cases}
\]

используя таблицы, приведенные в приложении. - 2. На 10м переконайтеся,
що прямі та обернені перетворення Фур'є є лінійними перетвореннями. - 3.
Використовуючи перетворення Фур'є, розв'яжете задачу Коші

\[
\begin{aligned}
(\text{УЧП})\quad & u_t = \alpha^2 u_{xx}, && -\infty < x < \infty,\\
(\text{НУ})\quad & u(x,0)=e^{-x^2}, && -\infty < x < \infty.
\end{aligned}
\]

\begin{enumerate}
\def\labelenumi{\arabic{enumi}.}
\setcounter{enumi}{3}
\tightlist
\item
  Проверьте справедливость соотношений
\end{enumerate}

\[
\begin{aligned}
\mathcal{F}\left[u_{x}\right] &= i\xi \mathcal{F}\left[u\right], \\
\mathcal{F}\left[u_{xx}\right] &= -\xi^{2} \mathcal{F}\left[u\right].
\end{aligned}
\]

УКАЗАНИЕ, Воспользуйтесь интегрированием по частям. 5. Перевірити, чи
можна записати згортку двох функцій f і g у двох еквівалентних формах

\[
(f * g)(x) = \frac{1}{\sqrt{2\pi}} \int_{-\infty}^{+\infty} f(\xi)\, g(x - \xi)\, d\xi
\]

или

\[
(f * g)(x) = \frac{1}{\sqrt{2\pi}} \int_{-\infty}^{+\infty} f(x-\xi)\, g(\xi)\, d\xi.
\]

Лекция 13

\section{МЕТА ЛЕКЦІЇ З ПЕРЕТВОРЕННЯ
ЛАПЛАСА}\label{ux43cux435ux442ux430-ux43bux435ux43aux446ux456ux457-ux437-ux43fux435ux440ux435ux442ux432ux43eux440ux435ux43dux43dux44f-ux43bux430ux43fux43bux430ux441ux430}

МЕТА ЛЕКЦІЇ: Ввести перетворення Лапласа, ознайомити читача з формулою
оберненого перетворення та показати, як ці інструменти застосовуються до
диференціальних рівнянь.

Перетворення Лапласа має важливу перевагу над перетворенням Фур'є:
завдяки згасаючому множнику \(e^{-st}\) воно застосовується до значно
ширшого класу функцій. Крім того, перетворення Лапласа природно працює
на півосі \(0 \le t < \infty\), що особливо зручно для задач з
початковими умовами.

\section{Властивості перетворення
Лапласа}\label{ux432ux43bux430ux441ux442ux438ux432ux43eux441ux442ux456-ux43fux435ux440ux435ux442ux432ux43eux440ux435ux43dux43dux44f-ux43bux430ux43fux43bux430ux441ux430}

Свойство 1 (Прямое и обратное преобразования) Перетворення Лапласа та
формула його інверсії мають вигляд

\[
\mathscr{L}[f] = F(s) = \int_{0}^{\infty} f(t) \, e^{-st} \, dt
\quad \text{(перетворення Лапласа).}
\]

(13.2)

\[
\mathscr{L}^{-1}[F] = f(t) = \frac{1}{2\pi i} \int_{c-i\infty}^{c+i\infty} F(s) \, e^{st} \, ds
\quad \text{(формула інверсії).}
\]

Перетворення Лапласа має велику перевагу над перетворенням Фур'є:
множник \(e^{-st}\) забезпечує збіжність для ширшого класу функцій.
Більш точний опис цього класу дає наступна теорема.

\section{Достатні умови існування перетворення
Лапласа}\label{ux434ux43eux441ux442ux430ux442ux43dux456-ux443ux43cux43eux432ux438-ux456ux441ux43dux443ux432ux430ux43dux43dux44f-ux43fux435ux440ux435ux442ux432ux43eux440ux435ux43dux43dux44f-ux43bux430ux43fux43bux430ux441ux430}

Нехай 1. для будь-якого додатного числа \(A\) функція \(f(t)\) є
кусково-неперервною на інтервалі \(0 \le t \le A\); 2. існують такі
константи \(M\), \(a\) і \(T\), що \(|f(t)| \le Me^{at}\) для всіх
\(t > T\). Тоді перетворення Лапласа

\[
\mathscr{L}[f] = F(s) = \int_0^\infty f(t)e^{-st}\,dt
\]

існує для всіх \(s > a\). Приведем примеры некоторых функций и их
преобразований Лапласа У визначенні перетворення Лапласа зазвичай
вважають, що \(s\) --- дійсна змінна. Однак часто корисно розширити це
визначення на комплексні значення \(s\), а обернене перетворення
записати як

\[
\mathscr{L}^{-1}[F] = \frac{1}{2\pi i} \int_{\epsilon - i\infty}^{\epsilon + i\infty} F(s)e^{st}\,ds.
\]

Використання цієї формули передбачає знайомство з технікою інтегрування
вздовж контурів у комплексній площині та теорією дедукцій. Ми не будемо
відволікатися на вивчення цих питань тут, і коли йдеться про
перетворення Лапласа, ми використаємо таблицю, наведену в додатку.

\section{Властивість 2 (Перетворення часткових
похідних)}\label{ux432ux43bux430ux441ux442ux438ux432ux456ux441ux442ux44c-2-ux43fux435ux440ux435ux442ux432ux43eux440ux435ux43dux43dux44f-ux447ux430ux441ux442ux43aux43eux432ux438ux445-ux43fux43eux445ux456ux434ux43dux438ux445}

Нехай \(u(x,t)\) --- функція двох змінних, а

\[
U(x,s)=\mathcal{L}\!\left[u(x,t)\right].
\]

Оскільки перетворення Лапласа виконується за змінною \(t\), маємо
формули

\[
\begin{aligned}
\mathcal{L}\!\left[u_t\right]
&=\int_0^\infty u_t(x,t)e^{-st}\,dt
=sU(x,s)-u(x,0),\\
\mathcal{L}\!\left[u_{tt}\right]
&=\int_0^\infty u_{tt}(x,t)e^{-st}\,dt
=s^2U(x,s)-su(x,0)-u_t(x,0),\\
\mathcal{L}\!\left[u_x\right]
&=\int_0^\infty u_x(x,t)e^{-st}\,dt
=\frac{\partial U}{\partial x}(x,s),\\
\mathcal{L}\!\left[u_{xx}\right]
&=\int_0^\infty u_{xx}(x,t)e^{-st}\,dt
=\frac{\partial^2 U}{\partial x^2}(x,s).
\end{aligned}
\]

Нижче наведено кілька простих прикладів функцій та їхніх образів
Лапласа. 1. \(f(t)=1\), \(0<t<\infty\).
\pandocbounded{\includegraphics[keepaspectratio]{images/_page_99_Figure_3.jpeg}}

\[
F(s)=\int_0^\infty e^{-st}\,dt=\frac{1}{s}.
\]

\pandocbounded{\includegraphics[keepaspectratio]{images/_page_99_Figure_5.jpeg}}
2. \(f(t)=e^{2t}\), \(0<t<\infty\).
\pandocbounded{\includegraphics[keepaspectratio]{images/_page_99_Figure_7.jpeg}}

\[
F(s)=\frac{1}{s-2}, \qquad s>2.
\]

\pandocbounded{\includegraphics[keepaspectratio]{images/_page_99_Figure_9.jpeg}}
3. \(f(t)=\sin(\omega t)\).
\pandocbounded{\includegraphics[keepaspectratio]{images/_page_99_Figure_12.jpeg}}
\pandocbounded{\includegraphics[keepaspectratio]{images/_page_99_Figure_13.jpeg}}

\[
F(s)=\frac{\omega}{s^2+\omega^2}.
\]

\begin{enumerate}
\def\labelenumi{\arabic{enumi}.}
\setcounter{enumi}{3}
\tightlist
\item
  \(f(t)=e^{t^2}\) не має перетворення Лапласа. Рис. 13.1. Графіки
  деяких функцій та їхніх зображень за Лапласом. Правила перетворення
  просторових похідних випливають із теореми про диференціювання
  інтеграла за параметром:
\end{enumerate}

\[
\frac{d}{dx}\int_{a}^{b}f\left(x,\ y\right)dy=\int_{a}^{b}\frac{df}{dx}\left(x,\ y\right)dy,
\]

а формулы для преобразования производных \(u_t\) и \(u_{tt}\) можно
получить интегрированием по частям.

\section{Свойство 3 (Теорема о
свертке)}\label{ux441ux432ux43eux439ux441ux442ux432ux43e-3-ux442ux435ux43eux440ux435ux43cux430-ux43e-ux441ux432ux435ux440ux442ux43aux435}

Свертка играет здесь ту же роль, что и в теории преобразования Фурье, но
вводится несколько иначе.

\section{Определение конечной
свертки}\label{ux43eux43fux440ux435ux434ux435ux43bux435ux43dux438ux435-ux43aux43eux43dux435ux447ux43dux43eux439-ux441ux432ux435ux440ux442ux43aux438}

Конечная свертка двух функций f и g определяется формулами

\[
(f * g)(t)
= \int_0^t f(\tau)g(t-\tau)\,d\tau
= \int_0^t f(t-\tau)g(\tau)\,d\tau.
\]

(Можна довести, що ці інтеграли рівні). Тут інтегрування здійснюється
вздовж скінченного інтервалу від 0 до t, а не вздовж нескінченної (
\(\infty\) , \(+\infty\) ), як це було раніше з визначенням нескінченної
згортки. Ось приклад скінченної згортки двох функцій. Нехай

\[
f(t)=t,\qquad g(t)=t.
\]

тоді

\[
(f * g)(t)=\int_0^t \tau(t-\tau)\,d\tau=\frac{t^3}{6}.
\]

Як і у випадку нескінченної згортки, її найважливіша властивість
визначається формулою (13.4)

\[
\mathcal{L}[f*g]=\mathcal{L}[f]\mathcal{L}[g].
\]

або еквівалентна формула (13.5)

\[
\mathcal{L}^{-1}\!\left(\mathcal{L}[f]\mathcal{L}[g]\right)=f*g.
\]

Ці формули дозволяють застосовувати перетворення Лапласа у випадках,
коли образ можна подати як добуток \(\mathcal{L}[f]\) і
\(\mathcal{L}[g]\). Функції \(F\) і \(G\) легко знайти. Після
знаходження функцій \(f\) і \(g\) можна знайти їхню згортку. Наприклад,

\[
\mathcal{L}^{-1}\!\left(\frac{1}{s(s^2+1)}\right)=1-\cos t.
\]

Тут використано пари

\[
\mathcal{L}^{-1}\!\left(\frac{1}{s}\right)=1,\qquad
\mathcal{L}^{-1}\!\left(\frac{1}{s^2+1}\right)=\sin t.
\]

Тепер ми готові до важливого змішаного виклику.

\section{Теплопровідність у напівнескінченному
середовищі}\label{ux442ux435ux43fux43bux43eux43fux440ux43eux432ux456ux434ux43dux456ux441ux442ux44c-ux443-ux43dux430ux43fux456ux432ux43dux435ux441ux43aux456ux43dux447ux435ux43dux43dux43eux43cux443-ux441ux435ux440ux435ux434ux43eux432ux438ux449ux456}

Розглянемо глибокий резервуар з рідиною, і нехай бічна поверхня
резервуара буде теплоізольованою. Нехай \(u_0\) буде початковою
температурою рідини, а температура повітря над рідиною ---

\pandocbounded{\includegraphics[keepaspectratio]{images/_page_101_Picture_7.jpeg}}

Рис. 13.2. Схематичне представлення проблеми теплопровідності: a ---
тепло тече всередину, якщо ви(0, t) \textless{} 0, і назовні, якщо ви(0,
t) \textgreater{} 0; b --- резервуар настільки глибокий, що гранична
умова внизу не впливає на розв'язок при значеннях x.

кістка дорівнює нулю (орієнтир температур) (див. рис. 13.2). Наше
завдання --- знайти температурне поле в рідині на різних глибинах і в
різний час. Іншими словами, нам потрібно розв'язати проблему

(13.6)

\[
\begin{aligned}
(\text{УЧП})\quad & u_t = u_{xx}, && 0 < x < \infty,\; 0 < t < \infty, \\
(\Gamma \text{У})\quad & u_x(0,t) - u(0,t) = 0, && 0 < t < \infty, \\
(\text{НУ})\quad & u(x,0) = u_0, && 0 < x < \infty.
\end{aligned}
\]

Щоб розвинути це завдання, ми можемо відтворити Лапласа у мінливу t.
Відновити повагу, щоб можна було зупинити перетворення Лапласа на зміну
x, змінити з 0 на \(\infty\) . Обробка розмноження

Ми отримуємо звичайне диференціальне рівняння за змінною X.

(13.7)

\$\$

(ODY) \quad sU(x) - u\_0 = \frac{d^2U}{dx^2}, \quad 0 \textless{} x
\textless{} \infty,

\$\$

(\Gamma Y) \quad \frac{dU}{dx}(0) = U(0)

\$\$

(Ми перетворили лише рівняння та граничну умову, а не початкову умову.)
Маємо звичайне диференціальне рівняння другого порядку, але з однією
граничною умовою при x=0. Насправді існує ще одна гранична умова. Він
визначається фізичними міркуваннями і фіксується як \(U(x) \to 0\) ,
коли \(x \to +\infty\) . Зверніть увагу, що для спрощення позначення ми
опустили параметр s із аргументів функції U(x) всюди, оскільки рівняння
(13.7) є диференціальним рівнянням лише відносно x.

Щоб розв'язати задачу (13.7), запишіть загальне рішення (загальний
розв'язок однорідного + часткового розв'язку гетерогенного) TAC:

\$\$

U(x) = c\_1 e\^{}\{V\_\{5x\}\} + c\_2 e\^{}\{-V\_\{5x\}\} + u\_0/s.

\$\$

Підставляючи цей вираз у граничну умову (13.7), можна визначити
константи \(c_1\) та \(c_2\) (одразу зрозуміло, що \(c_3 = 0\) ,

\pandocbounded{\includegraphics[keepaspectratio]{images/_page_102_Figure_7.jpeg}}

Рис. 13.3. Граф ймовірнісного інтеграла (eri) та його комплементарної
функції (eric).

Інакше температура буде зростати нескінченно зі збільшенням координати
X). Визначаючи сталу \(c_2\) з граничної умови при x=0, отримуємо
фінальний вираз для \(U\left(x\right)\) (13.8)

\$\$

U(x) = u\_0 \left\{ \frac{e^{-V \widehat{s} x}}{s(V \widehat{s} + 1)}
\right\} + \frac{1}{s}.

\$\$

Залишається зробити останній крок. Щоб визначити температурне поле,
потрібно обчислити \(u\left(x,\,t\right)\) \$\$

u(x, t) = \mathcal{L}\^{}\{-1\}{[}U(x, s){]}

\$\$

(ми повернулися до написання всіх аргументів функції U(x, s)). Щоб
повернути перетворення Лапласа, ми використаємо таблиці, наведені в
додатку. У результаті ми отримуємо

(13.9)

\$\$

u(x, t) = u\_0 - u\_0
\left[ \text{eric} \left( \frac{x}{2} \sqrt{t} \right) + \text{eric} \left( \sqrt{t} + \frac{x}{2} \sqrt{t} \right) e^{x+t} \right],

\$\$

де

\$\$

\operatorname{eric}(x) = \frac{2}{\sqrt{\pi}} \int\_\{x\}\^{}\{\infty\}
e\textsuperscript{\{-\xi}\{2\}\} d\xi

\$\$

є функцією, доповнюючою до ймовірнісного інтеграла. Графіки цієї функції
показані на рис. 13.3. Аналіз цього рішення P H C

\pandocbounded{\includegraphics[keepaspectratio]{images/_page_103_Figure_7.jpeg}}

13.4. Температура всередині напівнескінченного середовища в різний час.

Потребує дуже малих зусиль. Фіг. (13.4) представлено результати
чисельного аналізу, отриманого на комп'ютері, оснащеному плоттером.

\section{ПРИМІТКИ}\label{ux43fux440ux438ux43cux456ux442ux43aux438-1}

\begin{itemize}
\tightlist
\item
  \begin{enumerate}
  \def\labelenumi{\arabic{enumi}.}
  \tightlist
  \item
    Перетворення Лапласа також можна звести до неоднорідного
    диференціального рівня з частковими множниками (для виконання методу
    диференціації різних рівнів можна звести до одного рівня), але можна
    звести його до рівня різних множників (метод диференціації можна
    звести до рівня коефіцієнтів заміщення). Таблиця 1. 13.2 У сфері
    захисту шкіри існує два методи.
  \end{enumerate}
\item
  \begin{enumerate}
  \def\labelenumi{\arabic{enumi}.}
  \setcounter{enumi}{1}
  \tightlist
  \item
    Для виконання цих і межових завдань також проводиться трансформація
    Ганкеля і Мелліна, а також трансформація Лапласа в одному аспекті.
    Відтворення Лапласа підстановки оперативного множення формули
  \end{enumerate}
\end{itemize}

\$\$

\mathscr{L}{[}y'{]} = s\mathscr{L}{[}y{]} - y(0),

\$\$

Таблиця 13.2 Порівняння методів розділення перетворення Лапласа та
змінних

{\def\LTcaptype{none} % do not increment counter
\begin{longtable}[]{@{}
  >{\raggedright\arraybackslash}p{(\linewidth - 4\tabcolsep) * \real{0.2907}}
  >{\raggedright\arraybackslash}p{(\linewidth - 4\tabcolsep) * \real{0.3140}}
  >{\raggedright\arraybackslash}p{(\linewidth - 4\tabcolsep) * \real{0.3953}}@{}}
\toprule\noalign{}
\begin{minipage}[b]{\linewidth}\raggedright
\end{minipage} & \begin{minipage}[b]{\linewidth}\raggedright
Метод
\end{minipage} & \begin{minipage}[b]{\linewidth}\raggedright
\end{minipage} \\
\midrule\noalign{}
\endhead
\bottomrule\noalign{}
\endlastfoot
& Перетворення Лапласа & Розділення \\
Гетерогенні PPP & Так & Так \\
Гетерогенні ГІ & Так & Так \\
Змінні ймовірності & Так & de' \\
Нелінійні рівняння & Так & Так \\
\end{longtable}
}

а перетворення Ганкеля і Мелліна замінюють диференціальний оператор на
множник. Наприклад, для перетворення Ганкеля, визначеного формулою

\$\$

H{[}y{]} = \int\emph{\{0\}\^{}\{\infty\} r J}\{o\}(\xi r) y(r) dr,

\$\$

справедливе співвідношення

\$\$

H\left[y''\left(r\right)+\frac{1}{r}\ y'\left(r\right)\right]=-\xi\^{}\{2\}H\left[y\right].

\$\$

Таким чином можна розв'язати спеціальне диференціальне рівняння з
змінними коефіцієнтами (рівняння Бесселя). 3. Перетворення Лапласа на
змінній t можна інтерпретувати як проекцію площини xt на осі x, що
призводить до У ц ь о м у в и

\pandocbounded{\includegraphics[keepaspectratio]{images/_page_104_Figure_9.jpeg}}

падку початкові рівняння, граничні та початкові умови перетворюються на
нове диференціальне рівняння та нові граничні умови (див. діаграму
вище).

\section{ЗАВДАННЯ}\label{ux437ux430ux432ux434ux430ux43dux43dux44f-10}

\begin{enumerate}
\def\labelenumi{\arabic{enumi}.}
\tightlist
\item
  Переконайтеся, що наступна формула для перетворення часткової похідної
  \(u_t\) є істинною:
\end{enumerate}

\$\$

\mathscr{L}\left[u_{t}(x, t)\right] = sU(x, s) - u(x, 0).

\$\$

\begin{enumerate}
\def\labelenumi{\arabic{enumi}.}
\setcounter{enumi}{1}
\tightlist
\item
  Використати перетворення Лапласа для розв'язання задачі Коші
\end{enumerate}

\$\$

\begin{aligned}
(\text{УЧП})\quad & u_t = \alpha^2 u_{xx}, && -\infty < x < \infty,\ 0 < t < \infty, \\
(\text{НУ})\quad & u(x,0) = \sin x, && -\infty < x < \infty.
\end{aligned}

\$\$

\begin{enumerate}
\def\labelenumi{\arabic{enumi}.}
\setcounter{enumi}{2}
\tightlist
\item
  За допомогою перетворення Лапласа за змінною \(t\) розв'язати задачу
\end{enumerate}

\$\$

\begin{aligned}
(\text{УЧП})\quad & u_t = u_{xx}, && 0 < x < \infty,\ 0 < t < \infty, \\
(\text{ГУ})\quad & u(0,t) = \sin t, && 0 < t < \infty, \\
(\text{НУ})\quad & u(x,0) = 0, && 0 \leqslant x < \infty.
\end{aligned}

\$\$

Дайте фізичну інтерпретацію цієї задачі.

\begin{enumerate}
\def\labelenumi{\arabic{enumi}.}
\setcounter{enumi}{3}
\tightlist
\item
  Розв'язати задачу з граничним значенням
\end{enumerate}

(ОДА)

\$\$

\frac{d^2U}{dx^2} - sU = A, \quad 0 \textless{} x \textless{} 1,

\$\$

(\Gamma Y) \quad

\$\$

\begin{cases} \frac{dU}{dx} (0) = 0 \\ U(1) = 0. \end{cases}

\$\$

Лекція 14

\section{ПРИНЦИП
ДЮАМЕЛЯ}\label{ux43fux440ux438ux43dux446ux438ux43f-ux434ux44eux430ux43cux435ux43bux44f}

МЕТА ЛЕКЦІЇ: Показати, як переклад Лапласа можна використовувати для
ідентифікації характеристик диференціальних рівнів. Зокрема, за
допомогою алгебранічного перетворення образів Лапласа, розвиток рівня
часткового оберненого можна відкинути принципом погляду Дюамеля.
Інтерпретація цього принципу в теорії найвищих диференціальних рівнів
світу використовується для розробки певних завдань УПП.

технології.

Давайте подивимось, як трансформація Лапласа здійснюється спробою
розвитку часткових диференціальних ліній, що дозволяє нам розкрити
фізичну сутність різних розвитків. За допомогою трансформації Лапласа ми
вже розробили принцип Дюамеля. Перше, що потрібно зробити --- звернутися
до принципу Дюамеля, ми бачимо проблему, яка часто є проблемою

\section{Теплопровідність у стрижні, кінці якого підтримуються при
заданій
температурі}\label{ux442ux435ux43fux43bux43eux43fux440ux43eux432ux456ux434ux43dux456ux441ux442ux44c-ux443-ux441ux442ux440ux438ux436ux43dux456-ux43aux456ux43dux446ux456-ux44fux43aux43eux433ux43e-ux43fux456ux434ux442ux440ux438ux43cux443ux44eux442ux44cux441ux44f-ux43fux440ux438-ux437ux430ux434ux430ux43dux456ux439-ux442ux435ux43cux43fux435ux440ux430ux442ux443ux440ux456}

Дуже часто необхідно знайти температуру всередині області, якщо відомо,
що на межі області вона змінюється з часом відповідно до певного закону
(граничні умови, залежні від часу).

На лівому кінці підтримується умова \(u(0,t)=0\), а на правому кінці
температура задається функцією \(f(t)\).

Рис. 14.1. Граничні умови, що залежать від часу.

Наприклад, розглянемо термоізольований стержень, правий кінець якого
підтримується у f(t) (див. рисунок 14.1). Потрібно вирішити цю проблему

(14.1)

\$\$

u\_\{t\} = u\_\{xx\}, \quad 0 \textless{} x \textless{} 1, \quad 0
\textless{} t \textless{} \infty,

\[ (\Gamma Y) \quad \begin{cases} u(0, t) = 0, \\ u(1, t) = f(t), \end{cases} \quad 0 < t < \infty,\]

(HY) \quad u(x, 0) = 0, \quad 0 \le x \le 1.

\$\$

Сподіваємося, що розв'язання цієї задачі (14.1) буде легким, оскільки ми
знаємо розв'язок простішої задачі, де температура в кінці є сталою:

Насправді, якщо обидві задачі ((14.1) і (14.2)) розв'язуються одночасно
за допомогою перетворення Лапласа, то отримаємо визначний результат:
розв'язок задачі (14.1) виражається через розв'язання задачі (14.2).

Отже, розв'язуючи обидві задачі одночасно, ми отримуємо

\section{Легке завдання (14.2) (Постійні
солдати)}\label{ux43bux435ux433ux43aux435-ux437ux430ux432ux434ux430ux43dux43dux44f-14.2-ux43fux43eux441ux442ux456ux439ux43dux456-ux441ux43eux43bux434ux430ux442ux438}

Застосування перетворення -Лапласа до задачі (14.2)

\$\$

\frac{d^2W}{dx^2} - sW(x) = 0,

\[ W(0) = 0,\]

W(1) = 1/s

\$\$.

Пемеріль ODU

\$\$

W(x, s) = \frac{1}{s}
\left[ \frac{\sinh(x \sqrt{s})}{\sinh(\sqrt{s})} \right].

\$\$

Пошук оберненого перетворення

\$\$

w(x, t) =

\[ = x + \frac{2}{\pi} \sum_{n=1}^{\infty} \frac{(-1)^n}{n} e^{-(n\pi)^2 t} \times\]
\(\times \sin(n\pi x)\) (розв'язання задачі з постійним ГІ).

\section{Важке завдання (14.1) (Солдати, що залежать від
часу)}\label{ux432ux430ux436ux43aux435-ux437ux430ux432ux434ux430ux43dux43dux44f-14.1-ux441ux43eux43bux434ux430ux442ux438-ux449ux43e-ux437ux430ux43bux435ux436ux430ux442ux44c-ux432ux456ux434-ux447ux430ux441ux443}

Застосування перетворення Лапласа до задачі (14.1)

\[\frac{d^2U}{dx^2} - sU(x) = 0,\]

U\left( \mathbf{0}\right) =0,

\[ U(1) = F(s)\].

\[U(x, s) = F(s) \left[ \frac{\sinh(x \sqrt{s})}{\sinh(\sqrt{s})} \right].\]

Розділіть і помножте на z

\[U(x, s) = F(s) \left\{ s \left[ \frac{\sinh(x \sqrt{s})}{s \sinh(s)} \right] \right\}.\]

Використовуючи співвідношення \(\mathscr{L}[w_t] = sW - w(x, 0),\)
отримуємо

\[U(x, s) = F(s) \mathcal{L}[w_t].\]

Відповідно,

\[u(x, t) = \mathcal{L}^{-1} \{ F(s) \mathcal{L}[w_t] \} =\]

= \mathcal{L}\^{}\{-1\} {[} F(s) {]} * \mathcal{L}\^{}\{-1\}{[}w\_t{]} =

\[ = f(t) * w_t(t) =\]

= \int\_0\^{}t w\_t(x, t-\tau) f(\tau) d\tau =

\$\$

(після інтегрування частинами)

\$\$

= \int\_\{0\}\^{}\{t\} w(x, t-\tau) f'(\tau) d\tau + f(0) w(x, t).

\$\$

(Розв'язок задачі з часово-залежними ГІ виражається через розв'язок
задачі з постійними ГІ

розв'язано у задачах з константою

Отже, ми висловили розв'язок задач із граничними умовами, що змінюються
в часі, через розв'язання задачі з постійними граничними умовами.
Відповідні формули такі:

(14.3)

плутанина у наступних формулах. Записати диференціальне рівняння в
похідних у нових координатах (~\$, т), мы воспользуемся следующими
формулами перехода;

\$\$

Наприклад, якщо dy/dx = B/2A = 3, то y - 3x = c, а отже
\(\xi(x, y) = y - 3x\) . Ця функція задовольняє умову
\([\xi_x/\xi_y] = -3\) і, відповідно, повертає коефіцієнт
\(\overline{A}\) до нуля.

Отже, половина роботи виконана: ми знайшли нову координатну
\(\xi = \xi(x, y)\) , яка забезпечила рівність \(\overline{A} = 0\) .
Залишається знайти \(\eta(x, y)\) таким чином, щоб перетворити
коефіцієнт \(\overline{B}\) на нуль.

Час зупинитися тут. Виявляється, що якщо вибрати \(B^2-4AC=0\) і \(\xi\)
так, щоб \(\overline{A}=0\) , то коефіцієнт \(\overline{B}\) автоматично
стане нулем 1). Тепер ми побачимо це.

Коефіцієнт \(\overline{B}\) задається виразом

\[ \overline{B}^2 - 4\overline{A}\overline{C} = (B^2 - 4AC) (\xi_x \eta_y - \xi_y \eta_x)^2 - \Pi \rho u u. \ \rho e \partial.\]

\overline{B} = 2 \left( V \overline{A} \xi\_x + V \overline{C} \xi\_y
\right) \left( V \overline{A} \eta\_x + V \overline{C} \eta\_y \right).

\$\$

Функциональная диаграмма на рис. 15.5 делает эти формулы совсем
очевидными. Эта диаграмма особенно полезна при вычислении частных
производных функций и и и по переменным Р И С

\pandocbounded{\includegraphics[keepaspectratio]{images/_page_115_Picture_14.jpeg}}

15.5 Диаграмма, иллюстрирующая функциональную зависимость переменных

\$ \xi \$ и \$ \tau \$ , поскольку \$ \xi \$ зависит и от x, и от t.
Переменная \$ \tau \$ зависит только от t. Это все, что касается
преобразований. Теперь подставим найденные выражения \$ u\_t \$ , \$
u\_x \$ и \$ u\_\{vx\} \$ в уравнение (15 2) и получим

\$\$

\overline{B} = 2 \sqrt{\overline{A}}
\left[ \sqrt{\overline{A}} \, \eta_x + \sqrt{\overline{C}} \, \eta_y \right].

\$\$

откуда

\$\$

u\_\{xy\} + 2u\_\{xy\} + u\_\{yy\} = 0.

\$\$.

Следовательно, новая задача с начальным условием в персменных \$ \xi \$
и \$ \tau \$ имеет вид

(УЧП)

\$\$

\frac{dy}{dx} = -{[}\xi\_x/\xi\_y{]} = B/2A = 1.

\$\$

(НУ) \$ u(\xi, 0) = 1 - H(\xi), \quad -\infty \textless{}
\xi \textless{} \infty. \$

Так как \$ \xi = x \$ при t = 0, то новые граничные условия совпадают со
старыми. Эта задача была решена в лекции 12 с помощью интегрального
преобразования Фурье, и ее решение записывается в виде

\$\$

y = x + c

\$\$

где \$ \phi(\beta) \$ --- начальное условие. Следовательно, в нашем
случае

\$\$

\eta = y

\$\$

После подстановки

\$\$

\begin{cases}
 \xi = y - x, \\
 \eta = y
 \end{cases}

\$\$

получаем окончательный результат

(15.3)

\$\$

\overline{C} = A \eta\_x\^{}2 + B \eta\_x \eta\_y + C \eta\_y\^{}2 = 1,

\[ \overline{E} = A \eta_{xx} + B \eta_{xy} + C \eta_{yy} + D \eta_x + E \eta_y = 0,
\]

\begin{cases} \frac{1}{2} \left[ 1 + \operatorname{erf} \left( \frac{-\xi}{2\sqrt{D\tau}} \right) \right], & \xi < 0, \\ \frac{1}{2} \operatorname{erfc} \left( \frac{\xi}{2\sqrt{D\tau}} \right), & 0 \leq \xi. \end{cases}

\[ \overline{G} = G = 0.\]

\overline{A}u\_\{\xi\xi\} + \overline{B}u\_\{\xi\eta\} +
\overline{C}u\_\{\eta\eta\} + \overline{D}u\_\{\xi\} +
\overline{E}u\_\{\eta\} + \overline{F}u = \overline{G}

\$\$

Графики этого решения для различных моментов времени изображены на рис.
15.6. Осталось сделать последний шаг --- записать решение нашей задачи в
координатах x и t. Подставляя

\$\$

u\_\{\eta\} = f(\xi),

\$\$, \$ \tau = t \$

в формулы (15.3), получаем

(15.4)

\$\$

u(\xi, \eta) = \eta f(\xi) + g(\xi),

\[ u(x, y) = yf(y-x) + g(y-x)\]

\begin{cases} \frac{1}{2} \left[ 1 + \operatorname{erf}\left(\frac{Vt - x}{2\sqrt{Dt}}\right) \right], & Vt > x, \\ \frac{1}{2} \operatorname{erfc}\left(\frac{x - Vt}{2\sqrt{Dt}}\right), & Vt \leq x. \end{cases}

\[u_{xx} + 2u_{xy} + u_{yy} = 0.\]

u(x, y) = y \sin(y-x) + (y-x)\^{}2

\[ Au_{xx} + Bu_{xy} + Cu_{yy} + Du_x + Eu_y + Fu = G\]

\begin{array}{ll} u_t = u_{xx} - 2u_x, & -\infty < x < \infty, & 0 < t < \infty, \\ u(x, 0) = \sin x, & -\infty < x < \infty. \end{array}

\[u_{\xi\xi} + u_{\eta\eta} = \varphi(\xi, \eta, u, u_{\xi}, u_{\eta}).\]

\overline{A}u\_\{\xi\xi\} + \overline{B}u\_\{\xi\eta\} +
\overline{C}u\_\{\eta\eta\} + \overline{D}u\_\{\xi\} +
\overline{E}u\_\{\eta\} + \overline{F}u = \overline{G}

\[ u_{\xi\eta} = \psi(\xi, \eta, u, u_{\xi}, u_{\eta}).\]

\begin{array}{ll} u_t = u_{xx} - 2u_x, & -\infty < x < \infty, \ 0 < t < \infty, \\ u(x, \ 0) = e^x \sin x, & -\infty < x < \infty. \end{array}

\[\frac{dy}{dx} = \frac{B - \sqrt{B^2 - 4AC}}{2A}\]

\xi(x, y) = \text{const}, \quad \eta(x, y) = \text{const}.

\[ \frac{dy}{dx} = -\sqrt{-4x^2} = -2ix,\]

\begin{array}{ll} (\text{YHI}) & u_t = -2u_x, \\ (\text{HY}) & u_t(x, 0) = e^{-x^2}. \end{array}

\$\$

також задаються комплексно-спряженим координатам

\$\$

\section{Конверсія
2}\label{ux43aux43eux43dux432ux435ux440ux441ux456ux44f-2}

Наступне перетворення з \((\xi,\eta)\) на \((\alpha,\beta)\) виконується
за формулами

\$\$

Проверьте ваш ответ.

\begin{enumerate}
\def\labelenumi{\arabic{enumi}.}
\setcounter{enumi}{3}
\tightlist
\item
  Решите задачу
\end{enumerate}

\$\$, \(\beta = \frac{\xi - \eta}{2i}\) .

Внаслідок застосування другого перетворення, рівняння

\$\$ набуває остаточної форми

\$\$

Перевірте своє рішення. Як виглядає рішення в різні моменти часу?

ПРИМІТКА. Перехід до рухомої системи відліку дозволяє негайно відкинути
член \$ Vu\_x \$ у рівнянні конвективної дифузії. Після розв'язання
нової проблеми

\$\$

де \(\phi\) та \(\phi\) використовуються для позначення правих сторін
рівнянь у загальному вигляді. Замість того, щоб доводити, що ці дві
трансформації дійсно ведуть до бажаного результату, давайте розглянемо
простий приклад.

\section{\texorpdfstring{Зведення рівняння \(y^2u_{xx} + x^2u_{yy} = 0\)
до його канонічної
форми}{Зведення рівняння y\^{}2u\_\{xx\} + x\^{}2u\_\{yy\} = 0 до його канонічної форми}}\label{ux437ux432ux435ux434ux435ux43dux43dux44f-ux440ux456ux432ux43dux44fux43dux43dux44f-y2u_xx-x2u_yy-0-ux434ux43e-ux439ux43eux433ux43e-ux43aux430ux43dux43eux43dux456ux447ux43dux43eux457-ux444ux43eux440ux43cux438}

Розглянемо рівняння

\$\$

у яких \(A=y^2\) , B=0, \(C=x^2\) , D=E=F=G=0. Дискримінантний
\(B^2-4AC\) дорівнює \(-4x^2y^2\) , і ми зведемо це рівняння до
канонічної форми в першому квадранті x\textgreater0, y\textgreater0.

КРОК 1 (Перше перетворення)

Запишемо характеристичні рівняння

\[ \frac{dy}{dx} = \frac{B + \sqrt{B^2 - 4AC}}{2A} = \frac{\sqrt{-4x^2y^2}}{2y^2} = i\frac{x}{y}\]

Замініть змінні \$ \xi = x - Vt \$ , \$ \tau = t \$ . У нашому випадку
інтеграл можна обчислити явно

\[y^2 + ix^2 = \text{const},\]

Такий інтеграл уже зустрічався у лекції 12. Можливо, читачу буде
зручніше звернутися до таблиць перетворення Фур'є.

\section{ОДНОВИМІРНЕ ХВИЛЬОВЕ РІВНЯННЯ (ГІПЕРБОЛІЧНІ
РІВНЯННЯ)}\label{ux43eux434ux43dux43eux432ux438ux43cux456ux440ux43dux435-ux445ux432ux438ux43bux44cux43eux432ux435-ux440ux456ux432ux43dux44fux43dux43dux44f-ux433ux456ux43fux435ux440ux431ux43eux43bux456ux447ux43dux456-ux440ux456ux432ux43dux44fux43dux43dux44f}

МЕТАЛЕКЦІЇ: Показати, що коливання навантаження описуються рівнянням

\$\$

Від

\$\$

і що це рівняння є наслідком законів Ньютона. Також розглядаються інші
типи хвильового рівняння:

\$\$ \(\eta(x, y) = y^2 - ix^2.\) (Неважливо, яка з двох функцій
називається \(\xi\) , а яка --- \(\eta\) ; ви можете поміняти їх
місцями, якщо хочете.) Після цього перетворення початкове рівняння
набуває вигляду

\$\$

Ми не звертатимемо уваги на це рівняння (яке є комплексним гіперболічним
рівнянням), а одразу перейдемо до другого перетворення.

КРОК 2. (Друга трансформація)

При здійсненні другої трансформації ми маємо

\$\$

\begin{split} u_{tt} &= \alpha^2 u_{xx} + F\left(x, \ t\right), \\ u_{tt} &= \alpha^2 u_{xx} - \beta u_t - \lambda u + F\left(x, \ t\right). \end{split}

\$\$

(реальна частина \(\xi\) і \(\eta\) ),
\(\beta = \frac{\xi - \eta}{2\iota} = x^2\) (уявна частина \(\xi\) і
\(\eta\) ).

Щоб спростити нотацію, найкраще перепозначити \((\alpha,\beta)\)
\((\xi,\eta)\) окремо і записати отримане \(\epsilon\) перетворення як

\$\$ КРОК 3 (Отримання нового рівняння)

Нову канонічну форму можна отримати, обчисливши всі коефіцієнти
\(\overline{A}, \overline{B}, \overline{C}, \overline{D}, \overline{E}, \overline{F}\)
і \(\overline{G}\) у рівнянні

\$\$

До цього часу ми розглядали фізичні процеси, які описуються одномірними
параболічними рівняннями (задачі дифузії). Тепер ми почнемо вивчати
інший клас диференціальних рівнянь у похідних --- гіперболічні рівняння.
Почнемо з одновимірного хвильового рівняння, яке описує, зокрема,
поперечні коливання струни.

\section{Вібрації
струн}\label{ux432ux456ux431ux440ux430ux446ux456ux457-ux441ux442ux440ux443ux43d}

Розглянемо невеликі коливання струни з нерухомими кінцями. Припустимо,
що струна щільно натягнута, зроблена з однорідного матеріалу і
коливається в одній площині (рис. 16.1).

\pandocbounded{\includegraphics[keepaspectratio]{images/_page_119_Figure_10.jpeg}}

Фіг. 16.1. Поперечні коливання струни: a --- зміщення струни з положення
рівноваги.

Ігнороване введення гравіювання на струну. Щоб створити математичну
модель проколювання струни, можна побачити всі сили, що діють на
невелику частину нитки (рис. 16.2). Щоб побачити, що рівень --- це не
більше ніж закон Ньютона (зміна імпульсу суми активних сил), який
зводиться до струни. Я вражений Фіг. На рисунку 16-1 можна уявити сили,
що діють на струну на прямих, перпендикулярних осі x.

\pandocbounded{\includegraphics[keepaspectratio]{images/_page_120_Figure_2.jpeg}}

Фіг. 16.2. Невелика частина \$ (x, x + \Delta x) \$ осцилюючою струною.

\begin{enumerate}
\def\labelenumi{\arabic{enumi}.}
\tightlist
\item
  Сила, що виникає через натяг струни ( \$ \alpha\^{}2 u\_\{xx\} \$ ).
  Значення поперечної компоненти сили натягу визначається відповідно до
  формули
\end{enumerate}

Сила =

\$\$

за формулами (41.5) з \(\xi = y^2\) та \(\eta = x^2\) . Зробивши це,
отримуємо канонічну форму еліптичного рівняння

\$\$ \$ \simeq T {[}u\_x(x + \Delta x, t) - u\_x(x, t){]}. \$ 2. Внешняя
сила F (x, t).

Зовнішня сила F(x, t), прикладена до струни, може довольно залежати від
x і t. Нижче наведено деякі приклади зовнішніх сил:

а) гравитационная сила F(x, t) = -mg;

\begin{enumerate}
\def\labelenumi{\alph{enumi})}
\setcounter{enumi}{1}
\item
  імпульси, розподілені вздовж струни та діючи на неї в різний час:
\item
  сили, що виникають внаслідок дії звукової хвилі на шкіру барабана
  (пізніше ми детально вивчаємо коливання мембрани, які описуються
  хвильовим рівнянням).
\end{enumerate}

\begin{enumerate}
\def\labelenumi{\arabic{enumi}.}
\setcounter{enumi}{2}
\item
  Сили тертя, що діють на струну (--- \$ \beta u\_t \$ ). Якщо струна
  вібрує у середовища, виникає сила опору, яка пропорційна швидкості \$
  u\_t \$ .
\item
  Возвращающая сила (- үи).
\end{enumerate}

Ця сила спрямована протилежно зміщенню струни. Якщо зміщення u додатне,
то сила від'ємна.

Якщо тепер застосувати рівняння руху Ньютона до невеликої ділянки
струни, отримаємо

\$\$

\section{ПРИМІТКИ}\label{ux43fux440ux438ux43cux456ux442ux43aux438-2}

\begin{enumerate}
\def\labelenumi{\arabic{enumi}.}
\tightlist
\item
  Аналогічно можна класифікувати лінійні рівняння другого порядку з
  \emph{трьох або більше змінних}, але для цього потрібен матричний
  аналіз. Наприклад, рівняння другого порядку з трьома змінними
\end{enumerate}

\$\$

де \$ \rho \$ --- лінійна густина струни. Поділимо обидві частини цього
співвідношення на \$ \Delta x \$ і поставимо \$ \Delta x \$ --- 0.
Результатом є добре відоме телеграфне рівняння

(16.1)

\$\$

згідно з цією класифікацією, буде параболічним, а рівняння

\$\$

У цьому рівнянні слід поділити \$ \alpha \$ , \$ \beta \$ , \$ \gamma \$
і F(x,t) на \$ \rho \$ , але для простоти залишимо старі позначки для
цих величин. Отримано бажане рівняння. Ось інтерпретація найпростішого
хвильового рівняння

\$\$

\begin{itemize}
\tightlist
\item
  Гіперболічно.
\end{itemize}

\begin{enumerate}
\def\labelenumi{\arabic{enumi}.}
\setcounter{enumi}{1}
\tightlist
\item
  Наш інтерес до класифікації диференціальних рівнянь у похідних
  пов'язаний із тим, що три основні типи рівнянь описують три різні типи
  фізичних явищ, і ми хотіли б отримати математичну класифікацію цих
  трьох класів фізичних явищ.
\end{enumerate}

\section{ЗАВДАННЯ}\label{ux437ux430ux432ux434ux430ux43dux43dux44f-11}

\begin{enumerate}
\def\labelenumi{\arabic{enumi}.}
\tightlist
\item
  Які з цих параболічних і еліптичних рівнянь записані у канонічній
  формі?
\end{enumerate}

\begin{enumerate}
\def\labelenumi{\alph{enumi})}
\tightlist
\item
  \(u_t = u_{xx} + hu\),
\end{enumerate}

б) \(u_{xy} + u_{xx} + 3u = \sin x\),

в) \(u_{xx} + 2u_{yy} = 0\),

г) \(u_{xx} = \sin^2 x\)?

\begin{enumerate}
\def\labelenumi{\arabic{enumi}.}
\setcounter{enumi}{1}
\item
  Перетворіть параболічне рівняння \(u_{xx} + 2u_{xy} + u_{yy} + u = 2\)
  до канонічної форми.
\item
  Перетворіть еліптичне рівняння
  \(u_{xx} + 2u_{yy} + x^2 u_x = e^{-x^2/2}\) до канонічної форми.
\end{enumerate}

\section{МЕТОД МОНТЕ-КАРЛО
(ВСТУП)}\label{ux43cux435ux442ux43eux434-ux43cux43eux43dux442ux435-ux43aux430ux440ux43bux43e-ux432ux441ux442ux443ux43f}

МЕТА ЛЕКЦІЇ: Пояснити основні ідени, що лежать в основі методу
Монте-Карло, і показати, як цей метод можна використовувати для
розв'язання різних задач. Виявляється, що за допомогою азартних ігор
можна змоделювати (загалом, на комп'ютері) наближені розв'язки реальних
задач. Як простий приклад можна вказати обчислення інтеграла

\$\$

основанную на интуитивных соображениях.

\section{Інтуїтивна інтерпретація хвильового
рівняння}\label{ux456ux43dux442ux443ux457ux442ux438ux432ux43dux430-ux456ux43dux442ux435ux440ux43fux440ux435ux442ux430ux446ux456ux44f-ux445ux432ux438ux43bux44cux43eux432ux43eux433ux43e-ux440ux456ux432ux43dux44fux43dux43dux44f}

Читач може запитати: чому рівняння типу (16.2) має описувати щось
подібне до вібрацій скрипкової струни? Щоб зрозуміти це, необхідно
зрозуміти Ф і г

\pandocbounded{\includegraphics[keepaspectratio]{images/_page_121_Picture_9.jpeg}}

. 16 3. Інтерпретація хвильового рівняння \$ u\_\{tt\} = \alpha\^{}2
u\_\{xx\} \$ \$ u\_\{tt\} \$ Equation (16.2) можна інтерпретувати так:
прискорення струни, викликане натягом, у кожній точці більше, чим більша
увігнутість струни \$ u\_\{xx\} \$ у заданій точці (коефіцієнт
пропорційності сталий \$ \alpha\^{}2 = T/\rho \$ ) (див. рис. 133). \#
ПОВАГА 1. Хвильове рівняння \$ u\_\{tt\} = \alpha\^{}2 u\_\{xx\} \$
також описує поздовжні та крутильні коливання стрижня. У поздовжніх
осцилях зміщення паралельні елементу, і значення поздовжнього зміщення
відносно положення рівноваги позначається через \$ u\left(x,,t\right) \$
. Такі коливання виникають, наприклад, коли молоток б'ють по кінці
стрижня. Торсіональний

ные колебания

\$\$

кидаючи дротик у одиничний квадрат \(\{(x, y): 0 < x < 1, 0 < y < 1\}\)
Після приблизно 100 кидків відносна частина дротиків, що опускається
нижче кривої \(y = x^2\) , вважається приблизним значенням інтеграла.
Узагальнення такої гри (з випадковим киданням дрогіка) завжди містить
процедуру генерації послідовності випадкових чисел.

У цій лекції ми коротко опишемо метод Монте-Карло, а в наступній лекції
покажемо, як його використовувати для розв'язання диференціальних
рівнянь у похідних

Методи Монте-Карло зазвичай називають групою методів розв'язання
детермінованих задач (тобто задач без випадковості), у яких суттєво
використовуються елементи випадковості. Загальна філософія етнічних
методів показана на рис. 42 1.

Ймовірнісна гра

Деторміноване завдання \(\rho\) --- розв'язання задачі (наприклад,
обчислення інтеграла, розв'язання UBP тощо).

Результат = \(\hat{p}\) Наближення p Рішення = p Рис. 42.1. Загальна
теорія методу Монте-Карло.

\section{Обчислення
інтеграла}\label{ux43eux431ux447ux438ux441ux43bux435ux43dux43dux44f-ux456ux43dux442ux435ux433ux440ux430ux43bux430}

Припустимо, ми хочемо обчислити інтеграл

\$\$

Тут \emph{k} --- це модуль Юнга, який характеризує еластичність
матеріалу. Матеріали з великим модулем Юнга коливаються на вищих
частотах. Звукові хвилі --- це поздовжні хвилі.

\begin{enumerate}
\def\labelenumi{\arabic{enumi}.}
\setcounter{enumi}{1}
\tightlist
\item
  Якщо лінійна густина струни \$ \rho(x) \$ залежить від координати, то
  хвильова аугура записується як
\end{enumerate}

\$\$

(детерміноване завдання). Щоб використати метод Монте-Карло, потрібно
винайти гру на удачу, яка дає приблизне значення цього інтеграла.

Звісно, можна придумати багато різних ігор. Остаточна версія залежить
від точності розрахунків, простоти гри тощо. Очевидний варіант гри для
обчислення інтегралу --- кинути дротик у прямокутний
\(R = \{(x,y): a \le x \le b, 0 \le y \le \max f(x)\}\) (див. рис. 42.2)

\pandocbounded{\includegraphics[keepaspectratio]{images/_page_126_Figure_14.jpeg}}

. 17.1. Волна \$ u(x, t) = e\textsuperscript{\{-(x-ct)}2\} \$ движется
слева направо

произвольной формы, движущихся в прогивоположных направлениях со
скоростью c.~Например, функции

\$\$

= \int\_\{0\}\^{}\{d\}
\left[\frac{\partial F}{\partial \overline{y}} \eta\left(x\right) + \frac{\partial F}{\partial \overline{y'}} \eta'\left(x\right)\right] dx
= 0.

\$\$

(волна, Лвижущаяся слева направо),

\$\$

\frac{d\varphi(e)}{de} \equiv \int\_\{a\}\^{}\{b\} \left\{
\frac{\partial F}{\partial \overline{y}} - \frac{\partial}{\partial x}
\left[ \frac{\partial F}{\partial \overline{y'}} \right] \right\} ,
\eta(x) , dx = 0.

\$\$

(волна, движущаяся справа налево),

\$\$

\frac{\partial F}{\partial \overline{y}} - \frac{d}{dx}
\left[ \frac{\partial F}{\partial \overline{y'}} \right] = 0

\$\$

(две волны, движущиеся в противоположных направлениях)

являются типичными решениями волиового уравнения. На рисунке 17.1
изображена простейшая бегущая волна.

ШАГ 4 (Учет начальных условий).

Напомним читателю общий метод решения задачи Коши в теории обыкновенных
дифференциальных уравнений. Сначала находят общее решение уравнения, а
затем оно подставляется в начальные условия для того, чтобы найти
конкретные значения произвольных постоянных. У нас ситуация аналогична.
Для решения задачи Коши для волнового уравнения мы подставим общее
решение

\$\$

\frac{\partial F}{\partial y} - \frac{d}{dx}
\left[ \frac{\partial F}{\partial y'} \right] = 0

\$\$

(содержащее две произвольные функции) в начальные условия

\$\$

J{[}y{]} = \int\_\{0\}\^{}\{1\} {[}y\^{}\{2\} + y'\^{}\{2\}{]} dx.

\$\$

\$ u\_1(x, 0) = g(x), \$ чтобы найти конкретные выражения для
произвольных функций ф и ф. После подстановки получаем

(17.5)

\$\$

F\_y - \frac{d}{dx} F\_\{y'\} = 0,

\$\$

Если теперь проинтегрировать второе уравнение (17.5), то получится
алгебранческая система двух уравнений относительно неизвестных функций
\$ \varphi(x) \$ и \$ \psi(x) \$ . В самом деле, после интегрирования в
пределах от \$ x\_0 \$ до x получаем

(17.6)

\$\$

\begin{aligned} F_y &= 2y, \ F_{y'} &= 2y'. \end{aligned}

\$\$

Решаем теперь уравнение (17.6) совместно с первым уравнением (17.5) и
получаем следующие выражения для функций \$ \varphi(x) \$ и \$ \psi(x)
\$ :

(17.7)

\$\$

2y - \frac{d}{dx}(2y') = 0,

\[ y'' - y = 0.
\]

Отже, загальне розв'язання задачі (17.1) задається формулою

(17.8)

\[ \overline{y}(x) = 0.42e^x - 0.42e^{-x}
\]

Мета досягнута. Ми отримали загальне розв'язання Вадачі (17.1).
Розв'язок у вигляді (17.8) зазвичай називають формулою д'Аламбера.
Читача запрошують самостійно слідкувати за розташуванням меж
інтегрування, щоб переконатися, що вони змінюються з x-ct на x+ct.
Проблема повністю вирішена.

Перед завершенням лекції ми наведемо кілька прикладів, які показують, як
застосовувати формулу д'Аламбера до конкретних валахів.

\section{Приклади застосувань формули
д'Аламбера}\label{ux43fux440ux438ux43aux43bux430ux434ux438-ux437ux430ux441ux442ux43eux441ux443ux432ux430ux43dux44c-ux444ux43eux440ux43cux443ux43bux438-ux434ux430ux43bux430ux43cux431ux435ux440ux430}

\section{1. Рух синузондалних
хвиль}\label{ux440ux443ux445-ux441ux438ux43dux443ux437ux43eux43dux434ux430ux43bux43dux438ux445-ux445ux432ux438ux43bux44c}

Рассмотрим начальные условия вида

\[ J[y] = \int_a^b y \sqrt{1 + y'^2} \, dx,
\] \$ u\_1(x, 0) = 0. \$ По формуле Даламбера получаем решение

\[ \overline{y}(x) = \alpha \operatorname{ch}[(x-\beta)/\alpha],
\]

Її можна інтерпретувати так. Початкове переміщення струни \$ u(x, 0) =
\sin x \$ поділяється на дві однакові частини

\[ \overline{y}(x) = 0.42e^x - 0.42e^{-x}
\]

H \$ \frac{\sin x}{2} \$ .

Кожна з частин поширюється зі швидкістю c у вигляді хвиль, що рухаються.
Одна з хвиль рухається зліва направо, а інша --- у протилежному
напрямку. Читача запрошують замислитися, як виглядає отримана хвиля у
цьому випадку.

\section{2. Распространение простейшего прямоугольного
импульса}\label{ux440ux430ux441ux43fux440ux43eux441ux442ux440ux430ux43dux435ux43dux438ux435-ux43fux440ux43eux441ux442ux435ux439ux448ux435ux433ux43e-ux43fux440ux44fux43cux43eux443ux433ux43eux43bux44cux43dux43eux433ux43e-ux438ux43cux43fux443ux43bux44cux441ux430}

В этом случае начальные условия задаются следующим образом:

\[ J[y] = \int_{0}^{1} [y^{2} + y'^{2}] dx.
\] \$ u\_t(x, 0) = 0 \$ (в остальных точках).

Оскільки початкове зміщення розділяється на дві півхвилі, що рухаються в
протилежних напрямках, отриманий рух буде таким, як показано на рис.
17.2.

\pandocbounded{\includegraphics[keepaspectratio]{images/_page_129_Figure_2.jpeg}}

Фіг. 17.2. Початкове збурення розділяється на дві хвилі, що рухаються.

\section{3. Задана начальная
скорость}\label{ux437ux430ux434ux430ux43dux430-ux43dux430ux447ux430ux43bux44cux43dux430ux44f-ux441ux43aux43eux440ux43eux441ux442ux44c}

Тепер припустимо, що в початковий момент струна перебуває в рівноваговій
позиції. Надамо струні початкову швидкість (як у фортепіано) форми \$
\sin x \$ , тобто

\[ \int\limits_{t_1}^{t_2}
\]

J\left[u\right] = \iint\limits\emph{\{D\}
F\left(x,~y,~u,~u}\{x\},~u\_\{y\}\right) , dx , dy.

\$\$

Решение находим по формуле Даламбера

\$\$

F\_\{u\} - \frac{\partial}{\partial x} F\_\{ux\} -
\frac{\partial}{\partial y} F\_\{uy\} = 0.

\$\$

Оно представляет сумму двух бегущих косинусоидальных воли. Читатель
должен спросить себя, кажется ли ему это решение разумным.

На этом лекция 17 заканчивается. В следующей лекции мы покажем, как
формула Даламбера помогает дать интерпретацию в плоскости переменных x и
t. \# ЗАУВАЖЕННЯ - 1. Общее решение уравнения с частными производными
второго порядка содержит две произвольные функции, а общее решение
обыкновенного дифференциального уравнения содержит две произвольные
постоянные. Значит, у уравнения с частными производными больше решений,
чем у обыкновенного дифференциального уравнения. - 2 Замена переменных
для приведения уравнения с частными производными к более простому виду
является одним из общих методов теории. Новые координаты (ξ, η) принято
называть каноническими координатами. Мы будем обращаться к ним и в
дальнейшем, особенно при изучении гиперболических уравнений. - 3 Метод
нахождения общего решения уравнения с частными производными и
подстановки этого решения в начальные условия не является общим методом
решения уравнений с частными производными. Решение, которое
рассматривается в данной лекции,--- единственный пример использования
этого метода. Обычно невозможно найти общее решение уравнения с частными
производными, но даже если удается это сделать, оказывается очень
сложным подставить это решение в краевые условия.

\section{ЗАДАЧН}\label{ux437ux430ux434ux430ux447ux43d}

\begin{enumerate}
\def\labelenumi{\arabic{enumi}.}
\item
  Проверьте, что решение Даламбера (17-8) удовлетворяет начальным
  условиям задачи (17.1).
\item
  Подставьте (17.7) в общее решение \$ u(x, t) = \psi(x - ct) + \psi(x +
  \epsilon t) \$ и получите формулу Даламбера.
\item
  Найдите решение задачи Коши:
\end{enumerate}

\$\$

J{[}y{]} = \int\_\{0\}\^{}\{1\} \sqrt{1 + y'^{2}} dx.

\$\$

Как выглядит это решение в различные моменты времени? 4. Найдите решение
задачи Коши:

(УЧП)

\$\$

\int\emph{\{t}\{1\}\}\^{}\{t\_\{2\}\} {[}KE - PE{]} dt = \frac{1}{2}
\int\emph{\{t}\{1\}\}\^{}\{t\_\{2\}\} {[}m\dot{y}\^{}\{2\} -
ky\^{}\{2\}{]} dt

\$\$

(НУ)

\$\$

J{[}y{]} = \int\_\{0\}\^{}\{\pi/2\} {[}y'\^{}\{2\} - y\^{}\{2\}{]} dx

\$\$

Постройте график решения в различные моменты времени. 5. Решая совместно
уравнения (17.6) и первое уравнение из (17.5),

получите формулы (17.7).

\section{ФОРМУЛА ДАЛАМБЕРА
(ПРОДОЛЖЕНИЕ)}\label{ux444ux43eux440ux43cux443ux43bux430-ux434ux430ux43bux430ux43cux431ux435ux440ux430-ux43fux440ux43eux434ux43eux43bux436ux435ux43dux438ux435}

ЦЕЛЬ ЛЕКЦИИ! Показать, как с помощью формулы Даламбера можно найти
решение волнового уравнения на полуограниченной прямой

(УЧП)

\$\$

J{[}u{]} = \iint\_D F{[}x, y, u, u\_x, u\_y) , dx , dy,

\$\$

(ГУ) \$ u(0, t) = 0, \quad 0 \textless{} t \textless{} \infty, \$ (НУ)

\$\$

F\_\{u\} - \frac{\partial}{\partial x} F\_\{uv\} -
\frac{\partial}{\partial y} F\_\{uv\} = 0.

\$\$

Дополнительно дается интерпретация формулы Даламбера в координатной
плоскости x, t.

На прошлой лекции было показано, что если заданы начальное смещение
струны u(x, 0) = f(x) и начальная скорость \$ u\_t(x, 0) = g(x) \$ , то
выражение

(18.i)

\$\$

u\_\{xx\} + u\_\{yy\} = f

\$\$

описывает смещение струны в произвольный момент времени. В этой лекции
мы познакомим читателя с одной интересной интерпретацией этой формулы в
плоскости переменных x и t (пространственновременная или фазовая
плоскость) и покажем, как изменить формулу Даламбера, чтобы получить
решение задачи о колебаниях полубесконечной струны.

Начнем с интерпретации формулы (18.1) в плоскости переменных x и t.

\section{Пространственно-временная интерпретация формулы
Даламбера}\label{ux43fux440ux43eux441ux442ux440ux430ux43dux441ux442ux432ux435ux43dux43dux43e-ux432ux440ux435ux43cux435ux43dux43dux430ux44f-ux438ux43dux442ux435ux440ux43fux440ux435ux442ux430ux446ux438ux44f-ux444ux43eux440ux43cux443ux43bux44b-ux434ux430ux43bux430ux43cux431ux435ux440ux430}

На предыдущей лекции мы показали, что решение задачи Коши

(18.2)

\$\$

J{[}u{]} = \iint\_\{D\} {[}u\_\{x\}\^{}\{2\} + u\_\{y\}\^{}\{2\} +
2uf{]} dx dy.

\[ J[u] = \int_{0}^{1} \int_{0}^{1} \left[ u_{1}^{2} + u_{y}^{2} \right] dx dy
\]

u(x, t) = \frac{1}{2} {[}f(x-ct) + f(x+ct){]} + \frac{1}{2c}
\int\_\{x-ct\}\^{}\{x+ct\} g(\xi) d\xi.

\[ u_{xx} + u_{yy} = 0
\]

u(x, 0) = f(x),

\[ u_{xx} + u_{yy} = 0
\]

u(x, t) = \frac{1}{2} {[}f(x-ct) + f(x+ct){]}.

\[ J[u] = \int_{0}^{1} \int_{0}^{1} \left[ u_{x}^{2} + u_{y}^{2} \right] dx dy
\]

x-ct=x\_0-ct\_0

\$\$

(PPP)

\$\$

(Y \text{ ЧП}) \qquad u\_\{tt\} = c\^{}2 u\_\{xx\},
\quad -\infty \textless{} x \textless{} \infty, \quad 0 \textless{} t
\textless{} \infty,

\[ J[u] = \iint_{D} [u] + u_{\eta} + 2uf ] dx dy.
\]

(HY) \qquad

\[ u_{xx} + u_{yy} = f
\]

J{[}u{]} = \int\emph{\{0\}\^{}\{1\} \int}\{0\}\^{}\{1\} {[}u' + u' +
2uf{]} dx dy,

\[ J[u] = \int_{0}^{1} \int_{0}^{1} \left[ u_{x}^{2} + u_{y}^{2} + 2uf \right] dx dy,
\]

Фіг. 18.2 Розв'язок зображено на площині змінних x, t, а не u, x, як у
попередній лекції.

\pandocbounded{\includegraphics[keepaspectratio]{images/_page_133_Figure_5.jpeg}}

Фіг. 18.2. Розв'язок задачі Коші (18.2) у xt-площині a --- це передній
край хвилі, що рухається вліво; b --- це задній край хвилі; b --- це
передній край хвилі, що рухається вправо

Тепер інтерпретуємо формулу д'Аламбера, де початкове зміщення дорівнює
нулю, а початкова швидкість довільна.

СЛУЧАЙ 2 (Начальное смещение равно нулю, начальная скорость
произвольна).

Рассмотрим начальные условия вида

\[ u_n(x, y) = a_1 \varphi_1(x, y) + a_2 \varphi_2(x, y) + \ldots + a_n \varphi_n(x, y),
\] \$ u\_t(x, 0) = g(x). \$ Решение в этом случае имеет вид

\[ \varphi_1(x, y) = xy(1-x)(1-y)
\]

Значення значення значення U в точці \$
(x\_\{\mathrm{ }e\},~t\_\{\mathrm{ }e\}) \$ можна інтерпретувати як
інтеграл початкової швидкості від \$
x\_\{\mathrm{ }e\}-ct\_\{\mathrm{ }e\} \$ до P і C

\pandocbounded{\includegraphics[keepaspectratio]{images/_page_134_Figure_2.jpeg}}

. 18.3. Інтерпретація задачі з початковою швидкістю в площині x1. \$
x\_\{\bullet\}+ct\_\{\bullet\} \$ (рис. 18.3). У заявці на рис. 18.4
Відображення завдань

(13.4)

Перші наближення в методі Рітца мають вигляд

\[
u_1(x,y)=a_1xy(1-x)(1-y),
\]

\[
u_3(x,y)=xy(1-x)(1-y)\,[a_1+a_2x+a_3y],
\]

і далі аналогічно для вищих наближень. Відповідний функціонал
записується як

\[
J[u_n] = \int_0^1 \int_0^1 \left[\left(\sum_{j=1}^n a_j \frac{\partial \varphi_j}{\partial x}\right)^2 + \left(\sum_{j=1}^n a_j \frac{\partial \varphi_j}{\partial y}\right)^2 + 2f\sum_{j=1}^n a_j\varphi_j\right] dx\,dy.
\]

Умови мінімуму отримують із рівностей

\[
\frac{\partial J[u_n]}{\partial a_k}=0, \qquad k=1,2,\ldots,n.
\]

Після цього задача зводиться до системи лінійних алгебраїчних рівнянь
для коефіцієнтів \(a_1,\ldots,a_n\).

\$\$

= \frac{1}{2c} \int\_\{x-ct\}\^{}\{-1\} d\xi, \quad (x, t)
\in \text{ области 3,}

\$\$ где \(A = (A_{IJ})\) матрица размером \(n \times n\) , элементы
которой вычисляются по формулам

(45.3) \$\$

= \frac{1}{2c} \int\_\{x-ct\}\^{}\{1\} d\xi, \quad (x, t)
\in \text{ области 4,}

\$\$ \(b = (b_t)\) --- вектор с компонентами

\$\$

= \frac{1}{2c} \int\_\{x-ct\}\^{}\{x+ct\} 0 d\xi, \quad (x, t)
\in \text{ области 5,}

\$\$ а \(a=(a_i)\) --- неизвестный вектор, компоненты которого являются
коэффициентами в приближенном решении
\(a_n(x,y)=a_1\varphi_1(x,y)+a_2\varphi_2(x,y)+\ldots+a_n\varphi_n(x,y)\)
.

ШАГ 3. Решаем линейную систему Aa = b относительно коэффициентов
\(a_1, a_2, \ldots, a_n\) и получаем приближенную минимизирующую функцию

\$\$

= \frac{1}{2c} \int\_\{x-ct\}\^{}\{x+ct\} d\xi, \quad (x, t)
\in \text{ области 6.}

\$\$ а эначит, и приближенное решение задачи Дирихле. \# ЗАУВАЖЕННЯ 1. В
этой лекции методом Ритца минимизировался функционал, вависящий от
функции двух переменных. Точно так же можно минимизировать функционалы
вида

\$\$

y(0) = 0, \quad y(1) = 1.

\$\$

\begin{array}{lll} (\text{Y}4\Pi) & u_{tt} = c^2 u_{xx}, & 0 < x < \infty, & 0 < t < \infty, \\ (\Gamma \text{Y}) & u(0, t) = 0, & 0 < t < \infty, \\ (\text{H}\text{Y}) & \begin{cases} u(x, 0) = f(x), \\ u_t(x, 0) = g(x). \end{cases} & 0 < x < \infty, \\ \end{array}

\$\$

\varphi\_i(0) = 0,

\[ J[u] = \int_{0}^{1} \int_{0}^{1} [u^{2} + u^{2}_{y}] dx dy,
\]

\section{Розв'язання задачі для напівнескінченного рядка за допомогою
формули
д'Аламбера}\label{ux440ux43eux437ux432ux44fux437ux430ux43dux43dux44f-ux437ux430ux434ux430ux447ux456-ux434ux43bux44f-ux43dux430ux43fux456ux432ux43dux435ux441ux43aux456ux43dux447ux435ux43dux43dux43eux433ux43e-ux440ux44fux434ux43aux430-ux437ux430-ux434ux43eux43fux43eux43cux43eux433ux43eux44e-ux444ux43eux440ux43cux443ux43bux438-ux434ux430ux43bux430ux43cux431ux435ux440ux430}

Наше завдання --- знайти хвильові рухи напівнескінченної струни, лівий
кінець якої фіксований, за заданих початкових умов. Для розв'язання
задачі (18.6) ми використаємо той самий метод, що й у пошуку загального
розв'язку хвильового рівняння на всій лінії. Заміна загального розв'язку

\[ u_{xx} + u_{yy} = 1
\]

в начальные условия, находим (как в предыд, щей лекции)

(18.7)

\[ J[y] = \int_{0}^{1} [y^{2} + {y'}^{2}] dx, \ y(0) = 0, \ y(1) = 1,
\]

z(x) = (1 - x) y(x),

\$\$

``Теперь перед нами проблема, которой не было в задаче о бесконечной
струне. Решение \$ u\left(x,\t\right) \$ должно быть определено всюду
только внутри первого квадранта \$ (x\textgreater0,,t\textgreater0) \$
плоскости переменных x и t. Значит, мы должны уметь вычислить функции

\$\$

u\_\{xx\} + u\_\{yy\} = 0

\$\$

при всех \$ -\infty \textless{} x-ct \textless{} \infty \$ , \$
\psi(x+ct) \$ при всех \$ 0 \textless{} x+ct \textless{} \infty \$ .

К сожалению, первая из формул (18.7) позволяет вычислять \$
\varphi(x-ct) \$ только при \$ x-ct\geqslant 0 \$ , поскольку в
начальных данных функции f(x) и g(x) определены только для положительных
значений аргумента.

Если \$ x-ct \geqslant 0 \$ , то после подстановки (18.7) в общее
решение

\$\$

J{[}u{]} = \int\emph{\{0\}\^{}\{1\} \int}\{0\}\^{}\{1\}
{[}u\_\{x\}\^{}\{2\} + u\_\{y\}\^{}\{2\}{]} dx dx.

\$\$

получаем

\$\$

u\_\{xx\} + u\_\{yy\} = \sin(\pi x)

Для півобмеженої струни розв'язок зручно записати у формі

\[
u(x,t)=\phi(x-ct)+\psi(x+ct)
=\frac12\bigl[f(x-ct)+f(x+ct)\bigr]+\frac{1}{2c}\int_{x-ct}^{x+ct} g(\xi)\,d\xi.
\]

Якщо \(x \ge ct\), ця формула збігається зі звичайною формулою Даламбера
для нескінченної струни. Якщо ж \(x<ct\), потрібно врахувати граничну
умову \(u(0,t)=0\), тобто відбиття хвилі від закріпленого кінця.

У результаті для півобмеженої струни маємо

\[
u(x,t)=
\begin{cases}
\dfrac12\bigl[f(x-ct)+f(x+ct)\bigr]+\dfrac{1}{2c}\displaystyle\int_{x-ct}^{x+ct} g(\xi)\,d\xi, & x \ge ct, \\[1.2ex]
\dfrac12\bigl[f(x+ct)-f(ct-x)\bigr]+\dfrac{1}{2c}\displaystyle\int_{ct-x}^{x+ct} g(\xi)\,d\xi, & x < ct.
\end{cases}
\]

Цим завершується побудова розв'язку. Його геометрична інтерпретація
подана на рис. 18.6.

\pandocbounded{\includegraphics[keepaspectratio]{images/_page_137_Figure_9.jpeg}}

Рис. 18.6. Интерпретация задачи Коши для полуограниченной струны в
xt-плоскости a---возмущение отражается от границы; \$ \delta \$
---возмущение распространяется вдоль характеристики из точки ( \$
x\_0+ct\_0 \$ , 0). \# ЗАУВАЖЕННЯ 1. Формула для півобмеженої струни
повністю узгоджується з фізичною картиною: при \(x \ge ct\) хвиля ще не
встигла відбитися від межі, а при \(x<ct\) у розв'язку з'являється
відбита хвиля зі зміненим знаком.

\begin{enumerate}
\def\labelenumi{\arabic{enumi}.}
\setcounter{enumi}{1}
\item
  Якщо замінити граничну умову \(u(0,t)=0\) на іншу, то явний вигляд
  розв'язку теж зміниться. Але сама ідея побудови - доозначення
  хвильових складових за допомогою граничної умови - залишається тією
  самою.
\item
  Прямі
\end{enumerate}

\[
x-ct=\text{const}, \qquad x+ct=\text{const},
\]

називаються характеристиками хвильового рівняння. Саме вздовж цих прямих
поширюється збурення.

\section{ЗАВДАННЯ}\label{ux437ux430ux432ux434ux430ux43dux43dux44f-12}

\begin{enumerate}
\def\labelenumi{\arabic{enumi}.}
\item
  Побудуйте графіки розв'язку для напівнескінченної струни в різні
  моменти часу та поясніть, як хвиля відбивається від закріпленої межі.
\item
  Опишіть алгоритм побудови такого розв'язку:

  \begin{itemize}
  \tightlist
  \item
    продовжити початкові дані на всю вісь,
  \item
    побудувати ліву і праву хвилі,
  \item
    врахувати відбиття від межі \(x=0\).
  \end{itemize}
\item
  Розв'яжіть змішану задачу для напівнескінченної струни, коли початкові
  дані задано графічно.
\end{enumerate}

\pandocbounded{\includegraphics[keepaspectratio]{images/_page_139_Figure_3.jpeg}}

\section{\texorpdfstring{Розв'язання нелінійного рівняння
\(\Delta u + u^2 = 0\) методом
збурень}{Розв'язання нелінійного рівняння \textbackslash Delta u + u\^{}2 = 0 методом збурень}}\label{ux440ux43eux437ux432ux44fux437ux430ux43dux43dux44f-ux43dux435ux43bux456ux43dux456ux439ux43dux43eux433ux43e-ux440ux456ux432ux43dux44fux43dux43dux44f-delta-u-u2-0-ux43cux435ux442ux43eux434ux43eux43c-ux437ux431ux443ux440ux435ux43dux44c}

Розглянемо нелінійну задачу Діріхле як збурення простішої лінійної
задачі. Для цього вводять сімейство задач

\[
\Delta u + \varepsilon u^2 = 0, \qquad 0 \le \varepsilon \le 1,
\]

і шукають розв'язок у вигляді ряду

\[
u = u_0 + \varepsilon u_1 + \varepsilon^2 u_2 + \cdots .
\]

Нульове наближення \(u_0\) обирають як розв'язок відповідної лінійної
задачі. Далі, підставляючи ряд у рівняння і прирівнюючи коефіцієнти при
однакових степенях \(\varepsilon\), послідовно знаходять поправки
\(u_1, u_2, \ldots\).

(46.4)

\$\$

Воспользуйтесь методом, аналогичным рассмотренному в лекции. Дайте
физическую интерпретацию этой задачи.

\begin{enumerate}
\def\labelenumi{\arabic{enumi}.}
\setcounter{enumi}{3}
\tightlist
\item
  Предположим, что колебания струны описываются уравнением \$ u\_\{tt\}
  = u\_\{xx\} \$ , начальное смещение изображено на рисунке, а начальная
  скорость равна нулю. Изобразите решение этой задачи в xt-плоскости.
  Обратите внимание на то, что начальные условия в этой задаче разрывны.
\end{enumerate}

\pandocbounded{\includegraphics[keepaspectratio]{images/_page_139_Figure_8.jpeg}}

\section{ХВИЛЬОВЕ РІВНЯННЯ І ГРАНИЧНІ
УМОВИ}\label{ux445ux432ux438ux43bux44cux43eux432ux435-ux440ux456ux432ux43dux44fux43dux43dux44f-ux456-ux433ux440ux430ux43dux438ux447ux43dux456-ux443ux43cux43eux432ux438}

\textbf{Мета лекції.} Ознайомити читача з основними типами граничних
умов для хвильового рівняння на скінченному відрізку.

Розглядатимемо одномірне хвильове рівняння

\[
u_{tt} = c^2 u_{xx}, \qquad 0 < x < L,\quad 0 < t < \infty.
\]

Для нього найчастіше виникають три типи граничних умов:

\begin{enumerate}
\def\labelenumi{\arabic{enumi}.}
\tightlist
\item
  \textbf{Умови першого роду}:
\end{enumerate}

\[
u(0,t)=g_1(t), \qquad u(L,t)=g_2(t).
\]

\begin{enumerate}
\def\labelenumi{\arabic{enumi}.}
\setcounter{enumi}{1}
\tightlist
\item
  \textbf{Умови другого роду}:
\end{enumerate}

\[
u_x(0,t)=g_1(t), \qquad u_x(L,t)=g_2(t).
\]

\begin{enumerate}
\def\labelenumi{\arabic{enumi}.}
\setcounter{enumi}{2}
\tightlist
\item
  \textbf{Умови третього роду}:
\end{enumerate}

\[
\alpha_1 u_x(0,t)+\beta_1 u(0,t)=g_1(t), \qquad
\alpha_2 u_x(L,t)+\beta_2 u(L,t)=g_2(t).
\]

У фізичному сенсі це відповідає заданому режиму на межі, заданій силі
або пружному закріпленню кінців струни.

Далі ми почнемо з умов першого роду і розглянемо задачу

\[
\begin{aligned}
u_{tt} &= c^2u_{xx}, && 0 < x < L,\quad 0 < t < \infty, \\
u(0,t) &= g_1(t), \quad u(L,t) = g_2(t), && 0 < t < \infty, \\
u(x,0) &= f(x), \quad u_t(x,0) = g(x), && 0 \leqslant x \leqslant L.
\end{aligned}
\]

Щоб розв'язати таку неоднорідну задачу, зручно подати розв'язок як суму
спеціально підібраної функції, що задовольняє граничні умови, і нової
невідомої з однорідними граничними умовами. Саме цей підхід і буде далі
використано.

Правильні частини форми \(r^n\) , \(r^n \cos{(n\theta)}\) ,
\(r^n \sin{(n\theta)}\) ведуть до рішень \(Ar^{n+2}\) ,
\(Br^{n+2} \cos{(n\theta)}\) \(Cr^{n+2} \sin{(n\theta)}\) відповідно.

Підставляючи функції \(u_p(r, \theta)\) у гетерогенне рівняння,
отримуємо A = -1/32, B = -1/24.

Отже,

\$\$

u\_\{x\}(1, t) = \frac{1}{k} v(t) - (k - \text{модуль Юнга}).

\$\$

Щоб розв'язати задачу, \(P_1\) має зробити останній крок --- підставити
загальне розв'язання \(u(r, \theta) = u_h(r, \theta) + u_p(r, \theta)\)
на граничні умови \(u(1, \theta) = 0\) і знайти коефіцієнти \(a_n\) і
\(b_n\) .

Після цього ми отримуємо

\$\$

u\_\{x\}(0, t) = \frac{h}{T} u(0, t),

\$\$

Отже, \(a_0=1/32\) , \(a_2=1/24\) , і всі інші \(a_n\) і \(b_n\)
дорівнюють нулю. Отже, рішення проблеми \(P_1\) знайдено:

\$\$

u\_\{x\}(L, t) = -\frac{h}{T} u(L, t).

\$\$

Функція \(u_1(r, \theta)\) називається пертурбацією першого порядку до
\(u_4(r, \theta)\) . Додавши \(u_0\) і \(u_1\) , отримуємо новий підхід
до розв'язання задачі (46.2):

(46.5)

\$\$

u\_\{x\}(0, t) - \frac{h}{T}u(0, t) = 0,

\$\$

Щоб знайти збурення наступного порядку \(u_{z}(r, \theta)\) , потрібно
розв'язати задачу

\$\$

u\_\{x\}(L, t) + \frac{h}{T}u(L, t) = 0.

\$\$
\(\begin{cases} \Delta u_{1} = -2u_{0}u_{1}, \\ u_{2}(1, \theta) = 0, \end{cases}\)
де \(u_{\phi}(r,\theta)\) і \(u_{1}(r,0)\) вже відомі, раніше знайдені
функції. Звісно, виконувати всі ці алгебраїчні перетворення без допомоги
комп'ютера надзвичайно складно. На щастя, вираз (46.5) є майже точним
рішенням нашої проблеми. Насправді, якщо підставити цей вираз зліва від
\(\Delta u + u^{2}\) нашого нелінійного рівняння, отримаємо вираз, який
майже всюди відрізняється від нуля в межах кола 0 \textless{} r
\textless{} 1.

Окрім розв'язання нелінійних задач, теорія збурень може
використовуватися для розв'язання задач у областях з неправильними
межами (звісно, якщо нерегулярність не дуже велика). Давайте розглянемо
простий приклад.

\section{Приклад задачі з порушеною
межею}\label{ux43fux440ux438ux43aux43bux430ux434-ux437ux430ux434ux430ux447ux456-ux437-ux43fux43eux440ux443ux448ux435ux43dux43eux44e-ux43cux435ux436ux435ux44e}

Осквернення можна знищити, не позбавляючи його права. Можна знайти
розвиток рівня Лаплас посередині деформованого стовпа, оскільки пробурив
рівень Лаплас у центрі круглої області. У застосуванні дозволено знайти
потенціал у квадраті за кривою \(r=1+\frac{1}{4}\sin\theta\) (за
полярними координатами). Нехай потенціал у на межі заданий. Тобто, нам
потрібно вирішити проблему

(46.6)

\[ u \left( 1 + \frac{1}{4} \sin \theta, \ \theta \right) = \cos \theta, \quad 0 \le \theta \le 2\pi.\]

\begin{aligned} u_x(0, T) &= \frac{h}{T} [u(0, t) - \theta_1(t)], \\ u_x(L, T) &= -\frac{h}{T} [u(L, t) - \theta_2(t)]. \end{aligned}

\[u\left(1+\frac{1}{4}\sin\theta,\ \theta\right)=\cos\theta\]

u(1+\epsilon\sin\theta, \theta) = \cos\theta, \quad 0
\le \epsilon \le 1/4.

\$\$

Этим мы завершаем обсуждение наиболее важных граничных условий,
связанных с гиперболическими задачами. Следующие несколько лекций мы
посвятим решению задач, содержащих полобные граничные условия. \#
ЗАУВАЖЕННЯ 1. В лекции не рассматривалось граничное условие еще одного
типа, когда струна подвергается действию силы, пропорциональной скорости
и направленной в противоположном направ-

лении. Граничное условие такого типа (на левом конце) записывается в
виде

\$\$

f(x+h) = f(x) + f'(x)h + f'\,'(x)\frac{h^2}{2!} + \dots,

\$\$

\begin{enumerate}
\def\labelenumi{\arabic{enumi}.}
\setcounter{enumi}{1}
\tightlist
\item
  Нелинейное упругое закрепление левого конца приводит к граничному
  условию вида
\end{enumerate}

\$\$

u(1+\varepsilon\sin\theta,~\theta)=u(1,~\theta)+u\_r(1,~\theta)(\varepsilon\sin\theta)+u\_\{rr\}(1,~\theta)\frac{(\varepsilon\sin\theta)^2}{21}+\ldots

\$\$

где \$ \varphi(u) \$ --- произвольная функция величины u. Например,
граничное условие

\$\$

\Delta u = 0, \quad 0 \textless{} r \textless{} 1 +
\varepsilon \sin \theta,

\$\$

говорит о том, что возвращающая сила на левом конце пропорциональна кубу
смещения (а не первой степени и, как это будет согласно закону Гука в
линейном случае).

\begin{enumerate}
\def\labelenumi{\arabic{enumi}.}
\setcounter{enumi}{2}
\tightlist
\item
  Если к нижнему концу пружины, совершающей продольные колебания,
  прикрепить массу \emph{m}, то граничное условие следует задавать в
  виде
\end{enumerate}

\$\$

Коли всі розрахунки будуть виконані, ми внесемо \(\varepsilon=1/4\) у
кінцеві результати. Звісно, здається, що цю задачу не можна легко
розв'язати, але її можна звести до послідовності задач, у кожній з яких
визначена одна з функцій \(u_{\mathrm{ }e}\) , \(u_{\mathrm{ }e}\) ,
\ldots, що включені у розклад

\$\$ \# ЗАДАЧІ Опираясь на интуитивные представления о граничных
условиях различных типов, изобразить в общих чертах решение задачи

\$\$

Після підстави ряду в (467) отримуємо таку послідовність задач для
визначення функцій \(u_a\) , \(u_1\) , \ldots;

\[ \begin{cases} \Delta u_{0} = 0, & 0 < r < 1 \text{ (внутри круга),} \\ u_{0}(1, \theta) = \cos \theta, & u_{0}(r, \theta) = r \cos \theta, \end{cases}
\]

\begin{array}{ll} (\mathrm{Y} \, \Pi) & u_{tt} = c^2 u_{xx}, \quad 0 < x < 1, \quad 0 < t < \infty, \\ (\mathrm{\Gamma} \, \mathrm{Y}) & \left\{ \begin{array}{ll} u \, (0, \, t) = 0, \\ u \, (1, \, t) = \sin t, \end{array}

$$ \begin{cases} \Delta u_{1} = 0, & 0 < r < 1 \text{ (внутри круга),} \\ u_{1}(1, \theta) = -\sin \theta \frac{\partial u_{0}(1, \theta)}{\partial r} = -\sin \theta \cos \theta \end{cases}
$$

u = u_0 + \frac{1}{4}u_1 + \left(\frac{1}{4}\right)^2 u_1 + \cdots

$$ u_u + \frac{1}{4} u_t$$

\begin{array}{ll} u \, (x, \, 0) = 0, \\ u_{t} \, (x, \, 0) = 0, \end{array}

$$u_t = (1+x) u_{xx}, \\ u(x, 0) = q(x), \quad -\infty < x < \infty,$$

\right. & 0 \leqslant x \leqslant 1, \end{array}

\$\$ Якщо розглянути рівняння з параметром

(46.11) \$\$

u\_t = (1 + \varepsilon x) u\_\{xx\},

\$\$

і шукати його розв'язок у вигляді

\$\$

u = u\_0 + \varepsilon u\_1 + \varepsilon\^{}2 u\_2 + \cdots,

\$\$

то після підстановки отримуємо послідовність задач

\$\$

P\_0:~

\begin{cases}
\dfrac{\partial u_0}{\partial t} = \dfrac{\partial^2 u_0}{\partial x^2}, \\
u_0(x,0) = \varphi(x),
\end{cases}

\qquad P\_1:~

\begin{cases}
\dfrac{\partial u_1}{\partial t} - \dfrac{\partial^2 u_1}{\partial x^2} = x\,\dfrac{\partial^2 u_0}{\partial x^2}, \\
u_1(x,0) = 0.
\end{cases}

\$\$

Зауважимо, що в перших двох задачах коефіцієнти сталі. Параметр
\(\varepsilon\) має бути достатньо малим, інакше нескінченний ряд може
виявитися розбіжним.

\subsection{ЗАДАЧИ}\label{ux437ux430ux434ux430ux447ux438}

\begin{enumerate}
\def\labelenumi{\arabic{enumi}.}
\item
  Подставьте разложение (46.8) в задачу (46.7) и получите
  последовательность задач \(P_0,\ P_1,\ P_2,\ \dots\).
\item
  Покажите, что нелинейную задачу
\end{enumerate}

\$\$

\begin{aligned}
\Delta u + u^2 &= 0, && 0 \leqslant r < 1,\ 0 \leqslant \theta < 2\pi, \\
u(1,\theta) &= \cos\theta, && 0 \leqslant \theta < 2\pi,
\end{aligned}

\$\$

можно свести к последовательности линейных задач \(P_0,\ P_1,\ P_2\).

\begin{enumerate}
\def\labelenumi{\arabic{enumi}.}
\setcounter{enumi}{2}
\item
  Подставьте приближенное решение (46.5) в нелинейную задачу 2 и оцените
  его точность.
\item
  Решите задачу \(P_1\) из примера с возмущенной границей и проверьте,
  насколько хорошо приближенное решение
\end{enumerate}

\$\$

u(r,\theta) = u\_0(r,\theta) + \frac{1}{4}u\_1(r,\theta)

\$\$

удовлетворяет соотношениям

\$\$

\Delta u = 0,

\[ u\left(1 + \frac{1}{4}\sin\theta,\ \theta\right) = \cos\theta.
\]

\section{Лекція 47}\label{ux43bux435ux43aux446ux456ux44f-47}

\section{РОЗВ'ЯЗАННЯ ДИФЕРЕНЦІАЛЬНИХ РІВНЯНЬ У ПОХІДНИХ ПОХІДНИХ ЗА
ДОПОМОГОЮ МЕТОДУ КОНФОРМНИХ
ВІДОБРАЖЕНЬ}\label{ux440ux43eux437ux432ux44fux437ux430ux43dux43dux44f-ux434ux438ux444ux435ux440ux435ux43dux446ux456ux430ux43bux44cux43dux438ux445-ux440ux456ux432ux43dux44fux43dux44c-ux443-ux43fux43eux445ux456ux434ux43dux438ux445-ux43fux43eux445ux456ux434ux43dux438ux445-ux437ux430-ux434ux43eux43fux43eux43cux43eux433ux43eux44e-ux43cux435ux442ux43eux434ux443-ux43aux43eux43dux444ux43eux440ux43cux43dux438ux445-ux432ux456ux434ux43eux431ux440ux430ux436ux435ux43dux44c}

МЕТА ЛЕКЦІЇ: Показати, як деякі двовимірні граничні задачі можна
перетворити на інші, простіші задачі за допомогою конформного
відображення. Припустимо, наприклад, що ми хочемо розв'язати рівняння
Лапласа

\$\$

\begin{cases} u(0, t) = 0, \\ u_x(1, t) = -u(1, t), \\ u(x, 0) = x, \\ u_t(x, 0) = 0. \end{cases}

\$\$

з певною граничною умовою в області комплексної форми (площина змінних x
і y). Цю задачу можна перетворити на нову, в якій потрібно знайти
розв'язок рівняння Лапласа

\$\$ у простішій області (площина змінних і і і).

\pandocbounded{\includegraphics[keepaspectratio]{images/_page_147_Picture_5.jpeg}}

. 20.1. Три типові стоячі хвилі: X(x) T(t).

Декілька прикладів стоячих хвиль показано на рис. 20.1. Якщо профілі цих
стоячих хвиль відомі \$ X\_n(x) \$ і природа коливань кожної хвилі \$
T\_n(t) \$ , то розв'язок проблеми коливань гітарних струн є
суперпозицією найпростіших коливань \$ X\_n(x),T\_n(t) \$ і

сывается в виде

\$\$

= \ln\left[\frac{z-1}{z+1}\right] + i \arg\left[\frac{z-1}{z+1}\right]

\$\$

Залишилося лише обрати шанси \$ c\_n \$ так, щоб при t = 0 початкові
умови \$ u\left(x,~0\right)=f\left(x\right) \$ і \$
u\_t\left(x,~0\right)=g\left(x\right) \$ були виконані.

Задачу о колебаниях гитарной струны мы решим методом разделения
переменных.

\section{Розв'язання проблеми коливань обмеженої струни методом
розділення
змінних}\label{ux440ux43eux437ux432ux44fux437ux430ux43dux43dux44f-ux43fux440ux43eux431ux43bux435ux43cux438-ux43aux43eux43bux438ux432ux430ux43dux44c-ux43eux431ux43cux435ux436ux435ux43dux43eux457-ux441ux442ux440ux443ux43dux438-ux43cux435ux442ux43eux434ux43eux43c-ux440ux43eux437ux434ux456ux43bux435ux43dux43dux44f-ux437ux43cux456ux43dux43dux438ux445}

Чтобы решить смешанную задачу

(УЧП)

\$\$

v = \arg\left[\frac{z-1}{z+1}\right] =
\arg\left[\frac{x+iy-1}{x+1+iy}\right] =

\$\$

= \arg\left[\frac{2y}{x^2+y^2-1}\right]

\$\$

\varphi(x, y) = \frac{1}{\pi}
\arctan\left[\frac{2y}{x^2 + y^2 - 1}\right].

\$\$

\begin{cases} u(0, t) = 0, \\ u(L, t) = 0, \end{cases}

\$\$

x\^{}\{2\} + y\^{}\{2\} = 1,

\$\$

якщо потенціал на внутрішньому колі дорівнює одинці, а на зовнішньому
колі дорівнює двом (рис. 47.5). Іншими словами, нам потрібно розв'язати
проблему

(47.3)

\$\$

Проблема полягає в тому, щоб знайти конформне відображення, яке
перетворює цю область на щось простіше, де нашу проблему легко
розв'язати. У цьому випадку очевидно, що ця проста ділянка має бути
кільцем. Відповідна експозиція Р И С У Н О К

\pandocbounded{\includegraphics[keepaspectratio]{images/_page_149_Figure_5.jpeg}}

20.2. Построение стоячих воли при раздичных значениях параметра λ.

Вода та граничні умови. Залишається лише знайти таку лінійну комбінацію
цих коливань, щоб сума задовольняла початкові умови при t = 0.

Подстановка (20.2) в граничные условия \(u(0, t) = u(L, t) = 0\) дает

\[
u(0, t) = X(0)T(t) = D\bigl[A\sin(\alpha\beta t) + B\cos(\alpha\beta t)\bigr] = 0 \Rightarrow D = 0,
\]

и

\[
u(L, t) = X(L)T(t) = C\sin(\beta L)\bigl[A\sin(\alpha\beta t) + B\cos(\alpha\beta t)\bigr] = 0 \Rightarrow \sin(\beta L) = 0.
\]

Другими словами, константа разделения \(\beta\) должна удовлетворять
уравнению \(\sin(\beta L) = 0\) , откуда

\[
\beta_n = \frac{n\pi}{L}, \qquad n = 0, 1, 2, \ldots
\]

Зверніть увагу, що якщо поставити \(C=0\) у друге рівняння (20,3),
отримаємо тривіальне розв'язання \(X(x)T(t) \equiv 0\) . Відповідно, ми
знайшли послідовність елементарних коливань струни:

\[
u_n(x, t) = X_n(x)T_n(t) = \sin\!\left(\frac{n\pi x}{L}\right)\left[a_n \sin\!\left(\frac{n\pi \alpha t}{L}\right) + b_n \cos\!\left(\frac{n\pi \alpha t}{L}\right)\right].
\]

(20.4) або

\[
u_n(x, t) = R_n \sin\!\left(\frac{n\pi x}{L}\right)\cos\!\left(\frac{n\pi \alpha (t - \delta_n)}{L}\right).
\]

Тут \(a_n\) , \(b_n\) , \(R_n\) і \(\delta_n\) --- довільні сталі. Кожне
таке елементарне коливання є стоячою хвилею.

\pandocbounded{\includegraphics[keepaspectratio]{images/_page_114_Figure_4.jpeg}}

Фіг. 15.3 Конвективна хвиля a --- попереду фронту домішок \$ \delta \$
--- одна одиниця попереду фронту хвиль

Зверніть увагу, що у задачі (15.2) ми перемістили межу до - \$ \infty \$
(тепер розв'язується задача Коші), щоб вона більше не впливала на оцінку
розв'язку. Для розв'язання задачі (15.2) можна використати перетворення
Лапласа P і c

\pandocbounded{\includegraphics[keepaspectratio]{images/_page_114_Figure_7.jpeg}}

. 15 4 Начальные условия в задаче конвективной диффузии.

Однак у цьому випадку набагато цікавіше зробити щось зовсім інше.
Давайте представимо нову систему координат, яка рухається вздовж старої
зі швидкістю V. Іншими словами, замість системи координат, прив'язаних
до берега річки, ми розглянемо систему координат, що рухається зі
швидкістю фронту домішок (звісно, якщо є дифузія, фронт буде
«розмитий»). Математично це означає, що ми замінюємо просторову
координату x на нову \$ \xi = x - Vt \$ .

Теперь ясно, что

когда \$ \xi = 0 \$ , мы находимся па фронте распространяющейся примеси.

когда \$ \xi = 1 \$ , мы находимся на одну единицу длины впереди фронта.

когда \$ \xi = -1 \$ , мы находимся на одну единицу длины позади фронта.

Наша задача --- преобразовать исходную задачу с НУ

\$\$ \overline{D} = A\xi\emph{\{xx\} + B\xi}\{xy\} + C\xi\emph{\{yy\} +
D\xi}\{x\} + E\xi\_\{y\},

\$\$

{[}\xi\_x/\xi\_y{]} = -B/2A

\$\$ \xi(x, y) = c

Для параболічного випадку достатньо однієї характеристики. Її записують
як

\[
\frac{dy}{dx} = -\frac{B}{2A}.
\]

Якщо характеристика побудована, то другу змінну можна вибрати довільно,
наприклад \(\eta = y\), за умови що вона не пропорційна \(\xi\). Після
цього початкове рівняння приводиться до канонічної форми, а коефіцієнти
при старших похідних спрощуються до одного ненульового члена.

Замість того щоб зберігати тут розірвані проміжні перетворення,
достатньо пам'ятати головний результат: правильний вибір координат
перетворює параболічну задачу на рівняння з однією канонічною похідною,
яке далі розв'язується стандартними методами.

\section{Приклад}\label{ux43fux440ux438ux43aux43bux430ux434}

Для рівняння

\[
u_{xx} + 2u_{xy} + u_{yy} = 0
\]

характеристичне рівняння має вигляд

\[
\frac{dy}{dx} = -1.
\]

Тому можна взяти нові координати

\[
\xi = y - x, \qquad \eta = y.
\]

У цих координатах початкове рівняння переходить у просту канонічну форму

\[
u_{\eta\eta} = 0.
\]

Це означає, що загальний розв'язок має вигляд

\[
u(\xi,\eta) = \eta f(\xi) + g(\xi),
\]

де \(f\) та \(g\) - довільні функції. \# ЗАВДАННЯ

\begin{enumerate}
\def\labelenumi{\arabic{enumi}.}
\item
  Розв'яжіть задачу Коші методом характеристик для лінійного рівняння
  першого порядку.
\item
  Зведіть рівняння
\end{enumerate}

\[
u_{xx} + 2u_{xy} + u_{yy} + u = 2
\]

до канонічної форми.

\begin{enumerate}
\def\labelenumi{\arabic{enumi}.}
\setcounter{enumi}{2}
\tightlist
\item
  Зведіть еліптичне рівняння
\end{enumerate}

\[
u_{xx} + 2u_{yy} + x^2u_x = e^{-x^2/2}
\]

до канонічної форми.

\begin{enumerate}
\def\labelenumi{\arabic{enumi}.}
\setcounter{enumi}{3}
\item
  Для хвильового рівняння запишіть граничні умови першого, другого і
  третього роду та поясніть, як вони впливають на вигляд розв'язку.
\item
  Знайдіть розв'язок задачі переносу методом переходу до рухомої
  координати.
\end{enumerate}

На цьому блок задач завершується. Далі в рукописі починається наступна
лекція, але в цьому chunk вона вже не потрібна для побудови коректного
PDF.

Этим завершается наша интерпретация формулы Даламбера в плоскости
переменных \$ \boldsymbol{x} \$ и t. В оставшейся части лекции мы
займемся решением смешанной задачи для полубесконечной струны:

(18.6)

\$\$

u = \gamma [(x-s)(x-t)-\gamma^2],

\$\$

Единственное отличие в том, что аппроксимирующие функции \$ \phi\_1,
\phi\_2, \phi\_3, \ldots, \phi\_n \$ должны удовлетворять граничным
условиям:

\$\$

\gamma = 2t/{[}(x-t)\^{}2 + y\^{}2{]}.

\$\$

является иравнение Лапласа, но это доказательство проводится аналогично
тому, как это делалось в лекции 44 для функционала, зависящего от
функции одной переменной.

\begin{enumerate}
\def\labelenumi{\arabic{enumi}.}
\setcounter{enumi}{2}
\item
  В книгах по вариационному исчислению показано, как строить
  энергетический функционал не только для уравнения Пуассона с граничным
  условием Дирихле u=0, но и для многих других дифференциальных
  уравнений и многих других типов граничных условий. Следовательно,
  большое число краевых задач можно решить, минимизируя соответствующие
  энергетические функционалы.
\item
  Чем большее значение п мы выбираем, тем меньшее значение \$
  J{[}\bar\{u\}\_n{]} \$ получается. Значит, с увеличением n растет
  точность л приближенного решения дифференциального уравнения. Один из
  методов определения и, основан на последовательном вычислении
  функционала \$ J{[}u\_n{]} \$ для возрастающих значений n.~Вычисления
  прекращаются, если с ростом п функционал перестает практически
  уменьшаться.
\item
  Если и велико, то проведение вычислений по методу Ритца требует
  привлечения ЭВМ. На рис. 45.2 (с. 345) приведена блок-
\end{enumerate}

схема решения краевой задачи методом Ритца.

Пользователь должен позаботиться о подпрограмме, которая должна
вычислять значения пробных функций. Такая подпрограмма может иметь вид:

SUBROUTINE BC (X.Y.PHI)

\begin{itemize}
\tightlist
\item
  SUBROUTINE PROVIDED BY THE USER TO C - EVALUATE THE - FUNCTIONS
  PHI(1), PHI(2), \ldots, PHI(N) DIMENSION PHI(20)
\end{itemize}

PHI(1) = X*Y*(1-X)*(1-Y)

\$ PHI(2) = X\^{}*PHI(1) \$ (используемые функции)

PHI(3) = Y*PHI(1)

PHI(N) = (что получится)RETURN END

\section{\texorpdfstring{\textbf{ЗА}ДАЧИ}{ЗАДАЧИ}}\label{ux437ux430ux434ux430ux447ux438-1}

Найдите энергетический функционал для задачи

(УЧП)

\$\$

Щоб завершити перетворення, потрібно так само перенести початкові умови
на нову невідому функцію. Після цього задача для хвильового рівняння з
неоднорідними граничними умовами зводиться до задачі з однорідними
граничними умовами, але вже з модифікованими правими частинами.

На практиці це означає, що складність переноситься з межі в саме
рівняння, зате далі можна застосовувати стандартні методи: розділення
змінних, розкладання за власними функціями або інтегральні перетворення.

\pandocbounded{\includegraphics[keepaspectratio]{images/_page_144_Picture_4.jpeg}}

Рис. 19.6. Виникнення неоднорідних граничних умов пружного закріплення:
\(h\) - коефіцієнт жорсткості пружини.

Для умов пружного закріплення однорідні граничні умови часто записують у
вигляді

\[
u_x(0,t) - h_1u(0,t) = 0, \qquad u_x(L,t) + h_2u(L,t) = 0.
\]

Саме з такого запису зазвичай починають побудову власних функцій і
відповідних спектральних розкладів для подальшого розв'язання задачі.

\pandocbounded{\includegraphics[keepaspectratio]{images/_page_150_Picture_5.jpeg}}

Фіг. 20.3. Стоячи \(u_n(x, t) = X_n(x)T_n(t)\) .

У нашій задачі рівняння і граничні умови є лінійними та однорідними,
тому довільна сума елементарних стоячих хвиль також є розв'язком
хвильового рівняння і задовольняє граничним умовам. Тепер ми повинні
знайти таку суперпозицію стоячих хвиль, щоб вона також задовольняла
початкові умови, і тоді отримаємо розв'язок нашої проблеми. Заміна
кількості

\$\$

Зверніть увагу на зображення.

\pandocbounded{\includegraphics[keepaspectratio]{images/_page_154_Figure_5.jpeg}}

Головна різниця між поперечними коливаннями тонкої балки та поперечними
коливаннями струни полягає в тому, що промінь чинить опір згину. Не
заглиблюючись у механіку тонких балок, припустимо, що, враховуючи опір
згинання (замість рівняння хвилі) веде до рівняння четвертого порядку

(21.1)

\$\$

\section{Шарнірно оперта
балка}\label{ux448ux430ux440ux43dux456ux440ux43dux43e-ux43eux43fux435ux440ux442ux430-ux431ux430ux43bux43aux430}

Розглянемо малі коливання тонкої балки, кінці якої вільно спираються на
дві опори. Це означає, що кінці балки не зміщуються, але можуть вільно
повертатися.

Тому на кінцях балки мають виконуватися граничні умови

\[
u(0, t) = 0, \qquad u(1, t) = 0,
\]

а також умови нульового згинального моменту

\[
u_{xx}(0, t) = 0, \qquad u_{xx}(1, t) = 0.
\]

Із теорії тонких балок випливає, що згинальний момент пропорційний
\(u_{xx}\), тому для вільно опертих кінців саме ці умови є природними.

\pandocbounded{\includegraphics[keepaspectratio]{images/_page_155_Figure_8.jpeg}}

Рис. 21.1. Шарнірно оперта балка.

Отже, коливання балки описуються такою змішаною задачею (для простоти
відповідний коефіцієнт вважаємо рівним одиниці):

(21.2)

\[
\begin{aligned}
(\text{УЧП})\quad & u_{tt} = -u_{xxxx}, && 0 < x < 1,\ 0 < t < \infty, \\
(\text{ГУ})\quad & u(0,t) = 0, \quad u_{xx}(0,t) = 0, \quad u(1,t) = 0, \quad u_{xx}(1,t) = 0, && 0 < t < \infty, \\
(\text{НУ})\quad & u(x,0) = f(x), \quad u_t(x,0) = g(x), && 0 \leqslant x \leqslant 1.
\end{aligned}
\]

\$\$

Щоб розв'язати цю задачу, ми використаємо метод розділення змінних. Ми
будемо шукати лише періодичні розв'язки, тобто флуктуації

(21.3)

\[ [2/\pi\omega]^{1/2}\]

Зверніть увагу, що вибір розв'язку вигляду (21.3) фактично означає, що
ми взяли константу розділення у методі розділення змінних як від'ємну.

Подставим теперь (21.3) в уравнение колебаний балки и получим уравнение
для \$ X\left(x\right) \$ \[\frac{2}{\pi\omega} [1-\cos(\omega a)]\]

общее решение которого записывается в виде

\[1\]

Щоб визначити коефіцієнти C, D, E та F, підставляємо загальний розв'язок
у граничні умови

\[e^{-1}a + \omega\]

\frac{1}{2}e\^{}\{-\alpha\}\sin\omega

\$\$

\begin{split} u\left(0,\ t\right) &= 0 \Rightarrow X\left(0\right)T\left(t\right) = 0 \Rightarrow X\left(0\right) = 0 \Rightarrow C + E = 0 \\ u_{xx}\left(0,\ t\right) &= 0 \Rightarrow X''\left(0\right)T\left(t\right) = 0 \Rightarrow X''\left(0\right) = 0 \Rightarrow -C + E = 0 \\ u\left(1,\ t\right) &= 0 \Rightarrow D\sin\sqrt{\omega} + F\sin\sqrt{\omega} = 0, \\ u_{xx}\left(1,\ t\right) &= 0 \Rightarrow -D\sin\sqrt{\omega} + F\sin\sqrt{\omega} = 0. \end{split}

\$\$

\frac{1-e^{-a\omega}}{\omega}

\[ \frac{2}{\pi}F''(\omega)\]

F \sinh \sqrt{\omega} = 0,

\[f(x) = \int_{0}^{\infty} F(\omega) \cos(\omega x) d\omega \qquad \qquad F(x) = \frac{2}{\pi} \int_{0}^{\infty} f(x) \cos(\omega x) dx\]

D \sin \sqrt{\omega} = 0.

\$\$

\begin{enumerate}
\def\labelenumi{\arabic{enumi}.}
\tightlist
\item
\end{enumerate}

\$\$

F = 0

\$\$

\section{ПЕРЕХІД ДО БЕЗРОЗМІРНИХ
ЗМІННИХ}\label{ux43fux435ux440ux435ux445ux456ux434-ux434ux43e-ux431ux435ux437ux440ux43eux437ux43cux456ux440ux43dux438ux445-ux437ux43cux456ux43dux43dux438ux445}

МЕТА ЛЕКЦІЇ: Показати, яким чином крайові, початкові та інші фізичні
задачі можна записувати у безрозмірному вигляді. У різних галузях знань
рівняння формулюються по-різному, але після переходу до безрозмірних
змінних їх часто можна звести до спільної математичної форми.

Переход к безразмерным переменным позволяет привести уравнения к одному
и тому же виду. Именно по этой причине при изучении уравнений с частными
производными в математике стараются отвлекаться от физического смысла
входящих в них параметров. Вот почему химик, физик или биолог должен
преобразовать свое уравнение к той форме, которая принята в нашей книге.

Основная идея введения новых (безразмерных) переменных состоит в том,
что после перехода к безразмерным переменным вадача становится чисто
математической и уже не содержит

карактерных физических констант. Именно таким способом многие различные
уравнения физики, биологии и химии, содержащие всякие нюансы, связанные
с физическими параметрами, приводятся к одной и той же простой форме
(см. рис. 22.1).

Познакомимся с этим преобразованием на простом примере.

\pandocbounded{\includegraphics[keepaspectratio]{images/_page_161_Figure_4.jpeg}}

Рис. 22.1. Некоторые задачи, приводящиеся к одной и той же безразмерной
форме.

\section{Приведення дифузійної задачі до безрозмірного
вигляду}\label{ux43fux440ux438ux432ux435ux434ux435ux43dux43dux44f-ux434ux438ux444ux443ux437ux456ux439ux43dux43eux457-ux437ux430ux434ux430ux447ux456-ux434ux43e-ux431ux435ux437ux440ux43eux437ux43cux456ux440ux43dux43eux433ux43e-ux432ux438ux433ux43bux44fux434ux443}

Начнем со смешанной задачи для отрезка с начальным условием на
температуру \(u\left(x,\,0\right)=\sin\left(\pi x/L\right)\) и
граничными температурами на концах \(T_1\) и \(T_2\) . Другими словами,
мы решаем задачу

(22.1) \$\$

\begin{array}{ccccccccccccccccccccccccccccccccccc

$$

(см. рис. 22.2).

![](images/_page_161_Figure_10.jpeg)

Рис. 22.2. Задача теплопроводности в плоскости переменных (х, т).

Наша цель — дать новую эквивалентную (22.1) постановку задачи таким образом, чтобы

в новой задаче не было никаких физических параметров (вроде α);

2) начальные и граничные условия стали проще.

Для осуществления этой цели, мы введем три новых безразмерных переменных U,  $ \xi $  и  $ \tau $ , чтобы заменить прежние u, x и t по схеме

$$

u \to U

$$

(безразмерная температура),  $ x \to \xi $  (безразмерная длина),  $ t \to \tau $  (безразмерное время).

Мы проведем все три преобразования одновременно.

# Перетворення залежної змінної $ u \rightarrow U $

Определим функцию U(x, t) по формуле

$$

U(x, t) = \frac{u(x, t) - T_1}{T_2 - T_1}

$$.

Ясно, что новая температура  $ U\left(x,\,t\right) $  безразмерна, поскольку мы делим градусы Цельсия на градусы Цельсия. Новые граничные условия для  $ U\left(x,\,t\right) $  при x=0 и x=L примут вид  $ U\left(0,\,t\right)=0 $  и  $ U\left(L,\,t\right)=1 $ . Давайте поставим задачу для новой функции  $ U\left(x,\,t\right) $ . После несложных преобразований задача (22.1) переходит в новую задачу

задачу

$$

\begin{array}{ll} \text{(УЧП)} & U_t = \alpha^2 U_{xx}, \quad 0 < x < L, \quad 0 < t < \infty, \\ U_t = 0, \quad 0 < t < \infty, \\ U_t = 0, \quad 0 < t < \infty, \\ U_t = 0, \quad 0 < t < \infty, \\ U_t = 0, \quad 0 < t < \infty, \\ U_t = 0, \quad 0 < t < \infty, \\ U_t = 0, \quad 0 < t < \infty, \\ U_t = 0, \quad 0 < t < \infty, \\ U_t = 0, \quad 0 < t < \infty, \\ U_t = 0, \quad 0 < t < \infty, \\ U_t = 0, \quad 0 < t < \infty, \\ U_t = 0, \quad 0 < t < \infty, \\ U_t = 0, \quad 0 < t < \infty, \\ U_t = 0, \quad 0 < t < \infty, \\ U_t = 0, \quad 0 < t < \infty, \\ U_t = 0, \quad 0 < t < \infty, \\ U_t = 0, \quad 0 < t < \infty, \\ U_t = 0, \quad 0 < t < \infty, \\ U_t = 0, \quad 0 < t < \infty, \\ U_t = 0, \quad 0 < t < \infty, \\ U_t = 0, \quad 0 < t < \infty, \\ U_t = 0, \quad 0 < t < \infty, \\ U_t = 0, \quad 0 < t < \infty, \\ U_t = 0, \quad 0 < t < \infty, \\ U_t = 0, \quad 0 < t < \infty, \\ U_t = 0, \quad 0 < t < \infty, \\ U_t = 0, \quad 0 < t < \infty, \\ U_t = 0, \quad 0 < t < \infty, \\ U_t = 0, \quad 0 < t < \infty, \\ U_t = 0, \quad 0 < t < \infty, \\ U_t = 0, \quad 0 < t < \infty, \\ U_t = 0, \quad 0 < t < \infty, \\ U_t = 0, \quad 0 < t < \infty, \\ U_t = 0, \quad 0 < t < \infty, \\ U_t = 0, \quad 0 < t < \infty, \\ U_t = 0, \quad 0 < t < \infty, \\ U_t = 0, \quad 0 < t < \infty, \\ U_t = 0, \quad 0 < t < \infty, \\ U_t = 0, \quad 0 < t < \infty, \\ U_t = 0, \quad 0 < t < \infty, \\ U_t = 0, \quad 0 < t < \infty, \\ U_t = 0, \quad 0 < t < \infty, \\ U_t = 0, \quad 0 < t < \infty, \\ U_t = 0, \quad 0 < t < \infty, \\ U_t = 0, \quad 0 < t < \infty, \\ U_t = 0, \quad 0 < t < \infty, \\ U_t = 0, \quad 0 < t < \infty, \\ U_t = 0, \quad 0 < t < \infty, \\ U_t = 0, \quad 0 < t < \infty, \\ U_t = 0, \quad 0 < t < \infty, \\ U_t = 0, \quad 0 < t < \infty, \\ U_t = 0, \quad 0 < t < \infty, \\ U_t = 0, \quad 0 < t < \infty, \\ U_t = 0, \quad 0 < t < \infty, \\ U_t = 0, \quad 0 < t < \infty, \\ U_t = 0, \quad 0 < t < \infty, \\ U_t = 0, \quad 0 < t < \infty, \\ U_t = 0, \quad 0 < t < \infty, \\ U_t = 0, \quad 0 < t < \infty, \\ U_t = 0, \quad 0 < t < \infty, \\ U_t = 0, \quad 0 < t < \infty, \\ U_t = 0, \quad 0 < t < \infty, \\ U_t = 0, \quad 0 < t < \infty, \\ U_t = 0, \quad 0 < t < \infty, \\ U_t = 0, \quad 0 < t < \infty, \\ U_t = 0, \quad 0 < t < \infty, \\ U_t = 0, \quad 0 < t < \infty, \\ U_t = 0, \quad 0 < t < \infty, \\ U_t = 0, \quad 0 < t < \infty, \\ U_t = 0, \quad 0 < t < \infty, \\ U_t = 0, \quad 0 < t < \infty, \\ U_t = 0, \quad 0 < t < \infty, \\ U_t = 0, \quad 0 < t < \infty, \\ U_t = 0, \quad 0 < t < \infty, \\ U_t = 0, \quad 0 < t < \infty, \\ U_t = 0, \quad 0 < t < \infty, \\ U_t = 0, \quad 0 < t < \infty, \\ U_t = 0, \quad 0 < t < \infty, \\ U_t = 0, \quad 0 < t < \infty, \\ U_t = 0, \quad 0 < t < \infty, \\ U_t = 0, \quad 0 < t < \infty, \\ U_t = 0, \quad 0 < t < \infty, \\ U_t = 0, \quad 0 < t < \infty, \\ U_t

$$

Если угодно, то здесь можно остановиться, решить задачу для  $ U\left(x,\, t\right) $  и затем найти  $ u\left(x,\,t\right) $  по формуле

$$

u(x, t) = T_1 + (T_2 - T_1) U(x, t).

$$

Однако мы пойдем дальше и займемся преобразованием независимых переменных x и t.

# Перетворення просторової координати $ x \rightarrow \xi $

Представляется совершенно очевидным, как следует выбрать безразмерную переменную  $ \xi $ . Поскольку  $ 0 \leqslant x \leqslant L $ , полагаем

$$

\xi = x/L

$$.

После вычисления производных

$$

U_x = U_{\xi} \xi_x = \frac{1}{L} U_{\xi},

$$ U_{xx} = \frac{1}{L^2} U_{\xi\xi}$$

наступна задача стає очевидною (у змінних $U$, $\xi$, $t$):

$$
\begin{aligned}
(\text{УЧП})\quad & U_t = (\alpha/L)^2 U_{\xi\xi}, && 0 < \xi < 1,\quad 0 < t < \infty, \\
(\text{ГУ})\quad & U(0,t)=0,\quad U(1,t)=1, && 0 < t < \infty, \\
(\text{НУ})\quad & U(\xi,0)=\frac{\sin(\pi\xi)-T_1}{T_2-T_1}, && 0 \leqslant \xi \leqslant 1.
\end{aligned}
$$

Ми вже пройшли дві третини шляху. Останній крок пов'язаний із введенням безрозмірного часу $ \tau $ для виключення коефіцієнта $ [\alpha/L]^2 $ з диференціального рівняння.

# $ t \rightarrow \tau $ Перетворення часу

Як обрати час без розмірів? Це не так очевидно, як у випадку перших двох змінних. Однак, оскільки наша мета — виключити сталу $ [\alpha/L]^2 $ з рівняння (22.3), ми діятимемо наступним чином.

- 1. Спробуємо замінити вигляд $ \tau = ct $ , де c — невідома константа.
  - 2. Давайте розрахуємо $ u_t = u_{\tau} \tau_t = c u_{\tau} $ .
  - 3. Підставимо цю похідну у рівняння (22.3) і отримаємо

$$cu_{\mathbf{t}} = [\alpha/L]^2 u_{\xi\xi}.$$

Тому потрібно обрати $ c = [\alpha/L]^2 $ , тобто новий час слід вводити за формулою

$$\tau = [\alpha/L]^2 t.$$

Застосовуючи це перетворення до попередньої задачі (22.3), ми повністю переходимо до безрозмірних змінних $U$, $\xi$ та $\tau$:

$$
\begin{aligned}
(\text{УЧП})\quad & U_{\tau} = U_{\xi\xi}, && 0 < \xi < 1,\quad 0 < \tau < \infty, \\
(\text{ГУ})\quad & U(0,\tau) = 0,\quad U(1,\tau) = 1, && 0 < \tau < \infty, \\
(\text{НУ})\quad & U(\xi,0) = \frac{\sin(\pi \xi) - T_1}{T_2 - T_1}, && 0 \le \xi \le 1.
\end{aligned}
$$

Нова безрозмірна система має такі властивості:

1. Рівняння не містить фізичних параметрів.
2. Граничні умови мають простий вигляд.
3. Початкова умова суттєво не ускладнюється і містить лише відомі функції.
4. Якщо знайдено розв'язок $U(\xi,\tau)$, то розв'язок початкової задачі відновлюється за формулою

$$
u(x,t) = T_1 + (T_2 - T_1)\,U\left(x/L,\alpha^2 t/L^2\right).
$$

На цьому завершується обговорення безрозмірного запису задач. Універсального рецепта введення нових змінних не існує: зазвичай доводиться спиратися на фізичну інтуїцію і перевіряти кілька природних варіантів.

Завершимо лекцію простим прикладом.

# Приклад перетворення гіперболічної задачі на безрозмірний вигляд

Розглянемо коливну струну.

(22.5)

$$
\begin{aligned}
(\text{УЧП})\quad & u_{tt} = \alpha^2 u_{xx}, && 0 < x < L,\quad 0 < t < \infty, \\
(\text{ГУ})\quad & u(0,t) = 0,\quad u(L,t) = 0, && 0 < t < \infty, \\
(\text{НУ})\quad & u(x,0) = \sin(\pi x/L) + 0.5\sin(3\pi x/L),\quad u_t(x,0) = 0, && 0 \le x \le L.
\end{aligned}
$$

Переходячи до безрозмірних змінних

$$
\xi = x/L, \qquad \tau = (\alpha/L)t,
$$

одержуємо нову задачу

$$
\begin{aligned}
(\text{УЧП})\quad & u_{\tau\tau} = u_{\xi\xi}, && 0 < \xi < 1,\quad 0 < \tau < \infty, \\
(\text{ГУ})\quad & u(0,\tau) = 0,\quad u(1,\tau) = 0, && 0 < \tau < \infty, \\
(\text{НУ})\quad & u(\xi,0) = \sin(\pi \xi) + 0.5\sin(3\pi \xi),\quad u_{\tau}(\xi,0) = 0, && 0 \le \xi \le 1.
\end{aligned}
$$

Розв'язком цієї задачі є функція

$$
u(\xi,\tau) = \cos(\pi \tau)\sin(\pi \xi) + 0.5\cos(3\pi \tau)\sin(3\pi \xi).
$$

Повертаючись до координат $x$ і $t$, дістаємо

$$
u(x,t) = \cos(\pi \alpha t/L)\sin(\pi x/L) + 0.5\cos(3\pi \alpha t/L)\sin(3\pi x/L).
$$

## Зауваження

1. Аналіз розмірностей особливо важливий у чисельних методах, оскільки більшість програм орієнтована на загальні математичні моделі, а не на конкретний набір фізичних параметрів.
2. Не завжди потрібно робити безрозмірними всі змінні; інколи достатньо перетворити лише одну або дві з них.
3. Безрозмірний запис часто дозволяє побачити головні параметри задачі значно чіткіше.

# ЗАВДАННЯ

1. За допомогою заміни

$$
\xi = x/L, \qquad \tau = (\alpha/L)t
$$

перетворіть задачу (22.5) до безрозмірного вигляду.

2. Знайдіть безрозмірну форму задачі

$$
\begin{aligned}
(\text{УЧП})\quad & u_t = \alpha^2 u_{xx}, && 0 < x < L, \\
(\text{ГУ})\quad & u(0,t) = T_1,\quad u(L,t) = 0, && 0 < t < \infty, \\
(\text{НУ})\quad & u(x,0) = T_2, && 0 \le x \le L.
\end{aligned}
$$

3. За допомогою заміни

$$
U(x,t) = \frac{u(x,t) - T_1}{T_2 - T_1}
$$

перетворіть задачу (22.1) на задачу (22.2).

4. Поясніть, чому заміна змінної $ \tau = \alpha t $ усуває параметр $\alpha$ з хвильового рівняння $u_{tt} = \alpha^2 u_{xx}$ і який фізичний зміст має ця заміна.

5. Як слід вибрати нову просторову координату $\xi$, щоб усунути параметр $v$ з рівняння

$$
u_t + vu_x = 0?
$$

## Лекція 23

# КЛАСИФІКАЦІЯ РІВНЯНЬ У ЧАСТИННИХ ПОХІДНИХ (КАНОНІЧНА ФОРМА ГІПЕРБОЛІЧНОГО РІВНЯННЯ)

МЕТА ЛЕКЦІЇ: Показати, що лінійне рівняння другого порядку з двома незалежними змінними

$$
Au_{xx}+Bu_{xy}+Cu_{yy}+Du_x+Eu_y+Fu=G
$$

залежно від знака виразу $B^2-4AC$ належить до одного з трьох типів:

1. гіперболічного, якщо $B^2-4AC>0$;
2. параболічного, якщо $B^2-4AC=0$;
3. еліптичного, якщо $B^2-4AC<0$.

Щоб спростити таке рівняння, вводять нові змінні $\xi=\xi(x,y)$ і $\eta=\eta(x,y)$, у яких воно переходить до канонічної форми. Типовими є такі варіанти:

1. у гіперболічному випадку
   $$
   u_{\xi\eta}=\Psi(\xi,\eta,u,u_\xi,u_\eta)
   \quad\text{або}\quad
   u_{\xi\xi}-u_{\eta\eta}=\Psi(\xi,\eta,u,u_\xi,u_\eta);
   $$
2. у параболічному випадку
   $$
   u_{\eta\eta}=\Phi(\xi,\eta,u,u_\xi,u_\eta);
   $$
3. в еліптичному випадку
   $$
   u_{\xi\xi}+u_{\eta\eta}=\Phi(\xi,\eta,u,u_\xi,u_\eta).
   $$

Тут $\Phi$ і $\Psi$ — деякі функції від $u$, її похідних першого порядку та нових незалежних змінних.

У читателя может возникнуть мысль, что глава, посвященная классификации уравнений, должна размещаться в самом начале книги. Вероятно, это правильно, и многие книги действительно начинаются с этого материала. Однако правда и то, что многие студенты не могут с энтузиазмом изучать нечто такое, о чем у них нет никакого представления. Именно по этой причине мы отложили изучение темы, посвященной классификации уравнений с частными производными, до настоящего момента.

Наша задача - классифицировать все уравнения вида

(23 1)

$$

Au_{xx} + Bu_{xy} + Cu_{yy} + Du_x + Eu_y + Fu = G,

$$

де $A$, $B$, $C$, $D$, $E$, $F$ і $G$ є, взагалі кажучи, функціями $x$ та $y$. Згідно з цією схемою рівняння називається

1. гіперболічним у точці $(x_0, y_0)$, якщо $B^2(x_0, y_0) - 4A(x_0, y_0)C(x_0, y_0) > 0$;
2. параболічним у точці $(x_0, y_0)$, якщо $B^2(x_0, y_0) - 4A(x_0, y_0)C(x_0, y_0) = 0$;
3. еліптичним у точці $(x_0, y_0)$, якщо $B^2(x_0, y_0) - 4A(x_0, y_0)C(x_0, y_0) < 0$.

Залежно від типу його можна звести до відповідної канонічної форми. Щоб краще зрозуміти класифікаційну схему, наведемо кілька прикладів.

# Приклади гіперболічних, параболічних та еліптичних рівнянь

1. У равнение теплопроводности  $ u_t = u_{xx} $  — линейное уравнение второго порядка вида (23.1) с коэффициентами

$$

A=1

$$,  $ B=0 $ ,  $ C=0 $ ,  $ D=0 $ ,  $ E=-1 $ ,  $ F=0 $ ,  $ G=0 $ ,

так что  $ B^2-4AC=0 $  при всех x и t. Следовательно, это уравнение будет параболическим при всех x и t. В общем уравнении (23 1) переменная y играет роль времени t. Результат не изменится, если в общем уравнении поменять местами x и y.

2. Волновое уравнение  $ u_{1t} = u_{xx} $  также принадлежит классу уравнений вида (23.1) с коэффициентами

$$

A=1

$$,  $ B=0 $ ,  $ C=-1 $ ,  $ D=E=F=G=0 $ .

Так как  $ B^2 - 4AC = 4 $  для всех x и t, то это уравнение является еиперболическим при всех x и t.

<sup>1)</sup> Следует указать, что приведенные в примерах уравнения являются частными случаями более общих параболических, гиперболических в эллиптических уравнений, а основные их свойства характерны и для общего случая.

3. У равнение Лапласа  $ u_{xx} + u_{yy} = 0 $  является эллиптическим при всех x и y, так как  $ B^2 - 4AC = - 4 < 0 $ .

4. Линейное уравнение  $ xu_{xx} + u_{yy} = \sin x $  с переменными коэффициентами также принадлежит классу уравнений вида (23.1), но для него  $ B^2-4AC=-4x $  и, следовательно, уравнение будет эллиптическим при x > 0,

лараболическим при x=0, гиперболическим при x < 0.

Этот пример иллюстрирует тот факт, что тип уравнения о переменными коэффициентами может изменяться от точки к точке.

Читатель должен обратить внимание на следующее: тип уравнения (23.1) определяется только коэффициентами при вторых првизводных и никак не зависит ни от коэффициентов при первых производных и самой функции, ни от свободного члена.

Обратимся теперь к главной теме нашей лекции — приведению гиперболических уравнений к каноническому виду. Известно, что если в данной области пространства уравнение является гиперболическим, то вместо x и y можно ввести координаты  $ \xi $  и  $ \eta $ (характеристические переменные) таким образом, что уравнение примет простейшую форму

(23.2)

$$

u_{\xi\eta} = \Phi(\xi, \eta, u, u_{\xi}, u_{\eta}).

$$

Это уравнение содержит только одну производную второго порядка  $ u_{\xi\eta} $ , а функция  $ \Phi(\xi,) \eta,u,u_{\xi},u_{\eta}) $  зависит только от новых независимых переменных § и η, зависимой переменной и и первых производных из и ип. Конкретный вид функции Ф зависит, конечно, от исходного уравнения и формул перехода к новым координатам § и п. Научимся отыскивать эти функции.

# Каноническая форма гиперболического уравнения

Рассмотрим уравнение общего вида

(23.3)

$$

Au_{xx} + Bu_{xy} + Cu_{yy} + Du_x + Eu_y + Fu = G

$$

и пусть в интересующей нас области  $ B^2 - 4AC > 0 $ . Нам нужно ввести новые координаты 1)

$$

\xi = \xi(x, y),

$$ \eta = \eta(x, y)$$

тому в рівнянні (23.3) залишається лише одна часткова похідна другого порядку $ u_{\xi\eta} $ (не слід шукати заміну змінним, за допомогою якої гіперболічне вирівнювання

<sup>1)</sup> Необхідно, щоб це перетворення координат було локально оборотним і подвійно диференційованим. — Ед. Ед.

У цьому випадку це неможливо).

Вычислим сначала частные производные

$$ \begin{array}{c} u_x = u_\xi \xi_x + u_\eta \eta_x, \\ u_y = u_\xi \xi_y + u_\eta \eta_y, \\ (23.4) \ u_{xx} = u_{\xi\xi} \xi_x^2 + 2u_{\xi\eta} \xi_x \eta_x + u_{\eta\eta} \eta_x^2 + u_{\xi} \xi_{xx} + u_{\eta} \eta_{xx}, \\ u_{xy} = u_{\xi\xi} \xi_x^2 + u_{\xi\eta} (\xi_x \eta_y + \xi_y \eta_x) + u_{\eta\eta} \eta_x \eta_y + u_{\xi} \xi_{xy} + u_{\eta} \eta_{xy}, \\ u_{yy} = u_{\xi\xi} \xi_y^2 + 2u_{\xi\eta} \xi_y \eta_y + u_{\eta\eta} \eta_y^2 + u_{\xi} \xi_{yy} + u_{\eta} \eta_{yy}. \end{array}

\$\$

Підставляючи ці відношення до початкового рівняння (23.3), після
простих, але громіздких обчислень отримуємо

(23.5)

\[\overline{A}u_{\xi\xi} + \overline{B}u_{\xi\eta} + \overline{C}u_{\eta\eta} + \overline{D}u_{\xi} + \overline{E}u_{\eta} + \overline{F}u = \overline{G},\]

\overline{A} = A\xi\emph{\{x\}\^{}\{2\} + B\xi}\{x\}\xi\emph{\{y\} +
C\xi}\{y\}\^{}\{2\},

\[ \overline{B} = 2A\xi_{x}\eta_{y} + B(\xi_{x}\eta_{y} + \xi_{y}\eta_{x}) + 2C\xi_{y}\eta_{y},\]

\overline{C} = A\eta\emph{\{x\}\^{}\{2\} + B\eta}\{x\}\eta\emph{\{y\} +
C\eta}\{y\}\^{}\{2\},

\[ \overline{D} = A\xi_{xx} + B\xi_{xy} + C\xi_{yy} + D\xi_{x} + E\xi_{y},\]

\overline{E} = A\eta\emph{\{xx\} + B\eta}\{xy\} + C\eta\emph{\{yy\} +
D\eta}\{x\} + E\eta\_\{y\},

\[ \overline{F} = F,\]

\overline{G} = G.

\$\$

Читателю предоставится возможность проделать эти выкладки при решении
задач.

Наш наступний крок --- вибрати функції \(\xi=\xi(x,y)\) і
\(\eta=\eta(x,y)\) так, щоб коефіцієнти \(\overline A\) та
\(\overline C\) у новому рівнянні обернулися на нуль. Це і дає канонічну
форму гіперболічного рівняння. Отже, треба вимагати

\[
\overline A=A\xi_x^2+B\xi_x\xi_y+C\xi_y^2=0,
\qquad
\overline C=A\eta_x^2+B\eta_x\eta_y+C\eta_y^2=0.
\]

Після ділення на \(\xi_y^2\) та \(\eta_y^2\) відповідно одержуємо
характеристичні рівняння

\[
A\left(\frac{\xi_x}{\xi_y}\right)^2+B\left(\frac{\xi_x}{\xi_y}\right)+C=0,
\qquad
A\left(\frac{\eta_x}{\eta_y}\right)^2+B\left(\frac{\eta_x}{\eta_y}\right)+C=0.
\]

Отже,

\[
\frac{\xi_x}{\xi_y}=\frac{-B+\sqrt{B^2-4AC}}{2A},
\qquad
\frac{\eta_x}{\eta_y}=\frac{-B-\sqrt{B^2-4AC}}{2A}.
\]

Саме ці співвідношення визначають дві родини характеристик.

\pandocbounded{\includegraphics[keepaspectratio]{images/_page_170_Figure_3.jpeg}}

Фіг. 23.1. Характеристики \(\xi(x,y)=c\) та \(\eta(x,y)=c\).

Щоб побачити, як це працює, розглянемо просте рівняння

\[
u_{xx}-4u_{yy}+u_x=0.
\]

Для нього маємо дві родини характеристик:

\[
\frac{dy}{dx}=-2,
\qquad
\frac{dy}{dx}=2.
\]

Інтегруючи, одержуємо

\[
y+2x=c_1,
\qquad
y-2x=c_2.
\]

Тому природно покласти

\[
\xi=y+2x,
\qquad
\eta=y-2x.
\]

\pandocbounded{\includegraphics[keepaspectratio]{images/_page_171_Figure_3.jpeg}}

Фіг. 23.2. Характеристичні координати для рівняння
\(u_{xx}-4u_{yy}+u_x=0\).

Після підстановки нових координат у загальну формулу

\[
\overline{A}u_{\xi\xi}+\overline{B}u_{\xi\eta}+\overline{C}u_{\eta\eta}+\overline{D}u_{\xi}+\overline{E}u_{\eta}+\overline{F}u=\overline{G}
\]

одержуємо рівняння у канонічній формі.

Розглянемо тепер приклад

\[
y^2u_{xx}-x^2u_{yy}=0,
\qquad x>0,\; y>0,
\]

яке є гіперболічним у першому квадранті.

Характеристичні рівняння тут мають вигляд

\[
\frac{dy}{dx}=\frac{x}{y},
\qquad
\frac{dy}{dx}=-\frac{x}{y}.
\]

Інтегруючи, дістаємо дві родини кривих:

\[
y^2-x^2=\text{const},
\qquad
y^2+x^2=\text{const}.
\]

Тому беремо нові координати

\[
\xi=y^2-x^2,
\qquad
\eta=y^2+x^2.
\]

\pandocbounded{\includegraphics[keepaspectratio]{images/_page_172_Figure_2.jpeg}}

Фіг. 23.3. Нові характеристичні координати.

Після обчислення коефіцієнтів і підстановки у формулу (23.6) одержуємо

\[
u_{\xi\eta}
=
\frac{-(x^2+y^2)u_{\xi}+(y^2-x^2)u_{\eta}}{8x^2y^2}.
\]

Нарешті, виражаючи \(x\) і \(y\) через \(\xi\) та \(\eta\), дістаємо
канонічну форму

\[
u_{\xi\eta}=\frac{\eta u_{\xi}-\xi u_{\eta}}{2(\eta^2-\xi^2)}.
\]

Таким чином, задача зводиться до рівняння у характеристичних
координатах, яке вже можна досліджувати стандартними методами. \#
ЗАУВАЖЕННЯ 1. На самом деле для гиперболических уравнений существует две
канонические формы. Вторую можно получить из первов заменой переменных
вида

\$\$

\alpha = \xi + \eta,

\$\$

\beta = \xi - \eta.

Скористаймося цією заміною для перетворення рівняння (23.8). Тоді

\$\$

\begin{aligned}
u_{\xi}&=u_{\alpha}\alpha_{\xi}+u_{\beta}\beta_{\xi}=u_{\alpha}+u_{\beta},\\
u_{\eta}&=u_{\alpha}\alpha_{\eta}+u_{\beta}\beta_{\eta}=u_{\alpha}-u_{\beta},\\
u_{\xi\eta}&=u_{\alpha\alpha}-u_{\beta\beta}.
\end{aligned}

\$\$

і, отже,

\$\$

u\_\{\alpha\alpha\} - u\_\{\beta\beta\} =
\frac{-\beta u_{\alpha} - \alpha u_{\beta}}{2\alpha\beta}.

При бажанні можна виразити \(\alpha\) і \(\beta\) через початкові
координати \(x\) та \(y\):

\[
\alpha=\xi+\eta=2y^2,
\qquad
\beta=\xi-\eta=-2x^2.
\]

Потреба в канонічній формі має простий зміст:

\begin{enumerate}
\def\labelenumi{\arabic{enumi}.}
\tightlist
\item
  вона відокремлює хвильові, дифузійні та стаціонарні процеси;
\item
  саме для канонічних форм існує найбільше готових аналітичних методів;
\item
  чисельні схеми для гіперболічних задач зазвичай теж будують уже після
  такого перетворення.
\end{enumerate}

Як тільки розв'язок знайдено в нових координатах, його завжди можна
перетворити назад у початкові координати \((x,y)\).

\section{ЗАВДАННЯ}\label{ux437ux430ux432ux434ux430ux43dux43dux44f-13}

\begin{enumerate}
\def\labelenumi{\arabic{enumi}.}
\item
  Визначте, чи є наступні рівняння гіперболічними, параболічними чи
  еліптичними: \(u_{xx}u_{xy}=0\); \(u_{tt}=u_{xx}+u_x+hu\);
  \(u_{xx}+3u_{yy}=\sin x\); \(u_{xx}+u_{yy}=f(x,y)\);
  \(u_{rr}+\frac{1}{r}u_r+\frac{1}{r^2}u_{\theta\theta}=f(r,\theta)\).
\item
  Отримайте співвідношення (23.4), (23.5) та (23.6).
\item
  Переконайтеся, що рівняння
\end{enumerate}

\[3u_{xx} + 7u_{xy} + 2u_{yy} = 0\]

є гіперболічним при всіх \(x\) та \(y\), і знайдіть його характеристики.

\begin{enumerate}
\def\labelenumi{\arabic{enumi}.}
\setcounter{enumi}{3}
\tightlist
\item
  Зведіть рівняння з попередньої задачі до канонічної форми
\end{enumerate}

\[u_{\xi\eta} = \Phi(\xi, \eta, u, u_{\xi}, u_{\eta}).\]

\begin{enumerate}
\def\labelenumi{\arabic{enumi}.}
\setcounter{enumi}{4}
\tightlist
\item
  Продовжте задачу 4 і приведіть це рівняння до другої канонічної форми
\end{enumerate}

\[u_{\alpha\alpha} - u_{\beta\beta} = \Psi(\alpha, \beta, u, u_{\alpha}, u_{\beta}).\]

\begin{enumerate}
\def\labelenumi{\arabic{enumi}.}
\setcounter{enumi}{5}
\tightlist
\item
  Найдите характеристики уравнения
\end{enumerate}

\[u_{xx} + 4u_{xy} = 0.\]

Перетворіть рівняння на нові координати, розв'яжіть його і поверніться
до координат x і y, щоб отримати розв'язок початкової здогадки.

\section{ХВИЛЬОВЕ РІВНЯННЯ У ВІЛЬНОМУ ПРОСТОРІ (ДВОВИМІРНІ ТА ТРИВИМІРНІ
ЗАДАЧІ)}\label{ux445ux432ux438ux43bux44cux43eux432ux435-ux440ux456ux432ux43dux44fux43dux43dux44f-ux443-ux432ux456ux43bux44cux43dux43eux43cux443-ux43fux440ux43eux441ux442ux43eux440ux456-ux434ux432ux43eux432ux438ux43cux456ux440ux43dux456-ux442ux430-ux442ux440ux438ux432ux438ux43cux456ux440ux43dux456-ux437ux430ux434ux430ux447ux456}

МЕТА ЛЕКЦІЇ: Узагальнити формулу д'Аламбера на двовимірний і тривимірний
випадки та пояснити геометричний зміст відповідних розв'язків.

Розглянемо задачу Коші для хвильового рівняння у тривимірному просторі:

\[
u_{tt}=c^2\Delta u,
\qquad
(x,y,z)\in\mathbb{R}^3,
\qquad
t>0,
\]

\[
u(x,y,z,0)=\varphi(x,y,z),
\qquad
u_t(x,y,z,0)=\psi(x,y,z).
\]

Тут оператор Лапласа має вигляд

\[
\Delta u=u_{xx}+u_{yy}+u_{zz}.
\]

Одновимірні хвилі в просторі називають плоскими, двовимірні ---
циліндричними, а тривимірні --- сферичними. Для сферичних хвиль роль
формули д'Аламбера відіграє \textbf{формула Пуассона}.

\subsection{Формула
Пуассона}\label{ux444ux43eux440ux43cux443ux43bux430-ux43fux443ux430ux441ux441ux43eux43dux430}

Розв'язок задачі Коші у \(\mathbb{R}^3\) можна подати у вигляді

\[
u(x,y,z,t)=\frac{\partial}{\partial t}\!\left[t\,\overline{\varphi}(x,y,z,t)\right]+t\,\overline{\psi}(x,y,z,t),
\]

де \(\overline{\varphi}\) та \(\overline{\psi}\) позначають середні
значення початкових функцій по сфері радіуса \(ct\) з центром у точці
\((x,y,z)\).

\pandocbounded{\includegraphics[keepaspectratio]{images/_page_177_Picture_3.jpeg}}

Фіг. 24.1. Середнє значення початкових даних береться по сфері радіуса
\(ct\).

Геометричний зміст формули простий: значення розв'язку в точці
\((x,y,z)\) у момент часу \(t\) залежить лише від тих початкових даних,
що лежать на сфері радіуса \(ct\).

\pandocbounded{\includegraphics[keepaspectratio]{images/_page_177_Picture_7.jpeg}}

Фіг. 24.2. Початкове збурення поширюється сферично в усіх напрямках.

\subsection{Принцип
Гюйгенса}\label{ux43fux440ux438ux43dux446ux438ux43f-ux433ux44eux439ux433ux435ux43dux441ux430}

Якщо початкове збурення локалізоване в обмеженій області, то в
тривимірному випадку розв'язок має чітко визначений передній і задній
фронт. Це означає, що після проходження хвильового фронту через точку
простору вплив збурення зникає. Такий ефект називають \textbf{принципом
Гюйгенса}.

\pandocbounded{\includegraphics[keepaspectratio]{images/_page_179_Picture_3.jpeg}}

Фіг. 24.3. Точка починає реагувати лише тоді, коли сфера радіуса \(ct\)
перетинає область початкового збурення.

У двовимірному випадку аналогічна задача розв'язується методом спуску,
але там за хвильовим фронтом лишається «хвіст», тому принцип Гюйгенса
вже не виконується.

\subsection{Конформні
відображення}\label{ux43aux43eux43dux444ux43eux440ux43cux43dux456-ux432ux456ux434ux43eux431ux440ux430ux436ux435ux43dux43dux44f}

Тепер перейдемо до одного важливого класу функцій комплексної змінної,
які широко використовують у теорії потенціалу.

\subsubsection{Визначення}\label{ux432ux438ux437ux43dux430ux447ux435ux43dux43dux44f}

Відображення \(w=f(z)\) називається \textbf{конформним} у точці \(z_0\),
якщо похідна \(f'(z_0)\neq 0\). Якщо ця умова виконується в кожній точці
області \(D\), то відображення є конформним у всій області \(D\).

Прості приклади:

\begin{itemize}
\tightlist
\item
  \(f(z)=z^2\) є конформним всюди, крім точки \(z=0\);
\item
  \(f(z)=e^z\) є конформним у всій комплексній площині, бо
  \(f'(z)=e^z\neq 0\).
\end{itemize}

Головна причина, чому ці відображення важливі в задачах математичної
фізики, полягає в тому, що вони переносять складні області у простіші,
не руйнуючи гармонічну структуру. Саме тому рівняння Лапласа в нових
координатах часто набуває зручнішого вигляду.

\subsubsection{Приклад: задача Діріхле у верхній
півплощині}\label{ux43fux440ux438ux43aux43bux430ux434-ux437ux430ux434ux430ux447ux430-ux434ux456ux440ux456ux445ux43bux435-ux443-ux432ux435ux440ux445ux43dux456ux439-ux43fux456ux432ux43fux43bux43eux449ux438ux43dux456}

Нехай потрібно розв'язати задачу

\[
\varphi_{xx}+\varphi_{yy}=0,
\qquad -\infty<x<\infty,\quad y>0,
\]

із граничною умовою

\[
\varphi(x,0)=
\begin{cases}
1, & |x|\le 1,\\
0, & |x|>1.
\end{cases}
\]

Для переходу до простішої області використовують відображення

\[
w=\ln\!\left(\frac{z-1}{z+1}\right).
\]

Це відображення переводить верхню півплощину у смугу, для якої гранична
задача має значно простіший вигляд. Далі розв'язок знаходять у новій
області та повертають назад у змінні \(x\) і \(y\).

\subsection{Принцип
суперпозиції}\label{ux43fux440ux438ux43dux446ux438ux43f-ux441ux443ux43fux435ux440ux43fux43eux437ux438ux446ux456ux457}

Для лінійних рівнянь і лінійних граничних умов працює принцип
суперпозиції: якщо \(u_1\) та \(u_2\) --- розв'язки двох простіших
задач, то сума \(u=u_1+u_2\) розв'язує задачу з сумарним правим боком і
сумарними початковими/граничними даними.

\pandocbounded{\includegraphics[keepaspectratio]{images/_page_189_Picture_10.jpeg}}

Фіг. 26.1. Основна ідея принципу суперпозиції: оператор \(L\) діє
лінійно, тому сумі входів відповідає сума виходів.

Саме тому змішану задачу

\[
\begin{aligned}
u_t &= u_{xx}+f(x,t), && 0<x<1,\quad 0<t<\infty,\\
u(0,t)&=g_1(t),\qquad u(1,t)=g_2(t), && 0<t<\infty,\\
u(x,0)&=\varphi(x), && 0\le x\le 1
\end{aligned}
\]

зручно розкладати на кілька простіших підзадач: окремо враховувати
початкову умову, окремо неоднорідність \(f(x,t)\) і окремо кожну з
граничних умов.

\section{ЗАУВАЖЕННЯ}\label{ux437ux430ux443ux432ux430ux436ux435ux43dux43dux44f-8}

Використання скінченного синусного перетворення приводить до розкладу
розв'язку в ряд. Для інших інтегральних перетворень розв'язки часто
подаються у вигляді інтегралів, тобто неперервних спектральних
розкладів.

Типовий приклад задачі для такого методу:

\[
\begin{aligned}
u_{tt}&=u_{xx}+\sin(3\pi x), && 0<x<1,\quad 0<t<\infty,\\
u(0,t)&=0,\qquad u(1,t)=0, && 0<t<\infty,\\
u(x,0)&=\sin(\pi x),\qquad u_t(x,0)=0, && 0\le x\le 1.
\end{aligned}
\]

Після синусного розкладу по \(x\) початкова змішана задача зводиться до
набору звичайних диференціальних рівнянь для коефіцієнтів Фур'є.

\section{ЗАВДАННЯ}\label{ux437ux430ux432ux434ux430ux43dux43dux44f-14}

\begin{enumerate}
\def\labelenumi{\arabic{enumi}.}
\tightlist
\item
  Розв'яжіть наведену вище задачу для хвильового рівняння за допомогою
  синусного розкладу.
\item
  Поясніть, для яких рівнянь сума двох розв'язків знову є розв'язком, а
  для яких --- ні.
\item
  Розкладіть загальну змішану задачу на кілька простіших підзадач, щоб
  можна було застосувати принцип суперпозиції.
\end{enumerate}

\section{РІВНЯННЯ ПЕРШОГО ПОРЯДКУ (ХАРАКТЕРИСТИЧНИЙ
МЕТОД)}\label{ux440ux456ux432ux43dux44fux43dux43dux44f-ux43fux435ux440ux448ux43eux433ux43e-ux43fux43eux440ux44fux434ux43aux443-ux445ux430ux440ux430ux43aux442ux435ux440ux438ux441ux442ux438ux447ux43dux438ux439-ux43cux435ux442ux43eux434}

МЕТА ЛЕКЦІЇ: Ввести поняття рівнянь першого порядку і пояснити, як метод
характеристик перетворює задачу Коші для УЧП на систему звичайних
диференціальних рівнянь.

Розглянемо типову задачу Коші

\[
a(x,t)u_x+b(x,t)u_t=c(x,t),
\qquad
u(x,0)=\varphi(x).
\]

На відміну від рівняння теплопровідності, де початковий вплив миттєво
розмазується по всій області, для рівнянь першого порядку інформація
переноситься вздовж певних кривих у площині \((x,t)\). Ці криві
називають характеристиками.

\pandocbounded{\includegraphics[keepaspectratio]{images/_page_196_Figure_10.jpeg}}

Фіг. 27.1. Початкове значення в точці \(x\) впливає на розв'язок лише
вздовж характеристики в площині \((x,t)\).

Ідея методу така: треба знайти криві \(x=x(s)\) і \(t=t(s)\), для яких

\[
\frac{dx}{ds}=a(x,t),
\qquad
\frac{dt}{ds}=b(x,t).
\]

Тоді вздовж характеристики частинне рівняння перетворюється на звичайне
диференціальне рівняння для \(u(x(s),t(s))\).

Якщо коефіцієнти сталі, характеристики є прямими. Якщо вони залежать від
\(x\) і \(t\), характеристики можуть викривлятися, а швидкість переносу
змінюватися вздовж потоку.

\pandocbounded{\includegraphics[keepaspectratio]{images/_page_197_Figure_5.jpeg}}

Фіг. 27.2. Для хвильового рівняння початкове збурення в точці породжує
дві хвилі, що поширюються вздовж двох характеристик.

Отже, геометрія характеристик дає не лише спосіб розв'язку, а й наочну
фізичну інтерпретацію задачі.

\section{ПОРІВНЯННЯ АНАЛІТИЧНИХ РІШЕНЬ ІЗ
ЧИСЕЛЬНИМИ}\label{ux43fux43eux440ux456ux432ux43dux44fux43dux43dux44f-ux430ux43dux430ux43bux456ux442ux438ux447ux43dux438ux445-ux440ux456ux448ux435ux43dux44c-ux456ux437-ux447ux438ux441ux435ux43bux44cux43dux438ux43cux438}

МЕТА ЛЕКЦІЇ: Порівняти переваги та недоліки аналітичних і чисельних
методів розв'язання задач математичної фізики.

Під \textbf{аналітичним розв'язком} будемо розуміти формулу,
нескінченний ряд або інтегральне представлення, яке задає шукану функцію
через змінні задачі та параметри системи.

Під \textbf{чисельним розв'язком} будемо розуміти наближений розв'язок,
отриманий після заміни початкового рівняння дискретною схемою. У цьому
випадку результат зазвичай подається таблицею значень \(u(x_i,t_j)\) на
сітці або графіками цих значень.

У друкованій версії книги тут наводиться типова чисельна таблиця
значень. Для нашої мети важливіше не самі числа, а принцип: аналітичний
розв'язок описує поведінку системи в цілому, тоді як чисельний розв'язок
дає набір значень у вузлах сітки.

Саме тому порівняння цих двох підходів треба проводити не абстрактно, а
залежно від мети: чи хочемо ми зрозуміти загальну структуру розв'язку,
чи отримати швидкі числові значення для конкретних параметрів.

\section{Переваги аналітичного
розв'язку}\label{ux43fux435ux440ux435ux432ux430ux433ux438-ux430ux43dux430ux43bux456ux442ux438ux447ux43dux43eux433ux43e-ux440ux43eux437ux432ux44fux437ux43aux443}

\begin{enumerate}
\def\labelenumi{\arabic{enumi}.}
\tightlist
\item
  Аналітична формула дає не лише окремі числові значення, а й цілісне
  уявлення про поведінку розв'язку при зміні параметрів.
\item
  Якщо розв'язок подано рядом або інтегралом, його можна обчислювати з
  потрібною точністю, контролюючи похибку за кількістю взятих членів або
  точністю квадратур.
\item
  Аналітичний розв'язок дозволяє досліджувати асимптотику, стабільність,
  вплив граничних умов і характер поширення збурень.
\end{enumerate}

Разом з тим чисельні методи незамінні там, де явну формулу отримати не
вдається. На практиці обидва підходи доповнюють один одного: аналітичний
розв'язок служить орієнтиром і тестом для чисельної схеми, а чисельний
дозволяє працювати з геометрично складними або нелінійними задачами.

\section{Властивості
перетворень}\label{ux432ux43bux430ux441ux442ux438ux432ux43eux441ux442ux456-ux43fux435ux440ux435ux442ux432ux43eux440ux435ux43dux44c}

Перш ніж переходити до прикладів, корисно зафіксувати кілька загальних
фактів. Коли перетворення виконується лише за просторовою змінною \(x\),
отримані коефіцієнти залежать уже тільки від часу \(t\). Саме це робить
інтегральні та спектральні методи зручними для задач математичної
фізики.

Окремо слід уважно працювати з похідними: після перетворення вони
переходять у нові алгебраїчні або диференціальні співвідношення. Саме
тому перед обчисленнями завжди виписують кілька базових формул для
похідних і граничних членів.

\section{Варіаційні методи розв'язання рівнянь у частинних
похідних}\label{ux432ux430ux440ux456ux430ux446ux456ux439ux43dux456-ux43cux435ux442ux43eux434ux438-ux440ux43eux437ux432ux44fux437ux430ux43dux43dux44f-ux440ux456ux432ux43dux44fux43dux44c-ux443-ux447ux430ux441ux442ux438ux43dux43dux438ux445-ux43fux43eux445ux456ux434ux43dux438ux445}

\textbf{Мета лекції.} Пояснити, як деякі крайові задачі можна розглядати
як задачі мінімізації відповідного функціонала. Тоді функція, що
мінімізує цей функціонал, одночасно є розв'язком диференціального
рівняння.

Класичний приклад дає задача Діріхле для еліптичного рівняння

\[
\Delta u = f \quad \text{в області } D, \qquad u = 0 \quad \text{на межі } \partial D.
\]

Її можна пов'язати з енергетичним функціоналом

\[
J[u] = \iint_D \left[ \frac{1}{2}\left(u_x^2 + u_y^2\right) + fu \right] dx\,dy.
\]

За відповідних умов функція, що мінімізує \(J[u]\) серед усіх допустимих
функцій із заданою граничною умовою, задовольняє рівняння
Ейлера-Лагранжа і тим самим дає розв'язок початкової крайової задачі.

\subsection{Метод
Рітца}\label{ux43cux435ux442ux43eux434-ux440ux456ux442ux446ux430}

Один із найвідоміших варіаційних підходів --- метод Рітца. Його ідея
проста:

\begin{enumerate}
\def\labelenumi{\arabic{enumi}.}
\tightlist
\item
  Обрати набір пробних функцій \(\phi_1, \phi_2, \ldots, \phi_n\), які
  задовольняють граничні умови.
\item
  Шукати наближений розв'язок у вигляді
\end{enumerate}

\[
u_n(x,y) = c_1\phi_1(x,y) + c_2\phi_2(x,y) + \cdots + c_n\phi_n(x,y).
\]

\begin{enumerate}
\def\labelenumi{\arabic{enumi}.}
\setcounter{enumi}{2}
\tightlist
\item
  Підставити \(u_n\) у функціонал \(J[u]\) і мінімізувати отриману
  функцію багатьох змінних за коефіцієнтами \(c_1,\ldots,c_n\).
\end{enumerate}

У результаті задача для рівняння в частинних похідних зводиться до
скінченно-вимірної задачі мінімізації. Чим вдаліше підібрано пробні
функції, тим кращим буде наближення.

Для задачі Діріхле в одиничному квадраті природно брати пробні функції,
що обертаються в нуль на межі області. Типові приклади мають вигляд

\[
\phi_1(x,y) = x(1-x)y(1-y), \qquad
\phi_2(x,y) = x\,\phi_1(x,y), \qquad
\phi_3(x,y) = y\,\phi_1(x,y),
\]

і далі аналогічно будують розширений базис.

На цьому етапі головна ідея така: енергетичний підхід не підміняє
класичну теорію рівнянь, а дає інший, часто дуже ефективний спосіб
знаходити наближені розв'язки складних крайових задач.

На цьому завершується розділ про інтегральні перетворення, принцип
суперпозиції та варіаційні методи. Далі ми переходимо до ще одного
важливого класу задач --- рівнянь першого порядку, для яких ключову роль
відіграє геометрія характеристичних кривих.

\subsection{Лекція 27}\label{ux43bux435ux43aux446ux456ux44f-27}

\section{РІВНЯННЯ ПЕРШОГО ПОРЯДКУ (МЕТОД
ХАРАКТЕРИСТИК)}\label{ux440ux456ux432ux43dux44fux43dux43dux44f-ux43fux435ux440ux448ux43eux433ux43e-ux43fux43eux440ux44fux434ux43aux443-ux43cux435ux442ux43eux434-ux445ux430ux440ux430ux43aux442ux435ux440ux438ux441ux442ux438ux43a}

МЕТА ЛЕКЦІЇ: Ввести поняття диференціального рівняння в частинних
похідних першого порядку і ознайомитися з методом характеристик для
задач із початковими умовами.

Розглянемо типову задачу Коші для лінійного рівняння першого порядку:

\begin{itemize}
\tightlist
\item
  (УЧП): \(a(x,t)u_x + b(x,t)u_t = c(x,t)\), \(-\infty < x < \infty\),
  \(0 < t < \infty\).
\item
  (НУ): \(u(x,0) = \varphi(x)\), \(-\infty < x < \infty\).
\end{itemize}

На відміну від рівняння теплопровідності, де початкове збурення миттєво
впливає на весь простір, для рівнянь першого порядку інформація
поширюється вздовж спеціальних кривих у площині \((x,t)\). Такі криві
називаються характеристиками.

Ідея методу характеристик така: замість того, щоб розв'язувати рівняння
одразу в координатах \((x,t)\), ми переходимо до параметра \(s\), який
біжить уздовж характеристики. Уздовж правильно вибраної характеристики
частинні похідні поєднуються так, що задача зводиться до звичайного
диференціального рівняння.

Якщо характеристику записати у вигляді \(x=x(s)\), \(t=t(s)\), то
природно вимагати, щоб вона задовольняла співвідношення:

\begin{itemize}
\tightlist
\item
  \(\dfrac{dx}{ds} = a(x,t)\),
\item
  \(\dfrac{dt}{ds} = b(x,t)\).
\end{itemize}

Після цього значення \(u(x,t)\) уздовж характеристики визначається ще
одним звичайним диференціальним рівнянням. Саме тому метод характеристик
перетворює складну задачу для УЧП на послідовність простіших задач для
звичайних рівнянь.

Геометрично це означає, що початкові дані з лінії \(t=0\) не
розмазуються по всій площині, а переносяться вздовж характеристичних
кривих. У наступному фрагменті лекції ці криві буде побудовано явно, а
тоді з їхньою допомогою буде отримано розв'язок задачі Коші.

\pandocbounded{\includegraphics[keepaspectratio]{images/_page_205_Figure_5.jpeg}}

Фіг. 28.1. Типова залежність потоку від густини.

Тепер ми готові застосувати закон збереження до задачі дорожнього руху.
Якщо потік \(f(u)\) залежить від густини \(u\), то

\[
f_x=\frac{df}{du}\,u_x.
\]

Отже, рівняння збереження

\[
u_t + f_x = 0
\]

можна переписати у вигляді

\[
u_t + \frac{df}{du}u_x = u_t + g(u)u_x = 0,
\]

де \(g(u)=f'(u)\).

Наприклад, якщо залежність потоку від густини квадратична, тобто
\(f(u)=u^2\), то закон збереження набуває форми

\[
u_t + 2u\,u_x = 0.
\]

Якщо початкова густина розподілу автомобілів задається функцією
\(u(x,0)=\varphi(x)\), то потрібно розв'язати задачу Коші

\[
\begin{aligned}
(\text{УЧП})\quad & u_t + 2u\,u_x = 0, && -\infty < x < \infty,\ 0 < t < \infty, \\
(\text{НУ})\quad & u(x,0)=\varphi(x), && -\infty < x < \infty.
\end{aligned}
\]

Тепер розглянемо загальну нелінійну задачу Коші.

\section{Нелінійна задача з початковою умовою (задача
Коші)}\label{ux43dux435ux43bux456ux43dux456ux439ux43dux430-ux437ux430ux434ux430ux447ux430-ux437-ux43fux43eux447ux430ux442ux43aux43eux432ux43eux44e-ux443ux43cux43eux432ux43eux44e-ux437ux430ux434ux430ux447ux430-ux43aux43eux448ux456}

Розглянемо задачу

\[
\begin{aligned}
(\text{УЧП})\quad & u_t + g(u)u_x = 0, && -\infty < x < \infty,\ 0 < t < \infty, \\
(\text{НУ})\quad & u(x,0)=\varphi(x), && -\infty < x < \infty.
\end{aligned}
\]

Нагадаємо, що для рівняння переносу

\[
u_t + v u_x = 0
\]

розв'язок є сталим уздовж прямих \(x=x_0+vt\). У нелінійному випадку
відбувається те саме, але швидкість тепер залежить від самого розв'язку.
Якщо частинка середовища стартує з точки \((x_0,0)\) і рухається зі
швидкістю \(g(u)\), то її траєкторія має вигляд

\[
x = x_0 + g(u)\,t.
\]

Це і є рівняння характеристики. Вздовж характеристики значення
\(u(x,t)\) залишається сталим і дорівнює початковому значенню
\(u(x_0,0)\).

\pandocbounded{\includegraphics[keepaspectratio]{images/_page_206_Figure_8.jpeg}}

Фіг. 28.2. Характеристики рівняння \(u_t + g(u)u_x = 0\); вздовж кожної
характеристики розв'язок залишається сталим і дорівнює значенню в точці
\((x_0,0)\), з якої ця характеристика виходить.

Отже, характеристика, що починається в точці \((x_0,0)\), задається
формулою

\[
x = x_0 + g\!\bigl[u(x_0,0)\bigr]\,t.
\]

Наприклад, для рівняння

\[
u_t + 3u\,u_x = 0
\]

якщо \(u(2,0)=1\), то характеристика, що виходить із точки \((2,0)\),
визначається співвідношенням

\[
x = 2 + g[1]\,t = 2 + 3t.
\]

Отже, на прямій \(x=2+3t\) розв'язок зберігає значення \(u=1\). Саме так
читається зміст характеристик у нелінійній задачі Коші.

Тепер зрозуміло, що, знаючи характеристики в кожній точці і знаючи, що
розв'язок за характеристиками не змінюється, можна отримати розв'язок \$
u\left(x,~t\right) \$ у будь-який момент часу t. Ми не отримали явного
виразу для \$ u\left(x,~t\right) \$ через змінні x і t, але знання
характеристик рівняння дозволяє розв'язати деякі цікаві задачі.

\section{Задача дорожнього
руху}\label{ux437ux430ux434ux430ux447ux430-ux434ux43eux440ux43eux436ux43dux44cux43eux433ux43e-ux440ux443ux445ux443}

Розглянемо модель руху, у якій потік задається формулою \(f(u)=u^2\), а
початкова густина має вигляд, показаний на рис. 28.3.

\pandocbounded{\includegraphics[keepaspectratio]{images/_page_207_Figure_8.jpeg}}

. 28.3. Початкова густина автомобілів, що рухаються зліва направо.

Іншими словами, нам потрібно розв'язати задачу

\[
\begin{aligned}
(\text{УЧП})\quad & u_t + 2u\,u_x = 0, && -\infty < x < \infty,\ 0 < t < \infty, \\
(\text{НУ})\quad &
u(x,0)=
\begin{cases}
1, & x \le 0,\\
1-x, & 0 < x < 1,\\
0, & 1 \le x,
\end{cases}
&& -\infty < x < \infty.
\end{aligned}
\]

Почнемо з визначення характеристик, що виходять із точок \((x_0,0)\).

Якщо \(x_0<0\), то \(u(x_0,0)=1\), отже

\[
x = x_0 + g[1]\,t = x_0 + 2t,
\]

або, що те саме,

\[
t = \frac{x-x_0}{2}.
\]

Тепер розглянемо точки \(0<x_0<1\). Тут \(u(x_0,0)=1-x_0\), тому

\[
x = x_0 + g(1-x_0)t = x_0 + 2(1-x_0)t,
\]

звідки

\[
t = \frac{x-x_0}{2(1-x_0)}.
\]

Нарешті, якщо \(1\le x_0<\infty\), то \(u(x_0,0)=0\), і характеристики є
вертикальними:

\[
x = x_0.
\]

Усі ці характеристики показані на рис. 28.4.

\pandocbounded{\includegraphics[keepaspectratio]{images/_page_208_Figure_10.jpeg}}

Фіг. 28.4. Характеристики рівняння \$ u\_t + 2uu\_x = 0 \$ .

Ситуація на дорозі при \(0<t<1/2\) досить зрозуміла. Рис. 28.5 показує
розв'язок у різні моменти часу. Зверніть увагу, що характеристики
сходяться в одній точці при \(t=1/2\). Отже, для \(t>1/2\) виникає
ударна хвиля, тобто розрив розв'язку. Потрібно визначити швидкість руху
її фронту.

Виявляється, що швидкість поширення розриву задається формулою
\(S=\dfrac{f(u_R)-f(u_L)}{u_R-u_L}\), де \(u_R\) і \(u_L\) --- значення
розв'язку праворуч і ліворуч від фронту хвилі, а \(f(u_R)\) і \(f(u_L)\)
--- значення потоку при цих густинах.

\pandocbounded{\includegraphics[keepaspectratio]{images/_page_209_Figure_5.jpeg}}

Фіг. 28.5. Щільність руху в різний час.

У нашому прикладі швидкість поширення щілини виявляється рівною одиниці.
Справді, для \(f(u)=u^2\) маємо

\[
S=\frac{0-1}{0-1}=1.
\]

Це означає, що при \(t>1/2\) фронт хвилі рухатиметься зліва направо зі
швидкістю одиниці. Повне розв'язання цієї задачі показано на рис. 28.5.

\section{ЗАУВАЖЕННЯ}\label{ux437ux430ux443ux432ux430ux436ux435ux43dux43dux44f-9}

\begin{enumerate}
\def\labelenumi{\arabic{enumi}.}
\item
  Ударна хвиля в нашому прикладі виникає через те, що потік швидко
  зростає зі збільшенням густини \(u\). Якщо б потік задавався формулою
  \(f(u)=u\), то рівняння мало б вигляд \(u_t+u_x=0\), а розв'язок був
  би просто хвилею, що рухається вправо без спотворення форми.
\item
  Прямою підстановкою розв'язку вигляду
\end{enumerate}

\[
u=\varphi(x-g(u)t)
\]

можна перевірити, що ця формула неявно задає розв'язок задачі

\[
\begin{aligned}
(\text{УЧП})\quad & u_t + g(u)u_x = 0, && -\infty < x < \infty,\quad 0 < t < \infty,\\
(\text{НУ})\quad & u(x,0)=\varphi(x), && -\infty < x < \infty.
\end{aligned}
\]

Наприклад, для задачі Коші

\[
\begin{aligned}
(\text{УЧП})\quad & u_t + u\,u_x = 0, && -\infty < x < \infty,\quad 0 < t < \infty,\\
(\text{НУ})\quad & u(x,0)=x, && -\infty < x < \infty,
\end{aligned}
\]

неявний розв'язок задається співвідношенням

\[
u=\varphi[x-g(u)t]=x-ut.
\]

У нашому конкретному випадку можливо отримати рішення у явній формі (у
багатьох інших випадках це неможливо):

\[
u(x,t)=\frac{x}{1+t}.
\]

Переконайтеся, що ця функція справді є розв'язком нашої задачі.

\section{ЗАВДАННЯ}\label{ux437ux430ux432ux434ux430ux43dux43dux44f-15}

\begin{enumerate}
\def\labelenumi{\arabic{enumi}.}
\tightlist
\item
  Розв'яжіть задачу Коші
\end{enumerate}

\[
\begin{aligned}
(\text{УЧП})\quad & u_t + u\,u_x = 0, && -\infty < x < \infty,\quad 0 < t < \infty, \\
(\text{НУ})\quad &
u(x,0)=
\begin{cases}
0, & x < 0,\\
x, & 0 \le x.
\end{cases}
\end{aligned}
\]

Побудуйте графіки розв'язку для різних моментів часу і поясніть їх
фізичний зміст.

\begin{enumerate}
\def\labelenumi{\arabic{enumi}.}
\setcounter{enumi}{1}
\tightlist
\item
  Розв'яжіть задачу
\end{enumerate}

\[
\begin{aligned}
(\text{УЧП})\quad & u_t + u^2 u_x = 0, && -\infty < x < \infty,\quad 0 < t < \infty, \\
(\text{НУ})\quad &
u(x,0)=
\begin{cases}
0, & x < 0,\\
x, & 0 \le x.
\end{cases}
\end{aligned}
\]

Порівняйте поведінку цього розв'язку з результатом задачі 1.

\begin{enumerate}
\def\labelenumi{\arabic{enumi}.}
\setcounter{enumi}{2}
\tightlist
\item
  Припустимо, що нев'язка рідина тече по трубі і водночас просочується
  крізь стінки згідно із законом \(F(u)=ku\). Тоді закон збереження
  набуває вигляду
\end{enumerate}

\[
u_t + f_x = -F(u).
\]

Знайдіть розв'язок, якщо \(f(u)=u\), а початковий розподіл задається
довільною функцією \(u(x,0)=\varphi(x)\).

\begin{enumerate}
\def\labelenumi{\arabic{enumi}.}
\setcounter{enumi}{3}
\item
  Розв'яжіть попередню задачу, якщо втрати задаються функцією
  \(F(x,t)=1/x\).
\item
  Перевірте, що неявно задана функція
\end{enumerate}

\[
u = \varphi[x-g(u)t]
\]

є розв'язком нелінійної задачі

\[
\begin{aligned}
(\text{УЧП})\quad & u_t + g(u)u_x = 0,\\
(\text{НУ})\quad & u(x,0)=\varphi(x).
\end{aligned}
\]

\section{СИСТЕМИ РІВНЯНЬ З ЧАСТКОВИМИ
ПОХІДНИМИ}\label{ux441ux438ux441ux442ux435ux43cux438-ux440ux456ux432ux43dux44fux43dux44c-ux437-ux447ux430ux441ux442ux43aux43eux432ux438ux43cux438-ux43fux43eux445ux456ux434ux43dux438ux43cux438}

Мета лекції --- показати, що багато фізичних моделей, особливо в
гідродинаміці, не можна описати одним рівнянням. Для цього потрібні
системи пов'язаних рівнянь для тиску, густини, температури та компонент
швидкості.

Почнемо з лінійної системи

\[
u_t + A u_x = 0,
\]

яку зручно звести до незалежних рівнянь діагоналізацією матриці \(A\).

У гідродинаміці зустрічається, наприклад, система рівнянь Ейлера

\[
\begin{aligned}
u_t + uu_x + vu_y + \frac{1}{\rho}p_x &= 0,\\
v_t + uv_x + vv_y + \frac{1}{\rho}p_y &= 0,\\
\rho_t + (\rho u)_x + (\rho v)_y &= 0,\\
\left(\frac{p}{\rho^\gamma}\right)_t
+ u\left(\frac{p}{\rho^\gamma}\right)_x
+ v\left(\frac{p}{\rho^\gamma}\right)_y &= 0.
\end{aligned}
\]

Тут \(p(x,y,t)\) --- тиск, \(u(x,y,t)\) і \(v(x,y,t)\) --- компоненти
швидкості, а \(\rho(x,y,t)\) --- густина рідини.

Системи виникають і з іншої причини: одне рівняння вищого порядку часто
можна звести до системи рівнянь першого порядку. Наприклад, телеграфне
рівняння

\[
u_{tt} = c^2 u_{xx} + a u_t + b u
\]

після введення змінних

\[
u_1=u,\qquad u_2=u_x,\qquad u_3=u_t
\]

зводиться до системи

\[
\begin{cases}
\dfrac{\partial u_1}{\partial x} = u_2,\\[4pt]
\dfrac{\partial u_1}{\partial t} = u_3,\\[4pt]
\dfrac{\partial u_3}{\partial t} = c^2 \dfrac{\partial u_2}{\partial x} + a u_3 + b u_1.
\end{cases}
\]

Для розв'язання систем такого типу потрібно добре розуміти власні
значення та власні вектори матриць, тому зробимо короткий огляд лінійної
алгебри.

\section{Огляд лінійної
алгебри}\label{ux43eux433ux43bux44fux434-ux43bux456ux43dux456ux439ux43dux43eux457-ux430ux43bux433ux435ux431ux440ux438}

Визначення. Якщо \(A\) --- квадратна матриця порядку \(n\), то її власні
значення є коренями характеристичного рівняння

\[
\det(A-\lambda I)=0.
\]

Наприклад, для матриці

\[
A=\begin{bmatrix} 1 & 1 \\ 4 & 1 \end{bmatrix}
\]

маємо

\[
\det\!\begin{bmatrix} 1-\lambda & 1 \\ 4 & 1-\lambda \end{bmatrix}
= (1-\lambda)^2 - 4
= \lambda^2 - 2\lambda - 3 = 0.
\]

Звідси

\[
\lambda_1=-1,\qquad \lambda_2=3.
\]

Власний вектор, що відповідає власному значенню \(\lambda\),
визначається умовою

\[
Ax=\lambda x.
\]

Для \(\lambda_2=3\) одним із власних векторів є

\[
X=\begin{bmatrix} 1 \\ 2 \end{bmatrix},
\]

оскільки

\[
\begin{bmatrix} 1 & 1 \\ 4 & 1 \end{bmatrix}
\begin{bmatrix} 1 \\ 2 \end{bmatrix}
=
3\begin{bmatrix} 1 \\ 2 \end{bmatrix}.
\]

Матриця \(A^{-1}\) називається оберненою до \(A\), якщо

\[
AA^{-1}=A^{-1}A=I.
\]

Наприклад, для

\[
A=\begin{bmatrix} 2 & -2 \\ 1 & 1 \end{bmatrix}
\]

оберненою є матриця

\[
A^{-1}=
\begin{bmatrix}
0.25 & 0.5\\
-0.25 & 0.5
\end{bmatrix},
\]

бо

\[
AA^{-1}=A^{-1}A=
\begin{bmatrix}
1 & 0\\
0 & 1
\end{bmatrix}.
\]

\$\$

\section{Діагональна
матриця}\label{ux434ux456ux430ux433ux43eux43dux430ux43bux44cux43dux430-ux43cux430ux442ux440ux438ux446ux44f}

Нехай \(A\) --- квадратна матриця порядку \(n\) з різними власними
значеннями \(\lambda_1,\lambda_2,\dots,\lambda_n\). Побудуємо матрицю

\[
P=[X_1\ X_2\ \dots\ X_n],
\]

де \(X_k\) --- власний вектор, що відповідає власному значенню
\(\lambda_k\). Тоді

\[
\Lambda = P^{-1}AP
\]

є діагональною матрицею

\[
\Lambda =
\begin{bmatrix}
\lambda_1 & 0 & \dots & 0\\
0 & \lambda_2 & \dots & 0\\
\vdots & \vdots & \ddots & \vdots\\
0 & 0 & \dots & \lambda_n
\end{bmatrix}.
\]

Наприклад, для матриці

\[
A=
\begin{bmatrix}
0 & 8\\
2 & 0
\end{bmatrix}
\]

власні значення дорівнюють

\[
\lambda_1=4,\qquad \lambda_2=-4,
\]

а відповідні власні вектори можна взяти у вигляді

\[
X_1=
\begin{bmatrix}
2\\
1
\end{bmatrix},
\qquad
X_2=
\begin{bmatrix}
-2\\
1
\end{bmatrix}.
\]

\pandocbounded{\includegraphics[keepaspectratio]{images/_page_221_Picture_9.jpeg}}

Рис. 30.1. Кругла мембрана. Гіперболічна змішана задача.

Коливання круглої мембрани описуються рівнянням

\[
u_{tt}=c^2\Delta u,
\qquad
\Delta u=u_{rr}+\frac{1}{r}u_r+\frac{1}{r^2}u_{\theta\theta}.
\]

Оскільки мембрана закріплена на межі круга, маємо граничну умову

\[
u(1,\theta,t)=0.
\]

Шукаємо розв'язок у формі

\[
u(r,\theta,t)=U(r,\theta)T(t).
\]

Після відокремлення змінних отримуємо задачу на власні значення для
рівняння Гельмгольца:

\[
\Delta U+\lambda^2U=0,
\qquad
U(1,\theta)=0.
\]

При \(\mu\le 0\) відповідна крайова задача має лише тривіальний
розв'язок, тому нас цікавлять додатні власні значення.

\section{Розв'язання задачі на власні значення для рівняння
Гельмгольца}\label{ux440ux43eux437ux432ux44fux437ux430ux43dux43dux44f-ux437ux430ux434ux430ux447ux456-ux43dux430-ux432ux43bux430ux441ux43dux456-ux437ux43dux430ux447ux435ux43dux43dux44f-ux434ux43bux44f-ux440ux456ux432ux43dux44fux43dux43dux44f-ux433ux435ux43bux44cux43cux433ux43eux43bux44cux446ux430}

Покладемо

\[
U(r,\theta)=R(r)\Theta(\theta).
\]

Тоді одержуємо два звичайні диференціальні рівняння:

\[
\Theta''+n^2\Theta=0
\]

і

\[
r^2R''+rR'+(\lambda^2r^2-n^2)R=0,
\qquad
0<r<1,\quad R(0)<\infty,\quad R(1)=0.
\]

Перше з них описує кутову залежність, а друге є рівнянням Бесселя.

\section{Рівняння
Бесселя}\label{ux440ux456ux432ux43dux44fux43dux43dux44f-ux431ux435ux441ux441ux435ux43bux44f}

Рівняння

\[
r^2R''+rR'+(\lambda^2r^2-n^2)R=0
\]

має два лінійно незалежні розв'язки: \(J_n(\lambda r)\) --- функцію
Бесселя першого роду та \(Y_n(\lambda r)\) --- функцію Бесселя другого
роду. Тому загальний розв'язок має вигляд

\[
R(r)=A J_n(\lambda r)+B Y_n(\lambda r).
\]

\pandocbounded{\includegraphics[keepaspectratio]{images/_page_223_Figure_8.jpeg}}

\pandocbounded{\includegraphics[keepaspectratio]{images/_page_223_Figure_9.jpeg}}

Рис. 30.2. Функції Бесселя першого і другого роду.

Оскільки \(Y_n(\lambda r)\) необмежена при \(r=0\), у фізичній задачі
про мембрану залишається лише

\[
R(r)=A J_n(\lambda r).
\]

Гранична умова \(R(1)=0\) дає рівняння

\[
J_n(\lambda)=0.
\]

Отже, власні значення дорівнюють

\[
\lambda = k_{nm},
\]

де \(k_{nm}\) --- \(m\)-й корінь рівняння \(J_n(r)=0\).

Таблиця 30.1

{\def\LTcaptype{none} % do not increment counter
\begin{longtable}[]{@{}lrrrrr@{}}
\toprule\noalign{}
\(m\) & \(J_0(r)=0\) & \(J_1(r)=0\) & \(J_2(r)=0\) & \(J_3(r)=0\) &
\(J_4(r)=0\) \\
\midrule\noalign{}
\endhead
\bottomrule\noalign{}
\endlastfoot
1 & 2.40 & 3.83 & 5.13 & 6.38 & 7.59 \\
2 & 5.52 & 7.02 & 8.42 & 9.76 & 11.06 \\
3 & 8.65 & 10.17 & 11.62 & 13.02 & 14.37 \\
4 & 11.79 & 13.32 & 14.80 & 16.22 & 17.62 \\
5 & 14.93 & 16.47 & 17.96 & 19.41 & 20.83 \\
\end{longtable}
}

Тоді власні функції круглої мембрани мають вигляд

\[
U_{nm}(r,\theta)=J_n(k_{nm}r)\,[A\sin(n\theta)+B\cos(n\theta)].
\]

На рис. 30.3 зображено деякі з цих фундаментальних форм коливань.

Кожна така функція описує фундаментальну форму коливань круглої
мембрани. Їй відповідає частота

\[
f_{nm}=\frac{k_{nm}c}{2\pi},
\]

оскільки часовий множник визначається розв'язком рівняння

\[
T'' + k_{nm}^{2}c^{2}T = 0,
\]

тобто

\[
T_{nm}(t)=A\sin(k_{nm}ct)+B\cos(k_{nm}ct).
\]

Отже, всі можливі коливання круглої мембрани можна розглядати як
суперпозиції власних форм \(U_{nm}(r,\theta)\) з відповідними власними
частотами.

\paragraph{ЗАДАЧІ}\label{ux437ux430ux434ux430ux447ux456-2}

\begin{enumerate}
\def\labelenumi{\arabic{enumi}.}
\tightlist
\item
  Використовуючи тригонометричні тотожності, розкладіть функцію
\end{enumerate}

\[
u(1,\varphi)=\cos(3\varphi)
\]

за поліномами Лежандра.

\begin{enumerate}
\def\labelenumi{\arabic{enumi}.}
\setcounter{enumi}{1}
\tightlist
\item
  Розв'яжіть задачу Діріхле в кулі
\end{enumerate}

\[
\Delta u = 0, \qquad 0 < r < 1,\quad 0 \le \varphi \le \pi,\quad 0 \le \theta < 2\pi,
\]

якщо

\[
u(1,\varphi)=
\begin{cases}
1, & 0 \le \varphi \le \pi/2,\\
-1, & \pi/2 < \varphi \le \pi.
\end{cases}
\]

У цій задачі верхня півсфера нагріта до \(+1\), а нижня охолоджена до
\(-1\).

\begin{enumerate}
\def\labelenumi{\arabic{enumi}.}
\setcounter{enumi}{2}
\tightlist
\item
  Знайдіть розв'язок зовнішньої задачі Діріхле для кулі
\end{enumerate}

\[
\Delta u = 0, \qquad 1 < r < \infty,\quad 0 \le \varphi \le \pi,\quad 0 \le \theta < 2\pi,
\]

якщо

\[
u(1,\varphi)=1+\cos\varphi.
\]

Перевірте відповідь.

Лекция 36

\section{НЕОДНОРОДНАЯ ЗАДАЧА ДИРИХЛЕ (ФУНКЦИЯ
ГРИНА)}\label{ux43dux435ux43eux434ux43dux43eux440ux43eux434ux43dux430ux44f-ux437ux430ux434ux430ux447ux430-ux434ux438ux440ux438ux445ux43bux435-ux444ux443ux43dux43aux446ux438ux44f-ux433ux440ux438ux43dux430}

МЕТА ЛЕКЦІЇ: Показати, як гетерогенну задачу Діріхле можна розв'язати за
допомогою функції Гріна (вихідної функції). У цьому важливому методі
права частина рівняння розглядається як певна вхідна дія і розкладається
на неперервний набір дельтаподібних джерел, розподілених по заданій
площі. Потім визначається відповідь системи на кожне таке джерело
(функція Гріна), і всі відповіді підсумовуються (інтегровані). У
результаті ми отримуємо повне розв'язання проблеми.

Досить загальною задачею прикладної математики є визначення потенціалу в
деякій області простору, якщо всередині цієї області задано розподіл
джерел \(f(x,y)\). В електростатиці потенціал у області \(D\)
визначається розподілом зарядів із щільністю \(f(x,y)\), а в теорії
теплопровідності аналогічно задається джерелами тепла. Типова двовимірна
задача в крузі приводить до рівняння Пуассона.

\[
\begin{aligned}
(36.1)\quad (\text{УЧП})\quad &
u_{rr} + \frac{1}{r}u_r + \frac{1}{r^2}u_{\theta\theta} = f(r,\theta),
&& 0 \le r < 1,\quad 0 \le \theta < 2\pi, \\
(\text{ГУ})\quad &
u(1,\theta)=0,
&& 0 \le \theta \le 2\pi.
\end{aligned}
\]

Ми не випадково обрали нульові граничні умови. Якщо б потрібно було
розв'язати задачу в загальному випадку, з ненульовими граничними умовами
та неоднорідним рівнянням, то внесок граничних даних можна було б знайти
за допомогою інтеграла Пуассона з лекції 33.

Щоб отримати наочне уявлення про неоднорідні еліптичні рівняння,
розглянемо задачу

\[
\begin{aligned}
(\text{УЧП})\quad & \Delta u = -q, && 0 < r < 1,\quad 0 \le \theta \le 2\pi, \\
(\text{ГУ})\quad & u(1,\theta)=0, && 0 \le \theta \le 2\pi,
\end{aligned}
\]

де \(q>0\) --- стала. На межі потенціал підтримується рівним нулю, а
лапласіан є сталим у всій області. Рівняння Пуассона означає, що графік
\(u(r,\theta)\) має всюди однакову опуклість. Геометрично його можна
уявляти як тонку мембрану з нерухомим краєм, яку прогинає зовнішня сила.

Тепер перейдемо до головної частини лекції: визначення функції Гріна і
побудови розв'язку рівняння (36.1). Але спершу введемо потенціали
точкових джерел і стоків.

\section{Потенциалы точечных источников и
стоков}\label{ux43fux43eux442ux435ux43dux446ux438ux430ux43bux44b-ux442ux43eux447ux435ux447ux43dux44bux445-ux438ux441ux442ux43eux447ux43dux438ux43aux43eux432-ux438-ux441ux442ux43eux43aux43eux432}

Щоб розв'язати неоднорідне лінійне рівняння, достатньо зрозуміти
відповідь системи на одне точкове джерело. Тоді загальний розв'язок
можна отримати підсумовуванням таких елементарних внесків. У
теплопровідності точкове джерело є місцем виділення тепла, а в
електростатиці --- точковим зарядом.

Припустимо, що джерело інтенсивності \(+q\) розміщене в початку
координат. Тоді потік є радіальним, а повний потік через коло радіуса
\(r\) дорівнює

\[
-\int_0^{2\pi} u_r(r)\,r\,d\theta = -2\pi r\,u_r(r).
\]

За законом збереження цей потік повинен дорівнювати кількості тепла, що
генерується всередині кола, тобто

\[
-2\pi r\,u_r(r)=q.
\]

Розв'язуючи це рівняння відносно функції \(u(r)\), отримуємо двовимірний
потенціал точкового джерела.

\pandocbounded{\includegraphics[keepaspectratio]{images/_page_273_Picture_10.jpeg}}

36 1 Радіальний потік тепла з точкового джерела; A --- тепло виділяється
в початку.

\begin{equation*}
u(r) = \frac{-q}{2\pi} \ln r = \frac{q}{2\pi} \ln \frac{1}{r}.
\end{equation*}

(Розв'язок зображено на рис. 36.2.)

В електростатиці різниця потужності t(B) до u(A) дорівнює роботі, яку
потрібно виконати для переміщення одного позитивного заряду з точки A в
точку B (рисунок 36.2).

\pandocbounded{\includegraphics[keepaspectratio]{images/_page_273_Figure_15.jpeg}}

Рис. 36.2. Потенціал точкового заряду у двовимірному випадку. Тут
\(\alpha\) --- робота, потрібна для переміщення одиничного заряду з
точки A до точки B.

Стік можна інтерпретувати як від'ємне джерело, так що стік -q створює
потенціал

\begin{equation*}
u\left(r\right)=-\frac{q}{2\pi}\ln\frac{1}{r}.
\end{equation*}

На цьому завершується обговорення потенціалів точкових зарядів. Тепер ми
готові розв'язати неоднорідне рівняння за допомогою методу функції
Гріна.

\section{Рівняння Пуассона в
крузі}\label{ux440ux456ux432ux43dux44fux43dux43dux44f-ux43fux443ux430ux441ux441ux43eux43dux430-ux432-ux43aux440ux443ux437ux456}

Розглянемо важливу задачу.

(36.2)

\[
\begin{aligned}
(\text{УЧП})\quad & u_{rr} + \frac{1}{r} u_r + \frac{1}{r^2} u_{\theta\theta} = f(r, \theta), && 0 < r < 1,\ 0 \le \theta \le 2\pi, \\
(\text{ГУ})\quad & u(1, \theta) = 0, && 0 \le \theta \le 2\pi.
\end{aligned}
\]

Метод функції Гріна складається з двох кроків.

\begin{enumerate}
\def\labelenumi{\arabic{enumi}.}
\item
  Встановіть заряд точкової одиниці в точці \((\rho, \varphi)\) і
  знайдіть \(G(r, \theta, \rho, \varphi)\) потенціал, створений цим
  зарядом і який на межі дорівнює нулю.
\item
  Підсумуйте окремі відповіді, \(G(r, \theta, \rho, \phi)\), зважені
  правильною частиною (щільністю заряду) \(f(r, \theta)\), вздовж усього
  кола. У результаті ми отримуємо те рішення, яке шукаємо
\end{enumerate}

\begin{equation*}
u\left(r,\theta\right) = \int_{0}^{2\pi} \int_{0}^{1} G\left(r,\theta,\rho,\varphi\right) f\left(\rho,\varphi\right)\rho\, d\rho\, d\varphi.
\end{equation*}

Тепер знайдемо вихідну функцію для нашого завдання.

\section{\texorpdfstring{Означення функції джерела
\(G(r, \theta, \rho, \phi)\)}{Означення функції джерела G(r, \textbackslash theta, \textbackslash rho, \textbackslash phi)}}\label{ux43eux437ux43dux430ux447ux435ux43dux43dux44f-ux444ux443ux43dux43aux446ux456ux457-ux434ux436ux435ux440ux435ux43bux430-gr-theta-rho-phi}

Замінимо праву частину \(f(r,\theta)\) рівняння точковим джерелом
одиничної величини, розташованим у точці \((\rho, \varphi)\).
Математично таке джерело описується дельта-функцією. Наше завдання ---
знайти потенціал цього джерела за умови, що на межі кола потенціал
дорівнює нулю. Така функція і називається функцією Гріна, або функцією
джерела.

Див. В. С. Володимиров, Узагальнені функції в математичній фізиці,
Москва, Наука, 1976. -- Примітка редактора.

потенціал точкового заряду дорівнює

\begin{equation*}
\frac{1}{2\pi} \ln \left( \frac{1}{r} \right)
\end{equation*},

де \(r\) --- відстань від заряду до точки спостереження.

Шуканий потенціал має такий фізичний зміст:

\begin{enumerate}
\def\labelenumi{\arabic{enumi}.}
\tightlist
\item
  У теорії теплопровідності це стаціонарна температура всередині круга,
  якщо джерело тепла розташоване в точці \((\rho,\varphi)\), а на межі
  підтримується нульова температура.
\item
  Для мембрани це відхилення від положення рівноваги, якщо край
  закріплено на нулі, а в точці \((\rho,\varphi)\) прикладено
  зосереджене навантаження.
\item
  В електростатиці це розподіл потенціалу всередині круга, якщо точковий
  позитивний заряд розташований у точці \((\rho,\varphi)\), а коло
  заземлено.
\end{enumerate}

\pandocbounded{\includegraphics[keepaspectratio]{images/_page_275_Picture_11.jpeg}}

Рис. 36.3. Функція Гріна для точкового джерела, розташованого в точці
\((\rho, \phi)\), всередині одиничного кола.

Тепер побудуємо функцію Гріна; за формою вона буде подібна до графіка на
рис. 36.3.

\section{Побудова
розв'язку}\label{ux43fux43eux431ux443ux434ux43eux432ux430-ux440ux43eux437ux432ux44fux437ux43aux443}

КРОК 1. Оскільки

\begin{equation*}
\frac{1}{2\pi} \ln \frac{1}{R}
\end{equation*}

--- це потенціал поля в точці \(P=(r,\theta)\), який створюється
одиничним точковим зарядом, розміщеним у точці \(Q=(\rho,\phi)\).
Залишається лише модифікувати цю функцію так, щоб на межі кола вона
перетворювалася на нуль.

КРОК 2. З експерименту відомо, що рівнопотенціальні лінії точкових
зарядів є колами (рис. 36.4).

Основна ідея побудови функції Гріна така: потрібно розмістити інший,
від'ємний, заряд поза колом так, щоб сумарний потенціал на колі \(r=1\)
був сталим. Потім цю константу можна відняти й отримати нульовий
потенціал на межі.

\pandocbounded{\includegraphics[keepaspectratio]{images/_page_276_Picture_3.jpeg}}

Рис. 36.4. Потенціал поля, утвореного двома протилежно зарядженими
частинками; A і B --- лінії сталого потенціалу.

\pandocbounded{\includegraphics[keepaspectratio]{images/_page_276_Picture_5.jpeg}}

Рис. 36.5. Заряди \(Q\) і \(\overline{Q}\) створюють сталий потенціал на
колі \(r=1\).

Легко побачити, що якщо негативний заряд розміщений у точці
\(\bar{Q} = (\bar{\rho}, \bar{\phi}) = (1/\rho, \phi)\), то повний
потенціал

\begin{equation*}
u(r, \theta) = \frac{1}{2\pi} \ln \frac{1}{R} - \frac{1}{2\pi} \ln \frac{1}{\overline{R}}
\end{equation*}

буде сталою на колі \(r=1\). Тут відстані від зарядів до точки
спостереження позначено через \(R\) і \(\overline{R}\) (див. рис. 36.5).
Можна показати, що значення потенціалу на колі \(r=1\) дорівнює

\begin{equation*}
-\frac{1}{2\pi}\ln\rho
\end{equation*}

(додатна константа).

З урахуванням цього результату будуємо функцію Гріна

(36.3)

\begin{equation*}
G(r, \theta, \rho, \phi) =
\frac{1}{2\pi} \ln \frac{1}{R}
- \frac{1}{2\pi} \ln \frac{1}{\overline{R}}
+ \frac{1}{2\pi} \ln \rho
\end{equation*}

Отже, функція Гріна складається з потенціалу додатного заряду в точці
\(Q\), потенціалу від'ємного заряду в точці \(\overline{Q}\) та
коригувальної константи, яка забезпечує нульове граничне значення.

де

\begin{equation*}
\begin{aligned}
R &= \sqrt{r^2 - 2\rho r \cos(\theta - \phi) + \rho^2},\\
\overline{R} &= \sqrt{r^2 - 2\frac{r}{\rho}\cos(\theta - \phi) + \frac{1}{\rho^2}}.
\end{aligned}
\end{equation*}

Це формули для відстаней між відповідними точками в полярних
координатах. Щоб знайти розв'язок початкової задачі, потрібно
просумувати внесок усіх таких елементарних джерел. Переходимо до
останнього кроку.

КРОК 3. Суперпозиція імпульсних характеристик Цей крок досить простий. З
відношення

\begin{equation*}
u(r, \theta) = \int_{0}^{2\pi} \int_{0}^{1} G(r, \theta, \rho, \varphi) f(\rho, \varphi) \rho d\rho d\varphi
\end{equation*}

отримуємо

(36.4)

\begin{equation*}
u(r, \theta) = \frac{1}{2\pi} \int_{0}^{2\pi} \int_{0}^{1}
\ln \left(\frac{\rho \overline{R}}{R}\right) f(\rho, \varphi)\,\rho\, d\rho\, d\varphi.
\end{equation*}

Формула (36.4) виражає розв'язок задачі Діріхле для рівняння Пуассона
всередині кола через функцію Гріна цієї задачі. Якщо щільність заряду
\(f(r, \theta)\) відома, то інтеграл можна обчислити, наприклад,
чисельно. \# ЗАУВАЖЕННЯ

Зауваження 1. Задачу

\begin{equation*}
\begin{aligned}
(\text{УЧП}) \quad & \Delta u = 0, \quad 0 < r < 1, \quad 0 \leqslant \theta < 2\pi, \\
(\text{ГУ}) \quad & u(1, \theta) = g(\theta), \quad 0 \leqslant \theta \leqslant 2\pi.
\end{aligned}
\end{equation*}

також можна розв'язати за допомогою методу \emph{функції Гріна}. У цьому
випадку розв'язок має вигляд

\begin{equation*}
u(r, \theta) = \int_{0}^{2\pi} \frac{\partial G}{\partial r}(r, \theta, 1, \varphi)\, g(\varphi)\, d\varphi.
\end{equation*}

Останнє співвідношення можна переписати в більш обчислювально зручному
вигляді, явно підставивши вираз для \$ \partial G/\partial r \$. У
результаті отримуємо інтегральну формулу Пуассона, виведену раніше в
лекції 33.

(36.5)

\begin{equation*}
u(r, \theta) = \frac{1}{2\pi} \int_{0}^{2\pi}
\left[ \frac{1 - r^2}{1 - 2r\cos(\theta - \varphi) + r^2} \right] g(\varphi)\, d\varphi.
\end{equation*}

Зауваження 2. Розв'язок загальної задачі Діріхле

\begin{equation*}
\begin{aligned}
(\text{УЧП}) \quad & \Delta u = f(r, \theta), \quad 0 < r < 1, \quad 0 \le \theta < 2\pi, \\
(\text{ГУ}) \quad & u(1, \theta) = g(\theta), \quad 0 \le \theta \le 2\pi.
\end{aligned}
\end{equation*}

знаходиться як сума розв'язків (36.4) і (36.5).

Зауваження 3. Метод функцій Гріна дозволяє отримати розв'язки багатьох
задач в областях різної форми. Однак для кожної області, а точніше для
кожного оператора та граничної умови, потрібно знайти власну функцію
Гріна, і це не завжди легко зробити.

Зауваження 4. Щоб дійсно знайти розв'язок за формулою (36.4), у
більшості випадків необхідно обчислити інтеграли на комп'ютері. \#
ЗАВДАННЯ

\begin{enumerate}
\def\labelenumi{\arabic{enumi}.}
\item
  Знайти потенціал точкового джерела у тривимірному просторі.
\item
  Знайти функцію Гріна \(G(x,y,\xi,\eta)\) для рівняння Лапласа у
  верхній півплощині \(y>0\). Іншими словами, знайти потенціал у точці
  \((x,y)\), якщо точковий заряд розташований у \((\xi,\eta)\), а
  потенціал на межі \(y=0\) дорівнює нулю.
\end{enumerate}

\pandocbounded{\includegraphics[keepaspectratio]{images/_page_278_Figure_11.jpeg}}

ПРИМІТКА. Якщо від'ємний заряд розмістити в точці
\(\overline{Q}=(\xi,-\eta)\), то сумарний потенціал на прямій \(y=0\)
зануляється. Це і є основа методу дзеркальних відображень.

\begin{enumerate}
\def\labelenumi{\arabic{enumi}.}
\setcounter{enumi}{2}
\item
  Використовуючи результати задачі 2, знайдіть розв'язок рівняння
  Пуассона \(\Delta u=-k\) у верхній півплощині з нульовою граничною
  умовою.
\item
  Як би ви побудували функцію Гріна для першого квадранта \(x>0\),
  \(y>0\)?
\item
  Інший підхід до розв'язання рівняння Пуассона такий. Припустимо, що
  потрібно розв'язати задачу
\end{enumerate}

\[
\begin{aligned}
(\text{УЧП})\quad & \Delta u = 1, && 0 < r < 1,\quad 0 \le \theta < 2\pi,\\
(\text{ГУ})\quad & u(1,\theta)=\sin\theta, && 0 \le \theta \le 2\pi.
\end{aligned}
\]

Спочатку спробуйте знайти частинний розв'язок підстановкою

\[
u_p(r,\theta)=Ar^2.
\]

Після цього підставте вираз \(u_p\) в рівняння Пуассона і знайдіть
константу \(A\). Далі подайте шуканий розв'язок у вигляді \(u=w+u_p\),
де функція \(w(r,\theta)\) задовольняє відповідну однорідну задачу з
модифікованими граничними умовами. Знайдіть \(w(r,\theta)\) і запишіть
остаточний розв'язок \(u(r,\theta)\). Ретельно проаналізуйте отриману
відповідь.

\section{Частина 5. ЧИСЕЛЬНІ ТА ПРИБЛИЗНІ
МЕТОДИ}\label{ux447ux430ux441ux442ux438ux43dux430-5.-ux447ux438ux441ux435ux43bux44cux43dux456-ux442ux430-ux43fux440ux438ux431ux43bux438ux437ux43dux456-ux43cux435ux442ux43eux434ux438}

\subsection{Лекція 37}\label{ux43bux435ux43aux446ux456ux44f-37}

\section{ЧИСЕЛЬНІ РОЗВ'ЯЗКИ (ЕЛІПТИЧНІ
ЗАДАЧІ)}\label{ux447ux438ux441ux435ux43bux44cux43dux456-ux440ux43eux437ux432ux44fux437ux43aux438-ux435ux43bux456ux43fux442ux438ux447ux43dux456-ux437ux430ux434ux430ux447ux456}

МЕТА ЛЕКЦІЇ: Показати, як рівняння в частинних похідних можна звести до
систем алгебраїчних рівнянь, якщо замінити похідні наближеннями
скінченних різниць. Розв'язуючи ці системи ітеративними методами, можна
отримати наближені розв'язки еліптичних задач.

На цей момент ми вже бачили кілька точних методів розв'язання лінійних
рівнянь у частинних похідних. Проте більшість із них працює лише для
простих рівнянь і простих областей. У реальних прикладних задачах цього
часто недостатньо, тому доводиться вдаватися до чисельних методів.

У цій та наступних двох лекціях ми покажемо, як розв'язувати еліптичні,
гіперболічні та параболічні рівняння методом скінченних різниць. Спершу
ознайомимося з самим поняттям скінченних різниць, а потім покажемо, як
застосовувати їх до задачі Діріхле в квадраті.

\section{Апроксимації скінченних
різниць}\label{ux430ux43fux440ux43eux43aux441ux438ux43cux430ux446ux456ux457-ux441ux43aux456ux43dux447ux435ux43dux43dux438ux445-ux440ux456ux437ux43dux438ux446ux44c}

Згадаємо ряд Тейлора для функції \(f(x)\):

\[
f(x+h)=f(x)+f'(x)h+\frac{f''(x)}{2!}h^2+\dots
\]

Якщо обмежитися першими двома членами, отримаємо

\[
f(x+h)\cong f(x)+f'(x)h,
\]

звідки випливає права різницева похідна

\[
(37.1)\qquad f'(x)\cong \frac{f(x+h)-f(x)}{h}.
\]

У розкладі Тейлора можна замінити \(h\) на \(-h\) і отримати ліву
різницеву похідну:

\[
(37.2)\qquad f'(x)\cong \frac{f(x)-f(x-h)}{h}.
\]

Розкладаючи \(f(x-h)\) і \(f(x+h)\) в ряд Тейлора та відкидаючи члени
вищого порядку, отримуємо центральну різницеву похідну:

\[
(37.3)\qquad f'(x)\simeq \frac{f(x+h)-f(x-h)}{2h}.
\]

Якщо в ряду Тейлора залишити ще один член, то для другої похідної маємо

\[
(37.4)\qquad f''(x)\cong \frac{f(x+h)-2f(x)+f(x-h)}{h^2}.
\]

Тепер поширимо поняття скінченних різниць на часткові похідні. Виходячи
з розкладів Тейлора для функції двох змінних,

\[
\begin{aligned}
u(x+h,y)&=u(x,y)+u_x(x,y)h+u_{xx}(x,y)\frac{h^2}{2!}+\dots,\\
u(x-h,y)&=u(x,y)-u_x(x,y)h+u_{xx}(x,y)\frac{h^2}{2!}-\dots.
\end{aligned}
\]

отримуємо такі наближення:

\[
\begin{aligned}
u_x(x,y) &\cong \frac{u(x+h,y)-u(x,y)}{h},\\
u_{xx}(x,y) &\cong \frac{u(x+h,y)-2u(x,y)+u(x-h,y)}{h^2},\\
u_y(x,y) &\cong \frac{u(x,y+k)-u(x,y)}{k},\\
u_{yy}(x,y) &\cong \frac{u(x,y+k)-2u(x,y)+u(x,y-k)}{k^2}.
\end{aligned}
\]

У цих формулах часткові похідні апроксимуються правою, центральною та
лівою різницею, але в цій лекції ми будемо використовувати лише
центральні наближення.

\pandocbounded{\includegraphics[keepaspectratio]{images/_page_282_Figure_5.jpeg}}

Фіг. 37.1 Сітка для задачі Діріхле всередині квадрата.

Щоб ознайомитися з основними правилами використання апроксимації
скінченних різниць, розв'яжемо просту задачу Діріхле.

Розв'язання задачі Діріхле методом скінченної різниці

Нехай потрібно розв'язати задачу Діріхле

\[
\begin{aligned}
(\text{УЧП})\quad & u_{xx}+u_{yy}=0, && 0<x<1,\quad 0<y<1,\\
(\text{ГУ})\quad & u=0, && \text{на верхній та бічних сторонах квадрата},\\
& u(x,0)=\sin(\pi x), && 0\le x\le 1.
\end{aligned}
\]

Намалюємо сітку на площині \((x,y)\), як показано на рис. 37.1. Зручно
(особливо якщо ми збираємося користуватися комп'ютером) використовувати
такі позначення:

\[
\begin{aligned}
u(x,y)&=u_{ij},\\
u(x,y+k)&=u_{i+1,j},\\
u(x,y-k)&=u_{i-1,j},\\
u(x-h,y)&=u_{i,j-1},\\
u(x+h,y)&=u_{i,j+1},\\
u_x(x,y)&\cong \frac{u_{i,j+1}-u_{i,j-1}}{2h},\\
u_y(x,y)&\cong \frac{u_{i+1,j}-u_{i-1,j}}{2k},\\
u_{xx}(x,y)&\cong \frac{u_{i,j+1}-2u_{i,j}+u_{i,j-1}}{h^2},\\
u_{yy}(x,y)&\cong \frac{u_{i+1,j}-2u_{i,j}+u_{i-1,j}}{k^2}.
\end{aligned}
\]

Основна ідея розв'язання задачі Діріхле базується на заміні часткових
похідних у рівнянні Лапласа їхніми скінченнорізницевими апроксимаціями.
Зробивши це і використавши компактну нотацію, отримуємо різницеве
рівняння

\[
\Delta u \cong
\frac{u_{i,j+1}-2u_{i,j}+u_{i,j-1}}{h^2}
+
\frac{u_{i+1,j}-2u_{i,j}+u_{i-1,j}}{k^2}
=0.
\]

У випадку \(h=k\) рівняння Лапласа приводиться до вигляду

\[
u_{i+1,j}+u_{i-1,j}+u_{i,j+1}+u_{i,j-1}-4u_{ij}=0.
\]

Розв'язавши його відносно \(u_{ij}\), отримуємо

\[
u_{ij}=\frac{1}{4}\left(u_{i+1,j}+u_{i-1,j}+u_{i,j+1}+u_{i,j-1}\right).
\]

Отже, значення розв'язку у внутрішньому вузлі апроксимується середнім
значенням розв'язку в чотирьох сусідніх вузлах. На цьому і ґрунтується
метод Лібмана.

\section{Алгоритм чисельного розв'язання задачі Діріхле (метод
Лібмана)}\label{ux430ux43bux433ux43eux440ux438ux442ux43c-ux447ux438ux441ux435ux43bux44cux43dux43eux433ux43e-ux440ux43eux437ux432ux44fux437ux430ux43dux43dux44f-ux437ux430ux434ux430ux447ux456-ux434ux456ux440ux456ux445ux43bux435-ux43cux435ux442ux43eux434-ux43bux456ux431ux43cux430ux43dux430}

КРОК 1. Присвоїмо значенням \(u_{ij}\) у внутрішніх вузлах початкове
наближення, наприклад середнє значення всіх граничних умов.

КРОК 2. Перераховуємо значення у всіх внутрішніх вузлах, замінюючи старе
значення середнім значенням чотирьох сусідніх точок. Після кількох
ітерацій процес збігається до наближеного розв'язку задачі.

\section{ЗАУВАЖЕННЯ}\label{ux437ux430ux443ux432ux430ux436ux435ux43dux43dux44f-10}

\begin{enumerate}
\def\labelenumi{\arabic{enumi}.}
\tightlist
\item
  Якщо взяти сітку з чотирма внутрішніми точками \((m=n=4)\), то система
  набуває вигляду
\end{enumerate}

\[
\begin{aligned}
-4u_{22} + u_{23} + u_{32} &= -\sin(\pi/3),\\
u_{22} - 4u_{23} + u_{33} &= -\sin(2\pi/3),\\
u_{22} - 4u_{32} + u_{33} &= 0,\\
u_{23} + u_{32} - 4u_{33} &= 0.
\end{aligned}
\]

Цю систему потрібно розв'язувати відносно \(u_{22}\), \(u_{23}\),
\(u_{32}\) та \(u_{33}\). Розв'язок такої системи можна знайти
ітеративним методом. Метод Лібмана також є одним з ітеративних методів.

\begin{enumerate}
\def\labelenumi{\arabic{enumi}.}
\setcounter{enumi}{1}
\item
  Якщо зменшити кроки сітки \(h\) і \(k\), то отримаємо систему рівнянь
  того самого типу, але більшого розміру.
\item
  Систему рівнянь можна записати у матричній формі:
\end{enumerate}

\[
\begin{bmatrix}
-4 & 1 & 1 & 0 \\
1 & -4 & 0 & 1 \\
1 & 0 & -4 & 1 \\
0 & 1 & 1 & -4
\end{bmatrix}
\begin{bmatrix}
u_{22}\\
u_{23}\\
u_{32}\\
u_{33}
\end{bmatrix}
=
\begin{bmatrix}
-\sin(\pi/3)\\
-\sin(2\pi/3)\\
0\\
0
\end{bmatrix}
\approx
\begin{bmatrix}
-0.86\\
-0.86\\
0\\
0
\end{bmatrix}.
\]

Зі збільшенням кількості рівнянь матриця коефіцієнтів стає розрідженою,
тобто містить багато нулів. Такі системи розв'язуються спеціальними
методами, серед яких особливо важливі метод Якобі, метод Гаусса-Зайделя
та послідовна верхня релаксація.

\begin{enumerate}
\def\labelenumi{\arabic{enumi}.}
\setcounter{enumi}{3}
\item
  При розв'язанні задачі Неймана ті \emph{похідні}, що входять до
  граничної умови, також слід замінити різницею апроксимаціями.
\item
  Таким чином ви зможете розв'язувати задачі такого типу.
\end{enumerate}

а) \(u_{xx} + u_{yy} = f(x, y)\)

б) \(x u_{xx} + u_{yy} + 2u = \sin(x - y)\)

в) \(\sin x\,u_{xx} + u_{xy} + 3u = 0\)

(гетерогенні рівняння зі змінними коефіцієнтами).

\begin{enumerate}
\def\labelenumi{\arabic{enumi}.}
\setcounter{enumi}{5}
\tightlist
\item
  Якщо область, у якій розв'язується задача, має неправильну форму, її
  можна покрити сіткою і наближати розв'язок у вузлах, найближчих до
  межі, інтерполюючи граничні умови. Після цього задачу розв'язують
  звичайним способом.
\end{enumerate}

\pandocbounded{\includegraphics[keepaspectratio]{images/_page_285_Figure_8.jpeg}}

Фіг. 37.2. Поблизу криволінійної межі нові граничні умови отримують
інтерполяцією.

\begin{enumerate}
\def\labelenumi{\arabic{enumi}.}
\setcounter{enumi}{6}
\tightlist
\item
  Існують спеціальні програмні пакети для чисельного розв'язання
  еліптичних задач. Один з відомих прикладів --- пакет ELLPACK, який
  дозволяє працювати з дво- та тривимірними задачами в різних
  координатних системах і для досить широкого класу граничних умов.
\end{enumerate}

\section{ЗАВДАННЯ}\label{ux437ux430ux432ux434ux430ux43dux43dux44f-16}

\begin{enumerate}
\def\labelenumi{\arabic{enumi}.}
\tightlist
\item
  Отримайте наближення для другої похідної
\end{enumerate}

\[
f''(x)\approx \frac{f(x+h)-2f(x)+f(x-h)}{h^2}.
\]

\begin{enumerate}
\def\labelenumi{\arabic{enumi}.}
\setcounter{enumi}{1}
\item
  Виконайте дві ітерації для задачі Діріхле (37.5) методом Лібмана.
\item
  До якої алгебраїчної системи зводиться задача Діріхле для рівняння
  Пуассона в квадраті
\end{enumerate}

\[
\begin{aligned}
(\text{УЧП})\quad & u_{xx}+u_{yy}=f(x,y), && 0<x<1,\quad 0<y<1,\\
(\text{ГУ})\quad & u(x,y)=g(x,y), && \text{на межі квадрата},
\end{aligned}
\]

якщо розв'язувати її методом скінченних різниць?

\begin{enumerate}
\def\labelenumi{\arabic{enumi}.}
\setcounter{enumi}{3}
\tightlist
\item
  Розв'яжіть задачу 3 для рівняння
\end{enumerate}

\[
u_{xx}+u_{yy}+2u=0,\qquad 0<x<1,\quad 0<y<1,
\]

за умови \(u(x,y)=g(x,y)\) на межі квадрата.

\begin{enumerate}
\def\labelenumi{\arabic{enumi}.}
\setcounter{enumi}{4}
\tightlist
\item
  Як би ви розв'язували задачу Неймана в квадраті
\end{enumerate}

\[
\begin{aligned}
(\text{УЧП})\quad & u_{xx}+u_{yy}=0, && 0<x<1,\quad 0<y<1,\\
(\text{ГУ})\quad & u=0 && \text{на трьох сторонах квадрата},\\
& \frac{\partial u}{\partial x}(1,y)=1, && 0\le y\le 1,
\end{aligned}
\]

методом скінченних різниць?

\begin{enumerate}
\def\labelenumi{\arabic{enumi}.}
\setcounter{enumi}{5}
\tightlist
\item
  Побудуйте блок-схему алгоритму розв'язання задачі Діріхле в квадраті
\end{enumerate}

\[
\begin{aligned}
(\text{УЧП})\quad & u_{xx}+u_{yy}=f(x,y), && 0<x<1,\quad 0<y<1,\\
(\text{ГУ})\quad & u(x,y)=g(x,y), && \text{на межі квадрата},
\end{aligned}
\]

для довільної кількості вузлів сітки. Якщо знаєте мову програмування,
напишіть програму для виконання цих обчислень.

\section{ЯВНІ РІЗНИЦЕВІ
СХЕМИ}\label{ux44fux432ux43dux456-ux440ux456ux437ux43dux438ux446ux435ux432ux456-ux441ux445ux435ux43cux438}

МЕТА ЛЕКЦІЇ: Ввести поняття явної різницевої схеми і показати, як її
можна застосовувати до задач теплопровідності та хвильових задач.
Основна ідея полягає в тому, щоб замінити похідні скінченними різницями
і отримати формули, які безпосередньо виражають розв'язок на новому
часовому шарі через значення на попередньому шарі.

На відміну від стаціонарних еліптичних задач, де значення у внутрішніх
вузлах сітки визначаються одночасно, для задач, що залежать від часу,
можна обчислювати розв'язок послідовно:
\(t=\Delta t,\ 2\Delta t,\ 3\Delta t,\dots\)

\section{Явна схема для рівняння
теплопровідності}\label{ux44fux432ux43dux430-ux441ux445ux435ux43cux430-ux434ux43bux44f-ux440ux456ux432ux43dux44fux43dux43dux44f-ux442ux435ux43fux43bux43eux43fux440ux43eux432ux456ux434ux43dux43eux441ux442ux456}

Розглянемо задачу теплопровідності для стрижня, початкова температура
якого дорівнює нулю. Температура лівого кінця фіксована, а на правому
кінці виконується умова теплообміну з навколишнім середовищем:

\pandocbounded{\includegraphics[keepaspectratio]{images/_page_288_Figure_3.jpeg}}

Фіг. 38.1. Сітка для рівняння теплопровідності.

Нехай температура середовища задається функцією \(g(t)\). Тоді маємо
задачу

\[
\begin{aligned}
(38.1)\quad (\text{УЧП})\quad & u_t=u_{xx}, && 0<x<1,\quad 0<t<\infty,\\
(\text{ГУ})\quad & u(0,t)=1,\quad u_x(1,t)=-\left[u(1,t)-g(t)\right], && 0<t<\infty,\\
(\text{НУ})\quad & u(x,0)=0, && 0\le x\le 1.
\end{aligned}
\]

Побудуємо прямокутну сітку з вузлами

\[
\begin{aligned}
x_j&=jh,\qquad j=0,1,2,\dots,n,\\
t_i&=ik,\qquad i=0,1,2,\dots,m.
\end{aligned}
\]

Позначимо через \(u_{i,j}\) наближене значення функції \(u(x_j,t_i)\).
Значення на лівій межі та на початковому шарі відомі з умов задачі. Щоб
знайти решту значень, замінимо похідні наближеннями

\[
\begin{aligned}
u_t(x_j,t_i)&\approx \frac{u_{i+1,j}-u_{i,j}}{k},\\
u_{xx}(x_j,t_i)&\approx \frac{u_{i,j+1}-2u_{i,j}+u_{i,j-1}}{h^2}.
\end{aligned}
\]

Підставляючи ці вирази в рівняння теплопровідності, дістаємо явну
формулу

\[
(38.2)\qquad
u_{i+1,j}=u_{i,j}+\frac{k}{h^2}\left[u_{i,j+1}-2u_{i,j}+u_{i,j-1}\right].
\]

Щоб врахувати праву граничну умову, замінимо похідну \(u_x(1,t)\) лівою
різницею:

\[
\frac{u_{i,n}-u_{i,n-1}}{h}=-\left[u_{i,n}-g_i\right],
\qquad g_i=g(ik).
\]

Звідси отримуємо

\[
(38.4)\qquad
u_{i,n}=\frac{u_{i,n-1}+h g_i}{1+h}.
\]

Формули (38.2) і (38.4) дають повну явну схему для послідовного переходу
від одного часового шару до наступного.

\section{Алгоритм обчислень за явною
схемою}\label{ux430ux43bux433ux43eux440ux438ux442ux43c-ux43eux431ux447ux438ux441ux43bux435ux43dux44c-ux437ux430-ux44fux432ux43dux43eux44e-ux441ux445ux435ux43cux43eux44e}

КРОК 1. Обчислити значення розв'язку на шарі \(t=\Delta t\) за формулою

\[
u_{2,j}=u_{1,j}+\frac{k}{h^2}\left[u_{1,j+1}-2u_{1,j}+u_{1,j-1}\right],
\qquad j=2,3,\dots,n-1.
\]

КРОК 2. Значення на правій межі знайти з (38.4):

\[
u_{2,n}=\frac{u_{2,n-1}+h g_2}{1+h}.
\]

КРОК 3. Для переходу до шарів \(t=2\Delta t,\ 3\Delta t,\dots\)
повторювати кроки 1 і 2, щоразу використовуючи значення з попереднього
часового шару.

\pandocbounded{\includegraphics[keepaspectratio]{images/_page_290_Figure_3.jpeg}}

Фіг. 38.2. Ілюстрація явної різницевої схеми.

\section{ЗАУВАЖЕННЯ}\label{ux437ux430ux443ux432ux430ux436ux435ux43dux43dux44f-11}

\begin{enumerate}
\def\labelenumi{\arabic{enumi}.}
\tightlist
\item
  Явна схема має істотний недолік: якщо часовий крок занадто великий
  порівняно з кроком за просторовою координатою, похибки округлення
  можуть швидко зрости і зруйнувати чисельний розв'язок. Для рівняння
  теплопровідності зазвичай потрібно, щоб
\end{enumerate}

\[
\frac{k}{h^2}\le 0.5.
\]

\begin{enumerate}
\def\labelenumi{\arabic{enumi}.}
\setcounter{enumi}{1}
\item
  Зменшення кроків \(\Delta t\) і \(\Delta x\) зазвичай підвищує
  точність, але водночас збільшує обсяг обчислень. Залежність повної
  похибки від кроку сітки показано на рис. 38.5.
\item
  Для гіперболічної задачі
\end{enumerate}

\[
\begin{aligned}
(\text{УЧП})\quad & u_{tt}=u_{xx}, && 0<x<1,\quad 0<t<\infty,\\
(\text{ГУ})\quad & u(0,t)=g_1(t),\quad u(1,t)=g_2(t), && 0<t<\infty,\\
(\text{НУ})\quad & u(x,0)=\varphi(x),\quad u_t(x,0)=\psi(x), && 0\le x\le 1
\end{aligned}
\]

також можна побудувати явну різницеву схему. Для цього використовують
центральні різниці

\[
\begin{aligned}
u_{tt}(x,t)&\approx \frac{u(x,t+k)-2u(x,t)+u(x,t-k)}{k^2},\\
u_{xx}(x,t)&\approx \frac{u(x+h,t)-2u(x,t)+u(x-h,t)}{h^2}.
\end{aligned}
\]

а початкову умову для швидкості наближують рівністю

\[
u_t(x,0)\approx \frac{u(x,k)-u(x,0)}{k}
=\frac{u(x,k)-\varphi(x)}{k}.
\]

Звідси отримуємо явну схему

\[
(38.5)\qquad
u(x,t+k)=2u(x,t)-u(x,t-k)+\left(\frac{k}{h}\right)^2
\left[u(x+h,t)-2u(x,t)+u(x-h,t)\right].
\]

Щоб почати обчислення, потрібно спочатку знайти значення на шарі
\(t=\Delta t\):

\[
u(x,k)=\varphi(x)+k\psi(x).
\]

\pandocbounded{\includegraphics[keepaspectratio]{images/_page_291_Figure_4.jpeg}}

Фіг. 38.3. Алгоритм розв'язання гіперболічної задачі за явною схемою.

\pandocbounded{\includegraphics[keepaspectratio]{images/_page_292_Figure_6.jpeg}}

Фіг. 38.4. Використання масивів при реалізації явної схеми.

\section{ЗАВДАННЯ}\label{ux437ux430ux432ux434ux430ux43dux43dux44f-17}

\begin{enumerate}
\def\labelenumi{\arabic{enumi}.}
\tightlist
\item
  Побудуйте явну скінченнорізницеву схему для задачі
\end{enumerate}

\[
\begin{aligned}
(\text{УЧП})\quad & u_t=u_{xx}, && 0<x<1,\quad 0<t<\infty,\\
(\text{ГУ})\quad & u(0,t)=0,\quad u(1,t)=0, && 0<t<\infty,\\
(\text{НУ})\quad & u(x,0)=\sin(\pi x), && 0\le x\le 1.
\end{aligned}
\]

Вважаючи \(h=\Delta x=0.1\), знайдіть чисельний розв'язок при
\(t_1=0.005\), \(t_2=0.010\), \(t_3=0.015\) і побудуйте графік розв'язку
для \(t=0.015\).

\pandocbounded{\includegraphics[keepaspectratio]{images/_page_293_Figure_6.jpeg}}

Фіг. 38.5. Залежність повної похибки від кроку сітки.

\begin{enumerate}
\def\labelenumi{\arabic{enumi}.}
\setcounter{enumi}{1}
\item
  Розв'яжіть задачу 1 аналітично методом розділення змінних і порівняйте
  значення у вузлах сітки з чисельним розв'язком.
\item
  Побудуйте блок-схему для чисельного розв'язання наведеної вище
  гіперболічної задачі.
\item
  Розв'яжіть задачу 1 з новою граничною умовою при \(x=1\):
\end{enumerate}

\[
u_x(1,t)=-\left[u(1,t)-1\right].
\]

\section{СХЕМИ НЕЯВНИХ РІЗНИЦЬ (СХЕМА
КРАНКА-НІКОЛСОНА)}\label{ux441ux445ux435ux43cux438-ux43dux435ux44fux432ux43dux438ux445-ux440ux456ux437ux43dux438ux446ux44c-ux441ux445ux435ux43cux430-ux43aux440ux430ux43dux43aux430-ux43dux456ux43aux43eux43bux441ux43eux43dux430}

МЕТА ЛЕКЦІЇ: Показати, як задачі, що залежать від часу, можна
розв'язувати за допомогою неявної різницевої схеми. У такій схемі
значення на новому часовому шарі вже не виражаються явно через
попередній шар; замість цього на кожному кроці потрібно розв'язувати
систему алгебраїчних рівнянь.

Перевага неявних схем над явними полягає в тому, що вони дозволяють
брати значно більший часовий крок без руйнування розв'язку через
накопичення похибок округлення. Для параболічних задач ми використаємо
відому схему Кранка-Ніколсона.

Як уже зазначалося в попередній лекції, явна схема для рівняння
теплопровідності працює надійно лише тоді, коли кроки сітки \(\Delta t\)
і \(\Delta x\) задовольняють умову

\[
\frac{\Delta t}{(\Delta x)^2}\le 0.5.
\]

Інакше схема стає чисельно нестабільною. Наприклад, якщо
\(\Delta x=0.1\), то часовий крок не повинен перевищувати
\(0.5\Delta x^2=0.005\), і для переходу від \(t=0\) до \(t=1\) потрібно
зробити 200 кроків.

Неявні схеми дають можливість працювати з більшими часовими кроками.
Ціна за це --- необхідність на кожному часовому кроці розв'язувати
систему алгебраїчних рівнянь.

\section{Неявна схема для рівняння
теплопровідності}\label{ux43dux435ux44fux432ux43dux430-ux441ux445ux435ux43cux430-ux434ux43bux44f-ux440ux456ux432ux43dux44fux43dux43dux44f-ux442ux435ux43fux43bux43eux43fux440ux43eux432ux456ux434ux43dux43eux441ux442ux456}

Розглянемо задачу

\[
\begin{aligned}
(39.2)\quad (\text{УЧП})\quad & u_t=u_{xx}, && 0<x<1,\quad 0<t<\infty,\\
(\text{ГУ})\quad & u(0,t)=0,\quad u(1,t)=0, && 0<t<\infty,\\
(\text{НУ})\quad & u(x,0)=1, && 0<x<1.
\end{aligned}
\]

Використаємо такі наближення для часткових похідних:

\[
\begin{aligned}
u_t(x,t) &\cong \frac{u(x,t+k)-u(x,t)}{k},\\
u_{xx}(x,t) &\cong
\frac{\lambda}{h^2}\bigl[u(x+h,t+k)-2u(x,t+k)+u(x-h,t+k)\bigr]\\
&\qquad +\frac{1-\lambda}{h^2}\bigl[u(x+h,t)-2u(x,t)+u(x-h,t)\bigr].
\end{aligned}
\]

Тут параметр \(\lambda\) належить відрізку \([0,1]\). При
\(\lambda=0.5\) отримуємо схему Кранка-Ніколсона, а при \(\lambda=0\)
--- звичайну явну схему.

Після підстановки цих апроксимацій у рівняння теплопровідності отримуємо

\[
\frac{u_{i+1,j}-u_{i,j}}{k}
=
\frac{\lambda}{h^2}\left(u_{i+1,j+1}-2u_{i+1,j}+u_{i+1,j-1}\right)
+
\frac{1-\lambda}{h^2}\left(u_{i,j+1}-2u_{i,j}+u_{i,j-1}\right).
\]

Позначаючи \(r=k/h^2\) і переносячи всі невідомі з верхнього часового
шару в ліву частину, одержуємо різницеву схему

\[
(39.4)\qquad
-\lambda r\,u_{i+1,j+1}
+(1+2\lambda r)u_{i+1,j}
-\lambda r\,u_{i+1,j-1}
=
r(1-\lambda)u_{i,j+1}
+[1-2r(1-\lambda)]u_{i,j}
+r(1-\lambda)u_{i,j-1}.
\]

Якщо індекс \(i\) зафіксовано, а \(j\) пробігає значення
\(2,3,\dots,n-1\), то співвідношення (39.4) задають систему з \(n-2\)
рівнянь для \(n-2\) невідомих \(u_{i+1,2},u_{i+1,3},\dots,u_{i+1,n-1}\).

\pandocbounded{\includegraphics[keepaspectratio]{images/_page_296_Figure_5.jpeg}}

Фіг. 39.1. Сітка для схеми неявних різниць.

\pandocbounded{\includegraphics[keepaspectratio]{images/_page_296_Figure_7.jpeg}}

Фіг. 39.2. Шаблон неявної схеми \(u_{i+1,4},\ldots,u_{i+1,n-1}\).

Фіг. 39.2 дає візуальне уявлення про структуру кожного рівняння системи
(39.4).

\section{Алгоритм розв'язання задачі
(39.2)}\label{ux430ux43bux433ux43eux440ux438ux442ux43c-ux440ux43eux437ux432ux44fux437ux430ux43dux43dux44f-ux437ux430ux434ux430ux447ux456-39.2}

КРОК 1. Виберіть значення параметра \(\lambda\), де \(0\le\lambda\le1\).

КРОК 2. Нехай, наприклад, \(h=\Delta x=0.2\) і \(k=\Delta t=0.08\). Тоді
\(r=k/h^2=2\). Візьмемо \(\lambda=0.5\); це і є схема Кранка-Ніколсона.
Для перших чотирьох внутрішніх вузлів нового часового шару одержуємо
систему

\[
\begin{aligned}
3u_{22}-u_{23} &= 1,\\
-u_{22}+3u_{23}-u_{24} &= 1,\\
-u_{23}+3u_{24}-u_{25} &= 1,\\
-u_{24}+3u_{25} &= 1.
\end{aligned}
\]

У матричній формі це записується так:

\[
(39.5)\qquad
\begin{bmatrix}
3 & -1 & 0 & 0\\
-1 & 3 & -1 & 0\\
0 & -1 & 3 & -1\\
0 & 0 & -1 & 3
\end{bmatrix}
\begin{bmatrix}
u_{22}\\
u_{23}\\
u_{24}\\
u_{25}
\end{bmatrix}
=
\begin{bmatrix}
1\\
1\\
1\\
1
\end{bmatrix}.
\]

Матриця цієї системи є тридіагональною. Такі системи зручно розв'язувати
методом прогонки. Для загальної системи

\[
\begin{bmatrix}
b_1 & c_1 & 0 & \cdots & 0\\
a_1 & b_2 & c_2 & \cdots & 0\\
0 & a_2 & b_3 & \ddots & \vdots\\
\vdots & \ddots & \ddots & \ddots & c_{n-1}\\
0 & \cdots & 0 & a_{n-1} & b_n
\end{bmatrix}
\begin{bmatrix}
x_1\\
x_2\\
x_3\\
\vdots\\
x_n
\end{bmatrix}
=
\begin{bmatrix}
d_1\\
d_2\\
d_3\\
\vdots\\
d_n
\end{bmatrix}
\]

вводимо коефіцієнти

\[
c_1^*=\frac{c_1}{b_1},\qquad d_1^*=\frac{d_1}{b_1},
\]

а далі рекурентно обчислюємо

\[
c_j^*=\frac{c_j}{b_j-a_{j-1}c_{j-1}^*},\qquad
d_j^*=\frac{d_j-a_{j-1}d_{j-1}^*}{b_j-a_{j-1}c_{j-1}^*},
\qquad j=2,3,\dots,n.
\]

Після цього розв'язок знаходиться зворотною підстановкою:

\[
x_n=d_n^*,\qquad x_j=d_j^*-c_j^*x_{j+1},\qquad j=n-1,n-2,\dots,1.
\]

Застосувавши цей метод до системи (39.5), отримуємо

\[
u_{22}=0.60,\qquad
u_{23}=0.80,\qquad
u_{24}=0.80,\qquad
u_{25}=0.60.
\]

Це наближені значення розв'язку у внутрішніх точках сітки на шарі
\(t=\Delta t\). Щоб перейти до наступного моменту часу, потрібно знову
розв'язати аналогічну тридіагональну систему.

У неявній схемі обсяг обчислень на кожному кроці більший, ніж у явній,
зате хорошу точність можна отримати навіть при значно більшому часовому
кроці.

\section{ЗАВДАННЯ}\label{ux437ux430ux432ux434ux430ux43dux43dux44f-18}

\begin{enumerate}
\def\labelenumi{\arabic{enumi}.}
\item
  Виведіть рівняння (39.4) з різницевої схеми, що передує йому.
\item
  Побудуйте неявну скінченнорізницеву схему для задачі
\end{enumerate}

\[
\begin{aligned}
(\text{УЧП})\quad & u_t=u_{xx}, && 0<x<1,\quad 0<t<\infty,\\
(\text{ГУ})\quad & u(0,t)=1,\quad u_x(1,t)+u(1,t)=g(t), && 0<t<\infty,\\
(\text{НУ})\quad & u(x,0)=0, && 0\le x\le 1.
\end{aligned}
\]

\begin{enumerate}
\def\labelenumi{\arabic{enumi}.}
\setcounter{enumi}{2}
\tightlist
\item
  Як за допомогою неявної схеми розв'язати задачу
\end{enumerate}

\[
\begin{aligned}
(\text{УЧП})\quad & u_t=u_{xx}+u, && 0<x<1,\quad 0<t<\infty,\\
(\text{ГУ})\quad & u(0,t)=0,\quad u(1,t)=0, && 0<t<\infty,\\
(\text{НУ})\quad & u(x,0)=\sin(\pi x), && 0\le x\le 1?
\end{aligned}
\]

\begin{enumerate}
\def\labelenumi{\arabic{enumi}.}
\setcounter{enumi}{3}
\item
  Як виглядає шаблон рівняння (39.4), коли \(\lambda=1\)?
\item
  Побудуйте блок-схему для розв'язання рівняння теплопровідності (39.2).
  Якщо можливо, напишіть комп'ютерну програму. Корисно також розв'язати
  цю задачу чисельно для різних значень \(\lambda\) і порівняти
  результати з аналітичним розв'язком.
\end{enumerate}

\section{Лекція 40. АНАЛІТИЧНІ РОЗВ'ЯЗКИ VERSUS ЧИСЛОВІ
РІШЕННЯ}\label{ux43bux435ux43aux446ux456ux44f-40.-ux430ux43dux430ux43bux456ux442ux438ux447ux43dux456-ux440ux43eux437ux432ux44fux437ux43aux438-versus-ux447ux438ux441ux43bux43eux432ux456-ux440ux456ux448ux435ux43dux43dux44f}

\pandocbounded{\includegraphics[keepaspectratio]{images/_page_301_Figure_4.jpeg}}

Рис. 40.1. Розв'язок рівняння теплопровідності (40.1) у різні моменти
часу.

\subsection{Переваги аналітичних
розв'язків}\label{ux43fux435ux440ux435ux432ux430ux433ux438-ux430ux43dux430ux43bux456ux442ux438ux447ux43dux438ux445-ux440ux43eux437ux432ux44fux437ux43aux456ux432}

\begin{enumerate}
\def\labelenumi{\arabic{enumi}.}
\item
  Аналітичний розв'язок дозволяє обчислювати точне значення функції в
  будь-якій точці, а не лише чисельно наближене.
\item
  Аналітичне розв'язання завжди дозволяє обчислити значення розв'язку в
  одній точці \((x, t)\) без необхідності обчислювати значення
  розв'язків в інших точках, як це відбувається при розв'язанні задач із
  явними або неявними різницевими схемами.
\item
  Аналітичне рішення дозволяє визначити розв'язок у будь-якій точці, а
  не лише на вузлах сітки.
\item
  Найважливішою перевагою аналітичного розв'язку для нас є можливість
  простежити вплив фізичних параметрів, початкових і граничних умов на
  природу розв'язку.
\end{enumerate}

Чисельні методи не виявляють цих регулярностей, оскільки дозволяють
знаходити розв'язок лише за заданих параметрів --- початкових і
граничних умовах. Іноді важливо знати взаємозв'язок між параметрами
моделі та розв'язком, особливо коли йдеться про оцінку фізичних
параметрів за типом розв'язку. Наприклад, припустимо, що розв'язок \(u\)
визначається експериментально, і ми знаємо аналітичний розв'язок:

\[
u = \text{функція від параметрів}.
\]

Тоді можна поставити питання про визначення параметрів як функцій
вхідних даних за схемою:

\[
\text{Параметри} = \text{функція від } u = \text{функція вхідних даних}.
\]

Така задача називається \textbf{параметричною ідентифікацією}. Для
параметричної ідентифікації необхідно вміти розв'язувати рівняння з
частковими похідними. Трохи пізніше в цій лекції наведемо приклад
параметричної ідентифікації в біології, але спочатку розглянемо переваги
числових розв'язань.

\subsection{Переваги числових
розв'язань}\label{ux43fux435ux440ux435ux432ux430ux433ux438-ux447ux438ux441ux43bux43eux432ux438ux445-ux440ux43eux437ux432ux44fux437ux430ux43dux44c}

Головна перевага числових розв'язків полягає в тому, що їх можна
отримати навіть якщо аналітичне рішення неможливо. Майже всі нелінійні
диференціальні рівняння в частинних похідних потрібно розв'язувати
числовими методами, і більшість реальних моделей фізики, хімії, біології
тощо є нелінійними. Як правило, лінійні моделі наближаються до
нелінійних лише якщо в них відкинути нелінійні члени. Наприклад, ось
кілька нелінійних рівнянь:

\begin{enumerate}
\def\labelenumi{\arabic{enumi}.}
\item
  нелінійне хвильове рівняння \(u_{tt} = u_{xx} + f(u)\);
\item
  рівняння реакції з нелінійною дифузією \(u_t = u_{xx} + f(u)\);
\item
  система Ходжкіна--Хакслі (складна нелінійна система, що описує
  генерацію нервових імпульсів).
\end{enumerate}

Аналітичні розв'язки для цих рівнянь невідомі ні при яких нелінійних
функціях \(f\) і \(g\). Тому загальний підхід до розв'язання нелінійних
(а в ряді випадків і лінійних) задач зазвичай базується на числових
розв'язаннях.

Тепер давайте розглянемо приклад того, як можна використати аналітичне
рішення для визначення фізичних параметрів.

\subsection{Параметрична ідентифікація (в
біології)}\label{ux43fux430ux440ux430ux43cux435ux442ux440ux438ux447ux43dux430-ux456ux434ux435ux43dux442ux438ux444ux456ux43aux430ux446ux456ux44f-ux432-ux431ux456ux43eux43bux43eux433ux456ux457}

Припустимо, що біолог намагається визначити швидкість, з якою іони калію
\((K^+)\) дифундують у екзоплазмовому розчині. Якщо знати цей коефіцієнт
дифузії, можна багато сказати про те, як нервові імпульси передаються
вздовж аксонів. Проблема в тому, що це співвідношення майже неможливо
знайти прямим вимірюванням. Однак можливо знайти математичний зв'язок
між концентрацією поташу \(u(x,t)\) і коефіцієнтом дифузії \(D\), а
потім за вимірюваннями величини \(u(x,t)\) визначити величину \(D\).

Біологи Ходжкін і Кейз виявили, що після розміщення гігантського аксону
у спеціальному соляному розчині концентрація радіоактивного калію
(\({}^{42}K\)) вздовж аксону приблизно описується кривою в початковий
момент:

\[
u(x, 0) = A \, e^{-x^2/a},
\]

де \(A\) і \(a\) --- параметри, що підбираються так, щоб крива
відповідала експериментальним даним.

\pandocbounded{\includegraphics[keepaspectratio]{images/_page_303_Figure_6.jpeg}}

Рис. 40.2. Початкова концентрація іонів (\({}^{42}K\)): а --- теоретична
крива; б --- експериментальні точки.

З часом іони поташу під дією дифузії та конвекції розтікаються вздовж
аксону. Ходжкін і Кейз припустили, що концентрація \(u\) є розв'язком
наступної задачі з початковими даними для рівняння дифузії:

(40.3) \[
\begin{cases}
u_t = D u_{xx} - V u_x, \quad -\infty < x < \infty,\ t > 0, \\
u(x, 0) = A\, e^{-x^2/a}.
\end{cases}
\]

Цю проблему вирішили в лекції 15, переключившись на рухому систему
координат. Загальна схема методу виглядає так: спочатку відкидаємо
конвективний член і отримуємо розв'язок суто дифузійної задачі (при
\(V = 0\)):

\[
u_0(x, t) = \frac{A\sqrt{a}}{\sqrt{a+4Dt}}\, e^{-x^2/(a+4Dt)},
\]

а потім враховуємо конвективний член і отримуємо:

\[
u(x, t) = \frac{A\sqrt{a}}{\sqrt{a+4Dt}}\, e^{-[(x-Vt)^{2}/(a+4Dt)]}.
\]

Параметри \(A\), \(a\), \(D\) і \(V\) визначаються з умови найкращого
узгодження з експериментом (наприклад, методом найменших квадратів):

\[
SS = \sum_{i=1}^{n} [y_i - (a + bx_i)]^2.
\]

\begin{center}\rule{0.5\linewidth}{0.5pt}\end{center}

\section{Лекція 41. КЛАСИФІКАЦІЯ РІВНЯНЬ З ЧАСТКОВИМИ ПОХІДНИМИ ДРУГОГО
ПОРЯДКУ ТА ЗВЕДЕННЯ ДО КАНОНІЧНОЇ
ФОРМИ}\label{ux43bux435ux43aux446ux456ux44f-41.-ux43aux43bux430ux441ux438ux444ux456ux43aux430ux446ux456ux44f-ux440ux456ux432ux43dux44fux43dux44c-ux437-ux447ux430ux441ux442ux43aux43eux432ux438ux43cux438-ux43fux43eux445ux456ux434ux43dux438ux43cux438-ux434ux440ux443ux433ux43eux433ux43e-ux43fux43eux440ux44fux434ux43aux443-ux442ux430-ux437ux432ux435ux434ux435ux43dux43dux44f-ux434ux43e-ux43aux430ux43dux43eux43dux456ux447ux43dux43eux457-ux444ux43eux440ux43cux438}

Загальне лінійне рівняння другого порядку з двома незалежними змінними
\(x\) і \(y\) має вигляд:

\[
A u_{xx} + B u_{xy} + C u_{yy} + D u_x + E u_y + F u = G,
\]

де коефіцієнти \(A, B, C, D, E, F, G\) можуть залежати від \(x\) і
\(y\). Класифікація рівняння визначається знаком дискримінанта
\(B^2 - 4AC\):

\begin{itemize}
\tightlist
\item
  \(B^2 - 4AC < 0\) → \textbf{еліптичне} (наприклад, рівняння Лапласа);
\item
  \(B^2 - 4AC = 0\) → \textbf{параболічне} (наприклад, рівняння
  теплопровідності);
\item
  \(B^2 - 4AC > 0\) → \textbf{гіперболічне} (наприклад, хвильове
  рівняння).
\end{itemize}

\subsection{Зведення до канонічної форми (загальна
схема)}\label{ux437ux432ux435ux434ux435ux43dux43dux44f-ux434ux43e-ux43aux430ux43dux43eux43dux456ux447ux43dux43eux457-ux444ux43eux440ux43cux438-ux437ux430ux433ux430ux43bux44cux43dux430-ux441ux445ux435ux43cux430}

Щоб привести рівняння до канонічної форми, вводять нові змінні
\(\xi = \xi(x, y)\), \(\eta = \eta(x, y)\). Після заміни змінних
рівняння набирає вигляду:

(41.2) \[
\overline{A} u_{\xi\xi} + \overline{B} u_{\xi\eta} + \overline{C} u_{\eta\eta} + \overline{D} u_{\xi} + \overline{E} u_{\eta} + \overline{F} u = \overline{G},
\]

де нові коефіцієнти виражаються через старі та часткові похідні заміни:

\[
\overline{A} = A\xi_x^2 + B\xi_x\xi_y + C\xi_y^2,
\]

\[
\overline{B} = 2A\xi_x\eta_x + B(\xi_x\eta_y + \xi_y\eta_x) + 2C\xi_y\eta_y,
\]

\[
\overline{C} = A\eta_x^2 + B\eta_x\eta_y + C\eta_y^2,
\]

\[
\overline{D} = A\xi_{xx} + B\xi_{xy} + C\xi_{yy} + D\xi_x + E\xi_y,
\]

\[
\overline{E} = A\eta_{xx} + B\eta_{xy} + C\eta_{yy} + D\eta_x + E\eta_y,
\]

\[
\overline{F} = F, \quad \overline{G} = G.
\]

\subsection{Приклад зведення до канонічної форми (гіперболічний
випадок)}\label{ux43fux440ux438ux43aux43bux430ux434-ux437ux432ux435ux434ux435ux43dux43dux44f-ux434ux43e-ux43aux430ux43dux43eux43dux456ux447ux43dux43eux457-ux444ux43eux440ux43cux438-ux433ux456ux43fux435ux440ux431ux43eux43bux456ux447ux43dux438ux439-ux432ux438ux43fux430ux434ux43eux43a}

Розглянемо рівняння

\[
u_{xx} - x^2 u_{yy} = 0.
\]

Знаходимо: \(A = 1\), \(B = 0\), \(C = -x^2\), значить
\(B^2 - 4AC = 4x^2 > 0\) --- рівняння гіперболічне. Характеристичні
рівняння:

\[
\frac{dy}{dx} = \pm x \quad \Rightarrow \quad y = \pm \frac{x^2}{2} + c.
\]

Нові координати:

\[
\xi = y - \frac{x^2}{2}, \quad \eta = y + \frac{x^2}{2}.
\]

Після підстановки в формули для \(\overline{A}\), \(\overline{B}\),
\(\overline{C}\):

\[
\overline{A} = 0, \quad \overline{B} \neq 0, \quad \overline{C} = 0,
\]

і рівняння набирає канонічної гіперболічної форми:

\[
u_{\xi\eta} = \frac{-(x^2+y^2)\,u_\xi + (y^2-x^2)\,u_\eta}{8x^2 y^3}.
\]

\subsection{Новa система координат для параболічного
прикладу}\label{ux43dux43eux432a-ux441ux438ux441ux442ux435ux43cux430-ux43aux43eux43eux440ux434ux438ux43dux430ux442-ux434ux43bux44f-ux43fux430ux440ux430ux431ux43eux43bux456ux447ux43dux43eux433ux43e-ux43fux440ux438ux43aux43bux430ux434ux443}

Нова система координат \(\xi = y - x\), \(\eta = y\) (рис. 41.1):

\pandocbounded{\includegraphics[keepaspectratio]{images/_page_309_Figure_3.jpeg}}

Рис. 41.1. Нова система координат \(\xi = y - x\), \(\eta = y\).

\textbf{ШАГ 2.} Залишилось знайти канонічну форму рівняння. Підставляючи
\(\xi\) і \(\eta\) у формули для коефіцієнтів \(\overline{A}\),
\(\overline{B}\), \(\overline{C}\), \(\overline{D}\), \(\overline{E}\),
\(\overline{F}\) і \(\overline{G}\), отримуємо:

\[\overline{A} = 0 \quad (\text{мало бути нулем; ми знаходили } \xi \text{ з умови рівності нулю цього коефіцієнта}),\]

\[\overline{B} = 0 \quad (\text{ми вже раніше показали, що цей коефіцієнт теж обертається в нуль}),\]

\[\overline{C} \neq 0.\]

Нові допоміжні коефіцієнти (41.5):

\[
\overline{D} = A\xi_{xx} + B\xi_{xy} + C\eta_{yy} + D\xi_x + E\xi_y = 0,
\]

\[
\overline{E} = A\eta_{xx} + B\eta_{xy} + C\eta_{yy} + D\eta_x + E\eta_y = 2(y^2 - x^2),
\]

\[
\widetilde{E} = A\eta_{xx} + B\eta_{xy} + C\eta_{yy} + D\eta_x + E\eta_y = 2(y^2 - x^2),
\]

\[
\tilde{F} = F = 0, \quad \overline{G} = G = 0.
\]

Отже, нове рівняння

\[
u_{\xi\eta} = \frac{\eta\, u_\xi - \xi\, u_\eta}{2(\xi^2 - \eta^2)}.
\]

є загальним розв'язком відповідного рівняння.

\subsection{Конвективна дифузія --- приклад застосування канонічних
координат}\label{ux43aux43eux43dux432ux435ux43aux442ux438ux432ux43dux430-ux434ux438ux444ux443ux437ux456ux44f-ux43fux440ux438ux43aux43bux430ux434-ux437ux430ux441ux442ux43eux441ux443ux432ux430ux43dux43dux44f-ux43aux430ux43dux43eux43dux456ux447ux43dux438ux445-ux43aux43eux43eux440ux434ux438ux43dux430ux442}

Задача (15.2):

\[
\begin{cases}
u_t = D u_{xx} - V u_x, \quad -\infty < x < \infty,\ t > 0, \\
u(x, 0) = H(x),
\end{cases}
\]

де \(H(x)\) --- функція Хевісайда. Замінимо змінні:

\[
\xi = x - Vt, \quad \eta = t.
\]

Розв'язок задачі конвективної дифузії:

\[
u(x, t) = \begin{cases} 0, & t < x/V, \\ P, & t \ge x/V. \end{cases}
\]

Початковий розподіл концентрації зображено на рис. 15.4. Цей результат
легко інтерпретувати, якщо уявити, що ми рухаємося відносно графіка,
показаного на рис. 15.6.

\pandocbounded{\includegraphics[keepaspectratio]{images/_page_319_Figure_14.jpeg}}

\begin{center}\rule{0.5\linewidth}{0.5pt}\end{center}

\section{Лекція 42. ОБЧИСЛЕННЯ ЗА МЕТОДОМ
МОНТЕ-КАРЛО}\label{ux43bux435ux43aux446ux456ux44f-42.-ux43eux431ux447ux438ux441ux43bux435ux43dux43dux44f-ux437ux430-ux43cux435ux442ux43eux434ux43eux43c-ux43cux43eux43dux442ux435-ux43aux430ux440ux43bux43e}

\textbf{МЕТА ЛЕКЦІЇ:} Показати, що деякі інтеграли можна оцінити
імовірнісним шляхом --- методом Монте-Карло --- і що ця ідея
переноситься на розв'язок рівнянь в частинних похідних.

\pandocbounded{\includegraphics[keepaspectratio]{images/_page_315_Figure_7.jpeg}}

Фіг. 42.2. Обчислення інтеграції Монте-Карло.

Очевидно, якщо випадковим чином кинути близько 100 дротиків у
прямокутник \(R\), що містить графік функції, можна отримати оцінку
значення інтеграла, помноживши площу прямокутника \(R\) на відносну
кількість дротиків під графіком функції. Результат нашої гри:

\[
I \approx \text{(площа } R) \times \frac{\text{кількість дротиків під графіком}}{\text{загальна кількість дротиків}}.
\]

Щоб виконати ці обчислення на комп'ютері, нам потрібно якимось чином
отримати послідовність випадкових точок, тобто наказати комп'ютеру
«кидати дротики». Блок-схема на рис. 42.3 показує, як комп'ютер вирішить
цю проблему.

\pandocbounded{\includegraphics[keepaspectratio]{images/_page_316_Figure_2.jpeg}}

\textbf{РИС. 42.3.} Блок-схема розрахунку інтегралу
\(\int_{a}^{b} f(x)\,dx\) методом Монте-Карло (100 кидків);
\(M = \max_a f(x)\).

\subsection{Випадкові
числа}\label{ux432ux438ux43fux430ux434ux43aux43eux432ux456-ux447ux438ux441ux43bux430}

Перш ніж перейти до методу Монте-Карло для рівнянь з частковими
похідними, зупинимося на важливому питанні --- випадкових числах. Щойно,
обчислюючи інтеграл, нам потрібно було побудувати послідовність
випадкових точок \(P_i = (x_i, y_i)\) всередині прямокутника \(R\).
Іншими словами, \(x\)-координата точки \(P\) має бути випадковим числом
з інтервалу \([a, b]\), а \(y_i\) --- з інтервалу \([0, M]\). Щоб знайти
випадкові числа з цих інтервалів, звертаємося до послідовності
випадкових чисел, рівномірно розподілених на відрізку \([0, 1]\). Тоді
числа, рівномірно розподілені на відрізку \([a, b]\), можна обчислити за
формулою

\[
x = a + (b - a)\,r_i,
\]

де \(r_i \in [0, 1]\) --- рівномірно розподілене випадкове число.

Тепер, звичайно, виникає питання: як отримати послідовність випадкових
чисел \(\{r_i: i=1,2,\ldots\}\), рівномірно розподілених на \([0,1]\)?
Найпоширенішим методом сьогодні є \textbf{метод конгруентних
відрахувань}. Використовуючи цей метод, ми отримуємо послідовність
випадкових цілих чисел, наприклад: \(2120,\ 1401,\ 177,\ 3013, \ldots\),
а потім ліворуч від кожного числа розміщуємо десяткову крапку:
\(0{,}212;\ 0{,}1401;\ 0{,}0177;\ 0{,}3013; \ldots\)

\subsection{Отримання випадкових чисел за допомогою методу
відрахування}\label{ux43eux442ux440ux438ux43cux430ux43dux43dux44f-ux432ux438ux43fux430ux434ux43aux43eux432ux438ux445-ux447ux438ux441ux435ux43b-ux437ux430-ux434ux43eux43fux43eux43cux43eux433ux43eux44e-ux43cux435ux442ux43eux434ux443-ux432ux456ux434ux440ux430ux445ux443ux432ux430ux43dux43dux44f}

\textbf{1.} Як перше випадкове ціле число, оберіть будь-яке число між
\(0\) і \(P\) (число \(P\) обирається наперед).

\textbf{2.} Помножте це випадкове число на деяке фіксоване число ---
попередньо обраний множник \(M\).

\textbf{3.} Додайте до добутку деяке фіксоване ціле число \(K\), обране
заздалегідь.

\textbf{4.} Поділіть отриману суму на \(P\) і візьміть залишок --- це і
є нове випадкове число.

Метод відрахування записується у вигляді формули:

\[
r_{i+1} \equiv (M r_i + K) \bmod P.
\]

Це означає, що для отримання наступного \(r_{i+1}\) потрібно взяти
попередній \(r_i\), помножити його на \(M\), додати \(K\), поділити суму
на \(P\) і взяти решту ділення.

\textbf{Наприклад}, якщо \(P = 100\), \(M = 37\), \(K = 16\),
\(r_0 = 15\):

\[
r_1 \equiv (37 \cdot 15 + 16) \bmod 100 = 571 \bmod 100 = 71,
\]

\[
M r_1 + K = 37 \cdot 71 + 16 = 2643, \quad r_2 = 2643 \bmod 100 = 43, \quad \ldots
\]

\subsection{Приклади інтегралів для методу
Монте-Карло}\label{ux43fux440ux438ux43aux43bux430ux434ux438-ux456ux43dux442ux435ux433ux440ux430ux43bux456ux432-ux434ux43bux44f-ux43cux435ux442ux43eux434ux443-ux43cux43eux43dux442ux435-ux43aux430ux440ux43bux43e}

\[
I = \int_{0}^{1} e^{\sin x}\, dx.
\]

\[
I = \int_{0}^{1}\int_{0}^{1}\int_{0}^{1} e^{-(x^{2}+y^{2}+z^{2})}\, dx\, dy\, dz.
\]

\[
I = \int_{0}^{1}\int_{0}^{1}\int_{0}^{1}\int_{0}^{1} e^{-(x^{2}+y^{2}+z^{2}+\omega^{2})}\, dx\, dy\, dz\, d\omega.
\]

\subsection{ЗАВДАННЯ}\label{ux437ux430ux432ux434ux430ux43dux43dux44f-19}

\begin{enumerate}
\def\labelenumi{\arabic{enumi}.}
\item
  Яку послідовність випадкових чисел дає формула
  \(r_{i+1} \equiv (3r_i + 4) \bmod 7\), де \(r_0 = 0\)?
\item
  Як отримати послідовність випадкових чисел, розподіл яких показаний на
  рисунку?
\end{enumerate}

\pandocbounded{\includegraphics[keepaspectratio]{images/_page_319_Figure_16.jpeg}}

Іншими словами, як отримати послідовність чисел \(\{0, 1, 2\}\), якщо
ймовірність виникнення \(0\) і \(2\) дорівнює \(0{,}25\), а ймовірність
появи \(1\) --- \(0{,}5\)?

\begin{center}\rule{0.5\linewidth}{0.5pt}\end{center}

\section{Лекція 43. РОЗВ'ЯЗАННЯ РІВНЯНЬ З ЧАСТКОВИМИ ПОХІДНИМИ МЕТОДОМ
МОНТЕ-КАРЛО}\label{ux43bux435ux43aux446ux456ux44f-43.-ux440ux43eux437ux432ux44fux437ux430ux43dux43dux44f-ux440ux456ux432ux43dux44fux43dux44c-ux437-ux447ux430ux441ux442ux43aux43eux432ux438ux43cux438-ux43fux43eux445ux456ux434ux43dux438ux43cux438-ux43cux435ux442ux43eux434ux43eux43c-ux43cux43eux43dux442ux435-ux43aux430ux440ux43bux43e}

\textbf{МЕТА ЛЕКЦІЇ:} Показати, як побудувати таку гру, результатом якої
є приблизне розв'язання диференціального рівняння. Одна з цих ігор
(модель випадкової ходьби) веде до скінченно-різницевої апроксимації
задачі Діріхле у квадраті. Цю модель можна узагальнити для отримання
розв'язків інших задач.

У попередній лекції було сказано, що можливо винайти такі ігри на удачу,
результатом яких є розв'язки (можливо, приблизні) задач для рівнянь з
частковими похідними. У цій лекції ми покажемо, як побудувати таку гру
для наближеного розв'язання задачі Діріхле:

\[
\begin{cases}
u_{xx} + u_{yy} = 0, & 0 < x < 1,\ 0 < y < 1, \\
u(x,y) = g(x,y), & \text{на межі квадрата.}
\end{cases}
\]

Після того, як цю конкретну задачу буде розв'язано методом Монте-Карло,
ми розглянемо, як розв'язувати більш загальні задачі за допомогою цього
методу.

Щоб проілюструвати метод Монте-Карло, розглянемо гру під назвою
\textbf{«Блукаючий п'яниця»}. Щоб грати в цю гру, потрібна таблиця з
сіткою ліній, як показано на рис. 43.1.

Перейдемо до опису правил гри.

\subsection{Як грати в «Блукаючого
п'яницю»?}\label{ux44fux43a-ux433ux440ux430ux442ux438-ux432-ux431ux43bux443ux43aux430ux44eux447ux43eux433ux43e-ux43fux44fux43dux438ux446ux44e}

\textbf{КРОК 1.} Блукання п'яниці починаються з довільної точки сітки (у
нашому випадку --- точки \(A\)).

\pandocbounded{\includegraphics[keepaspectratio]{images/_page_321_Figure_4.jpeg}}

\begin{itemize}
\tightlist
\item
  Початкова точка --- внутрішні точки сітки --- □ --- Кінцева точка
  \(p_i\);
\item
  \(g_i\) --- винагороди в т. \(p_i\). Фіг. 43.1. Ігрове поле для гри
  «Блукаючий п'яниця».
\end{itemize}

\textbf{КРОК 2.} У кожному ході гри п'яний випадково переходить до
однієї з чотирьох суміжних точок сітки (у нашому випадку сусідні точки
для \(A\) --- це точки \(B\), \(C\), \(D\) і \(E\), і ймовірність
досягти кожної з цих точок становить \(0{,}25\)).

\textbf{Таблиця 43.1}

Ймовірність закінчення випадкового блукання в граничній точці \(p_i\) і
величина винагороди \(g_i\)

{\def\LTcaptype{none} % do not increment counter
\begin{longtable}[]{@{}
  >{\raggedright\arraybackslash}p{(\linewidth - 4\tabcolsep) * \real{0.3333}}
  >{\raggedright\arraybackslash}p{(\linewidth - 4\tabcolsep) * \real{0.3333}}
  >{\raggedright\arraybackslash}p{(\linewidth - 4\tabcolsep) * \real{0.3333}}@{}}
\toprule\noalign{}
\begin{minipage}[b]{\linewidth}\raggedright
Гранична точка \(p_i\)
\end{minipage} & \begin{minipage}[b]{\linewidth}\raggedright
\(P_A(p_i)\) --- відносна частка блукань, що закінчились у точці \(p_i\)
\end{minipage} & \begin{minipage}[b]{\linewidth}\raggedright
\(g_i\) --- величина винагороди при досягненні точки \(p_i\)
\end{minipage} \\
\midrule\noalign{}
\endhead
\bottomrule\noalign{}
\endlastfoot
1 & 0,04 & 1 \\
2 & 0,15 & 1 \\
3 & 0,03 & 1 \\
4 & 0,06 & 0 \\
5 & 0,17 & 0 \\
6 & 0,05 & 0 \\
7 & 0,06 & 0 \\
8 & 0,15 & 0 \\
9 & 0,03 & 0 \\
10 & 0,06 & 0 \\
11 & 0,16 & 0 \\
12 & 0,04 & 0 \\
\end{longtable}
}

\textbf{КРОК 3.} Після переходу до наступної точки процес блукання
відновлюється. Таким чином, п'яниця ходить від точки до точки, поки
випадково не опиниться на межі \(p_t\). Тут він зупиняється, і ми
записуємо номер точки \(p_t\). Це кінець однієї випадкової прогулянки.

\textbf{КРОК 4.} Повторимо кроки 1--3 достатньо разів і знайдемо
відносну кількість прибуттів п'яниці в кожній з граничних точок \(p_i\).
Таблиця 43.1 показує типові результати, отримані внаслідок 100
випадкових прогулянок.

\textbf{КРОК 5.} Припустимо, що п'яниця отримує винагороду
\(g_i = g(p_i)\) (значення граничної умови в точці \(p_i\), якщо він
досягає точки \(p_i\)), і припустимо, що мета гри --- обчислити середню
\(R(A)\) винагороду за всі випадкові прогулянки, починаючи з точки
\(A\). Цей середній виграш визначається формулою:

\[
R(A) = g_1 P_A(p_1) + g_2 P_A(p_2) + \dots + g_{12} P_A(p_{12}).
\]

Гра закінчується, коли знайдено значення \(R(A)\). Відповідно до даних
таблиці 43.1:

\[
R(A) = 1\cdot(0{,}04) + 1\cdot(0{,}15) + 1\cdot(0{,}03) + 0\cdot(0{,}06) + \dots + 0\cdot(0{,}04).
\]

\subsection{Для чого грати в «Блукаючого
п'яницю»?}\label{ux434ux43bux44f-ux447ux43eux433ux43e-ux433ux440ux430ux442ux438-ux432-ux431ux43bux443ux43aux430ux44eux447ux43eux433ux43e-ux43fux44fux43dux438ux446ux44e}

Виявляється, що середня винагорода \(R(A)\) є приблизним розв'язком
задачі Діріхле в точці \(A\). Це цікаве спостереження базується на двох
фактах.

\textbf{1.} Припустимо, що п'яниця почав свій шлях з точки, що лежить на
межі. Кожна така прогулянка одразу закінчується в одному й тому ж місці,
і він отримує винагороду \(g_i\). Отже, середня винагорода за кожну
граничну точку дорівнює \(g_i\).

\textbf{2.} Тепер припустимо, що прогулянка починається з внутрішньої
точки. Тоді середня винагорода \(R(A)\) буде арифметичним середнім
середніх нагород за чотири сусідні бали:

\[
R(A) = \frac{1}{4}\bigl[R(B) + R(C) + R(D) + R(E)\bigr].
\]

Давайте запитаємо себе: чому середнє значення винагороди \(R(A)\)
наближається до розв'язання задачі Діріхле в точці \(A\)? Ми бачили, що
значення \(R(A)\) задовольняє двом рівнянням:

\[
R(A) = \frac{1}{4}\bigl[R(B) + R(C) + R(D) + R(E)\bigr] \quad \text{(у внутрішніх точках)},
\]

\[
R(A) = g_i \quad \text{(у граничних точках)}.
\]

Якщо \(g_i\) --- значення функції \(g(x, y)\) з граничної умови у
граничних точках \(p_i\), то наші два рівняння точно збігаються з двома
рівняннями, отриманими при розв'язанні задачі Діріхле методом скінченних
різниць:

\[
u_{i,j} = \frac{1}{4}\bigl[u_{i-1,j} + u_{i+1,j} + u_{i,j-1} + u_{i,j+1}\bigr], \quad (i,j) \text{ — внутрішня точка},
\]

\[
u_{ij} = g_{ij}, \quad (i,j) \text{ — гранична точка}.
\]

Отже, значення \(R(A)\) дійсно наближає розв'язок рівняння Лапласа в
точці \(A\).

\subsection{Розв'язання рівняння Лапласа за методом
Монте-Карло}\label{ux440ux43eux437ux432ux44fux437ux430ux43dux43dux44f-ux440ux456ux432ux43dux44fux43dux43dux44f-ux43bux430ux43fux43bux430ux441ux430-ux437ux430-ux43cux435ux442ux43eux434ux43eux43c-ux43cux43eux43dux442ux435-ux43aux430ux440ux43bux43e}

Наступні правила дозволяють отримати розв'язок в одній внутрішній точці
квадрата.

\begin{itemize}
\tightlist
\item
  \textbf{КРОК 1.} Зробимо кілька випадкових прогулянок, починаючи з
  точки \(A\) і закінчуючи однією з граничних точок. Відстежуємо,
  скільки разів прогулянки закінчувалися на кожній межі.
\item
  \textbf{КРОК 2.} Після завершення всіх прогулянок для кожної граничної
  точки обчислюємо відносну кількість прогулянок, які закінчилися в цій
  точці. Позначимо ці значення \(P_A(p_i)\).
\item
  \textbf{КРОК 3.} Обчислюємо наближений розв'язок \(u(A)\) за формулою:
\end{itemize}

\[
u(A) = g_1 P_A(p_1) + g_2 P_A(p_2) + \dots + g_N P_A(p_N),
\]

де \(g_i\) --- значення функції \(g\) у точці \(p_i\), а \(N\) ---
кількість граничних точок.

\subsection{Зауваження}\label{ux437ux430ux443ux432ux430ux436ux435ux43dux43dux44f-12}

\begin{enumerate}
\def\labelenumi{\arabic{enumi}.}
\item
  Зверніть увагу, що якщо відносні частоти \(P_A(p_i)\) відомі, то
  значення \(u(A)\) за інших граничних умов \(g_i\) легко знайти. Для
  цього достатньо підставити нові \(g_i\) у формулу --- не потрібно
  наново розіграти випадкові прогулянки.
\item
  У багатьох випадках дослідник зацікавлений у розв'язанні
  диференціального рівняння в певній точці. Якщо межа області є складною
  і рівняння залежить від трьох або чотирьох змінних, метод Монте-Карло
  може бути зручним. Методи Монте-Карло виникли для розв'язання дуже
  складних задач дифузії нейтронів. Ці задачі не можна розв'язати
  аналітично.
\end{enumerate}

\subsection{Розв'язання задачі Діріхле зі змінними
коефіцієнтами}\label{ux440ux43eux437ux432ux44fux437ux430ux43dux43dux44f-ux437ux430ux434ux430ux447ux456-ux434ux456ux440ux456ux445ux43bux435-ux437ux456-ux437ux43cux456ux43dux43dux438ux43cux438-ux43aux43eux435ux444ux456ux446ux456ux454ux43dux442ux430ux43cux438}

Розглянемо еліптичну крайову задачу в квадраті:

\[
u_{xx} + (\sin x)\, u_{yy} = 0, \quad 0 < x < \pi,\ 0 < y < \pi,
\]

\[
u(x, y) = g(x, y) \quad \text{на межі квадрата.}
\]

Щоб розв'язати цю задачу, замінимо \(u_{xx}\), \(u_{yy}\) центральними
різницевими похідними:

\[
u_{xx} = u_{i,\,j+1} - 2u_{i,\,j} + u_{i,\,j-1}, \quad u_{yy} = u_{i+1,\,j} - 2u_{i,\,j} + u_{i-1,\,j}.
\]

Підstavляємо у рівняння і розв'язуємо за \(u_{i,j}\):

\[
u_{i,j} = \frac{u_{i,j+1} + u_{i,j-1} + \sin x_j\,(u_{i+1,j} + u_{i-1,j})}{2(1+\sin x_j)}.
\]

Коефіцієнти при \(u_{i+1,j}\), \(u_{i,j+1}\), \(u_{i,j-1}\),
\(u_{i-1,j}\) є додатними і в сумі дорівнюють одиниці. Отже, \(u_{i,j}\)
є зваженим середнім значень у чотирьох сусідніх точках.

\pandocbounded{\includegraphics[keepaspectratio]{images/_page_324_Picture_9.jpeg}}

Рис. 43.2. Вузли сітки для випадкового маршруту.

Можна \textbf{модифікувати Блукаючого П'яницю} так, щоб ймовірності
переміщення до сусідніх точок дорівнювали коефіцієнтам відповідних
членів. Якщо п'яниця знаходиться в точці \((i, j)\), то він переходить у
точку:

\[
(i, j+1) \text{ з імовірністю } \frac{1}{2(1+\sin x_j)},
\]

\[
(i, j-1) \text{ з імовірністю } \frac{1}{2(1+\sin x_j)},
\]

\[
(i+1, j) \text{ з імовірністю } \frac{\sin x_j}{2(1+\sin x_j)},
\]

\[
(i-1, j) \text{ з імовірністю } \frac{\sin x_j}{2(1+\sin x_j)}.
\]

Інакше гра не змінюється. Модифікація для інших задач може бути більш
винахідливою, але основна ідея залишається незмінною.

\subsection{ЗАВДАННЯ}\label{ux437ux430ux432ux434ux430ux43dux43dux44f-20}

\begin{enumerate}
\def\labelenumi{\arabic{enumi}.}
\item
  Побудуйте блок-схему гри «Блукаючий п'яниця», щоб розв'язати задачу
  \(u_{xx} + u_{yy} = 0\) у вузлах внутрішньої сітки.
\item
  Напишіть програму для розв'язання задачі відповідно до блок-схеми з
  завдання 1.
\item
  Як модифікувати гру для вирішення задачі:
\end{enumerate}

\[
\begin{cases}
u_{xx} + u_{yy} = 0, & 0 < x < 1,\ 0 < y < 1, \\
u(x, y) = g(x, y) & \text{на межі?}
\end{cases}
\]

Як виглядає в цій грі випадкова прогулянка?

\begin{enumerate}
\def\labelenumi{\arabic{enumi}.}
\setcounter{enumi}{3}
\tightlist
\item
  Чи можна так модифікувати гру, щоб розв'язати задачу:
\end{enumerate}

\[
\begin{cases}
u_{xx} + u_{yy} = 0, & 0 < x < 1,\ 0 < y < 1, \\
u(x, 1) = 0, & u(x, 0) = 0, \\
u(0, y) = 1, & \tfrac{\partial u}{\partial x}(1, y) = 0?
\end{cases}
\]

\begin{enumerate}
\def\labelenumi{\arabic{enumi}.}
\setcounter{enumi}{4}
\tightlist
\item
  Як модифікувати гру для розв'язання задачі:
\end{enumerate}

\[
\begin{aligned}
u_t &= \alpha^2 u_{xx} - \beta u_x - \gamma u, & 0 < x < 1,\ 0 < t < \infty, \\
u(0,t) &= f(t),\quad u(1,t) = g(t), & 0 < t < \infty, \\
u(x,0) &= \varphi(x), & 0 \le x \le 1?
\end{aligned}
\]

(ПРИМІТКА. Замінимо це рівняння скінченною різницею Кранка--Ніколсона,
виразимо \(u_{i+1,j}\) через п'ять значень у суміжних точках.)

\begin{center}\rule{0.5\linewidth}{0.5pt}\end{center}

\section{Лекція 44. ВАРІАЦІЙНЕ ЧИСЛЕННЯ (РІВНЯННЯ
ЕЙЛЕРА--ЛАГРАНЖА)}\label{ux43bux435ux43aux446ux456ux44f-44.-ux432ux430ux440ux456ux430ux446ux456ux439ux43dux435-ux447ux438ux441ux43bux435ux43dux43dux44f-ux440ux456ux432ux43dux44fux43dux43dux44f-ux435ux439ux43bux435ux440ux430ux43bux430ux433ux440ux430ux43dux436ux430}

\textbf{МЕТА ЛЕКЦІЇ:} Ввести поняття функціоналу (функції від функції) і
пояснити, як функціонали природно виникають у фізиці. Типовий функціонал
є інтегралом

\[
J[y] = \int_a^b F(x,\, y,\, y')\, dx,
\]

що є функцією від функції \(y\) (припускається, що підінтегральна
функція \(F(x, y, y')\) задана). Ми покажемо, як знайти функцію
\(\overline{y}(x)\), яка мінімізує функціонал \(J[y]\). Виявляється, що
мінімізуюча функція \(\overline{y}\) має задовольняти так зване
\textbf{рівняння Ейлера--Лагранжа}.

\subsection{Вступ. Задача
брахістохрони}\label{ux432ux441ux442ux443ux43f.-ux437ux430ux434ux430ux447ux430-ux431ux440ux430ux445ux456ux441ux442ux43eux445ux440ux43eux43dux438}

Варіаційне числення виникло одночасно з математичним аналізом у зв'язку
з розв'язанням задач максимізації та мінімізації функцій функцій
(функціоналів). Першою задачею варіаційного числення була задача
брахістохрони, сформульована Йоганном Бернуллі у 1696 році.

Бернуллі показав, що час сходження записується у вигляді:

\[
T\{y\} = \int_a^b \sqrt{\frac{1+y'^2}{2g\,y}}\, dx,
\]

і, отже, є функцією від функції. Оскільки багато функціоналів мають
такий вигляд, ми зосередимося на вивченні загального функціоналу

(44.1) \[
J[y] = \int_a^b F(x,\, y,\, y')\, dx.
\]

\pandocbounded{\includegraphics[keepaspectratio]{images/_page_329_Figure_5.jpeg}}

Фіг. 44.1. Задача брахістохрони (час сходження треба мінімізувати).

\subsection{Мінімізація
функціоналу}\label{ux43cux456ux43dux456ux43cux456ux437ux430ux446ux456ux44f-ux444ux443ux43dux43aux446ux456ux43eux43dux430ux43bux443}

Розглянемо задачу пошуку функції \(y(x)\), яка мінімізує функціонал
\(J[y]\) у класі гладких функцій, що задовольняють граничні умови
\(y(a) = A\), \(y(b) = B\) (рис. 44.2).

Нехай існує \(y = \overline{y}(x)\) --- потрібна функція, і розглянемо
малу варіацію функції \(\overline{y}\), тобто функцію
\(\overline{y} + \varepsilon\,\eta(x)\), де \(\varepsilon\) --- мале
число, а \(\eta(x)\) --- гладка функція, яка задовольняє граничні умови
\(\eta(a) = \eta(b) = 0\) (рис. 44.2).

\pandocbounded{\includegraphics[keepaspectratio]{images/_page_329_Figure_7.jpeg}}

Рис. 44.2. Варіація функції: а --- варіація функції
\(\overline{y} + \varepsilon\,\eta(x)\); б --- мінімізуюча функція
\(\overline{y}(x)\).

Оскільки \(\overline{y}\) мінімізує \(J\), то

\[
\phi(\varepsilon) = J[\overline{y} + \varepsilon\,\eta] \ge J[\overline{y}] = \phi(0) \quad \text{для всіх } \varepsilon.
\]

Іншими словами, графік
\(\phi(\varepsilon) = J[\overline{y}+\varepsilon\eta]\) як функція
\(\varepsilon\) буде виглядати як рис. 44.3.

\pandocbounded{\includegraphics[keepaspectratio]{images/_page_330_Figure_6.jpeg}}

Рис. 44.3. Графік функції \(J[\overline{y}+\varepsilon\eta]\) в околиці
\(\varepsilon = 0\).

З рис. 44.3 очевидно, що в точці мінімуму:

\[
\phi'(0) = \frac{d}{d\varepsilon} J[\overline{y} + \varepsilon\,\eta] \bigg|_{\varepsilon=0} = 0.
\]

Диференціюючи під знаком інтеграла й інтегруючи за частинами, приходимо
до висновку, що друга частина підінтегрального виразу має дорівнювати
нулю.

(44.2) \[
\frac{\partial F}{\partial y} - \frac{d}{dx}\!\left(\frac{\partial F}{\partial y'}\right) = 0 \qquad \textbf{(Рівняння Ейлера–Лагранжа)}.
\]

Рівняння (44.2) називається \textbf{рівнянням Ейлера--Лагранжа}, і хоча
загалом воно виглядає складно, коли в нього підставляється конкретна
функція \(F(x, y, y')\), воно стає звичайним диференціальним рівнянням
другого порядку відносно невідомої функції \(\overline{y}(x)\).

\subsection{Теорема (необхідна умова мінімуму
функціоналу)}\label{ux442ux435ux43eux440ux435ux43cux430-ux43dux435ux43eux431ux445ux456ux434ux43dux430-ux443ux43cux43eux432ux430-ux43cux456ux43dux456ux43cux443ux43cux443-ux444ux443ux43dux43aux446ux456ux43eux43dux430ux43bux443}

Якщо функція \(y(x)\) мінімізує функціонал

\[
J[y] = \int_a^b F(x, y, y')\, dx
\]

у класі гладких функцій з граничними умовами \(y(a) = A\), \(y(b) = B\),
то вона повинна задовольняти рівняння Ейлера--Лагранжа:

\[
\frac{\partial F}{\partial y} - \frac{d}{dx}\!\left(\frac{\partial F}{\partial y'}\right) = 0, \quad y(a) = A,\ y(b) = B.
\]

\subsection{\texorpdfstring{Приклад: мінімальний функціонал
\(J[y] = \int_{0}^{1}[y^2 + y'^2]\,dx\)}{Приклад: мінімальний функціонал J{[}y{]} = \textbackslash int\_\{0\}\^{}\{1\}{[}y\^{}2 + y\textquotesingle\^{}2{]}\textbackslash,dx}}\label{ux43fux440ux438ux43aux43bux430ux434-ux43cux456ux43dux456ux43cux430ux43bux44cux43dux438ux439-ux444ux443ux43dux43aux446ux456ux43eux43dux430ux43b-jy-int_01y2-y2dx}

Спробуємо знайти функцію \(y(x)\), яка проходить через точки \((0,0)\) і
\((1,1)\) і мінімізує функціонал:

\[
J[y] = \int_{0}^{1}[y^2 + y'^2]\, dx, \qquad y(0) = 0,\ y(1) = 1.
\]

Тут \(F(x, y, y') = y^2 + y'^2\), отже:

\[
\frac{\partial F}{\partial y} = 2y, \quad \frac{\partial F}{\partial y'} = 2y'.
\]

Рівняння Ейлера--Лагранжа набирає вигляду:

\[
2y - \frac{d}{dx}(2y') = 0 \quad \Rightarrow \quad y'' - y = 0.
\]

Загальний розв'язок --- \(y = C_1 e^x + C_2 e^{-x}\). Розв'язуючи з
граничними умовами \(y(0)=0\) та \(y(1)=1\), отримуємо:

\[
\overline{y}(x) = \frac{\sinh x}{\sinh 1} = 0{,}42\,e^x - 0{,}42\,e^{-x}.
\]

\pandocbounded{\includegraphics[keepaspectratio]{images/_page_332_Figure_5.jpeg}}

Фіг. 44.4. Усі допустимі гладкі криві задовольняють граничні умови
\(y(0)=0\) та \(y(1)=1\): а --- допустима крива; б --- мінімізуюча
функція \(y(x) = 0{,}42\,e^x - 0{,}42\,e^{-x}\).

\subsection{Зауваження}\label{ux437ux430ux443ux432ux430ux436ux435ux43dux43dux44f-13}

\begin{enumerate}
\def\labelenumi{\arabic{enumi}.}
\item
  Рівняння Ейлера--Лагранжа аналогічне рівності нульової похідної в
  диференціальному численні. Рівність нульової похідної є лише
  \textbf{необхідною}, але не достатньою умовою екстремуму. Те саме
  стосується рівняння Ейлера--Лагранжа. Однак дуже часто розв'язок
  рівняння Ейлера--Лагранжа фактично дає локальний або навіть глобальний
  мінімум.
\item
  Якщо \(y(x)\) мінімізує функціонал
\end{enumerate}

\[
J[y] = \int_a^b 2\pi y \sqrt{1 + y'^2}\, dx,
\]

тоді площа поверхні, утворена обертанням кривої \(y(x)\) навколо осі
\(x\), буде мінімальною (рис. 44.5). Розв'язок рівняння Ейлера--Лагранжа
в цьому випадку:

\[
y(x) = \alpha \cosh\!\left(\frac{x - \beta}{\alpha}\right) \quad \textbf{(ланцюгова лінія)},
\]

де константи \(\alpha\) і \(\beta\) визначаються з умов \(y(a) = A\) та
\(y(b) = B\).

\begin{enumerate}
\def\labelenumi{\arabic{enumi}.}
\setcounter{enumi}{2}
\tightlist
\item
  У цій лекції ми виявили, що функція
\end{enumerate}

\[
\Delta u = 0
\]

доставляє мінімум функціоналу \(\Phi[u]\).

\pandocbounded{\includegraphics[keepaspectratio]{images/_page_333_Picture_8.jpeg}}

Рис. 44.5. Мінімальна поверхня обертання.

\begin{enumerate}
\def\labelenumi{\arabic{enumi}.}
\setcounter{enumi}{3}
\tightlist
\item
  Основні закони фізики найчастіше формулюються мовою варіаційних
  принципів, а не диференціальних рівнянь. Прикладами є:
\end{enumerate}

\begin{itemize}
\tightlist
\item
  \textbf{Принцип Ферма}: світло, поширюючись з однієї точки в іншу,
  обирає шлях, якому відповідає найкоротший час поширення.
\item
  \textbf{Принцип Гамільтона}: у консервативному полі частинка рухається
  так, що інтеграл дії
\end{itemize}

\[
\int_{t_1}^{t_2} (E_k - E_p)\, dt \quad \text{(кінетична енергія мінус потенціальна)}
\]

буде стаціонарним. Звідси отримують диференціальне рівняння руху
матеріальної точки.

Основні ідеї варіаційного числення можна розширити на функціонали, які
залежать від функцій кількох незалежних змінних.

\subsection{ЗАВДАННЯ}\label{ux437ux430ux432ux434ux430ux43dux43dux44f-21}

\begin{enumerate}
\def\labelenumi{\arabic{enumi}.}
\tightlist
\item
  Показати, що в класі гладких функцій, що задовольняють граничні умови
  \(y(0)=0\) та \(y(\pi/2)=1\), мінімальний функціонал має вигляд:
\end{enumerate}

\[
J[y] = \int_0^{\pi/2} [y'^2 - y^2]\, dx.
\]

\begin{enumerate}
\def\labelenumi{\arabic{enumi}.}
\setcounter{enumi}{1}
\item
  Якщо \(y(x)\) мінімізує функціонал
  \(J[y] = \int_a^b \sqrt{1+y'^2}\,dx\), то крива \(y(x)\) між двома
  точками є відрізком прямої лінії.
\item
  Покажіть, що ланцюгова лінія \(y(x) = \alpha\cosh((x-\beta)/\alpha)\)
  задовольняє рівняння Ейлера--Лагранжа для функціоналу
  \(J[y] = \int_a^b y\sqrt{1+y'^2}\,dx\).
\item
  Принцип Гамільтона стверджує, що матеріальна точка рухатиметься так,
  що інтеграл дії буде стаціонарним. Отримайте диференціальне рівняння
  руху матеріальної точки.
\end{enumerate}

\begin{center}\rule{0.5\linewidth}{0.5pt}\end{center}

\section{Лекція 45. МЕТОД
РІТЦА}\label{ux43bux435ux43aux446ux456ux44f-45.-ux43cux435ux442ux43eux434-ux440ux456ux442ux446ux430}

\textbf{Мета лекції:} Показати, як задачі на мінімум функціоналів можна
розв'язувати наближено, замінюючи нескінченновимірний простір допустимих
функцій скінченновимірним підпростором.

\textbf{Наближений розв'язок} шукається у вигляді:

\[
u_n(x,y) = a_1\varphi_1 + a_2\varphi_2 + \dots + a_n\varphi_n.
\]

де \(\{\varphi_k\}\) --- вибрані базисні функції (наприклад,
тригонометричні поліноми або поліноми Лежандра). Коефіцієнти \(a_k\)
обираються так, щоб значення функціоналу \(J[u_n]\) було мінімальним.

Зверніть увагу, що в перших двох задачах коефіцієнти є сталими. Читач
повинен розуміти, що параметр \(b\) має бути достатньо малим, а
необережний ряд може відрізнятися від точного розв'язку.

\pandocbounded{\includegraphics[keepaspectratio]{images/_page_343_Figure_2.jpeg}}

Рис. 45.2. Блок-схема методу Рітца.

Чи знаєте ви, як виглядає графік цієї функції на відрізку \([0, 1]\)?
Зверніть увагу, що синусоїдальне перетворення функції \(f(x)\)
визначається лише для додатних цілих чисел \(n\).

\subsection{ЗАВДАННЯ}\label{ux437ux430ux432ux434ux430ux43dux43dux44f-22}

\begin{enumerate}
\def\labelenumi{\arabic{enumi}.}
\item
  Підставте розклад (46.8) у задачу (46.7) і отримайте послідовність
  задач \(P_0,\ P_1,\ P_2,\ldots\)
\item
  Покажіть, що нелінійну задачу можна наближено розв'язати методом
  збурень.
\end{enumerate}

\begin{center}\rule{0.5\linewidth}{0.5pt}\end{center}

\section{Лекція 46. РОЗВ'ЯЗАННЯ РІВНЯНЬ З ЧАСТКОВИМИ ПОХІДНИМИ МЕТОДАМИ
ТЕОРІЇ
ЗБУРЕНЬ}\label{ux43bux435ux43aux446ux456ux44f-46.-ux440ux43eux437ux432ux44fux437ux430ux43dux43dux44f-ux440ux456ux432ux43dux44fux43dux44c-ux437-ux447ux430ux441ux442ux43aux43eux432ux438ux43cux438-ux43fux43eux445ux456ux434ux43dux438ux43cux438-ux43cux435ux442ux43eux434ux430ux43cux438-ux442ux435ux43eux440ux456ux457-ux437ux431ux443ux440ux435ux43dux44c}

\textbf{Мета лекції} --- показати, як складні задачі з нелінійностями,
змінними коефіцієнтами або нерегулярними областями можна розв'язувати як
збурення простіших задач.

Наприклад, можна розглядати задачу Діріхле в деформованій області:

\[
\Delta u = 0, \qquad 0 < r < 1 + \tfrac{1}{4}\sin\theta,
\]

з граничною умовою:

\[
u\!\left(1+\tfrac{1}{4}\sin\theta,\,\theta\right) = \cos\theta, \qquad 0 \le \theta \le 2\pi.
\]

Дуже часто задачу, розв'язок якої невідомий, можна шляхом неперервної
зміни деяких параметрів звести до близької, але легко розв'язуваної
задачі. Тоді і розв'язок вихідної задачі буде близьким до розв'язку
модифікованої задачі.

Наприклад, рівняння Лапласа можна неперервно модифікувати в нелінійне
рівняння \(\Delta u + \varepsilon u^2 = 0\) за допомогою сімейства
задач, параметризованих числом \(\varepsilon\), де
\(0 \le \varepsilon \le 1\).

Розглядається область \(0 < r < 1\), \(0 \le \theta < 2\pi\) з граничною
умовою \(u(1,\theta) = g(\theta)\).

Розв'язоk у вигляді ряду за степенями \(\varepsilon\):

\[
u = u_0 + \varepsilon u_1 + \varepsilon^2 u_2 + \dots
\]

Підстановка цього розкладу дає послідовність лінійних задач
\(P_0, P_1, P_2, \ldots\), кожна з яких розв'язується аналітично.

Перше наближення:

\[
u(r,\theta) = u_0(r,\theta) + \tfrac{1}{4}\,u_1(r,\theta),
\]

де:

\[
u_0(r,\theta) = r\cos\theta,
\]

\[
u = u_0 + u_1 = r\cos\theta - \frac{r^4-1}{32} - \frac{(r^4-r^2)}{24}\cos(2\theta).
\]

Фіг. 46.1 показує схему розв'язання нелінійної задачі шляхом її зведення
до розв'язку рівняння Лапласа.

за деяких граничних умов (наприклад, \(g(\theta) = \cos\theta\)) задача
має два різних матеріальні розв'язки, тоді як за інших (наприклад,
\(g(\theta) = A\cos\theta\) з \(A > 20{,}65\)) задача не має реального
розв'язку.

\subsection{ЗАВДАННЯ}\label{ux437ux430ux432ux434ux430ux43dux43dux44f-23}

\begin{enumerate}
\def\labelenumi{\arabic{enumi}.}
\tightlist
\item
  Отримайте перший доданок розкладу для задачі:
\end{enumerate}

\[
\Delta u = 0, \qquad 0 < r < 1 + \epsilon\sin\theta,
\]

\[
u(1+\epsilon\sin\theta, \theta) = \cos\theta, \quad 0 \le \epsilon \le 1/4.
\]

\begin{enumerate}
\def\labelenumi{\arabic{enumi}.}
\setcounter{enumi}{1}
\item
  Розгляньте задачу методом збурень при малому \(\epsilon\) і знайдіть
  наближення \(u_0\) і \(u_1\).
\item
  Покажіть, що наближений розв'язок задачі збурення задовольняє формулу:
\end{enumerate}

\[
f(x+h) = f(x) + f'(x)\,h + f''(x)\frac{h^2}{2!} + \dots
\]

(формула Тейлора, яку використовують для виведення умов збурень).

\begin{enumerate}
\def\labelenumi{\arabic{enumi}.}
\setcounter{enumi}{3}
\tightlist
\item
  Задача параметрів збурення для теплового рівняння:
\end{enumerate}

\[
\begin{cases}
u_t = \alpha^2 u_{xx}, & 0 < x < 1,\ 0 < t < \infty, \\
\begin{cases} u(0,t) = 0, & 0 < t < \infty, \\ u(1,t) = 0, & 0 < t < \infty, \end{cases} \\
\begin{cases} u(x,0) = \sin(\pi x), & 0 \le x \le 1, \\ u_t(x,0) = 0. & 0 \le x \le 1. \end{cases}
\end{cases}
\]

Дайте фізичну інтерпретацію цієї задачі.

\subsection{Лекція 47}\label{ux43bux435ux43aux446ux456ux44f-47-1}

\section{РОЗВ'ЯЗАННЯ ДВОВИМІРНИХ ЗАДАЧ МЕТОДОМ КОНФОРМНИХ
ВІДОБРАЖЕНЬ}\label{ux440ux43eux437ux432ux44fux437ux430ux43dux43dux44f-ux434ux432ux43eux432ux438ux43cux456ux440ux43dux438ux445-ux437ux430ux434ux430ux447-ux43cux435ux442ux43eux434ux43eux43c-ux43aux43eux43dux444ux43eux440ux43cux43dux438ux445-ux432ux456ux434ux43eux431ux440ux430ux436ux435ux43dux44c}

МЕТА ЛЕКЦІЇ: Показати, як деякі двовимірні крайові задачі для рівняння
Лапласа можна звести до простіших задач за допомогою конформного
відображення.

\pandocbounded{\includegraphics[keepaspectratio]{images/_page_353_Figure_9.jpeg}}

Ідея методу полягає в тому, що складну область на площині \((x,y)\)
замінюють простішою областю на площині \((u,v)\) за допомогою
комплексного відображення \(w=f(z)\), де

\[
z = x + iy, \qquad w = u + iv.
\]

Для конформних відображень гармонічність зберігається: якщо
\(\varphi(x,y)\) задовольняє рівняння Лапласа в початковій області, то
після переходу до нових координат ми знову дістаємо рівняння Лапласа,
але вже в геометрично простішій області.

У цій лекції розглянуто:

\begin{enumerate}
\def\labelenumi{\arabic{enumi}.}
\tightlist
\item
  основні властивості конформних відображень;
\item
  приклади побудови таких відображень;
\item
  застосування методу до задач Діріхле.
\end{enumerate}

\section{Комплексні функції та конформні
відображення}\label{ux43aux43eux43cux43fux43bux435ux43aux441ux43dux456-ux444ux443ux43dux43aux446ux456ux457-ux442ux430-ux43aux43eux43dux444ux43eux440ux43cux43dux456-ux432ux456ux434ux43eux431ux440ux430ux436ux435ux43dux43dux44f}

Комплексну функцію записують у вигляді

\[
w = f(z).
\]

Вона переводить точки площини \(z\) у точки площини \(w\). Наприклад,
відображення

\[
w = z^2
\]

переносить перший квадрант площини \(z\) у верхню півплощину площини
\(w\).

\pandocbounded{\includegraphics[keepaspectratio]{images/_page_355_Figure_5.jpeg}}

Рис. 47.1. Площина комплексної змінної \(z\).

\pandocbounded{\includegraphics[keepaspectratio]{images/_page_355_Figure_7.jpeg}}

Рис. 47.2. Відображення першого квадранта площини \(z\) у верхню
півплощину площини \(w\).

Щоб побачити це в дійсних координатах, підставимо \(z=x+iy\) у формулу
\(w=z^2\):

\[
w = (x+iy)^2 = x^2 - y^2 + 2ixy.
\]

Отже,

\[
u = x^2 - y^2, \qquad v = 2xy.
\]

Це і є дійсна форма того самого відображення. Вона показує, як криві у
площині \((x,y)\) переходять у криві в площині \((u,v)\).

\section{Означення конформного
відображення}\label{ux43eux437ux43dux430ux447ux435ux43dux43dux44f-ux43aux43eux43dux444ux43eux440ux43cux43dux43eux433ux43e-ux432ux456ux434ux43eux431ux440ux430ux436ux435ux43dux43dux44f}

Відображення \(w=f(z)\) називається конформним у точці \(z_0\), якщо
\(f'(z_0)\neq 0\). Воно називається конформним в області \(D\), якщо ця
умова виконується в кожній точці області.

Геометрично це означає, що відображення зберігає кути між кривими. Саме
ця властивість робить його корисним для розв'язування крайових задач:
складну межу можна перенести в простішу, не зруйнувавши структуру
задачі.

\section{Рівняння Лапласа у верхній
півплощині}\label{ux440ux456ux432ux43dux44fux43dux43dux44f-ux43bux430ux43fux43bux430ux441ux430-ux443-ux432ux435ux440ux445ux43dux456ux439-ux43fux456ux432ux43fux43bux43eux449ux438ux43dux456}

Розглянемо задачу Діріхле у верхній півплощині. Основна ідея полягає в
тому, щоб підібрати відображення, яке переведе верхню півплощину у
горизонтальну смугу. Для цього використовується класичне перетворення

\[
w = \ln\left(\frac{z-1}{z+1}\right).
\]

Це відображення переносить верхню півплощину у смугу

\[
-\infty < u < \infty, \qquad 0 < v < \pi.
\]

\pandocbounded{\includegraphics[keepaspectratio]{images/_page_357_Figure_11.jpeg}}

Рис. 47.3. Крайова задача у верхній півплощині.

\pandocbounded{\includegraphics[keepaspectratio]{images/_page_358_Figure_2.jpeg}}

Рис. 47.4. Перетворення складної задачі на простішу за допомогою
конформного відображення.

У смузі розв'язок уже знаходиться елементарно. Якщо на одній межі задано
нульовий потенціал, а на іншій одиничний, то

\[
\varphi(u,v) = \frac{v}{\pi}.
\]

Далі потрібно виразити \(v\) через початкові координати \(x\) і \(y\) та
підставити цей вираз у знайдену формулу.

\section{Область між двома
колами}\label{ux43eux431ux43bux430ux441ux442ux44c-ux43cux456ux436-ux434ux432ux43eux43cux430-ux43aux43eux43bux430ux43cux438}

Другий типовий приклад полягає у відображенні області між двома
неконцентричними колами на кільце. Після такого перетворення задача
Діріхле зводиться до стандартної задачі в кільці, яку вже можна
розв'язувати у полярних координатах.

\pandocbounded{\includegraphics[keepaspectratio]{images/_page_359_Figure_10.jpeg}}

Рис. 47.5. Конформне відображення області на кільце.

Суть методу лишається тією самою:

\begin{enumerate}
\def\labelenumi{\arabic{enumi}.}
\tightlist
\item
  знайти зручне конформне відображення;
\item
  перенести задачу у простішу область;
\item
  розв'язати її в нових координатах;
\item
  повернутися до початкових змінних.
\end{enumerate}

Саме тому конформні відображення є одним із найсильніших інструментів
для двовимірних гармонічних задач.

\section{ПІСЛЯМОВА
ВИДАННЯ}\label{ux43fux456ux441ux43bux44fux43cux43eux432ux430-ux432ux438ux434ux430ux43dux43dux44f}

\section{Дорогий
читачу!}\label{ux434ux43eux440ux43eux433ux438ux439-ux447ux438ux442ux430ux447ux443}

Будь ласка, надсилайте свої коментарі щодо змісту книги, якості
перекладу та оформлення за адресою видавництва «Мир».

\section{Фарлоу
Стейпл}\label{ux444ux430ux440ux43bux43eux443-ux441ux442ux435ux439ux43fux43b}

\section{ДИФЕРЕНЦІАЛЬНІ РІВНЯННЯ В ПОХІДНИХ ПОХІДНИХ ДЛЯ НАУКОВЦІВ І
ІНЖЕНЕРІВ}\label{ux434ux438ux444ux435ux440ux435ux43dux446ux456ux430ux43bux44cux43dux456-ux440ux456ux432ux43dux44fux43dux43dux44f-ux432-ux43fux43eux445ux456ux434ux43dux438ux445-ux43fux43eux445ux456ux434ux43dux438ux445-ux434ux43bux44f-ux43dux430ux443ux43aux43eux432ux446ux456ux432-ux456-ux456ux43dux436ux435ux43dux435ux440ux456ux432}

Науковий редактор М. В. Хатунцева. Молодший науковий редактор Е. І.
Качуліна. Художник В. Е. Карпов. Художній редактор В. І. Шаповалов.
Технічний редактор Т. А. Максимова. Коректор Т. І. Стіфсева.

UB No 5289. Подано до публікації 01.07.85. Видавництво «Мир», Москва.

\end{document}
